% This file was converted to LaTeX by Writer2LaTeX ver. 0.5
% see http://www.hj-gym.dk/~hj/writer2latex for more info
\documentclass{article}
\usepackage[ascii]{inputenc}
\usepackage[T1]{fontenc}
\usepackage[english,ngerman]{babel}
\usepackage{amsmath,amssymb,amsfonts,textcomp}
\usepackage{color}
\usepackage{makeidx}
\usepackage{multicol}
\usepackage{array}
\usepackage{supertabular}
\usepackage{hhline}
\usepackage{hyperref}
\hypersetup{colorlinks=true, linkcolor=blue, citecolor=blue, filecolor=blue, pagecolor=blue, urlcolor=blue}
\usepackage{graphicx}
\usepackage{ooomath}
\newcommand\gSub[1]{\ensuremath{{}_{\text{#1}}}}
% Text styles
\newcommand\gEs[1]{\textbf{\textcolor[rgb]{0.0,0.0,0.5019608}{#1}}}
\newcommand\gEs[1]{\textbf{\textcolor[rgb]{0.0,0.0,0.5019608}{#1}}}
\newcommand\gT[1]{\textbf{\textup{\textcolor[rgb]{0.0,0.5019608,0.0}{#1}}}}
\newcommand\gE[1]{\textit{#1}}
\newcommand\gZ[1]{\textcolor[rgb]{0.5019608,0.5019608,0.5019608}{#1}}
\newcommand\gC[1]{\textit{#1}}
\newcommand\gH[1]{\textmd{\textcolor[rgb]{0.0,0.0,0.5019608}{#1}}}
\newcommand\gTt[1]{#1}
\newcommand\gCa[1]{\textbf{\textcolor{black}{#1}}}
\newcommand\gNw[1]{#1}
\makeindex
\setlength\tabcolsep{1mm}
\renewcommand\arraystretch{1.3}
% Non-floating captions
\makeatletter
\newcommand\captionof[1]{\def\@captype{#1}\caption}
\makeatother
\newcounter{Drawing}
\renewcommand\theDrawing{\arabic{Drawing}}
\title{Arbeitsstandards f\"ur den Sanit\"atsdienst, Version f\"ur \"Osterreich, Stand November 2008 (1)}
\begin{document}
\clearpage\setcounter{page}{1}\begin{center}
\tablehead{}\begin{supertabular}{m{3.0279999cm}m{10.0130005cm}m{3.372cm}}
\multicolumn{3}{m{16.813cm}}{\centering
{\color[rgb]{0.7882353,0.0,0.08627451} Der Sanit\"atsdienst}\par

\begin{center}
\begin{minipage}{16.97cm}
\centering {\color[rgb]{0.7882353,0.0,0.08627451} Arbeits- und
Ausbildungsstandards}\par

\centering {\itshape\color[rgb]{0.0,0.0,0.5019608} und eine Einf\"uhrung
in die Welt der Medizin und Sanit\"atshilfe}\par

\end{minipage}
\end{center}
\begin{center}
 [Warning: Image ignored] % Unhandled or unsupported graphics:
%\includegraphics[width=15.432cm,height=11.903cm]{AASdev20090228-img1.png}

\end{center}
}\\\hline
{\mdseries Wien, {}--.-{}-.2008}

 &
 &
\raggedleft Ausgabe 0\par

\\\hline
{\mdseries    [Warning: Image ignored]
% Unhandled or unsupported graphics:
%\includegraphics[width=1.36cm,height=3.962cm]{AASdev20090228-img2.png}
 }

 &
{\color[rgb]{0.29803923,0.29803923,0.29803923} \textbf{Mission Statement
{\textperiodcentered}} Das Ausbildungszentrum (ABZ) ist eine
Dienst\-leistungs\-ein\-richtung des ASB Floridsdorf-Donau\-stadt und
orientiert sein Angebot an den Be\-d\"urf\-nissen unserer Partner und
an den strategischen Zielen des ASB. Wir sind sowohl f\"ur die 
medizinische und sanit\"atsdienstliche Ausbildung unserer Mitarbeiter
als auch f\"ur die Durchf\"uhrung von Schulungen und Lehrg\"angen zur
Schulung der Bev\"olkerung verantwortlich. In Zusammenarbeit mit der
\"Arztlichen Leitung und den \"ubergeordneten Gremien  entwickeln wir
Arbeitsstandards und Lehrmeinungen. }

{\color[rgb]{0.29803923,0.29803923,0.29803923} Um die Qualit\"at der
Arbeit des ASB auf hohem Niveau zu halten verfolgt das ABZ eine aktive
Informations- und Ausbildungspolitik. }

 &
{\centering   [Warning: Image ignored]
% Unhandled or unsupported graphics:
%\includegraphics[width=2.182cm,height=2.208cm]{AASdev20090228-img3.png}
 \par}

\centering {\bfseries\itshape\color[rgb]{0.5019608,0.5019608,0.5019608}
Ausbildungszentrum}\par

\centering {\bfseries\itshape Floridsdorf-Donaustadt}\par

\\\hline
 &
\centering {\bfseries\color{red} DRAFT -- ENTWURF -- DRAFT}\par

 &
\\
\end{supertabular}
\end{center}
\clearpage\setcounter{page}{1}\clearpage\setcounter{page}{1}\setcounter{tocdepth}{1}
\renewcommand\contentsname{Inhalt - Schnell\"ubersicht}
\tableofcontents
\setcounter{tocdepth}{10}
\renewcommand\contentsname{Inhalt - Detail\"ubersicht}
\tableofcontents
{\bfseries
Autoren und Mitwirkende}

{\mdseries\itshape
(in alphabetischer Reihenfolge)}

\begin{multicols}{2}
\begin{center}
\begin{tabular}{m{6.813cm}}
{\mdseries\itshape Josef Emhofer}

{\mdseries\itshape Mag. phil.}

{\mdseries Ausbildungszentrum des ASB Floridsdorf-Donaustadt}

{\mdseries Wallenberggasse 2, A-1220 Wien}

\\
{\mdseries\itshape Gerhard Gabriel}

{\mdseries\itshape Dr. med. univ. }

{\mdseries Facharzt f\"ur Innere Medizin}

{\mdseries A-2103, Langenzersdorf}

\\
{\mdseries\itshape Sebastian Gabriel}

{\mdseries Medizinische Universit\"at Wien -- Besondere Einrichtung zur
Medizinischen Aus- und Weiterbildung}

{\mdseries Spitalgasse 23, A-1090 Wien}

\\
{\mdseries\itshape Lena Hirtler}

{\mdseries Medizinische Universit\"at Wien -- Zentrum f\"ur Anatomie und
Zellbiologie, Abt. f\"ur systematische Anatomie}

{\mdseries W\"ahringerstra{\ss}e 13, A-1090 Wien }

\\
{\mdseries\itshape Thomas Hochreiter}

{\mdseries\itshape DGKP}

{\mdseries Ausbildungszentrum des ASB Floridsdorf-Donaustadt}

{\mdseries Wallenberggasse 2, A-1220 Wien}

\\
{\mdseries\itshape Roman Koch}

{\mdseries\itshape Ing.}

{\mdseries Leiter Ausbildungszentrum des ASB Floridsdorf-Donaustadt}

{\mdseries Wallenberggasse 2, A-1220 Wien}

\\
{\mdseries\itshape Christof Koller}

{\mdseries Ausbildungszentrum des ASB Floridsdorf-Donaustadt}

{\mdseries Wallenberggasse 2, A-1220 Wien}

\\
{\mdseries\itshape Michael Motal}

{\mdseries Ausbildungszentrum des ASB Floridsdorf-Donaustadt}

{\mdseries Wallenberggasse 2, A-1220 Wien}

\\
{\mdseries\itshape Johannes Steuer}

{\mdseries\itshape Dr. med. univ., M.Sc.}

{\mdseries \"Arztlicher Leiter und Medizinisch-wissenschaftlicher Leiter
der Curricula}

{\mdseries ASB Floridsdorf-Donaustadt}

{\mdseries Wallenberggasse 2, A-1220 Wien}

{\mdseries Sozialmedizinisches Zentrum Ost der Stadt Wien --
Donauspital}

{\mdseries Abteilung f\"ur An\"asthesiologie und Intensivmedizin}

{\mdseries Langobardenstra{\ss}e 122, A-1220 Wien}

\\
{\mdseries\itshape Michael Witthofner}

{\mdseries\itshape Dr. med, univ. }

{\mdseries Wien }

\\
\end{tabular}
\end{center}
\end{multicols}
{\bfseries
Danksagungen}

{\mdseries
New Somerset Hospital, Kapstadt, S\"udafrika}

{\mdseries
Western Cape METRO Emergency Medical Service, Kapstadt, S\"udafrika}

{\bfseries
Sponsoren}

{\bfseries
Software}

\begin{center}
\tablehead{}\begin{supertabular}{m{5.0740004cm}m{5.0740004cm}m{5.0740004cm}}
{\mdseries\itshape JabRef 2.4.2 }

 &
\centering Literaturverzeichnis\par

 &
\\
{\mdseries\itshape OpenOffice 3.0}

 &
\centering Textverarbeitung\par

 &
\raggedleft http://www.openoffice.com\par

\\
{\mdseries\itshape Ubuntu Linux 8.04, 8.10}

 &
\centering Betriebssystem\par

 &
\raggedleft http://www.ubuntu.com\par

\\
 &
 &
\\
 &
 &
\\
\end{supertabular}
\end{center}
\section{T\"atigkeitsbild und Orientierung}
{\mdseries\itshape
Autor fehlt}

{\mdseries
Derzeit in Vorbereitung.}

{\mdseries
Der Sanit\"ater in \"Osterreich -- zwischen Sch\"utzengraben und
Flusanit\"ater}

{\mdseries
Der Rettungsdienst in der gro{\ss}en weiten Welt }

{\mdseries
Ein funktionierendes Sozialversicherungssystem ist Grundlage f\"ur
optimale Leistungen}

\section{Betriebsinterne Regelungen}
{\mdseries
Angepasst an: Arbeiter-Samariter-Bund Gruppe Floridsdorf-Donaustadt und 
Arbeiter-Samariter-Bund Floridsdorf-Donaustadt Rettungs-
Krankentransport und soziale Dieste gem. GesmbH.}

{\mdseries
Derzeit in Vorbereitung.}

{\mdseries
\gEs{Siehe externe Beilage. }}


\section{[RS I] Erste Hilfe und erweiterte Erste Hilfe}
{\mdseries\itshape
Autor fehlt}

\begin{itemize}
\item {\mdseries
\"Uberpr\"ufung der Atemt\"atigkeit (bei  \"uberstreckter HWS) durch}

\end{itemize}
\begin{itemize}
\item {\mdseries
SEHEN }

\item {\mdseries
H\"OREN}

\item {\mdseries
F\"UHLEN}

\item {\mdseries
F\"ur mindestens 10  Sekunden!! }

\end{itemize}
\paragraph{Freimachen der Atemwege}
\begin{itemize}
\item {\mdseries
Freimachen der Atemwege}

\end{itemize}
\begin{itemize}
\item {\mdseries
\"Offnen des Mundes -{\textgreater} ESMARCH-Handgriff}

\item {\mdseries
Inspektion}

\item {\mdseries
Ausr\"aumen -{\textgreater}Fremdk\"orper/Erbrochenes durch
Seitw\"artsdrehen des Kopfes}

\item {\mdseries
(falls n\"otig) Absaugen}

\item {\mdseries
einzuf\"uhrende Katheterl\"ange: Distanz Ohrl\"appchen/Mundwinkel}

\item {\mdseries
\"Uberstrecken der HWS}

\item {\mdseries
Kinn/Stirn-Griff}

\item {\mdseries
Zungengrund wird angehoben -{\textgreater} Atemwege  sind frei!}

\end{itemize}
\subsection{AED}
\section{[RS IX] Ger\"atelehre und Sanit\"atstechnik}
{\mdseries\itshape
Thomas Hochreiter [R?]; BEMAW, Sebastian Gabriel}

{\mdseries
Derzeit in Vorbereitung. }

{\mdseries
Kapitelliste nicht vollst\"andig}

{\mdseries\color[rgb]{0.0,0.0,0.5019608}
Blutdruckmessung}

{\mdseries\color{black}
Grundlagen}

{\mdseries\color{black}
Fallstricke}

{\mdseries
Die Blutdruckmanschette soll nach M\"oglichkeit nicht auf der Seite}

\begin{itemize}
\item {\mdseries
eines Shuntarmes bei Dialysepatienten}

\end{itemize}
\begin{itemize}
\item {\mdseries
Ein Dialyse-Shunt ist ein k\"unstlicher Kurzschluss zwischen einer
Armarterie und einer Armvene. Dadurch wird die Vene erweitert und sie
pulsiert, mann kann viel mehr Durchfluss f\"ur die Dialyse
erreichen.\newline
Die Nahtstelle zwischen Vene und Arterie ist jedoch sehr sensibel, wenn
gestaut wird kann sie aufplatzen ${\rightarrow}$ }

\item {\mdseries
${\rightarrow}$ Gefahr einer massiven Shuntblutung}

\end{itemize}
\begin{itemize}
\item {\mdseries
eines Arm\"odems}

\item {\mdseries
einer Brust-OP mit Lymphknoten-Entfernung (\"Odemgefahr)}

\item {\mdseries
einer Halbseitenschw\"ache oder -l\"ahmung }

\end{itemize}
{\mdseries
angelegt werden. }

{\mdseries\color{black}
\gEs{Shunt-Arme} sind \gEs{tabu}
f\"ur die Blutdruckmessung!}

{\mdseries
Bitten Sie im Praktikum am KTW Dialysepatienten h\"oflich Ihnen ihren
Shunt-Arm zu zeigen!}

{\mdseries\color{black}
Ma{\ss}nahmen}

\begin{itemize}
\item {\mdseries
Funktion des Blutdruckmessger\"ates und des Stethoskops pr\"ufen}

\end{itemize}
\begin{itemize}
\item {\mdseries
Kontrolle der Dichte des Ventils, der Verbindungsschl\"auche und der
Manschette (Ventil schlie{\ss}en, Manschette etwas aufpumpen und
Zusammenpressen, Manschette ganz entl\"uften)}

\item {\mdseries
Kontrolle der G\"angigkeit des Ventils (es soll sich leicht auf und
zuschrauben lassen)}

\item {\mdseries
Kontrolle des Manometers auf korrekte Eichung: Die Nadel muss sich im
Ruhezustand genau bei 0mmHg befinden}

\item {\mdseries
Ohrst\"opsel des Stethoskops zeigen schr\"ag nach vorne}

\item {\mdseries
Trichterfunktion durch Ber\"uhren der Membran pr\"ufen}

\end{itemize}
\begin{itemize}
\item {\mdseries
Kontaktaufnahme mit dem Patienten}

\end{itemize}
\begin{itemize}
\item {\mdseries
Standardisierte Blutdruckmessung in Ruhe = konstante Bedingungen + vor
der Messung: kein Nikotin ( 30 min) -- Koffein ( 2 h) -- Alkohol (8 h) 
{}-- k\"orperliche Anstrengung (1h)}

\end{itemize}
\begin{itemize}
\item {\mdseries
Patientenlagerung}

\end{itemize}
\begin{itemize}
\item {\mdseries
Arm freimachen und entspannt auf einen Tisch legen in leichter Flexion
und Supination}

\item {\mdseries
Oberarm auf Herzh\"ohe}

\end{itemize}
\begin{itemize}
\item {\mdseries
Auswahl der geeigneten Manschettengr\"o{\ss}e}

\end{itemize}
\begin{itemize}
\item {\mdseries
Manschette laut Herstellerangaben w\"ahlen z.B Standardmanschette:
Oberarmumfang max. 31 cm}

\item {\mdseries
Messwertverf\"alschung vermeiden! Der aufblasbare Teil der Manschette
soll 80\% des Oberarmumfangs umfassen, die Breite der Manschette
mindestens 1/3 des Oberarmumfangs betragen}

\item {\mdseries
Zu kleine Manschetten f\"uhren zum Messen zu hoher Werte, zu weite
Manschetten k\"onnen nicht am Arm fixiert werden}

\end{itemize}
\begin{itemize}
\item {\mdseries
Handhabung des Stethoskops}

\item {\mdseries
Bestimmen der Anlegestelle f\"ur Manschette und der Auskultationsstelle
durch Palpation}

\end{itemize}
\begin{itemize}
\item {\mdseries
Arteria brachialis medial in der Ellenbeuge aufsuchen}

\end{itemize}
\begin{itemize}
\item {\mdseries
Anlegen der Manschette}

\end{itemize}
\begin{itemize}
\item {\mdseries
Manschette muss auf unbekleideter Haut aufliegen}

\item {\mdseries
m\"a{\ss}ig straff}

\item {\mdseries
Abstand zur Ellenbeuge 2,5 cm}

\item {\mdseries
aufblasbaren Teil \"uber A. brachialis zentrieren}

\item {\mdseries
Schl\"auche nicht im Weg,}

\item {\mdseries
CAVE: Halbseitenschw\"ache/-l\"ahmung, Shuntarm bei Dialysepatienten,
Arm\"odem}

\end{itemize}
\begin{itemize}
\item {\mdseries
Puls f\"uhlen beim Aufpumpen}

\end{itemize}
\begin{itemize}
\item {\mdseries
Schnelles Aufpumpen}

\item {\mdseries
Tasten des Radialispulses mit Finger 2-4}

\item {\mdseries
Aufpumpen bis 30 mmHg \"uber jenen Druck, ab dem kein Radialispuls mehr
getastet wird}

\end{itemize}
\begin{itemize}
\item {\mdseries
Messvorgang}

\end{itemize}
\begin{itemize}
\item {\mdseries
Kopf des Stethoskops auf a. brachialis legen}

\item {\mdseries
Druck langsam ablassen 2 mmHg/sec}

\item {\mdseries
Systolischer Wert: Beginn der turbulenten Korotkoffschen
Str\"omungsger\"ausche}

\item {\mdseries
diastolischer Wert: Verschwinden der Korotkoffschen Ger\"ausche}

\item {\mdseries
Wenn Ger\"ausche \"uberhaupt nicht verschwinden (z.B. bei Kindern,
Schwangeren etc.), so gilt als diastolischer Wert der Messpunkt, an dem
die Ger\"ausche deutlich leiser werden}

\item {\mdseries
Wenn Ger\"ausche verschwunden sind, noch weitere 10mmHg langsam ablassen
-- wenn dann keine Ger\"ausche mehr h\"orbar sind, kann die Manschette
z\"ugig entleert werden}

\item {\mdseries
CAVE: auskultatorische L\"ucke}

\end{itemize}
\begin{itemize}
\item {\mdseries
Protokollierung der Blutdruckwerte}

\end{itemize}
\begin{itemize}
\item {\mdseries
Immer sofort nach Messung protokollieren}

\end{itemize}
{\mdseries\color[rgb]{0.0,0.0,0.5019608}
Das EKG}

{\mdseries
Das Elektrokardiogramm (EKG) zeigt mittels eines entsprechenden
Ger\"ates die elektrische Herzaktivit\"at des Reizleitungssystems an.
Dazu werden auf den K\"orper nach einem bestimmten Muster Elektroden
aufgeklebt und mittels eines Kabels mit dem EKG-Ger\"at verbunden. }

{\mdseries
\gZ{EKG -- Ablauf einer Erregung}}

\begin{itemize}
\item \gZ{P-Welle \~{} Vorhoferregung}

\item \gZ{PQ-Zeit \~{} AV-\"Uberleitung}

\item \gZ{QRS-Komplex \~{} Kammererregung}

\item \gZ{ST-Strecke + T-Welle \~{}
Erregungsr\"uckbildung}

\end{itemize}
{\mdseries\color{black}
Das EKG gibt nur die elektrische Leitung, nicht aber die tats\"achliche
Muskelarbeit wieder! Es erlaubt \gEs{keine Aussage
\"uber die Herzleistung}! }

{\mdseries\color{black}
Grundsatz: \gEs{Selbst ein
}\gEs{toter}\gEs{ Patient kann
ein {\quotedblbase}sch\"ones{\textquotedblleft} EKG haben. }(allerdings
nicht lange)}

\subparagraph{EKG -- Ableitungen}
{\mdseries
Grunds\"atzlich unterscheidet man zwischen den Extremit\"atenableitungen
und den Brustwandableitungen. }

\begin{itemize}
\item {\mdseries
Extremit\"atenableitungen: I, II, III, aVL, aVR, aVF}

\item {\mdseries
Brustwandableitungen: V1 -- V6}

\end{itemize}
{\mdseries\color{black}
Extremit\"atenableitungen}

{\mdseries
Die Extremit\"atenableitungen werden mittels 4 Elektroden abgeleitet. Je
eine Elektrobe wird an den Enden der Extremit\"aten angebracht. Zur
Vereinfachung k\"onnen die Elektroden auch am \"Ubergang vom Rumpf zu
den Extremit\"aten angebracht werden (z.B. statt dem rechten Arm an der
rechten Schulter).}

{\mdseries
Aus den Informationen der 4 Elektroden erzeugt das EKG-Ger\"at
\gZ{6} Ableitungen:  \gZ{\textbf{I,
II, III, aVL, aVR, aVF}} }

  [Warning: Image ignored] % Unhandled or unsupported graphics:
%\includegraphics[width=12.304cm,height=8.975cm]{AASdev20090228-img4.png}
 

\begin{center}
\begin{tabular}{|m{4.367cm}|m{4.367cm}|m{4.367cm}|}
\hline
\centering {\bfseries\upshape\color[rgb]{0.0,0.0,0.5019608}
Elektrode}\par

 &
\centering {\bfseries\upshape\color[rgb]{0.0,0.0,0.5019608}
Position}\par

 &
\centering {\bfseries\upshape\color[rgb]{0.0,0.0,0.5019608} Alternative
Postion}\par

\\\hline
{\bfseries\color{red} Rot}

 &
{\mdseries Rechte Hand}

 &
{\mdseries Rechte Schulter}

\\\hline
{\bfseries\color[rgb]{1.0,0.70980394,0.08235294} Gelb}

 &
{\mdseries Linke Hand}

 &
{\mdseries Linke Schulter}

\\\hline
{\bfseries\color[rgb]{0.0,0.5019608,0.0} Gr\"un}

 &
{\mdseries Linker Fu{\ss}}

 &
{\mdseries Linke Leiste}

\\\hline
{\bfseries Schwarz}

 &
{\mdseries Rechter Fu{\ss}}

 &
{\mdseries Rechte Leiste}

\\\hline
\end{tabular}
\end{center}
{\mdseries\color{black}
Brustwandableitungen}

{\mdseries
Die Brustwandableitungen werden mittels 6 Elektroden abgeleitet und
erlauben eine genauere Zuordnung von St\"orungen zu der jeweiligen
Region des Herzens. Aus den Informationen der 6 Elektroden erzeugt das
EKG-Ger\"at 6 Ableitungen: \gZ{\textbf{V1 -- V6}}. }

\begin{center}
\begin{tabular}{|m{6.649cm}|m{6.651cm}}
\hline
\centering {\bfseries\upshape\color[rgb]{0.0,0.0,0.5019608}
Elektrode}\par

 &
\centering {\bfseries\upshape\color[rgb]{0.0,0.0,0.5019608}
Position}\par

\\\hline
{\bfseries\color{red} V1}

 &
{\mdseries 4. Zwischenrippenraum rechts vom Sternum}

\\\hline
{\bfseries\color[rgb]{1.0,0.70980394,0.08235294} V2}

 &
{\mdseries 4. Zwischenrippenraum links vom Sternum}

\\\hline
{\bfseries\color[rgb]{0.0,0.5019608,0.0} V3}

 &
{\mdseries Mitte zwischen V2 und V4}

\\\hline
{\bfseries V4}

 &
{\mdseries 5. Zwischenrippenraum links in der Verl\"angerung der
Schl\"usselbein-Mittellinie}

\\\hline
{\bfseries V5}

 &
{\mdseries Mitte zwischen V4 und V6}

\\\hline
{\bfseries V6}

 &
{\mdseries 6. Zwischenrippenraum in der mittleren Achsellinie}

\\\hline
\end{tabular}
\end{center}
\subsection{Traumafertigkeiten}
{\mdseries\itshape
Michael Motal}

{\mdseries\color[rgb]{0.0,0.0,0.5019608}
Helmabnahme}

{\selectlanguage{ngerman}
\textit{immer.}}

{\mdseries\color{black}
Material}

\begin{itemize}
\item {\mdseries
2 Helfer}

\item {\mdseries
1 stifneck f\"ur sp\"ater!}

\end{itemize}
{\mdseries\color{black}
Vorgehen}

\begin{itemize}
\item {\mdseries
1 Helfer fixiert den Helm von hinten und l\"asst nicht mehr los, bis der
Helm unten ist. Zwischen Helm und Schritt eine Helmh\"ohe Abstand
lassen, damit der Helm in einer Bewegung abgezogen werden kann}

\item {\mdseries
1 Helfer spricht den Patienten von vorne an, \"offnet Visier und
Kinnriemen}

\item {\mdseries
der seitliche Helfer stabilisiert den Kopf: eine Hand fasst mit Daumen
und Zeigefinger das Unterkiefer, die andere Hand st\"utzt den Kopf von
unten. W\"ahrend der Helm abgezogen wird rutscht diese Hand etwas nach}

\item {\mdseries
der hintere Helfer zieht den Helm in einer S-f\"ormigen Bewegung ab: den
Helm zuerst kippen bis die Nasenspitze sichtbar ist, dann die
Unterkante in der gegengleichen Bewegung zum Helfer bewegen}

\item {\mdseries
der hinere Helfer zieht den Helm sanft fertig ab und legt ihn sicher ab}

\item {\mdseries
der hintere Helfer \"ubernimmt den Kopf im doppelten C-Griff: Daumen und
Zeigefinger um ein Ohr, Kopf leicht auf Zug halten.}

\item {\mdseries
der seitliche Helfer legt ein stifneck an (siehe unten)}

\end{itemize}
{\mdseries\color[rgb]{0.0,0.0,0.5019608}
Stifneck}

{\mdseries
zur Stabilisierung der Halswirbels\"aule bei jedem Verdacht auf eine
Verletzung; auf jeden Fall bei jeder Helmabnahme!}

{\mdseries\color{black}
Material}

\begin{itemize}
\item {\mdseries
2 Helfer}

\item {\mdseries
1 stifneck}

\end{itemize}
{\mdseries\color{black}
Vorgehen}

\begin{itemize}
\item {\mdseries
Pat. auffordern, den Kopf nicht zu bewegen}

\item {\mdseries
nach Helmabnahme: Kopf weiterhin leicht auf Zug halten}

\item {\mdseries
der hintere Helfer h\"alt Kopf/HWS in Neutralstellung}

\item {\mdseries
der seitliche Helfer misst die Entfernug von Schulter bis Unterkiefer in
Fingerbreiten aus}

\item {\mdseries
stifneck entsprechend einstellen: gemessene Entfernung auftragen:
Hartplastikrand, der an der Schulter aufliegt unten, Markierung am
verschiebbaren Teil bzw H\"ohe der Kinnst\"utze oben)}

\item {\mdseries
Verschiebemechanismus des stifneck blockieren, ab nun darf kein
Verstellen mehr m\"oglich  sein }

\item {\mdseries
stifneck zuerst am Kinn anlegen/anpassen und festhalten}

\item {\mdseries
Klettverschlu{\ss}band so einfalten, dass es nicht am Boden scheifen
kann (h\"alt sonst nicht mehr)}

\item {\mdseries
Nackenteil unter Pat. durchf\"uhren (Vorsicht: Klettverschlussband)}

\item {\mdseries
fest schlie{\ss}en}

\end{itemize}
\begin{itemize}
\item[] Kinnspitze sollte ein kurzes St\"uck \"uber Kante hinausschauen
(kein Durchrutschen) aber nicht allzu weit (falsch eingestellt)

\end{itemize}
{\mdseries\color[rgb]{0.0,0.0,0.5019608}
Schaufeltrage}

{\mdseries
zusammengesetzt aus 2 L\"angsh\"alften; Fu{\ss}teil jeder H\"alfte
l\"angenverstellbar, Kopfteil fixe L\"ange.}

{\mdseries
Zur beinahe bewegungsfreien Rettung bei Verdacht auf Wirbels\"aulen-,
Becken- oder Oberschenkelverletzungen. Letztendlich soll der Patient
auf der Vakuummatratze gelagert werden. Sollte der Patient gr\"o{\ss}er
als die Trage sein m\"ussen Kopf und Oberk\"orper auf jeden Fall
trotzdem auf der Trage zu liegen kommen. Ausnahme: Beinverletzung, bei
der eine Wirbels\"aulenverletzung ausgeschlossen werden kann.}

{\mdseries\color{black}
Material}

\begin{itemize}
\item {\mdseries
2 Helfer}

\item {\mdseries
Schaufeltrage}

\item {\mdseries
Gurte f\"ur Schaufeltrage}

\end{itemize}
{\mdseries\color{black}
Vorgehen}

\begin{itemize}
\item {\mdseries
Schaufeltrage neben Patienten in zusammengebautem Zustand auf die
richtige L\"ange einstellen. Sollte ein Patient zu gro{\ss} sein
m\"ussen Kopf und Oberk\"orper auf jeden Fall trotzdem auf der Trage
liegen! Vorsicht: Markierung der max. L\"ange bei Schraubmodell}

\item {\mdseries
\"Offnen der Trage: Schnappverschl\"usse; je nach Methode nur den
unteren oder beide \"offnen.}

\item {\mdseries
Aufschaufeln, 1. Methode: Fu{\ss}seite offen, Pat. wird von oben kommend
wie auf eine Schere aufgeladen, die Trage unter ihm geschlossen.
Kontrolle des sicheren Schlu{\ss}es!}

\item {\mdseries
Aufschaufeln: 2. Methode: beide Teile getrennt, Pat. leicht an Schulter
und H\"ufte zur Seite gedreht und jew. Tragenteil daruntergeschoben.
H\"alften zuerst am Kopf-, dann am Fu{\ss}ende schliessen. Kontrolle
des sicheren Schlu{\ss}es!}

\item {\mdseries
Angurten: 2 Gurte; jeden Gurt zuerst um den Handgriff der einen, dann
der anderen H\"alfte f\"uhren, dann \"uber dem Patienten schlie{\ss}en
und festziehen.}

\item {\mdseries
Umlagern auf Vakuummatratze: Trage mit vorbeireiteter Matratze (s.u.)
neben Patienten stellen um diesen einen m\"oglichst geringen Weg
tragen/bewegen zu m\"ussen}

\item {\mdseries
Gurte \"offnen und entfernen.}

\item {\mdseries
H\"alften trennen: beide Schnappverschl\"usse \"offnen}

\item {\mdseries
H\"alften parallel zum Patienten flach wegziehen. Der Patient soll
hierbei nicht bewegt werden!}

\item {\mdseries
Sollten die Schnappschl\"osser klemmen kann man die Trage wie bei der
Scherentechnik an nur einer Seite \"offnen und in V-Form (je ein Helfer
auf einer Seite) entfernen.}

\end{itemize}
{\mdseries\color[rgb]{0.0,0.0,0.5019608}
Vakuummatratze:}

{\selectlanguage{ngerman}
zur Ruhigstellung des gesamten K\"orpers, bei Verdacht auf m\"ogliche
Wirbels\"aulenverletzungen, Becken- oder Oberschenkelfrakturen (lockere
Indikationsstellung!)}

{\selectlanguage{ngerman}
Am Ende soll die Matratze an die K\"orperform der Patienten eng
angepasst sein und somit ein Bewegen verhindern.}

{\mdseries\color{black}
Material}

\begin{itemize}
\item {\mdseries
2 Helfer}

\item {\mdseries
Fahrtrage}

\item {\mdseries
Vakuummatratze}

\item {\mdseries
Accuvac}

\item {\mdseries
Leintuch}

\item {\mdseries
ev. Adapter, je nach Matratzentyp}

\end{itemize}
{\mdseries\color{black}
Vorgehen}

\begin{itemize}
\item {\mdseries
Fahrtrage auf tiefste Stufe stellen (ganz unten)}

\item {\mdseries
Vorbereiten: Vakuummatratze auf Fahrtrage legen. Die weichere Seite ist
die Patientenseite, Ventil beim Kopf. Kugeln so verteilen, dass der
gesamte K\"orper gut sowie die betroffene Stelle besonders gut fixiert
werden kann.}

\item {\mdseries
Vorabsaugen mit Accuvac, bis gut formbar, aber nicht allzu fest.}

\item {\mdseries
Vorformen, Ventil verschlie{\ss}en. Wenn vorhanden Leintuch auf die
Matratze legen.}

\item {\mdseries
Angegurteten Patienten mit Schaufeltrage auf die Vakuummatratze legen,
Schaufeltrage entfernen (siehe
{\quotedblbase}Schaufeltrage{\textquotedblleft})}

\item {\mdseries
Absaugen mit Accuvac, dabei z\"ugig Matratze an den Patienten anpassen.
Besonderes Augenmerk gilt der Fixierung der verletzten Regionen sowie
dem Kopf ({\quotedblbase}Ohren{\textquotedblleft} formen: Ecken neben
dem Kopf hochformen).}

\item {\mdseries
Ist kein weiteres Absaugen mehr m\"oglich und die Matratze gut angepasst
(Matratze muss jetzt bretthart sein) Ventil schlie{\ss}en, Accuvac
abdrehen.}

\item {\mdseries
Pat. an Fahrtrage angurten.}

\end{itemize}
{\mdseries\color[rgb]{0.0,0.0,0.5019608}
Vakuum-Beinschiene}

{\selectlanguage{ngerman}
Schienung einer Fu{\ss}-, Knie- oder Unterschenkelverletzung. Achtung:
Oberschenkel: Vakuummatratze}

{\selectlanguage{ngerman}
1 Helfer seitlich, 1 Helfer beim Fu{\ss}ende}

{\mdseries\color{black}
Material}

\begin{itemize}
\item {\mdseries
2 Helfer}

\item {\mdseries
Vakuumbeinschiene}

\item {\mdseries
Accuvac}

\item {\mdseries
ev. Adapter}

\end{itemize}
{\mdseries\color{black}
Vorgehen}

\begin{itemize}
\item {\mdseries
L\"ange anpassen, gegebenenfalls oben umfalten.
Klettverschlu{\ss}b\"ander herrichten.}

\item {\mdseries
Kugeln gut verteilen. Auch am Fu{\ss}ende m\"ussen genug Kugeln sein, um
dem Fu{\ss} sowohl unter der Sohle als auch am Fu{\ss}r\"ucken gut Halt
geben zu k\"onnen}

\item {\mdseries
Schiene leicht vorabsaugen, damit Kugeln nicht mehr verrutschen
k\"onnen}

\item {\mdseries
Schuh am gesunden Bein ausziehen: Finden der schonensten Methode. Ein
Helfer h\"alt das Bein auf Zug, der andere \"offnet den Schuh (B\"ander
ausf\"adeln, zur Not aufschneiden)}

\item {\mdseries
vorsichtiges Wiederholen zu zweit am verletztem Bein, Socken ausziehen
(DMS)}

\item {\mdseries
Stiefelgriff: gleich nach Ausziehen von Schuh/Socken nimmt der Helfer am
Fu{\ss}ende durch das Loch in der Schiene die Ferse des verletzten
Fu{\ss}es, die andere Hand h\"alt den Fu{\ss}r\"ucken. Das Bein wird so
auf dem Boden liegend auf Zug gehalten. Der Helfer verbleibt so, bis
die Schiene fest sitzt!}

\item {\mdseries
Der seitliche Helfer richtet die Schiene vorsichtig unter dem Patienten
aus und fixiert die Schiene mit den Klettverschlussb\"andern. Um den
Fu{\ss} werden Ohren gebogen und kreuzweise festgeklettet: der Fu{\ss}
soll von allen Seiten so gut wie m\"oglich umschlossen sein!}

\item {\mdseries
Der seitliche Helfer saugt ab und zieht dabei die Klettb\"ander immer
wieder fest nach. Der andere Helfer h\"alt nach wie vor den Fu{\ss} im
Stiefelgriff.}

\item {\mdseries
Fertig abgesaugt: Ventil schliessen, Accuvac abschalten. Der
Stiefelgriff kann jetzt losgelassen werden.}

\end{itemize}
{\selectlanguage{ngerman}
Die Schiene soll bretthart sein, das Bein leicht auf Zug fixiert. Der
Fu{\ss} soll weder seitlich noch zum K\"orper hin beweglich sein.}

{\mdseries\color[rgb]{0.0,0.0,0.5019608}
Sam-Splint}

{\selectlanguage{ngerman}
Unterarmschienung}

{\mdseries\color{black}
Material}

\begin{itemize}
\item {\mdseries
2 Helfer}

\item {\mdseries
Sam-Splint}

\item {\mdseries
Peha-Haft, Mullbinde + Leukoplast oder 2-3 Dreiecktuchkravatten}

\item {\mdseries
ein weiteres Dreiecktuch}

\end{itemize}
{\mdseries\color{black}
Vorgehen}

\begin{itemize}
\item {\mdseries
Sam-Splint in der H\"alfte falten, an gesundem Arm der L\"ange nach
anpassen, Knick beim Ellbogen}

\item {\mdseries
Anlegen an verletztem Arm, sanft anpassen. Fixierung mit Peha-Haft,
Mullbinden oder Dreiecktuchkravatten}

\item {\mdseries
Fingerspitzen m\"ussen sichtbar sein (DMS), gegebenenfalls Enden
umfalten.}

\item {\mdseries
Lagerung des geschienten Armes im Armtragetuch}

\end{itemize}
{\mdseries\color[rgb]{0.0,0.0,0.5019608}
Druckverband}

{\selectlanguage{ngerman}
zur Stillung einer starken, d.h. potentiell lebensbedrohlichen Blutung}

{\mdseries\color{black}
Material}

\begin{itemize}
\item {\mdseries
1-2 Helfer}

\item {\mdseries
sterile Wundauflage je nach Wundgr\"o{\ss}e}

\item {\mdseries
saugf\"ahiger Druckk\"orper: zB Mullbinde}

\item {\mdseries
Dreiecktuchkravatte}

\end{itemize}
{\mdseries\color{black}
Vorgehen}

\begin{itemize}
\item {\mdseries
sofort bei Erkennen einer starken Blutung durch direkten Druck auf die
Wunde und Hochhalten die Blutung verlangsamen (Pat. kann dies je nach
Zustand selbst tun!)}

\item {\mdseries
sterile Wundauflage auf die Wunde legen}

\item {\mdseries
Druckk\"orper dem Wundverlauf folgend \"uber die sterile Wundauflage
legen, weiterhin dr\"ucken.}

\item {\mdseries
Druckk\"orper mit einer Dreiecktuchkravatte fest und m\"oglichst
vollst\"andig umschlie{\ss}en}

\item {\mdseries
Dreiecktuchkravatte festziehen und mit einem festen Knoten bei
ausreichendem Druck direkt \"uber der Wunde fixieren}

\item {\mdseries
2. Knoten zur Sicherung des ersten}

\end{itemize}
{\selectlanguage{ngerman}
sollte die Blutung zu einem sp\"ateren Zeitpunkt wieder st\"arker werden
oder der Druckverband zu locker angelegt sein einen weiteren
Druckk\"orper mit einer weiteren Dreiecktuchkravatte noch fester
fixieren. Einen schon angelegten Druckverband nie \"offnen!}

\subsection{Assistenzt\"atigkeiten}
\begin{center}
\tablehead{}\begin{supertabular}{|m{6.649cm}|m{6.651cm}|}
\hline
{\mdseries\itshape Teile des Textes und der Handlungsanleitungen
enstammen dem Skriptum {\quotedblbase}\"Arztliche
Grundfertigkeiten{\textquotedblleft} der  Besonderen Einrichtung zur
medizinischen Aus- und Weiterbildung -- Medizinische Universit\"at
Wien, Abt. Methodik und Entwicklung, welches in Zusammenarbeit  mit dem
Klinischem Institut f\"ur Hygiene und Medizinische Mikrobiologie
entstanden ist. }

{\mdseries\itshape Wir danken Hrn. Univ.Prof. Dr. Michael Schmidts und
seinem Team f\"ur die Unterst\"utzung!}

 &
{\bfseries\itshape Powered by}

{\centering   [Warning: Image ignored]
% Unhandled or unsupported graphics:
%\includegraphics[width=4.551cm,height=2.17cm]{AASdev20090228-img5.gif}
 \par}

\\\hline
\end{supertabular}
\end{center}
{\mdseries\color[rgb]{0.0,0.0,0.5019608}
Aufziehen eines Medikamentes aus einer Ampulle}

{\mdseries
Ampulle mit rotem Punkt: Sollbruchstelle, Daumen muss beim aufbrechen
auf dem roten Punkt sein. }

{\mdseries
Ampulle mit {\quotedblbase}Halsring{\textquotedblleft}: Ampulle kann
nach alle Richtungen aufgebrochen werden. }

\begin{center}
\begin{minipage}{6.556cm}
{\mdseries
\gH{\textit{\textcolor{blue}{Warum
mach{\textquotesingle} cih mir wegen der Verletzungsgefahr Sorgen? Ist
ja nur ein kleiner Schnitt, ist ja meine Sache {\dots}}}}}

\gH{Nein! Unter Normalbedingungen, wenn ich in aller
Ruhe mein Medi aufziehen kann ist es tats\"achlich mein eigenes Bier.
Was aber wenn ich mich bei einer Reanimation an der Ampulle schneide?
Die anderen Kollegen sind besch\"aftigt (Dr\"ucken, blasen, Defi
{\dots}) und dann steh ich da mit meinem kleinen Schnitt und
}\gH{blute alles an}\gH{. Und der
Notarzt braucht dringend das Adrenalin {\dots} Was jetzt?}

{\centering
\gH{\textcolor{red}{Inakzeptabel!}}
\par}

\end{minipage}
\end{center}
{\mdseries\color{black}
Material }

\begin{itemize}
\item {\mdseries
Spritze}

\item {\mdseries
Kan\"ule(n)}

\item {\mdseries
Ampulle mit Medikament}

\item {\mdseries
Trockene Tupfer}

\end{itemize}
{\mdseries\color{black}
Aufziehen des Medikamentes}

\begin{itemize}
\item {\mdseries
Kontrolle Ablaufdatum}

\item {\mdseries
Kontrolle Inhalt: Verf\"arbung, Ausflockung?}

\item {\mdseries
\"Offnen der Ampulle}

\end{itemize}
\begin{itemize}
\item {\mdseries
 \textbf{Glasampulle}}

\end{itemize}
\begin{itemize}
\item {\mdseries
Ampullenspitze beklopfen\foreignlanguage{ngerman}{, darauf achten dass
sich keine Fl\"ussigkeit im Apullenkopf mehr befindet}}

\item {\mdseries
\"Offnen der Medikamenten-Ampulle mit Tupfer, roter Punkt weist zum
K\"orper}

\end{itemize}
[Warning: Draw object ignored]

\begin{center}
 [Warning: Image ignored] % Unhandled or unsupported graphics:
%\includegraphics[width=4.764cm,height=3.145cm]{AASdev20090228-img6.png}

\end{center}
\begin{itemize}
\item {\mdseries
\textbf{Plastikampulle},
{\quotedblbase}Dreh-und-Drink{\textquotedblleft}-Verschlu{\ss}}

\end{itemize}
\begin{itemize}
\item {\mdseries
Den Verschlussknopf ohne Kontamination des Ampullenhalses
(Sollbruchstelle) abdrehen}

\end{itemize}
\begin{itemize}
\item {\mdseries
Ampulle abstellen}

\end{itemize}
\begin{itemize}
\item {\mdseries
Verpackung der Spritze aufsch\"alen, Spritze in Verpackung lassen}

\item {\mdseries
Verpackung der Nadel aufsch\"alen, Spritze auf Konus aufstecken}

\item {\mdseries
Entfernen der Kan\"uleabdeckung (nicht abdrehen, sondern abziehen!)}

\item {\mdseries
Einf\"uhren der Nadel in die Ampulle}

\end{itemize}
\begin{itemize}
\item {\mdseries
Beim Einf\"uhren der Kan\"ule in die Ampulle den Ampullenrand nicht
ber\"uhren}

\end{itemize}
\begin{itemize}
\item {\mdseries
Aufziehen des Medikamentes}

\item {\mdseries
Vorsichtig entl\"uften}

\item {\mdseries
Variante {\quotedblbase}Verabreichung \"uber Venflon{\textquotedblleft}
}

\end{itemize}
\begin{itemize}
\item {\mdseries
Nadel h\"andisch abziehen und in den Kontamed verwerfen. }

\item {\mdseries
Arzt das Medikament \"ubergeben, Ampulle zeigen.}

\item {\mdseries
Wird das Medikamen nicht sofort verwendet spritze mit roten Stoppel
verschliessen und Ampulle ankleben.}

\end{itemize}
\begin{itemize}
\item {\mdseries
Variante {\quotedblbase}Zuspritzen in eine
Infusionsflasche{\textquotedblleft}}

\end{itemize}
\begin{itemize}
\item {\mdseries
Gummimembran desinfizieren}

\item {\mdseries
Kan\"ule nicht bis zum Schaft einstechen}

\item {\mdseries
Medikament zuspritzen}

\item {\mdseries
Beschriften der Flasche}

\end{itemize}
\begin{itemize}
\item {\mdseries
Variante {\quotedblbase}direktes Spritzen{\textquotedblleft}}

\end{itemize}
\begin{itemize}
\item {\mdseries
Arzt nach Nadelgr\"o{\ss}e fragen}

\item {\mdseries
Aufziehnadel h\"andisch entfernen}

\item {\mdseries
Spritznadel aufstecken}

\item {\mdseries
Arzt Spritze mit Verschlusskappe reichen, Ampulle zeigen.}

\end{itemize}
{\mdseries\color[rgb]{0.0,0.0,0.5019608}
Vorbereitung einer Infusion}

{\mdseries\color{black}
Material}

\begin{itemize}
\item \item {\mdseries
Infusionsflasche}

\item {\mdseries
Infusionsbesteck}

\item {\mdseries
Aufh\"angevorrichtung}

\item {\mdseries
alkoholhaltiges Desinfektionsmittel}

\item {\mdseries
einige saubere Tupfer}

\item {\mdseries
Permanentmarker}

\end{itemize}
{\mdseries\color{black}
Handlungsschritte}

\begin{itemize}
\item {\mdseries
H\"andedesinfektion         }

\item {\mdseries
Aufh\"angevorrichtung auf Flasche aufsetzen}

\item {\mdseries
Er\"offnen der Flaschenabdeckung (Membranschutz)}

\item {\mdseries
Wischdesinfektion der Gummimembran}

\end{itemize}
\begin{itemize}
\item {\mdseries
Alkoholgetr\"ankter Tupfer kann w\"ahrend weiterer Vorbereitung auf
Membran belassen werden}

\end{itemize}
\begin{itemize}
\item {\mdseries
Montage des Infusionsbestecks}

\end{itemize}
\begin{itemize}
\item {\mdseries
\gEs{Radklemme schlie{\ss}en} }

\item {\mdseries
Einstich des Tropfkammerdorns in Infusionsflasche mit geschlossenem
L\"uftungsventil und geschlossener Radklemme}

\item {\mdseries
{\quotedblbase}No touch technique{\textquotedblleft} -- Dorn der
Tropfkammer und Gummimembran der Infusionsflasche d\"urfen nicht
ber\"uhrt werden}

\end{itemize}
\begin{itemize}
\item {\mdseries
Infusionsschlauch f\"ullen}

\end{itemize}
\begin{itemize}
\item {\mdseries
Flasche wenden und hochhalten}

\item {\mdseries
Tropfkammerluftventil \"offnen}

\item {\mdseries
Tropfkammer durch Zusammendr\"ucken bis zur H\"alfte f\"ullen}

\item {\mdseries
Konusabdeckung des Infusionsschlauchs am anderen Ende muss bei der
Bel\"uftung nicht entfernt werden}

\item {\mdseries
Infusionsflasche und Ende des Infusionsschlauchs mit nicht dominanter
Hand halten \newline
\gZ{(Wenn sich das Infusionsschlauchende auf gleicher
H\"ohe wie der Spiegel in der Tropfkammer befindet, entl\"uftet sich
der Schlauch problemlos -- \newline
Prinzip der kommunizierenden Gef\"a{\ss}e, an kommunizierenden
Gef\"a{\ss}en stellt sich immer dieselbe Fl\"ussigkeitsspiegelh\"ohe
ein.)}}

\item {\mdseries
Radklemme \"offnen}

\item {\mdseries
Infusionsschlauch {\quotedblbase}ohne vergie{\ss}en{\textquotedblleft}
entl\"uften %
%
\gZ{(Fl\"ussigkeit bis ca. 5cm vor Konnektor rinnen
lassen, endg\"ultig entl\"uften bei Patienten vor Verabreichung)}}

\item {\mdseries
Radklemme schlie{\ss}en}

\end{itemize}
{\mdseries\color[rgb]{0.0,0.0,0.5019608}
Venenverweilkan\"ule ({\quotedblbase}Venflon{\textquotedblleft}):
Vorbereitung und Assistenz }

  [Warning: Image ignored] % Unhandled or unsupported graphics:
%\includegraphics[width=13.467cm,height=4.658cm]{AASdev20090228-img7.png}
 

{\mdseries\upshape
Bild: BEMAW}

{\color{black}
Farbcodierung der peripheren Venenverweilkan\"ulen}

\begin{flushleft}
\begin{tabular}{|m{1.5619999cm}|m{1.7229999cm}|m{1.7229999cm}|m{2.0509999cm}|m{1.7839999cm}|m{2.102cm}|}
\hline
\multicolumn{3}{|m{5.408cm}|}{\centering {\bfseries Nur zur
Illustration!}\par

} &
\centering Durchfluss [ml/min]\par

 &
\centering Farbcode\par

 &
\\\hline
\centering \gZ{{\O} mm}\par

 &
\centering \gZ{G}\par

 &
\centering \gZ{Lmm}\par

 &
 &
 &
\\\hline
\centering \gZ{0,5 x 0,8}\par

 &
\centering \gZ{22}\par

 &
\centering \gZ{25}\par

 &
\centering \gZ{25}\par

 &
\centering blau\par

 &
\centering Sehr d\"unn\par

\\\hline
\centering \gZ{0,7 x 1,0}\par

 &
\centering \gZ{20}\par

 &
\centering \gZ{32}\par

 &
\centering \gZ{55}\par

 &
\centering rosa\par

 &
\centering d\"unn\par

\\\hline
\centering \gZ{0,9 x 1,2}\par

 &
\centering \gZ{18}\par

 &
\centering \gZ{40}\par

 &
\centering \gZ{90}\par

 &
\centering gr\"un\par

 &
\centering normal\par

\\\hline
\centering \gZ{1,1 x 1,4}\par

 &
\centering \gZ{17}\par

 &
\centering \gZ{42}\par

 &
\centering \gZ{135}\par

 &
\centering wei{\ss}\par

 &
\centering dick\par

\\\hline
\centering \gZ{1,3 x 1,7}\par

 &
\centering \gZ{16}\par

 &
\centering \gZ{42}\par

 &
\centering \gZ{170}\par

 &
\centering grau\par

 &
\centering Sehr dick\par

\\\hline
\centering \gZ{1,6 x 2,1}\par

 &
\centering \gZ{14}\par

 &
\centering \gZ{42}\par

 &
\centering \gZ{265}\par

 &
\centering orange\par

 &
\centering Pferd\par

\\\hline
\end{tabular}
\end{flushleft}
{\mdseries\upshape
Tabelle: BEMAW (modifiziert)}

{\mdseries\color{black}
Material}

\begin{itemize}
\item {\mdseries
Nierentasse }

\item {\mdseries
Staubinde}

\item {\mdseries
Venflon: }

\end{itemize}
\begin{itemize}
\item {\mdseries
Die Gr\"o{\ss}e (= Farbe) ist vom Arzt vorher zu erfragen!)}

\item {\mdseries
Die Fl\"ugel werden hinunter geklappt}

\item {\mdseries
Die farbige Lasche des Zuspritzventils in rechten Winkel zur
Venflonachse stellen.}

\end{itemize}
\begin{itemize}
\item {\mdseries
Einmalverschlusskappe ({\quotedblbase}roter Stoppel{\textquotedblleft})}

\item {\mdseries
Fixierungspflaster}

\item {\mdseries
einige saubere trockene Tupfer}

\item {\mdseries
Alko-Tupfer}

\item {\mdseries
(Einwegspritze mit 5ml Sp\"ulfl\"ussigkeit)}

\item {\mdseries
(Faserschreiber)}

\item {\mdseries
durchstichsicherer Entsorgungsbeh\"alter
({\quotedblbase}sharp{\textquotedblleft})}

\end{itemize}
{\mdseries\color{black}
Assistenz}

\begin{itemize}
\item {\mdseries
Stauschlauch zureichen}

\item {\mdseries
ge\"offnete Verpackung des Alko-Tupfer hinreichen sodass dieser
entnommen werden kann}

\item {\mdseries
Kontamed in Reichweite des  Arztes / NKV bereitstellen sodass er
sp\"ater die Nadel abwerfen kann.}

\item {\mdseries
Venflon zureichen}

\end{itemize}
\begin{itemize}
\item {\mdseries
Venflon an der Schutzkappe halten sodass der Arzt / NKV den Venflon
abziehen kann.}

\end{itemize}
\begin{itemize}
\item {\mdseries
Trockene Tupfer zureichen.}

\item {\mdseries
Je nach Situation Infusionsschlauch, roten Stoppel oder gar nichts
zureichen.  }

\item {\mdseries
Fixierpflaster zureichen}

\end{itemize}
{\mdseries\color[rgb]{0.0,0.0,0.5019608}
Intubation: Vorbereitung und Assistenz }

\subsection[O2 \ - Sauerstoff]{O\gSub{2}  {}- Sauerstoff}
{\mdseries
\gT{Sauerstoff} (chem. Symbol
\gT{O}\gT{\gSub{2}}) ist ein
lebenswichtiges Gas welches in unserer Atemluft vorkommt und
au{\ss}erdem ein essentielles
Notfall{\textquotedblright}medikament{\textquotedblleft}, die
fr\"uhzeitige Verabreichung ist wichtig f\"ur die Prognose des
Patienten!}

{\mdseries
Sauerstoff darf und muss bei Bedarf vom RS verabreicht werden.}

\begin{center}
 [Warning: Image ignored] % Unhandled or unsupported graphics:
%\includegraphics[width=4.522cm,height=5.628cm]{AASdev20090228-img8.png}
\captionof{figure}[: NEIN!]{: \textbf{NEIN!}}

\end{center}
\paragraph{Sauerstoff -- technische Grundlagen}
{\mdseries
Sauerstoff wird in Druckflaschen  gelagert. Laut EU-Norm sind diese
Flaschen wei{\ss} und mit einem
{\quotedblbase}\gTt{N}{\textquotedblleft} gekennzeichnet.
Das {\quotedblbase}\gTt{N}{\textquotedblleft} steht f\"ur
{\quotedblbase}neue Kennzeichnung{\textquotedblleft}, da fr\"uher die
Flaschen blau gekennzeichnet wurden. }

{\mdseries
Wichtige Begriffe}

\begin{itemize}
\item {\mdseries
F\"ullgr\"o{\ss}e }

\item {\mdseries
F\"ullvolumen }

\item {\mdseries
Durchflu{\ss} }

\end{itemize}
{\mdseries
Die Druckflaschen haben unterschiedliche F\"ullgr\"o{\ss}en (Volumina),
g\"angige Flaschengr\"o{\ss}en sind 0,5~l, 1~l, 1,5~l, 2~l, 5~l, 10~l.
Der Sauerstoff wird unter Druck gelagert, durch diesen Trick kann die
Menge des zur Verf\"ugung stehenden Sauerstoffs um ein Vielfaches
gesteigert werden. Der Druck einer gef\"ullten Flasche betr\"agt
normalerweise ca. 200~bar. }

{\mdseries
Jede Druckflasche verf\"ugt \"uber ein \gT{Hauptventil}
und ein Standardgewinde. An diesem Gewinde ist enteder ein
Druckschlauch oder ein \gE{Druckminderer} angeschlossen.
Der  \gT{Druckminderer} hat ein
\gE{Abgabeventil} welches den Druck drosselt und eine
Abgabe an den Patienten erm\"oglicht. Je nach Bauart kann der
Druckminderer \"uber einen \gE{Regler}, ein
\gE{Flowmeter} und ein \gE{Barometer}
verf\"ugen. Mit dem Regler kann man die Durchlu{\ss}- bzw. Abgaberate
ein\-stellen, das \gT{Flowmeter} (Durchflu{\ss}anzeige)
zeigt die aktuelle Durchflu{\ss}rate in Liter pro Minute (~l~/~min~)an.
Das \gT{Barometer} zeugt den Druck in der
Sauerstofflasche an wenn das Haupt\-ventil ge\"offnet ist. }

\begin{center}
 [Warning: Image ignored] % Unhandled or unsupported graphics:
%\includegraphics[width=7.189cm,height=4.94cm]{AASdev20090228-img9.png}

\end{center}
{\mdseries
Das effektive F\"ullvolumen erh\"alt man wenn man die F\"ullgr\"o{\ss}e
(in~l) mit dem Flaschendruck (in bar) multipliziert.  }

{\mdseries
Wenn man wissen m\"ochte wie lange die vorhandene Sauerstoffmenge bei
einer bestimmten Abgabemenge reicht dividiert man einfach die
vorhandene Menge durch die Abgaberate. Achtung: Es ist eine
Ausreichende Reserve einzuplanen, und es ist zu Bedenken dass in der
Flasche ein Restdruck verbleiben muss! }

{\mdseries\color{black}
Mit Sauerstoff l\"asst sich gut rechnen ...}

\[\begin{gathered}\boldsubformula{\text{Enthaltene
}}\boldsubformula{O}_{2}\boldsubformula{\text{{}-Menge }}=\text{
Volumen der Flasche }\times \text{
Druck}\\\boldsubformula{\text{Verf\"ugbare Zeit }}=\frac{\text{
Enthaltene Menge }}{\text{ Durchflu{\ss} }}\end{gathered}\]
Beispiel:

Flasche 10l mit 150bar: 1500l  O\gSub{2}

Patient ben\"otigt 5l O\gSub{2} pro Minute: 

Die F\"ullung reicht f\"ur 300 Minuten, abz\"uglich eines
Sicherheitsvolumens.

{\mdseries\color{black}
Die Sauerstoffflaschen sind pfleglich zu behandeln}

Sauerstoff ist zusammen mit Fett selbstentz\"undlich! Die Armaturen und
Gewinde sind fettfrei zu halten! Das Hantieren mit Feuer und offenem
Licht ist in der N\"ahe von Sauerstoffflaschen verboten! 

Der gro{\ss}e Druck birgt weitere Gefahren: Sollte ein Ventil abbrechen
kann die Flasche eine enorme Kraft entwickeln und zu einer Rakete
werden. Die Kraft einer mit 200~bar gef\"ullten Flasche reicht aus um
Mauern zu druchbrechen!  

{\mdseries\color{black}
\gCa{Achtung!} Vorsicht beim Hantieren -  Explosionsgefahr!}

{\mdseries\color{black}
Flaschen stehen unter Druck, ein Abbrechen der Ventile ist
katastrophal!}

{\mdseries\color{black}
Sauerstoff ist zusammen mit Fett
\gEs{selbstentz\"undlich} ${\rightarrow}$ Ventile
nicht mit fetten/\"oligen H\"anden bedienen.}

\paragraph{Verabreichungsarten}
\subparagraph{Sauerstoffbrille}
{\mdseries
Nur bei vorhandener Spontanatmung!}

{\mdseries
zwei kurze, in die Nasen\"offnungen  eingef\"uhrte
Kunst\-stoff\-stutzen}

{\mdseries
bei ca. 4 l O\gSub{2} /min bis 40\% O\gSub{2}  in
Atemgas  m\"oglich}

\begin{center}
\tablehead{}\begin{supertabular}{|m{6.649cm}|m{6.651cm}|}
\hline
\centering {\bfseries\upshape\color[rgb]{0.0,0.0,0.5019608}
Vorteile}\par

 &
\centering {\bfseries\upshape\color[rgb]{0.0,0.0,0.5019608}
Nachteile}\par

\\\hline
{\mdseries gute Toleranz}

{\mdseries fehlendes Engegef\"uhl}

{\mdseries Sprechen\&Husten m\"oglich}

 &
{\mdseries ungenaue Dosierung}

{\mdseries Austrocknung  der Schleimh\"aute}

\\\hline
\end{supertabular}
\end{center}
Bei behinderter Nasenatmung kann die  Sauerstoffbrille im Mundwinkel
plaziert werden. 

\begin{center}
\tablehead{}\begin{supertabular}{|m{3.225cm}|m{3.225cm}|m{3.225cm}|m{3.225cm}|}
\hline
\centering {\bfseries\upshape\color[rgb]{0.0,0.0,0.5019608} Was}\par

 &
\centering {\bfseries\upshape\color[rgb]{0.0,0.0,0.5019608} Wann}\par

 &
\centering {\bfseries\upshape\color[rgb]{0.0,0.0,0.5019608}
Vorteile}\par

 &
\centering {\bfseries\upshape\color[rgb]{0.0,0.0,0.5019608}
Nachteile}\par

\\\hline
\centering {\bfseries\upshape\color[rgb]{0.0,0.0,0.5019608}
Sauerstoffbrille}\par



\begin{center}
 [Warning: Image ignored] % Unhandled or unsupported graphics:
%\includegraphics[width=2.74cm,height=2.74cm]{AASdev20090228-img10.png}

\end{center}
 &
{\mdseries Spontanatmung}

{\mdseries O2-Flow {\textless} 5l/min}

 &
{\mdseries gute Toleranz}

{\mdseries fehlendes Engegef\"uhl}

{\mdseries Sprechen\&Husten m\"oglich}

 &
{\mdseries ungenaue Dosierung}

{\mdseries Austrocknung  der Schleimh\"aute}

\\\hline
\centering {\bfseries\upshape\color[rgb]{0.0,0.0,0.5019608}
Inhalationsmaske}\par



\begin{center}
 [Warning: Image ignored] % Unhandled or unsupported graphics:
%\includegraphics[width=2.437cm,height=2.707cm]{AASdev20090228-img11.png}

\end{center}
 &
{\mdseries Spontanatmung}

{\mdseries O2-Flow ab 5l/min}

 &
{\mdseries h\"ohere O\gSub{2}{}-Konzentration}

 &
{\mdseries h\"ohere O\gSub{2}{}-Konzentration}

{\mdseries CO\gSub{2}{}-R\"uckatmung  wenn
O\gSub{2}{}-Durchflu{\ss} {\textless} 5 l/min!)}

\\\hline
  [Warning: Image ignored] % Unhandled or unsupported graphics:
%\includegraphics[width=2.652cm,height=2.476cm]{AASdev20090228-img12.png}
 

 &
 &
{\mdseries Noch h\"ohere O\gSub{2}{}-Konzentration}

 &
{\mdseries Wie oben.}

\\\hline
\centering {\bfseries\upshape\color[rgb]{0.0,0.0,0.5019608}
Beatmungsbeutel mit  Reservoir}\par

 &
{\mdseries Beatmung des Patienten (bei Reanimation oder im Ausnahmefall
bei Atemstillstand oder nicht ausreichender Eigenatmung)}

{\mdseries Beatmung eines intubierten Patienten}

 &
 &
\\\hline
\end{supertabular}
\end{center}
\paragraph{Beatmung mit Beatmungsbeutel}
{\mdseries
JEDE BEATMUNG SOLLTE MIT RESERVOIRBEUTEL UND 10-15 l O2/min
DURCHGEF\"UHRT WERDEN!}

{\mdseries
Beatmungsfrequenz beim Erwachsenen: 12-15 x / min (Eigenfrequenz)}

{\mdseries\color{black}
\gCa{Achtung!} Zuviel ist schlecht! Nicht zu schnell beatmen
${\rightarrow}$ Hyperventilationsgefahr!}

\paragraph{Sauerstofftherapie}
{\mdseries
Prinzipiell sollte sich die gew\"ahlte Therapie an der Diagnose und am
Schweregrad orientieren!}

{\mdseries
Bedenke: }

\begin{itemize}
\item {\mdseries
Diagnose!}

\item {\mdseries
Schockbek\"ampfung}

\item {\mdseries
Isch\"amie (bei Akutem Koronoarsyndrom, Insult, ..)}

\item {\mdseries
Atemst\"orungen (Fremdk\"orper, Atemnot, {\dots})}

\end{itemize}
\paragraph[Pulsoxymetrie]{Pulsoxymetrie%
%
}
\index{Pulsoxymetrie}\begin{center}
 [Warning: Image ignored] % Unhandled or unsupported graphics:
%\includegraphics[width=3.516cm,height=3.516cm]{AASdev20090228-img13.png}

\end{center}
{\mdseries
Das Pulsoxymeter (kurz: Pulsoxy) mi{\ss}t dden Puls und die
Sauerstoffs\"attigung
(\index{SpO2}\gT{SpO}\gT{\gSub{2}})des
Blutes (Verh\"altnis O2-beladene/ nichtbeladene Erythrozyten). }

{\mdseries
Normalwert 96-99\%, }

{\mdseries
\textbf{\textcolor[rgb]{0.7882353,0.0,0.08627451}{[F078?] }}Achtung: Ein
{\quotedblbase}guter{\textquotedblleft}
\index{SpO2}\gT{SpO}\gT{\gSub{2}}{}-Wert
bedeutet \gEs{nicht} {\quotedblbase}keinen
Sauerstoff geben{\textquotedblleft}. Die Entscheidung bez\"uglich
Sauerstoffgabe ber\"ucksichtigt die Diagnose und den Zustand des
Patienten, der
\index{SpO2}\gT{SpO}\gT{\gSub{2}}{}-Wert
hat wenig bis gar nichts zu sagen! }

{\mdseries\color{black}
Achtung: Das Pulsoxymeter verwechselt O2 und CO-Molek\"ule, bei einer
CO-Vergiftung zeigt das Pulsoxy falsch-hohe S\"attigungswerte an!}

{\mdseries
Weitere Fehlerquellen: kein Signal: Nagellack, Zentralisation}

{\mdseries
Das Pulsoxymeter misst die Herzfrequenz und die Sauerstoffs\"attigung
(=\gT{SpO2}) des Blutes, aber es gilt:}

{\mdseries\color{black}
\textbf{\textcolor[rgb]{0.7882353,0.0,0.08627451}{[F078?]}} Ein
{\quotedblbase}guter{\textquotedblleft} S\"attigungswert ist keine
Begr\"undung keinen Sauerstoff zu verabreichen!}

\section{[RS XIII] Untersuchung und Anamnese}
Hier lernen Sie wie Sie die bereits gelernten therotischen Grundlagen in
der \gEs{Praxis} einsetzen, wie und in welcher
\gEs{Reihenfolge} Sie im Einsatz vorgehen, wie Sie
\gEs{Dringlichkeiten} einsch\"atzen und
\gEs{Verdachtsdiagnosen} erarbeiten k\"onnen und wie
Sie dabei dann auch noch den \gEs{\"Uberblick}
behalten k\"onnen.

\subsection{{\quotedblbase}Hilfe, ein Patient!{\textquotedblleft} - Was
mach{\textquotesingle} ich jetzt?}
Blabla

\"Ubersicht 

Ersteinsch\"atzung ${\rightarrow}$  Erstma{\ss}nahmen ${\rightarrow}$ 
Plan zurechtlegen wie es weiter gehen soll

{\mdseries
blabla}

{\mdseries
mehr blabla}

{\mdseries\color[rgb]{0.0,0.0,0.5019608}
Die \gEs{Anamnese} erz\"ahlt uns die Geschichte des
Patienten}

\paragraph{S{\textquotesingle}KAMEL denkt an alles}
\subparagraph{Symptome und Schmerzen}
\subparagraph{Krankheiten}
\subparagraph{Allergien}
\subparagraph{Medikamente}
\subparagraph{Ereignisse}
\subparagraph{Das Letzte ...}
{\mdseries\color[rgb]{0.0,0.0,0.5019608}
\gEs{Untersuchungen} k\"onnen uns weiterhelfen}

\subparagraph{Sehen -- H\"oren -- F\"uhlen }
\subparagraph{Vitalwerte}
\subparagraph{Atmung}
\subparagraph{Pulsoxymetrie}
\subparagraph{Blutzuckermessung}
\subparagraph{K\"orpertemperatur}
\subparagraph{Neurocheck}
\subparagraph{Traumacheck}
{\selectlanguage{ngerman}\itshape
Suche nach Verletzungen oder Verletzungsanzeichen}

\begin{itemize}
\item {\selectlanguage{ngerman}
Patient orientiert? Unfallhergang erinnerlich? Schmerzen?}

\end{itemize}
\begin{itemize}
\item {\selectlanguage{ngerman}
Abtasten mit festem Griff: von Kopf bis Fu{\ss} (Sch\"adeldach,
Hinterkopf, Stirn, Wangenknochen, Nasenbein, Kiefer, Hals, Schultern,
Arme, H\"ande, Brustkorb vorne und seitlich (Prellmarken?
Atemmechanik?), Bauch in 4 Quadranten (Prellmarken?
Abwehrspannung/harte Bauchdecke?), Schambein, Becken (stabil?), Beine
(sichtbare Verletzungen?), F\"u{\ss}e: Schuhe und Socken ausziehen)}

\item {\selectlanguage{ngerman}
{\quotedblbase}DMS{\textquotedblleft}: Finger und Zehen: Kontrolle von
\textbf{Motorik} (bewegen lassen), \textbf{Sensibilit\"at}
({\quotedblbase}Sp\"uren Sie das?{\textquotedblleft}) und
\textbf{Durchblutung}: Nagelbettprobe/Rekapillarisierungszeit: leichter
Druck auf Nagelende bis Nagelbett wei{\ss} ist, loslassen. Zeit bis
Nagelbett wieder rosarot ist sollte nicht l\"anger als 1-2 Sekunden
sein. Seitenvergleich! Eine verl\"angerte Rekapzeit deutet auf eine
Durchblutungsst\"orung hin.}

\item {\selectlanguage{ngerman}
Kontrolle von Mund, Nase und Ohren auf Blutungen (mit einer Lampe).
Tupfertest: etwas Blut aus Ohren mit einem Tupfer auffangen. Sollte
Liquor im Blut enthalten sein bildet sich um den roten Blutfleck ein
heller, bernsteinfarbener Hof.}

\end{itemize}
{\selectlanguage{ngerman}
bei Sch\"adelverletzungen ist ein Neurocheck auf jeden Fall
erforderlich!}

{\bfseries\color[rgb]{0.5019608,0.0,0.0}
Tarumavortr\"age m\"ussen hierher verweisen}

\subparagraph{Abtasten des Abomens }
\clearpage\setcounter{page}{1}{\color[rgb]{0.0,0.0,0.5019608}
Vorgehen in der Praxis}

\paragraph{Szene\"uberblick}
\paragraph[Die ersten \ \ 30 Sekunden ]{Die ersten  30 Sekunden }
{\mdseries
Ersteinsch\"atzung ${\rightarrow}$  Erstma{\ss}nahmen ${\rightarrow}$ 
Plan zurechtlegen wie es weiter gehen soll}

{\mdseries\color{black}
Erster Schritt: Ersteinsch\"atzung}

\begin{center}
\tablehead{}\begin{supertabular}{|m{3.158cm}m{10.142cm}|}
\hline
\centering {\bfseries\upshape\color[rgb]{0.0,0.0,0.5019608}
Ersteindruck}\par

 &
{\mdseries Aussehen, Gesichtsausdruck, Haltung, Sprache}

{\mdseries Bl\"asse, Zyanose, Schwei{\ss}, Schohnhaltung, Atemnot,
verwaschene Sprache, Des\-orientiertheit, Kr\"ampfe }

\\\hline
\centering {\bfseries\upshape\color[rgb]{0.0,0.0,0.5019608}
Bewusstsein}\par

 &
{\mdseries Reaktion auf verbale Ansprache, Weckversuch, Schmerzreiz}

{\mdseries klar, somnolent, sopor\"os, komat\"os/bewusstlos}

\\\hline
\centering {\bfseries\upshape\color[rgb]{0.0,0.0,0.5019608} Atemweg}\par

 &
{\mdseries Atemger\"ausch, Gestik}

{\mdseries frei oder verlegt}

\\\hline
\centering {\bfseries\upshape\color[rgb]{0.0,0.0,0.5019608} Atmung}\par

 &
{\mdseries Frequenz, Atemmuster, Anstrengung beim Atmen,
Atemger\"ausche, Heimsauerstoff}

\\\hline
\centering {\bfseries\upshape\color[rgb]{0.0,0.0,0.5019608}
Kreislauf}\par

 &
{\mdseries Bewusstsein, Hautfarbe, Radialispuls, Rekap-Zeit}

\\\hline
\centering {\bfseries\upshape\color[rgb]{0.0,0.0,0.5019608} Traumacheck
1}\par

 &
{\mdseries Schwere Verletzungen: Backen-, Oberschenkel-,
Serienrippenfraktur, massive Blutung}

\\\hline
\centering {\bfseries\upshape\color[rgb]{0.0,0.0,0.5019608}
Hauptbeschwerde}\par

 &
{\mdseries Grund der Berufung, Schmerzen und andere f\"ur den Patienten
dringliche Symptome}

\\\hline
\end{supertabular}
\end{center}
{\mdseries
Nun stellen sich folgende Fragen:}

\begin{itemize}
\item {\mdseries
Was ist das \gE{Problem} des Patienten?}

\begin{center}
\begin{minipage}{3.863cm}
{\centering\bfseries\color{black}
Alarmzeichen
\par}

{\centering\bfseries\color{black}
Massive St\"orungen von 
\par}

Bewusstsein

Atemweg / Atmung 

Kreislauf 

{\centering\bfseries\color{black}
typische Symptome 
\par}

Atemnot die sich nicht bessert 

Brodelndes Atemger\"ausch 

Thoraxschmerz  

Schocksymptome 

entgleiste Vitalwerte

schwere Verletzungen

{\dots} 

\end{minipage}
\end{center}
\item {\mdseries
Ist der Patient \gE{akut vital bedroht}?}

\end{itemize}
\begin{itemize}
\item {\mdseries
\gE{Anzeichen} einer vitalen Bedrohung k\"onnen u.a.
sein: St\"orungen von Be\-wusstsein, Atemweg, Atmung oder Kreislauf
sowie typische Symptome wie sich nicht bessernde
\gE{Atemnot}, \gE{brodelndes
Atem\-ge\-r\"ausch}, \gE{Thoraxschmerz},
\gE{Schocksymptome} und \gE{entgleiste
Vital\-werte}.}

\end{itemize}
{\mdseries
W\"ahrend des weiteren Verlaufs (Anamnesegespr\"ach, etc.) kann man auf
\gE{weitere Befunde} sto{\ss}en welche auf eine vitale
Be\-drohung hinweisen (zB. massiv erh\"ohter Blutdruck bei bei
vor\-ge\-sch\"adigtem Herzen, pl\"otzliche oder schleichende
Zustands\-verschlechterung, etc.). }

{\mdseries
Die Ersteinsch\"atzung muss st\"andig \"uberpr\"uft werden. Neue Befunde
und der Zustand des Patienten k\"onnen die Situation \"andern.}

{\mdseries\color{black}
Zweiter Schritt: Erstma{\ss}nahmen}

{\mdseries
Welche Ma{\ss}nahmen sind f\"ur den Patienten am wichtigsten und
m\"ussen zuerst erfolgen? }

\begin{center}
\begin{minipage}{4.2cm}
{\mdseries
\gH{Bei der Pr\"ufung
{\quotedblbase}blind{\textquotedblleft} auf Nummer sicher gehen und
alle gleich f\"ur {\quotedblbase}vital bedroht{\textquotedblleft}
erkl\"aren gilt nicht! }}

{\mdseries
\gH{Viele Patienten bei der Pr\"ufung sind nicht vital
bedroht und brauchen keinen Notarzt. }}

{\mdseries
\gH{Erkl\"art man einen Patienten f\"ur vital bedroht
muss man das begr\"unden k\"onnen!}}

\end{minipage}
\end{center}
{\mdseries
Bei \gE{vital bedrohten} Patienten gibt es 5
Standardma{\ss}nahmen die sofort zu erfolgen haben. Bei
\gE{allen anderen} Patienten muss
\gE{individuell} \"uberlegt werden welche Ma{\ss}nahmen
zuerst zu treffen sind.}

\begin{center}
\tablehead{}\begin{supertabular}{|m{5.3770003cm}m{7.9230003cm}|}
\hline
\multicolumn{2}{|m{13.500001cm}|}{\centering
{\bfseries\upshape\color[rgb]{0.0,0.0,0.5019608} Erstma{\ss}nahmen bei
vital bedrohten Patienten}\par

}\\\hline
{\mdseries Situationsgerechte \gEs{Lagerung}}

 &
\\\hline
{\mdseries \gEs{Sauerstoffgabe}}

 &
{\mdseries je nach Indikation, Vorsicht bei COPD und Hyper\-ventilation}

\\\hline
{\mdseries \gEs{Notarzt}{}-Nachforderung}

 &
{\mdseries mit \gE{kurzer} Begr\"undung}

\\\hline
{\mdseries \gEs{Reanimationsbereitschaft}}

 &
{\mdseries Platz schaffen, Beatmungsbeutel zu\-sammen\-bauen, ggfs. Defi
und Absaugung holen und in Griffweite stellen}

\\\hline
{\mdseries \gEs{Vitalwerte} erheben und
\gEs{Monitoring}}

 &
{\mdseries RR, HF, Pulsoyxymetrie wenn vor\-handen}

\\\hline
\end{supertabular}
\end{center}
{\mdseries\color{black}
Dritter Schritt: Durchatmen, nachdenken, planen}

{\mdseries
Nach den Erstma{\ss}nahmen muss man nun planen welche Ma{\ss}nahmen,
Untersuchungen und Fragen am wichtigsten sind und zuerst durchgef\"uhrt
werden m\"ussen und was eher in der Abfolge nach hinten ver\-schoben
wird. Die folgende Checkliste gibt einen \"Uberblick \"uber die
Ma{\ss}nahmen und das Anamnese\-ge\-spr\"ach, die Reihenfolge ist bei
jedem Patienten anders!}

\clearpage\paragraph{Checkliste}
\begin{flushleft}
\tablehead{}\begin{supertabular}{m{2.467cm}|m{1.0209999cm}m{6.926cm}m{2.483cm}|}
\hline
\multicolumn{1}{|m{2.467cm}|}{} &
 &
\multicolumn{2}{m{9.609cm}|}{{\mdseries m\"ogliche Ma{\ss}nahmen, die
Reihenfolge ist je nach Patient unterschiedlich}

}\\\hline
\multicolumn{1}{|m{2.467cm}|}{\centering SELBSTSCHUTZ\par

\begin{center}
 [Warning: Image ignored] % Unhandled or unsupported graphics:
%\includegraphics[width=1.005cm,height=0.9cm]{AASdev20090228-img14.gif}

\end{center}
} &
 &
\multicolumn{2}{m{9.609cm}|}{\begin{itemize}
\item Schutzausr\"ustung (Handschuhe, Uniform, Helm)

\item Gefahrenbereich, Unfallstelle absichern, Hunde wegsperren lassen

\item ggfs. R\"uckzug, Spezialkr\"afte anfordern, {\dots} 

\item {\dots} 

\end{itemize}
}\\\hline
\multicolumn{1}{|m{2.467cm}|}{{\centering \par}

\begin{center}
 [Warning: Image ignored] % Unhandled or unsupported graphics:
%\includegraphics[width=1.12cm,height=0.63cm]{AASdev20090228-img15.gif}

\end{center}
\centering ERKENNEN\par

} &
 &
\multicolumn{2}{m{9.609cm}|}{\begin{itemize}
\item Situation/Milieu (zB Alter, soziales Umfeld, SMV, Alkohol,
Pensionistenheim, ...)

\item offensichtliche Symptome (zB Zyanose, Bewusstlosigkeit, starke
Blutung, ...)

\item Schockgefahr (Zyanose, kalter Schwei{\ss}, Radialispuls
tastbar/nicht tastbar, Puls beschleunigt)

\end{itemize}
}\\\hline
\multicolumn{1}{|m{2.467cm}|}{{\centering \par}

\begin{center}
 [Warning: Image ignored] % Unhandled or unsupported graphics:
%\includegraphics[width=1.12cm,height=1.12cm]{AASdev20090228-img16.png}

\end{center}
\centering ERSTE HILFE\par

} &
 &
\multicolumn{2}{m{9.609cm}|}{\begin{itemize}
\item Lagerung

\item Helmabnahme

\item Blutstillung

\item psychischer Beistand

\item W\"armeerhaltung

\end{itemize}
}\\\hline
\multicolumn{1}{|m{2.467cm}|}{{\centering   [Warning: Image ignored]
% Unhandled or unsupported graphics:
%\includegraphics[width=1.18cm,height=1.18cm]{AASdev20090228-img17.png}
 \par}

\centering ANAMNESE\par

} &
\centering {\bfseries\color[rgb]{0.7882353,0.0,0.08627451} S}\par

 &
\multicolumn{2}{m{9.609cm}|}{\gEs{Symptome }\&
\gEs{Schmerzen} 

(Zeit/Verlauf, St\"arke, Art/Qualit\"at, Lokalisation/Ausstrahlung,
Provokation)

}\\\hline
 &
\centering {\bfseries\color[rgb]{0.7882353,0.0,0.08627451} K}\par

 &
\multicolumn{2}{m{9.609cm}|}{\gEs{Krankheiten}
(Bestehen Grunderkrankungen, zB Diabetes?)

}\\\hhline{~---}
 &
\centering {\bfseries\color[rgb]{0.7882353,0.0,0.08627451} A}\par

 &
\multicolumn{2}{m{9.609cm}|}{\gEs{Allergien}/Allergische
Reaktion (auch Allergien auf Medikamente)

}\\\hhline{~---}
 &
\centering {\bfseries\color[rgb]{0.7882353,0.0,0.08627451} M}\par

 &
\multicolumn{2}{m{9.609cm}|}{\gEs{Medikamente}
(Nimmt Patient Medikamente? Wogegen sind diese?)

}\\\hhline{~---}
 &
\centering {\bfseries\color[rgb]{0.7882353,0.0,0.08627451} E}\par

 &
\multicolumn{2}{m{9.609cm}|}{\gEs{Ereignisse} vor
Notfalleintritt (zB Anstrengung vor Brustschmerz, Sturz, ...)

}\\\hhline{~---}
 &
\centering {\bfseries\color[rgb]{0.7882353,0.0,0.08627451} L}\par

 &
\multicolumn{2}{m{9.609cm}|}{\gEs{Letzte} ... (zB
letzte Mahlzeit, letzte Regelblutung, letzter Spitalsaufenthalt)

}\\\hhline{~---}
\multicolumn{1}{|m{2.467cm}|}{{\centering   [Warning: Image ignored]
% Unhandled or unsupported graphics:
%\includegraphics[width=1.679cm,height=1.82cm]{AASdev20090228-img18.jpg}
 \par}

\centering BEFUNDE\par

} &
\centering {\bfseries\color[rgb]{0.0,0.4,0.8} RR}\par

 &
\multicolumn{2}{m{9.609cm}|}{auskultatorische und palpatorische Messung

}\\\hline
 &
\centering {\bfseries\color[rgb]{0.0,0.4,0.8} HF}\par

 &
\multicolumn{2}{m{9.609cm}|}{Herzfrequenz (Puls), Puls rhythmisch oder
arrhytmisch

}\\\hhline{~---}
 &
\centering {\bfseries\color[rgb]{0.0,0.4,0.8} AF}\par

 &
\multicolumn{2}{m{9.609cm}|}{Atemfrequenz ausz\"ahlen

}\\\hhline{~---}
 &
\centering {\bfseries\color[rgb]{0.0,0.4,0.8} SpO\gSub{2}}\par

 &
\multicolumn{2}{m{9.609cm}|}{Sauerstoffs\"attigung im Blut, Monitoring

}\\\hhline{~---}
 &
\centering {\bfseries\color[rgb]{0.0,0.4,0.8} BZ}\par

 &
\multicolumn{2}{m{9.609cm}|}{Blutzucker (Wann war die letzte
Mahlzeit/Insulinspritze? Diabetes bekannt?)

}\\\hhline{~---}
 &
\centering {\bfseries\color[rgb]{0.0,0.4,0.8} Temp.}\par

 &
\multicolumn{2}{m{9.609cm}|}{K\"orpertemperatur/Fieber

}\\\hhline{~---}
 &
\centering {\bfseries\color[rgb]{0.0,0.4,0.8} Neuro\newline
status}\par

 &
\multicolumn{2}{m{9.609cm}|}{Bewusstsein (ev. GCS), Orientierung,
Pupillen, Halbseitenzeichen, Herdblick, Meningismus, retrograde Amnesie

}\\\hhline{~---}
 &
\centering {\bfseries\color[rgb]{0.0,0.4,0.8} Trauma\newline
check}\par

 &
\multicolumn{2}{m{9.609cm}|}{abnorme Beweglichkeit, DMS an allen
Extremit\"aten (Nagelbettprobe), Wunden, paradoxe Atmung,
Abwehrspannung des Bauches, ...

}\\\hhline{~---}
\multicolumn{1}{|m{2.467cm}|}{\centering SANIT\"ATSHILFE\par

\begin{center}
 [Warning: Image ignored] % Unhandled or unsupported graphics:
%\includegraphics[width=1.349cm,height=1.319cm]{AASdev20090228-img19.png}

\end{center}
} &
 &
\begin{itemize}
\item Entscheidung Notarzt-Nachforderung

\item Fortf\"uhrung der Ersten Hilfe

\item Schienung

\item Sauerstoffgabe

\item Reanimationsbereitschaft

\item Verlaufskontrolle

\item ...

\end{itemize}
 &
\\\hline
\multicolumn{1}{|m{2.467cm}|}{\centering NOTARZT-ASSISTENZ\par

\begin{center}
 [Warning: Image ignored] % Unhandled or unsupported graphics:
%\includegraphics[width=1.349cm,height=1.319cm]{AASdev20090228-img20.png}

\end{center}
} &
 &
\multicolumn{2}{m{9.609cm}|}{\begin{itemize}
\item Infusion/peripherven\"osen Zugang vorbereiten

\item Intubation vorbereiten

\item Medikamente vorbereiten

\item EKG anlegen

\item ...

\end{itemize}
}\\\hline
\multicolumn{1}{|m{2.467cm}|}{{\centering   [Warning: Image ignored]
% Unhandled or unsupported graphics:
%\includegraphics[width=1.88cm,height=1.88cm]{AASdev20090228-img21.png}
 \par}

\centering KRANKENHAUS\par

} &
 &
\multicolumn{2}{m{9.609cm}|}{\begin{itemize}
\item Welches Spital/Bett buche ich ab? (Arbeitsunfall? Welche Abteilung
fahre ich an? Ist der Patient bereits in einem Spital in Behandlung?)

\item Brauche ich eine Spezialabteilung? (Schockraum, Stroke Unit, PTCA,
Verbrennungsstation)

\item Arzt\"ubergabe (Anamnese, Befunde, Ma{\ss}nahmen, Hinweise)

\end{itemize}
}\\\hline
\end{supertabular}
\end{flushleft}
\paragraph[Ein Beispiel: Herr Sepp hat Atemnot]{Ein Beispiel: Herr Sepp
hat Atemnot}
{\mdseries
Berufung: Atemnot}

{\mdseries
Sie kommen in eine Wohnung. Die Umgebung erscheint sicher. Die Wohnung
sieht sauber und ordentlich aus. Auf der Couch liegt flach ein
50-j\"ahriger Mann mit weit aufgerissenen Augen und sichtlicher
Atemnot. Er wirkt bla{\ss} und schwei{\ss}ig und fasst ich mit der
rechten Hand an den Brustbereich. }

{\mdseries
Sie begr\"u{\ss}en den Patienten, stellen sich vor und fragen welche
Beschwerden er hat. Der Patient ringt nach Luft und klagt \"uber
Schmerzen in der Brust. Diese haben vor einer Viertelstunde angefangen.
}

{\mdseries\itshape
Welche Informationen haben Sie bis jetzt gewonnen?}

{\mdseries\itshape
\gH{Der Patient ist bei Bewusstsein. Er hat Atemnot
und einen pl\"otzlich eingetretenen Thoraxschmerz. Er ist bla{\ss} und
schwei{\ss}ig, vielleicht steckt ein Schockgeschehen dahinter. }}

{\mdseries\itshape
K\"onnen Sie schon etwas tun?}

{\mdseries\itshape
\gH{Ja.  Bei erhaltenem Bewusstsein, Atemnot und
Thoraxschmerz wird der Patient mit erh\"ohtem Oberk\"orper gelagert.
Weiters sprechen die Symptome f\"ur den Einsatz von Sauerstoff. Es muss
noch schnell gekl\"art werden ob eine Hyperventilationstetanie (nicht
wahrscheinlich) oder COPD-Erkrankung (erfragen) vorliegt. }}

{\mdseries\itshape
Ist der Patient vital bedroht?}

{\mdseries\itshape
\gH{Ja. Zwar wurden noch keien Vitalwerte erhoben,
aber bereits jetzt gilt der Patient als vital bedroht aufgrund der
pl\"otzlich einsetzenden Thoraxschmerzen und der Atemnot. }}

\begin{itemize}
\item {\mdseries\color[rgb]{0.0,0.0,0.5019608}
Der Oberk\"orper des Patienten wird hochgelagert }

\item {\mdseries\color[rgb]{0.0,0.0,0.5019608}
Der Notarzt wird nachgefordert (Berufungsgrund: akuter Thoraxschmerz mit
Atemnot und schlechtem AZ).}

\item {\mdseries\color[rgb]{0.0,0.0,0.5019608}
Der Patient wird gefragt ob er an COPD leidet, dies wird verneint. Eine
Sauerstoffmaske mit einem Flow von 10l/min wird angelegt. }

\item {\mdseries\color[rgb]{0.0,0.0,0.5019608}
Reanimationsbereitschaft wird hergestellt}

\end{itemize}
{\mdseries\itshape
Was sind die n\"achsten Schritte?}

{\mdseries\itshape
\gH{Wichtig sind nun die Vitalwerte und das fr\"uhe
Beginnen des Monitorings.Dann muss man genaueres \"uber das
Beschwerdebild in Erfahrung bringen. Der weitere Plan h\"angt davon
ab.}}

{\mdseries
Die Vitalwerte werden gemessen und ein Pulsoxy angeh\"angt (RR 100/60,
HF 110, SpO2 94\%). Der Patient gibt an dass er st\"arkste
stechende/brennende Schmerzen hinter dem Brustbein hat die ins Kiefer
ausstrahlen. Ausserdem f\"uhle es sich so an als w\"urde jemand auf dem
Patienten drauf stehen. Die Schmerzen sind nicht atem- oder
bewegungsabh\"angig  und haben vor einer Viertelstunde
{\quotedblbase}wie aus dem Nichts heraus{\textquotedblleft} begonnen.
Ausserdem bekomme er schlecht Luft, sonst h\"atte er keine Beschwerden.
}

{\mdseries\itshape
An welche Differentialdiagnosen denken Sie?}

{\mdseries\itshape
\gH{Leitsymptom: Thoraxschmerz }}

\begin{center}
\tablehead{}\begin{supertabular}{|m{3.225cm}|m{3.225cm}|m{3.225cm}|m{3.225cm}|}
\hline
\multicolumn{4}{|m{13.5cm}|}{\centering
{\bfseries\upshape\color[rgb]{0.0,0.0,0.5019608} Tabelle:
Differentialdiagnosen Thoraxschmerz unseres Patienten}\par

}\\\hline
 &
\centering {\bfseries daf\"ur spricht}\par

 &
\centering {\bfseries dagegen spricht}\par

 &
{\mdseries Bewertung}

\\\hline
{\bfseries Herzinfarkt, Angina pectoris}

 &
{\mdseries Schmerzen hinter dem Brustbein, Ausstrahlung, nicht
atemabh\"angig, Druck-/Engegef\"uhl, Atemnot}

 &
 &
{\mdseries sehr wahrscheinlich}

\\\hline
{\bfseries Lungenembolie}

 &
{\mdseries Schmerzen hinter dem Brustbein, Atemnot}

 &
{\mdseries untypische Ausstrahlung, }

 &
{\mdseries gut m\"oglich}

\\\hline
{\bfseries Pneumothorax}

 &
{\mdseries Schmerzen im Brustbereich, Atemnot}

 &
{\mdseries pl\"otzliche st\"arkste Schmerzen}

 &
{\mdseries eher unwahrscheinlich}

\\\hline
{\bfseries Trauma}

 &
{\mdseries Schmerzen}

 &
{\mdseries kein Hinweis}

 &
{\mdseries unwahrscheinlich}

\\\hline
{\bfseries Erkrankung des St\"utz- und Bewegungsapparats}

 &
{\mdseries Schmerzen }

 &
{\mdseries untypische Ausstrahlung, Atemnot, nicht bewegungsabh\"angig}

 &
{\mdseries sehr unwahrscheinlich}

\\\hline
\end{supertabular}
\end{center}
\gH{Die wahrscheinlichsten Diagnosen sind der
Herzinfarkt und die Lungenembolie.}

{\mdseries\itshape
Fallen Ihnen bei diesen Erkrankungen wichtige weitere Untersuchungen
ein?}

{\mdseries\itshape
\gH{Bei beiden Erkrankungen ist es wichtig beide Beine
zu begutachten. Im Falle einer Lungenembolie sucht man nach einer
Beinvenenthrombose (einseitig \"uberw\"armtes, bl\"aulich-r\"otliches
geschwollenes Bein). Bei einer Herzinsuffizienz bzw. Herzinfarkt
beurteilt man den R\"uckstau des Blutes, dh. ob die Beine geschwollen
sind.}}

{\mdseries\itshape
\gH{Ausserdem, hat der Patient so etwas gehabt? Oder
sind einschl\"agige Erkrankungen bekannt?  }}

{\mdseries
Sie begutachten beide Beine des Patienten und stellen mittels
Daumendruckprobe fest dass beide Beine leicht geschwollen sind. Auf
Nachfrage gibt der Patient an dass er sonst keine geschwollenen Beine
hat. }

{\mdseries\itshape
Was sind die n\"achsten Schritte?}

{\mdseries\itshape
\gH{Die wichtigsten Untersuchungen wurden gemacht, nun
muss man sich um die Anamnese k\"ummern. }}

{\mdseries
Symptome: siehe oben.}

{\mdseries
Krankheiten: Der Patient leidet an nicht-Insulinpflichtigem Diabetes
mellitus und Bluthochdruck. Sonst sind keine Krankheiten bekannt.
Spitalsbriefe gibt es nicht. }

{\mdseries\itshape
Sie merken sich dass Sie nachher den Blutzucker messen wollen!}

{\mdseries
Allergien: keine bekannt.}

{\mdseries
Medikamente: Glucophage 850mg f\"ur den Zucker, Tritace f\"ur den
Blutdruck. Medikamente wurden heute wie verschrieben eingenommen. }

{\mdseries
Ereignisse vor Notfalleintritt: Schmerzen begannen pl\"otzlich als der
Patient W\"asche waschen wollte. In den letzten Wochen immer wieder
Stiche in der HErzgegend, hat er aber nicht ernst genommen. }

{\mdseries
Letzte(r):}

{\mdseries\itshape
Spitalsaufenthalt: Vor 10 Jahren in der Rudolfstiftung wegen Rotlauf.}

{\mdseries\itshape
Arztbesuch: Vor einer Woche wegen Neuaustellung eines Rezeptes.}

{\mdseries\itshape
Mahlzeit: Mittagessen, Gulasch}

{\mdseries\itshape
Regel: Der Patient ist m\"annlich!}

pro:

Lungenembolie 

Wichtig sind nun die Vitalwerte und das fr\"uhe Beginnen des
Monitorings.Dann muss man genaueres \"uber das Beschwerdebild in
Erfahrung bringen. Der weitere Plan h\"angt davon ab.

\section{[RS IV] Anatomie, Physiologie}
{\mdseries\itshape
Sebastian Gabriel, Michael Withofner}

{\mdseries\itshape
mit Beitr\"agen von Gerhard Gabriel}

{\bfseries\color[rgb]{0.5019608,0.0,0.0}
\gEs{Changelog:}}

{\bfseries\color[rgb]{0.5019608,0.0,0.0}
2009-02-25: Start Changelog.}

\subsection{Das Skelett und der Bewegungsapparat}
{\mdseries
Das Skelettsystem besteht aus}

\begin{itemize}
\item R\"ohren knochen (z.B. Oberarmknochen)

\item platten  Knochen (z.B. Schulterblatt,  Sch\"adelknochen)

\end{itemize}
{\mdseries
und dient  haupts\"achlich als}

\begin{itemize}
\item St\"utzger\"ust  des K\"orpers

\item Ursprungs-  und Ansatz fl\"ache f\"ur die Muskulatur

\end{itemize}
{\mdseries
und weiters als}

\begin{itemize}
\item KalZiumspeicher (Ca)

\item \textbf{Blutbildungsorgan} (rotes Knochenmark!)

\end{itemize}
\paragraph{Aufbau eines Knochens}
\begin{center}
\tablehead{}\begin{supertabular}{|m{5.5150003cm}|m{7.7850003cm}|}
\hline
\index{Periost}\gT{Periost}: \"uberzieht den Knochen und
ist sehr gut mit Nerven versorgt.  Verantwortlich f\"ur Dickenwachstum
und Frakturheilung \gZ{(Kallusbildung)}

\index{Kortikalis}\gT{Kortikalis}: feste \"au{\ss}ere
Schicht, \"au{\ss}ere Form

\index{Spongiosa}\gT{Spongiosa}:
\index{Knochenb\"alkchen}Knochenb\"alkchen (Trabekel) ${\rightarrow}$
Zug- und Druckfestigkeit

 &
[Warning: Draw object ignored]

\\\hline
\end{supertabular}
\end{center}
\paragraph[Das Knochenmark]{Das Knochenmark}
{\mdseries\color{black}
\textcolor{red}{Rotes} Knochenmark ist f\"ur die Blutbildung
zust\"andig}

{\mdseries
Das rote Knochenmark ist das blutbildende Organ schlechthin.
\gT{Rote Blutk\"orperchen}
\gZ{(Erythrozyzen)}, \gT{wei{\ss}e
Blutk\"orperchen} \gZ{(Leukozyten)} sowie
\gT{Blutpl\"attchen}
\gZ{(Thrombozyten)} werden hier gebildet. Bei der
Geburt findet man es in s\"amtlichen Knochen, im Laufe des Lebens
bildet es sich zur\"uck und verwandelt sich in {\quotedblbase}Gelbes
Fettmark{\textquotedblleft}. Rotes Knochenmark findet man dann
lediglich in den platten Knochen wie Becken, Brustbein und Sch\"adel wo
es f\"ur die Blutbildung verantwortlich ist und normalerweise
zeitlebens erhalten bleibt. }

{\mdseries
Einige medizinische Behandlungsmethoden verursachen als Nebenwirkung
eine massive \gE{Sch\"adigung des Knochenmarks} (z.B.
Bestrahlung, Chemotherapien bei Krebspatienten
\gZ{({\quotedblbase}onkologsiche
Patienten{\textquotedblleft})}). Diese Patienten leiden neben einer
Blutarmut auch an einem schwachen Immunsystem, da nicht gen\"ugend
Wei{\ss}e Blutk\"orperchen gebildet werden. Sie k\"onnen daher sehr
leicht \gE{schwerwiegende Infektionen} bekommen. Diese
Pateinten sind im Krankentransport h\"aufig anzutreffen. }

{\mdseries\color{black}
\textcolor[rgb]{1.0,0.70980394,0.08235294}{Gelbes} Knochenmark besteht
aus Fett }

{\mdseries
Im Laufe des Lebens bildet es sich das rote Knochenmark in den
R\"ohrenknochen zu fettreichem, gelben Knochenmark um. }

{\mdseries
Das gelbe Fettmark kann bei schweren
\gE{Knochenverletzungen} von Bedeutung sein, da
verletzungs\-bedingt kleine Fettptr\"opchen in die Blutbahn gelangen
k\"onnen, die dann als \gE{Fettgerinnsel}  zu ernsten
Komplikationen f\"uhren k\"onnen. }

\paragraph{Knochenwachstum}
\begin{center}
\tablehead{}\begin{supertabular}{|m{6.649cm}|m{6.651cm}|}
\hline
Das Skelett eines Neugeborenen besteht bei der Geburt zum gr\"o{\ss}ten
Teil aus Knorpel. Erst mit zunehmendem Lebensalter beginnt dieser von
Knochenkernen ausgehend, zu verkn\"ochern. Eine
\gEs{Wachstumszone}
(\index{Epiphysenfuge}\gT{Epiphysenfuge}) bleibt dabei
allerdings erhalten.

{\mdseries Die \textit{Epiphysenfuge} enth\"alt beim Kind /Jugendlichen
bis etwa zum Ende der Pubert\"at
\gEs{teilungsf\"ahiges Knorpelgewebe}, daher ist ein
\gEs{L\"angenwachstum} m\"oglich. Beim Erwachsenen
sind diese knorpeligen Fugen hingegen vollst\"andig
\gEs{verkn\"ochert}, soda{\ss} keinerlei Wachstum
mehr stattfinden kann.}

 &
{\centering   [Warning: Image ignored]
% Unhandled or unsupported graphics:
%\includegraphics[width=3.557cm,height=4.948cm]{AASdev20090228-img22.png}
 \par}

\\\hline
\end{supertabular}
\end{center}
\paragraph[Gelenke - \ prinzipieller Aufbau]{Gelenke -  prinzipieller
Aufbau}
\begin{center}
\tablehead{}\begin{supertabular}{|m{6.649cm}|m{6.651cm}|}
\hline
{\mdseries Bewegliche Verbindung zwischen mindestens zwei Knochen}

\begin{itemize}
\item Ein Gelenk ist von einer bindegewebigen Kapsel umgeben und durch
Sehnen und B\"ander stabil verbunden.

\item Die Gelenksfl\"achen sind zur Reibungs\-ver\-minderung mit 
Knorpel \"uberzogen und der Gelenksspalt mit einer d\"unnen Schicht 
Schmierfl\"ussigkeit \gZ{(Synovialfl\"ussigkeit)}
gef\"ullt. 

\end{itemize}
 &
{\centering [Warning: Draw object ignored]\par}

\\\hline
\end{supertabular}
\end{center}
\subparagraph{Es gibt verschiedene Arten von Gelenken}
{\mdseries
Klassische Beispiele: \index{Kugelgelenk}\gT{Kugelgelenk}
(H\"uftgelenk),
\index{Scharniergelenk}\gT{Scharniergelenk}
(Ellenbogengelenk),
\index{Sattelgelenk}\gT{Sattelgelenk}
(Daumengrundgelenk).  [Warning: Image ignored]
% Unhandled or unsupported graphics:
%\includegraphics[width=13.701cm,height=8.482cm]{AASdev20090228-img23.png}
 }

{\mdseries\color[rgb]{0.0,0.0,0.5019608}
Gliederung des Skeletts}

\begin{itemize}
\item {\mdseries
Kopf  (\gT{Caput})}

\end{itemize}
\begin{itemize}
\item {\mdseries
Gehirnsch\"adel}

\item {\mdseries
Gesichtssch\"adel}

\end{itemize}
\begin{itemize}
\item {\mdseries
Rumpf  }

\end{itemize}
\begin{itemize}
\item {\mdseries
Hals }

\item {\mdseries
Brustkorb (\gEs{Thorax} )}

\item {\mdseries
Bauch (\gEs{Abdomen} )}

\item {\mdseries
Becken (}

\end{itemize}
\begin{itemize}
\item {\mdseries
obere Extremit\"aten }

\item {\mdseries
untere Extremit\"aten}

\end{itemize}
{\mdseries\color[rgb]{0.0,0.0,0.5019608}
Der Sch\"adel }

\begin{center}
\tablehead{}\begin{supertabular}{|m{6.649cm}|m{6.651cm}|}
\hline
{\mdseries Kopf (\gT{Caput}) / Sch\"adel
(\gT{Cranium})}

{\mdseries Der Sch\"adel gliedert sich in den
\gT{Hirnsch\"adel} und den Gesichtssch\"adel}

\captionof{figure}[Hirnsch\"adel vs. Gesichts\-sch\"adel]{Hirnsch\"adel
vs. Gesichts\-sch\"adel}
 &
{\centering [Warning: Draw object ignored]\par}

\\\hline
\end{supertabular}
\end{center}
\paragraph[Gesichtssch\"adel]{Gesichtssch\"adel}
\begin{center}
\tablehead{}\begin{supertabular}{|m{6.649cm}|m{6.651cm}|}
\hline
{\mdseries a ... Stirnbein \gZ{(Os frontale)}}

{\mdseries b ... Nasenbein \gZ{(Os nasale)}}

{\mdseries c ... Scheitelbein \gZ{(Os parietale)}}

{\mdseries d ... Schl\"afenbein \gZ{(Os temporale)}}

{\mdseries e ... Jochbein \gZ{(Os zygomaticum)}}

{\mdseries f ... Oberkiefer \gZ{(Maxilla)}}

{\mdseries g ... Unterkiefer \gZ{(Mandibula)}}

 &
[Warning: Draw object ignored]

\\\hline
\end{supertabular}
\end{center}
\paragraph{Gehirnsch\"adel}
\begin{center}
\tablehead{}\begin{supertabular}{|m{6.649cm}|m{6.651cm}|}
\hline
\begin{itemize}
\item {\mdseries Sch\"adelbasis}

\item {\mdseries Sch\"adeldach (Kalotte )}

\item {\mdseries Die platten Knochen sind durch Sch\"adeln\"ahte 
(Suturae) verzahnt miteinander verkn\"ochert}

\end{itemize}
{\mdseries a ... Stirnbein \gZ{(Os frontale)}}

{\mdseries b ... Scheitelbein \gZ{(Os parietale)}}

{\mdseries c ... Schl\"afenbein \gZ{(Os temporale)}}

{\mdseries d .... Hinterhauptsbein \gZ{(Os
occipitale)}}

{\mdseries e ... gro{\ss}er Keilbeinfl\"ugel \gZ{(Os
sphenoidale)}}

 &
[Warning: Draw object ignored]

\\\hline
\end{supertabular}
\end{center}
\subparagraph[Fontanelle -- Sutur]{Fontanelle -- Sutur}
\index{Fontanelle}\index{Sutur}\begin{center}
\tablehead{}\begin{supertabular}{|m{8.6640005cm}|m{4.642cm}|}
\hline
{\mdseries Die \gT{Fontanellen} sind
\gEs{knorpelige  Verbindungen}  zwischen 
Sch\"adelknochen beim Neu\-geborenen (der Sch\"adel mu{\ss} sich f\"ur 
Geburt verformen k\"onnen!). Sie schlie{\ss}en  sich bis zum 18. LM,
verkn\"ochern  aber erst vollst\"andig mit zunehmendem  Lebensalter
(ca. 20 Lj. abgeschlossen). }

\begin{itemize}
\item {\mdseries gro{\ss}e Fontanelle (rot): zw. Stirnbein u.
Scheitelbeinen}

\item {\mdseries kleine Fontanelle (gr\"un): zw. Scheitelbeinen u.
Hinter\-haupts\-bein}

\end{itemize}
 &
{\centering [Warning: Draw object ignored]\par}

\\\hline
\end{supertabular}
\end{center}
{\mdseries
Die \gT{Sutur} ist die endg\"ultig
\gEs{fest verkn\"ocherte Naht} zwischen den
Sch\"adelknochen beim Erwachsenen ${\rightarrow}$ starrer  stabiler 
Sch\"adel, Schutz  f\"ur Gehirn.}

\paragraph[Sch\"adelbasis]{Sch\"adelbasis}
\begin{center}
\tablehead{}\begin{supertabular}{|m{6.649cm}|m{6.651cm}|}
\hline
{\mdseries a ... Stirnbein \gZ{(Os frontale)}}

{\mdseries b ... Siebbein \gZ{(Os ethmoidale)}}

{\mdseries c ... Keilbein  \gZ{(Os sphenoidale)}}

{\mdseries d ... Schl\"afenbein(e)  \gZ{(Os
temporale)}}

{\mdseries e ... Felsenbein  \gZ{(Pars petrosa ossis
temporalis)}}

{\mdseries f ... Hinterhauptsbein  \gZ{(Os
occipitale)}}

{\mdseries g ... gro{\ss}es Hinterhauptsloch (Foramen magnum)}

 &
{\centering [Warning: Draw object ignored]\par}

\\\hline
\end{supertabular}
\end{center}
{\mdseries
Gehirn liegt auf der Sch\"adelbasis auf \gZ{(ausser
wenn man auf dem Kopf steht)}. }

{\mdseries
Hirnnerven treten durch div. L\"ocher aus dem Sch\"adel aus.}

{\mdseries
\textbf{Das gro{\ss}e Hinterhauptsloch (Foramen magnum) ist die
Durchtrittsstelle des R\"uckenmarks.} }

{\mdseries\color{black}
\index{Hirnstammeinklemmung}\gT{Hirnstammeinklemmung}}

{\mdseries\color{black}
Das \gEs{gro{\ss}e Hinterhauptsloch} (Foramen magnum)
ist die Durchtrittsstelle des R\"uckenmarks. Bei einer
\gEs{Hirndrucksteigerung}
(\gEs{Hirn\"odem},
\gEs{Hirnblutung}) ist es die einzige M\"oglichkeit
zur Ausdehnung des Gehirns! \gEs{Dabei dr\"uckt sich
das Gehirn allerdings selber ab}, es kommt zu einer t\"odlichen
\gEs{Hirnstammeinklemmung!} }

{\mdseries\color{black}
${\rightarrow}$ \gCa{ Lebensgefahr! } }

{\mdseries\color[rgb]{0.0,0.0,0.5019608}
Die Wirbels\"aule (Columna vertebralis)}

{\bfseries\color[rgb]{0.5019608,0.0,0.0}
Bild von Wirbels\"aule gest\"ort!}

{\mdseries
[Warning: Draw object ignored]%
%
Die Wirbels\"aule besteht aus:}

\begin{itemize}
\item \textcolor{red}{7 Halswirbeln}  (Cervical-Wirbel)

\item \textcolor[rgb]{0.0,0.0,0.5019608}{12 Brustwirbeln} 
(Thorakal-Wirbel)

\item \textcolor[rgb]{0.0,0.5019608,0.0}{5 Lendenwirbeln} 
(\gEs{Lumbal}{}-Wirbel)

\item 5 verschmolzene Kreuzbein wirbel (Os sacrum)

\item \textcolor[rgb]{0.5019608,0.0,0.0}{3-4 Stei{\ss}bein
{\textquotedblleft}wirbel{\textquotedblleft}}  \gZ{(Os
coccygis)} 

\end{itemize}
{\mdseries
Doppelt S-f\"ormig gekr\"ummtes Achsenskelett}

\begin{itemize}
\item {\mdseries
erm\"oglicht den aufrechten Gang}

\end{itemize}
{\mdseries\color{black}
besteht aus}

\begin{itemize}
\item {\mdseries
Wirbeln \gZ{(vertebrae)}}

\item {\mdseries
\index{Bandscheiben}\gT{Bandscheiben} (Discus,
\gZ{Pl: Disci intervertebrales}): dienen der
{\quotedblbase}Sto{\ss}d\"ampfung{\textquotedblleft}}

\item {\mdseries
\gEs{Im Wirbelkanal liegt das R\"uckenmark
eingebettet}.}

\end{itemize}
\subparagraph{Die Wirbel: Aufbau und Arten}
{\mdseries
Prinzipieller Aufbau:}

\begin{itemize}
\item Wirbelk\"orper

\item Wirbelbogen  (umschlie{\ss}t das R\"uckenmark)

\item 2 Querforts\"atze

\item 1 Dornfortsatz

\end{itemize}
{\mdseries
%
%
In den entsprechenden Abschnitten sehen die Wirbel unterschiedlich aus.
}

{\bfseries\color[rgb]{0.5019608,0.0,0.0}
Bild von Wirbelk\"orper hier einf\"ugen!}

{\mdseries\color[rgb]{0.0,0.0,0.5019608}
Der \index{Brustkorb}Brustkorb (\index{Thorax}Thorax)}

{\mdseries
Der  Brustkorb besteht vorne aus dem \gT{Brustbein}
(\gT{Sterrnum}) und 12 Rippenpaaren, welche hinten mit
der Wirbels\"aule in gelenkiger Verbindung stehen. Die obersten 7
Rippen-Paare werden als  {\quotedblbase}\gT{echte
Rippen}{\textquotedblleft} bezeichnet. Sie erreichen mit ihren Spitzen
das Brustbein und bilden praktisch
{\quotedblbase}Ringe{\textquotedblleft} um die Brusteingeweide. Die 5
unteren Rippenpaare bilden den bogenf\"ormigen unteren Rand des
Brustkorbes. Die Rippenpaar 8-10 sind dabei nur indirekt \"uber einen
Knorpel mit dem Brustbein verbunden ({\quotedblbase}unechte
Rippen{\textquotedblleft}), das 11. und das 12. Paar endet mit
knorpeligen Spitzen frei zwischen den
Bauchmuskeln({\quotedblbase}\gT{fliehende
Rippen}{\textquotedblleft}) bezeichnet werden.}

{\mdseries
Der Brustkorb enth\"alt und sch\"utzt die Brust\-organe (Lunge, Herz)
sowei die Oberbauchorgane (Leber, Magen). Zwischen den Rippen spannt
sich die \gT{Zwischenrippenmuskulatur} (wichtig f\"ur
\gEs{Atemmechanik}), in diese sind
\gE{segmentale Nerven- und
Gef\"a{\ss}b\"undel}\gEs{ }eingelagert.}

\begin{center}
\tablehead{}\begin{supertabular}{|m{6.649cm}|m{6.651cm}|}
\hline
  [Warning: Image ignored] % Unhandled or unsupported graphics:
%\includegraphics[width=4.388cm,height=5.019cm]{AASdev20090228-img24.png}
 

 &
{\mdseries Der Thorax besteht aus:}

\begin{itemize}
\item {\mdseries Brustbein 
(\index{Sternum}\gT{Sternum})}

\item {\mdseries 12  Rippenpaaren }

\end{itemize}
\begin{itemize}
\item {\mdseries 7 {\quotedblbase}echte{\textquotedblleft}  }

\item {\mdseries 3 {\quotedblbase}unechte{\textquotedblleft}}

\item {\mdseries 2 {\quotedblbase}fliegenden{\textquotedblleft}}

\end{itemize}
\begin{itemize}
\item {\mdseries Brustwirbels\"aule (12 Brust\-wirbel), sie Steht mit
den Rippen in gelenkiger Verbindung }

\end{itemize}
{\mdseries ({\quotedblbase}\textit{echt}{\textquotedblleft} bedeutet die
Rippe hat eine direkte Verbindung zum Sternum)}

\\\hline
\end{supertabular}
\end{center}
{\mdseries\color[rgb]{0.0,0.0,0.5019608}
Der \index{Schulterg\"urtel}Schulterg\"urtel}

\begin{center}
\tablehead{}\begin{supertabular}{|m{6.649cm}|m{6.651cm}|}
\hline
{\mdseries ... besteht aus (von vorne nach hinten) :}

\begin{itemize}
\item {\mdseries \index{Brustbein}\gT{Brustbein} 
(\index{Sternum}\gT{Sternum})}

\item {\mdseries
\index{Schl\"usselbein}\gT{Schl\"usselbein} 
\gZ{(Clavicula)}}

\item {\mdseries \index{Schulterblatt}\gT{Schulterblatt} 
\gZ{(Scapula)}}

\end{itemize}
{\mdseries Der Schulterg\"urtel bildet die Gelenkspfanne f\"ur das
Schultergelenk, jeweils auf beiden Seiten.}

{\mdseries Die Schulterbl\"atter sind nur durch Muskulatur mit dem
Brustkorb verbunden.}

 &
{\centering   [Warning: Image ignored]
% Unhandled or unsupported graphics:
%\includegraphics[width=4.379cm,height=4.453cm]{AASdev20090228-img25.png}
 \par}

\\\hline
\end{supertabular}
\end{center}
{\mdseries\color[rgb]{0.0,0.0,0.5019608}
Die obere Extremit\"at}

\begin{center}
\tablehead{}\begin{supertabular}{|m{6.649cm}|m{6.651cm}|}
\hline
\begin{itemize}
\item {\mdseries Schultergelenk}

\item {\mdseries Oberarmknochen  (Humerus)}

\item {\mdseries Ellenbogengelenk}

\item {\mdseries Unterarmknochen (2)}

\end{itemize}
\begin{itemize}
\item {\mdseries Elle  (Ulna) }

\item {\mdseries Speiche  (Radius), daumenseitig!}

\end{itemize}
\begin{itemize}
\item {\mdseries Handwurzelgelenk}

\item {\mdseries
\gZ{(}\gZ{8}\gZ{)}
Handwurzelknochen}

\item {\mdseries
\gZ{(}\gZ{5}\gZ{)}
Mittelhandknochen }

\item {\mdseries Finger: je 3  Endglieder}

\item {\mdseries Daumen: nur 2  Endglieder}

\end{itemize}
 &
[Warning: Draw object ignored]

\\\hline
\end{supertabular}
\end{center}
{\mdseries
Achtung: Unterscheide: Hand -- Arm!}

{\mdseries\color[rgb]{0.0,0.0,0.5019608}
Der Beckeng\"urtel}

\begin{center}
\tablehead{}\begin{supertabular}{|m{6.649cm}|m{6.651cm}|}
\hline
{\mdseries Besteht aus (von hinten nach vorne) :}

\begin{itemize}
\item {\mdseries Kreuzbein \gZ{(Os sacrum)}}

\item {\mdseries Beckenknochen, bestehend aus:}

\end{itemize}
\begin{itemize}
\item {\mdseries \index{Darmbein}\gT{Darmbein}
\gZ{(Os ilii)}}

\item {\mdseries \index{Sitzbein}\gT{Sitzbein}
\gZ{(Os ischii)}}

\item {\mdseries
\index{Schambein}\gT{Schambein}\gZ{
}\gZ{(Os pubis)}}

\end{itemize}
{\mdseries Die \index{Symphyse}\gT{Symphyse} stellt eine
straffe Knorpelverbindung vorne zwischen den beiden Beckenknochen dar
(Dehnungs\-fuge f\"ur Geburt).}

 &
{\centering   [Warning: Image ignored]
% Unhandled or unsupported graphics:
%\includegraphics[width=5.527cm,height=3.555cm]{AASdev20090228-img26.png}
 \par}

\\\hline
\end{supertabular}
\end{center}
{\mdseries\color[rgb]{0.0,0.0,0.5019608}
Die untere Extremit\"at}

\begin{center}
\tablehead{}\begin{supertabular}{|m{6.649cm}|m{6.651cm}|}
\hline
\begin{itemize}
\item {\mdseries Oberschenkel  (Femur)}

\item {\mdseries Unterschenkelknochen (2)}

\end{itemize}
\begin{itemize}
\item {\mdseries Schienbein  (Tibia) , dick}

\item {\mdseries Wadenbein  (Fibula), d\"unn}

\end{itemize}
\begin{itemize}
\item {\mdseries 8 Fu{\ss}wurzelknochen}

\item {\mdseries 5 Mittelfu{\ss}knochen  (Ossa metatarsalia)}

\item {\mdseries 5 Zehen:}

\end{itemize}
\begin{itemize}
\item {\mdseries 4 Zehen: je 3  Endglieder}

\item {\mdseries 1 Gro{\ss}zehe: nur 2 Endglieder}

\end{itemize}
 &
{\centering %
%
[Warning: Draw object ignored]\par}

\\\hline
\end{supertabular}
\end{center}
{\bfseries\color[rgb]{0.5019608,0.0,0.0}
Untere Extremit\"at Bild muss ausgetuascht werden.}

{\mdseries
Achtung: Unterscheide Fu{\ss} -- Bein!}

{\mdseries\color[rgb]{0.0,0.0,0.5019608}
Die Muskulatur}

{\mdseries
Prinzipiell gibt es 3  Arten von Muskelzellen:}

\begin{itemize}
\item Quergestreifte (Skelett-)Muskulatur: willk\"urlich steuerbar,
trainierbar, er\-m\"ud\-bar

\item Glatte  Muskulatur: Nicht willentlich kontrollierbar,
unerm\"udlich, ausdauernd 

\item Herz muskel: Sonderstellung zwischen glatt und quergestreift,
Eigenschaften von beiden

\end{itemize}
{\mdseries
Die quergestreifte Muskulatur heisst deswegen
{\quotedblbase}quergestreift{\textquotedblleft}, da man unter dem
Mikroskop tats\"achlich eine Streifung sehen kann, allerdings nicht mit
dem freien Auge) }

\begin{center}
\tablehead{}\begin{supertabular}{|m{6.649cm}|m{6.651cm}|}
\hline
{\mdseries Jeder Muskel hat eine vorgegebene  Zugrichtung. }

{\mdseries Muskeln k\"onnen sich nicht aktiv  entspannen, daher sind
Gegenspieler  (Agonist u. Antagonist) erforderlich (z.B. ein Muskel
streckt den Arm, ein anderer beugt ihn).}

 &
%
%
[Warning: Draw object ignored]

\\\hline
\end{supertabular}
\end{center}
{\bfseries\color[rgb]{0.5019608,0.0,0.0}
Ellbogen -- Muskelzug: Bild ersetzen (grau)}

\subsection{Das Herz}
{\mdseries
Das Herz ist ein etwa faustgro{\ss}er
\gEs{Hohlmuskel} und hat die Funktion einer
\gEs{Pumpe}. Es liegt \gEs{hinter
dem Sternum}\footnote{Brustbein} und teilweise im linken Brustkorb.
Unten liegt es auf dem \gEs{Zwerchfell} auf. }

\begin{center}
 [Warning: Image ignored] % Unhandled or unsupported graphics:
%\includegraphics[width=4.887cm,height=6.693cm]{AASdev20090228-img27.png}

\end{center}
{\mdseries
Das Herz ist durch eine \gEs{Scheidewand}
(\gT{Septum}) in eine rechte und eine linke Seite
geteilt. Umgangssprachlich spricht man dabei jeweils vom
{\quotedblbase}rechten{\textquotedblleft} oder {\quotedblbase}linken
Herz{\textquotedblleft}. Jede H\"alfte wird durch Klappen in einen
\gEs{Vorhof} (\gT{Atrium}) und eine
\gEs{Herzkammer} (\gT{Ventrikel})
geteilt. }

{\mdseries\color{black}
Herzklappen}

{\mdseries
Insgesamt \gEs{4 Herzklappen} sorgen als
\gEs{Ventile} daf\"ur dass das Blut nur in eine
Richtung str\"omt. Es gibt \gEs{2 Segelklappen} und
\gEs{2 Taschenklappen}. Die Segelklappen befinden
sich zwischen den Vorh\"ofen und den Herzkammern, die Taschenklappen
befinden sich an den Ausfluss\"offnungen der Kammern am \"Ubergang zur
Lungenarterie (rechte Seite) bzw. zur Hauptschlagader (linke Seite). }

{\mdseries\color{black}
Blutversorgung}

{\mdseries
Die Blutversorgung erfolgt \"uber die
\gEs{Herzkranzgef\"a{\ss}e}
(\gT{Koronargef\"a{\ss}e}) welche direkt aus der
Hauptschlagader abgehen. }

{\mdseries\color{black}
Schichten}

{\mdseries
3 Schichten (au{\ss}en nach innen):}

\begin{itemize}
\item {\mdseries
Epikard (aussen)}

\item {\mdseries
\gEs{Myokard} (Muskel)}

\item {\mdseries
Endokard (innen)}

\end{itemize}
\begin{minipage}{13.6cm}
 [Warning: Image ignored] % Unhandled or unsupported graphics:
%\includegraphics[width=13.6cm,height=6.168cm]{AASdev20090228-img28.jpg}
\captionof{figure}[Schema der Herz-Anatomie. ]{Schema der Herz-Anatomie.
}
{\mdseries\upshape
Quelle: Wikipedia,  J\"org Rittmeister. G-FDL, CCAL}

\end{minipage}

{\mdseries\color{black}
Herzfunktion}

{\mdseries
Der Herzmuskel kontrahiert und entspannt sich beim Erwachsenen 60 bis
100 mal in der Minute. Bei jeder Anspannung (Kontraktion) wird das Herz
ca. 70ml Blut in die Arterien aus. }

{\mdseries
Das Herz pumpt in Ruhe ca. 4l Blut in der Minute. }

{\mdseries
Blutdruck: systolisch/diastolisch  (Anspannung/Entspannung)}

{\mdseries\color[rgb]{0.0,0.0,0.5019608}
%
%
Reizleitungssystem}

{\bfseries\color[rgb]{0.5019608,0.0,0.0}
@ Johannes: Reizleitungssystem -- wie ausf\"uhrlich? Pr\"ufungsfrage?}

\begin{center}
\tablehead{}\begin{supertabular}{|m{6.649cm}|m{6.651cm}|}
\hline
{\mdseries Eigenst\"andiges  Erregungs- und  Reizleitungssystem, streng
hierarchisch, aber jede  Station zur Erregungsbildung  f\"ahig }

{\mdseries bestehend aus:}

\begin{itemize}
\item {\mdseries Sinusknoten 
({\quotedblbase}Schrittmacher{\textquotedblleft}, feuert mit einer
Eigenfrequenz v. ca. 60/min. elektrische Impulse ab)}

\item {\mdseries AV-(AtrioVentricular)-Knoten: Leitung vom Vorhof in die
Kammern, filtert zu hohe Frequenzen}

\item {\mdseries \gZ{His-B\"undel}}

\item {\mdseries \gZ{Tawara-Schenkel (linker/rechter)}}

\item {\mdseries \gZ{Purkinje-Fasern
({\quotedblbase}Endstrecke{\textquotedblleft})}}

\end{itemize}
 &
{\centering   [Warning: Image ignored]
% Unhandled or unsupported graphics:
%\includegraphics[width=5.673cm,height=6.587cm]{AASdev20090228-img29.png}
 \par}

\\\hline
\end{supertabular}
\end{center}
{\mdseries
Der normale Schrittmacher ist der Sinusknoten, welcher die elektrischen
Impulse f\"ur die Vorh\"ofe und die Kammern erzeugt. Man spricht dann
vom \gEs{Sinusrhythmus}.}

{\mdseries
\gZ{Herzrhythmus}}

{\mdseries
\gZ{m\"ogliche Rhythmen:}}

{\mdseries
\gZ{Sinusrhythmus: ({\textless}60/min ${\rightarrow}$
Sinusbradykardie, {\textgreater}100/min ${\rightarrow}$
Sinustachy\-kardie): Sinus\-knoten ist Schrittmacher, vor jedem QRS ein
P, Abstand P-QRS konstant, Abstand QRS }\gZ{{}-- QRS
konstant, Abstand P -- P konstant}}

{\mdseries
\gZ{Arrhythmie: ({\textless}60/min ${\rightarrow}$
Brady arrhythmie, {\textgreater}100/min ${\rightarrow}$
Tachyarrhythmie): ev. wechselnde Schrittmacherzentren, nicht vor jedem
QRS ein P, Abst\"ande QRS -- QRS variabel}}

\subsection{Der Kreislauf }
{\mdseries
\gEs{Herz}\gEs{ +
}\gEs{Gef\"a{\ss}system}\gEs{ +
}\gEs{Blut}\gEs{ = Kreislauf}}

{\mdseries
Es wird der K\"orperkreislauf ({\quotedblbase}Gro{\ss}er
Kreislauf{\textquotedblleft},
{\quotedblbase}Hochdrucksystem{\textquotedblleft}) vom Lungenkreislauf
({\quotedblbase}Kleiner Kreislauf{\textquotedblleft},
{\quotedblbase}Niederdrucksystem{\textquotedblleft}) unterschieden. Das
Herz ist die f\"ur beide Systeme die gemeinsame Pumpe.}

{\mdseries\color{black}
Kreislaufsteuerung}

{\mdseries
Das Stammhirn  steuert den Kreislauf  (RR + HF )*, Blutdruckrezeptoren
in den  Gef\"a{\ss}w\"anden (z.B.Carotis)}

\begin{center}
\begin{tabular}{|m{5.328cm}m{7.977cm}}
\hline
\multicolumn{1}{|m{5.328cm}|}{\subparagraph{Kleiner (Lungen-)Kreislauf}
{\mdseries Die Hohlvenen m\"unden in den
\gEs{rechten Vorhof} des Herzens und liefern
sauerstoffarmes Blut aus dem K\"orper an. Von hier wird das Blut in die
\gEs{rechte Herzkammer} gepumpt, um dann in die
\gEs{Lungenarterie} zu gelangen. In der Lunge
ver\"asteln sich die Lungenarterien zu haard\"unnen Gef\"a{\ss}en, den
\gEs{Kapillaren}. Diese umspannen die
Lungenbl\"aschen es findet ein \gEs{Gasaustausch}
statt: Sauerstoff wird aufgenommen und Kohlendioxid abgegeben. Das
jetzt mit Sauerstoff angereicherte Blut wird \"uber die
\gEs{Lungenvenen} in den
\gEs{linken Vorhof} geleitet, der wiederum das Blut
in die \gEs{linke Herzkammer} pumpt. }

} &
\multicolumn{1}{m{7.977cm}|}{

\begin{center}
 [Warning: Image ignored] % Unhandled or unsupported graphics:
%\includegraphics[width=8.026cm,height=7.416cm]{AASdev20090228-img30.png}

\end{center}
\begin{center}
 [Warning: Image ignored] % Unhandled or unsupported graphics:
%\includegraphics[width=7.55cm,height=6.82cm]{AASdev20090228-img31.png}

\end{center}
{\mdseries Rechter Vorhof ${\rightarrow}$ rechte Herzkammer
${\rightarrow}$ Lungenarterie(n) ${\rightarrow}$ Lungenkappilaren
(CO2-Abgabe, O2-Aufnahme) ${\rightarrow}$ Lungenvenen ${\rightarrow}$
linker Vorhof }

{\mdseries ${\rightarrow}$ linke Herzkammer ${\rightarrow}$ Aorta
${\rightarrow}$ abzweigende Arterien ${\rightarrow}$ Kapillaren
(O2-Abgabe, CO2-Aufnahme)  ${\rightarrow}$ Venen ${\rightarrow}$
Hohlvene ${\rightarrow}$ Rechter Vorhof}

{\mdseries  ${\rightarrow}$ alles beginnt wieder von vorne}

}\\\hline
\subparagraph{Gro{\ss}er Kreislauf}
{\mdseries Das Blut wird aus der linken Herzkammer ausgesto{\ss}en und
gelangt \"uber die die Hauptschlagader
({\quotedblbase}\gEs{Aorta}{\textquotedblleft}) in
den K\"orperkreislauf. Aus der Aorta entspringen
\gEs{Arterien}
({\quotedblbase}Schlagadern{\textquotedblleft}), die sich immer weiter
verzweigen und an Durchmesser abnehmen. Sie versorgen den gesamten
K\"orper mit Blut. Der durch die Pumpleistung des Herzens erzeugte Puls
ist an oberfl\"achlich gelegenen Arterien tastbar (Radialis-Puls,
Carotis-Puls, ..). Daher r\"uhrt auch der Name
{\quotedblbase}Schlagader{\textquotedblleft}. }

{\mdseries Sie verzweigen sich, weiter an Durchmesser abnehmend,
\"ahnlich dem Ge\"ast eines Baumes, bishin zu
\gEs{Arteriolen} und haard\"unnen Gef\"a{\ss}en,
den Haargef\"a{\ss}en
({\quotedblbase}\gEs{Kapillaren}{\textquotedblleft}).
Hier findet wieder ein \gEs{Gasaustausch} statt,
diesmal allerdings spiegelbildlich zum Lungenkreislauf:
\gE{Sauerstoff wird an das Gewebe abgegeben}, im
Gegenzug wird \gE{Kohlendioxid abtransportiert}. }

{\mdseries Das nun sauerstoffarme Blut gelangt \"uber kleine
\gEs{Venolen} in das Venensystem und flie{\ss}t in
den \gEs{Venen} herzw\"arts. Am Weg zum Herzen
nehmen die Venen an Umfang immer mehr zu, bis sie letztendlich in die
gro{\ss}en Venen ({\quotedblbase}\gEs{obere und
untere Hohlvene}{\textquotedblleft}) m\"unden. }

 &
\\\hhline{-~}
\end{tabular}
\end{center}
{\bfseries\color[rgb]{0.5019608,0.0,0.0}
Klappen fehlen in der Beschreibung!}

{\mdseries\color[rgb]{0.0,0.0,0.5019608}
Wichtige Gef\"a{\ss}e}

\subparagraph{Gef\"a{\ss}e}
{\mdseries
\gT{Arterien}: vom Herzen weg, dicke Muskelschicht}

{\mdseries
\gT{Venen}: zum Herzen hin, Venenklappen}

{\mdseries
\gT{Kapillaren}: d\"unne
{\quotedblbase}Haargef\"a{\ss}e{\textquotedblleft} in denen der
Gasaustausch stattfindet, Verbindung zwischen Arterien und Venen}

{\mdseries
\gZ{Arterien sind durch eine dicke Muskelschicht
gekennzeichnet. Venen haben dagegen eine deutlich d\"unnere Wand als
die muskul\"osen Arterien}%
%
}

\begin{center}
\tablehead{}\begin{supertabular}{|m{2.728cm}|m{2.728cm}|m{2.728cm}|m{4.721cm}|}
\hline
\multicolumn{4}{|m{13.505cm}|}{\centering
{\bfseries\upshape\color[rgb]{0.0,0.0,0.5019608} \"Ubersicht wichtiger
Blutgef\"a{\ss}e}\par

}\\\hline
{\mdseries \gEs{Hauptschlagader}}

 &
{\mdseries \gT{Aorta}}

 &
{\mdseries Arterie}

 &
{\mdseries Geht von der linken Herzkammer ab, von dort zweigen alle
anderen K\"orperarterien ab}

\\\hline
{\mdseries \gEs{Halsschlagader}}

 &
{\mdseries \gT{Karotis}}

 &
{\mdseries Arterie}

 &
{\mdseries 1 Paar (2 Stk.), versorgen den Sch\"adel, Puls tastbar
seitlich vom Kehlkopf}

\\\hline
{\mdseries \gEs{Herzkranzgef\"a{\ss}e}}

 &
{\mdseries \gT{Koronararterien}}

 &
{\mdseries Arterie}

 &
{\mdseries Versorgen das Herz mit Blut, gehen aus der Aorta ab}

\\\hline
{\mdseries \gEs{Lungenarterie}}

 &
{\mdseries \gT{Pulmonalarterie}}

 &
{\mdseries Arterie -- O\gSub{2}{}-arm!}

 &
{\mdseries Geht von der rechten Herzkammer weg in die Lunge zu den
Lungenkapillaren}

\\\hline
{\mdseries \gEs{Oberarmarterie}}

 &
{\mdseries Brachialarterie}

{\mdseries A. brachialis}

 &
{\mdseries Arterie}

 &
{\mdseries Taststelle zum Puls-f\"uhlen beim Baby, Blutdruckmessung}

\\\hline
{\mdseries \gEs{Radialarterie}}

 &
{\mdseries A. radialis}

 &
{\mdseries Arterie}

 &
{\mdseries Verl\"auft am Unterarm speichenseitig (dort wo der Daumen
ist). Puls f\"uhlbar.}

\\\hline
{\mdseries \gEs{Obere Hohlvene}}

 &
{\mdseries \gZ{Vena cava sup.}}

 &
{\mdseries Vene}

 &
{\mdseries Obere Vene die in den rechten Vorhof m\"undet}

\\\hline
{\mdseries \gEs{Untere Hohlvene}}

 &
{\mdseries \gT{Vena cava}\gZ{ inf.}}

 &
{\mdseries Vene}

 &
{\mdseries Untere Vene die in den rechten Vorhof m\"undet. }

\\\hline
{\mdseries \gEs{Lungenvene}}

 &
{\mdseries \gZ{V. pulmonalis}}

 &
{\mdseries Vene -- O\gSub{2}{}-reich! }

 &
{\mdseries Kommt von den Lungenkapillaren und m\"undet in den linken
Vorhof}

\\\hline
{\mdseries \gEs{Pfortader}}

 &
{\mdseries \gZ{Vena portae}}

 &
{\mdseries Vene}

 &
{\mdseries F\"uhrt O2-armes, aber n\"ahrstoffreiches Blut von Teilen des
Darms zur Leber (zur Entgiftung)}

\\\hline
{\mdseries \gEs{Haargef\"a{\ss}e}}

 &
{\mdseries \gT{Kapillare}}

 &
{\mdseries Kapillare}

 &
{\mdseries D\"unne, dicht verzwweigte Gef\"a{\ss}e, Verbindung zwischen
Arterien und Venen. Hierfindet der Gas- und N\"ahrstoffaustausch statt.
Es gibt K\"orperkapillaren (K\"orperkreislauf) und Lungenkapillaren
(Lungenkreislauf)}

\\\hline
\end{supertabular}
\end{center}
\paragraph{Blutdruck }
{\mdseries
Abk\"urzung: RR. Gemessen immer in Herzh\"ohe!}

{\mdseries
systolisch:  Maximaldruck in der Arterie am H\"ohepunkt der
Kammerkontraktion}

{\mdseries
diastolisch: Druck in der Arterie w\"ahrend die Kammer erschlafft ist,
beeinflu{\ss}t durch  Gef\"a{\ss}wandspannung}

{\mdseries
Normbereich : zwischen 90/60 mmHg und 140/90 mmHg  }

{\mdseries
Hypotonie:  {\textless} 90/60 mmHg}

{\mdseries
Hypertonie: {\textgreater} 140/90mmHg  }

{\mdseries
Achtung: Eine Hypertonie ist zwar eine Erkrankung, aber nicht immer im
Einsatz von Bedeutung. Viele Patienten sind an ihren erh\"ohten
Blutdruck gewohnt und haben keine Beschwerden. Es darf nie nur nach dem
Wert gehandelt werden, es muss auch immer bedacht werden: }

\begin{itemize}
\item {\mdseries
Zustand des Patienten.}

\item {\mdseries
Der sonst \"ubliche Blutdruck des Patienten (erfragen,
Blutdruckprotokoll).}

\item {\mdseries
Der Verlauf des Blutdruckes w\"ahrend der Behandlung (Traumapatienten
!!!).}

\item {\mdseries
Der Erregungszustand des Patienten.}

\end{itemize}
{\mdseries\color{black}
Beispiele}

Ein 65-j\"ahriger, m\"annlicher Patient klagt \"uber Schwindelgef\"uhl,
sie messen einen Blutdruck von 100/70. Im Blutdruckprotokoll finden sie
f\"ur die vergangenen vier Tage folgende Eintr\"age: 170/110, 180/115,
175/113, 183/112, ... .

Wie beurteilen Sie den Blutdruck? 

{\itshape
Der Blutdruck ist f\"ur den Patienten zu niedrig, der K\"orper ist auf
eine so pl\"otzliche Umstellung nicht vorbereitet gewesen. Im
Anamnesegespr\"ach erfahren sie dass er heute das erste Mal einen
{\quotedblbase}Nitro-Spray{\textquotedblleft} verwendet hat. Sie
vermerken dies am Einsatzprotokoll. Im Krankenhaus erfahren Sie dass
wahrscheinlich eine \"Uberdosierung dieses Medikaments die Ursache
dieses Blutdruckabfalls war. ${\rightarrow}$ Obwohl der Patient laut
Lehrbuch einen {\quotedblbase}sch\"onen{\textquotedblleft} Blutdruck
hatte, war dieser dennoch zu niedrig!}

Ein 56-j\"ahriger Patient klagt \"uber starke kolikartige Schmerzen im
linken Unterbauch. Sie messen einen Blutdruck von 155/95. Der Patient
kennt seinen {\quotedblbase}gew\"ohnlichen{\textquotedblleft} Blutdruck
nicht.  

Wie beurteilen Sie den Blutdruck? 

{\itshape
Der Blutdruck ist laut Lehrbuch erh\"oht, aber dies ist erkl\"arbar die
starken Schmerzen und den damit verbundenen Erregungszustand. Die
Beschwerden stehen wahrscheinlich in keinem urs\"achlichen Zusammenhang
mit den Schmerzen. }

Eine 50-j\"ahrige Patientin hat Nasenbluten. Sie messen einen Blutdruck
von 220/120. Auf Nachfrage gibt sie an etwas Kopfsschmerzen zu haben,
und etwas schwindlig sei sie auch. 

Wie beurteilen Sie den Blutdruck? 

{\itshape
Der Blutdruck ist sehr stark erh\"oht, ausserdem zeigt die Patientin
Zeichen einer hypertensiven Krise (Nasenbluten, Kopfweh, Schwindel).
Hier ist der erh\"ohte Blutdruck Ursache f\"ur die Beschwerden und muss
mittels Medikamente unbedingt gesenkt werden. }

{\mdseries\color[rgb]{0.0,0.0,0.5019608}
Das \index{Blut}Blut}

{\mdseries
ca. 7-8\% des K\"orpergewichts = \gEs{5-6l}}

{\mdseries
besteht aus}

\begin{itemize}
\item {\mdseries
festen Bestandteilen (Zellen, Blutk\"orperchen)}

\item {\mdseries
\index{Plasma}Plasma}

\end{itemize}
  [Warning: Image ignored] % Unhandled or unsupported graphics:
%\includegraphics[width=13.664cm,height=5.432cm]{AASdev20090228-img32.png}
 

{\mdseries\color{black}
Blut - Aufgaben}

\begin{itemize}
\item {\mdseries
Sauerstofftransport zu den K\"orperzellen}

\item {\mdseries
CO2-Abtransport}

\item {\mdseries
Transport von N\"ahrstoffen und Abfallprodukten}

\item {\mdseries
W\"armeverteilung}

\item {\mdseries
\gZ{pH-Pufferungssystem (S\"aure/Basen-Gleichgewicht)
(pH-Wert mu{\ss} konstant zwischen 7,35 und 7,45 liegen!)}}

\end{itemize}
{\mdseries\color{black}
Zellul\"are Bestandteile}

\begin{itemize}
\item {\color{red}
Rote Blutk\"orperchen (Erythrozyten)}

\item {\mdseries
Wei{\ss}e Blutk\"orperchen (Leukozyten)}

\item {\mdseries
\gZ{\textcolor{black}{Blutpl\"attchen (Thrombozyten)}}}

\end{itemize}
\begin{multicols}{2}
{\mdseries\color{black}
\index{Erythrozyten}Erythrozyten }

\begin{itemize}
\item {\color{red}
rote Blutk\"orperchen}

\item {\mdseries
O2 - Transport}

\item {\mdseries
tragen Blutgruppenmerkmale (Blutgruppe)}

\item {\mdseries
seeeeeehr viele, die h\"aufigsten Blutk\"orperchen
\gZ{(4.000.000/mm3 )}}

\end{itemize}
{\mdseries\color{black}
\index{Leukozyten}Leukozyten }

\begin{itemize}
\item {\mdseries
wei{\ss}e Blutk\"orperchen }

\item {\mdseries
K\"orperabwehr, Entz\"undungszellen}

\item {\mdseries
\gZ{(4.000 -- 8.000/mm3 )}}

\end{itemize}
{\mdseries\color{black}
\index{Thrombozyten}Thrombozyten }

\begin{itemize}
\item {\mdseries
\index{Blutpl\"attchen}Blutpl\"attchen}

\item {\mdseries
wichtige Rolle beim Wundverschlu{\ss}}

\item {\mdseries
\gZ{(150.000 -- 300.000/mm3 )}}

\item {\mdseries
Blutgerinnung}

\end{itemize}
{\mdseries\color{black}
Das \index{Plasma}Plasma}

\begin{itemize}
\item {\mdseries
10\% Proteine, Enzyme, Hormone, Elektrolyte, Glucose, Fette}

\item {\mdseries
90\% H2O}

\item {\mdseries
Aufgaben:}

\end{itemize}
\begin{itemize}
\item {\mdseries
N\"ahrstofftransport}

\item {\mdseries
W\"armeverteilung}

\end{itemize}
\begin{itemize}
\item {\mdseries
enth\"alt Proteine und Enzyme f\"ur  Blutgerinnung (Serum )}

\item {\mdseries
Pufferung}

\end{itemize}
\end{multicols}
\subsection{Das Nervensystem}
{\mdseries\color{black}
Unterteilung}

{\mdseries
Man kann das Nervensystem nach seinem Aufbau und nach seiner Funktion
unterteilen.}

\begin{center}
\begin{minipage}{4.001cm}
{\mdseries
\gNw{NEU }\gNw{03.12.08}}

\end{minipage}
\end{center}
{\mdseries
Vom Blickwinkel des Aufbaus wird das Nervensystem grob unterteilt in ein
}

\begin{itemize}
\item {\mdseries
\gEs{Zentrales Nervensystem}
(\gT{ZNS}) und ein }

\item {\mdseries
\gEs{Peripheres Nervensystem}.}

\end{itemize}
{\mdseries
Das ZNS ist im wesentlichen f\"ur Steuer-, Regel- und Denkfunktionen
zust\"andig sowie Teilweise f\"ur die Weiterleitung von Impulsen und
Informationen. Das periphere Nervensystem ist fast ausschliesslich nur
f\"ur die Weiterleitung von Informationen und Impulsen zu den
Endorganen (oder umgekehrt, von den Endorganen zum ZNS) zust\"andig). }

{\mdseries
Betrachtet man die Funktion kann man das periphere Nervensystem in ein }

\begin{itemize}
\item {\mdseries
\gEs{Motorisches} (willk\"urliche Bewegung),}

\item {\mdseries
\gEs{Sensorisches} und ein}

\item {\mdseries
\gEs{Vegetatives} (autonomes) Nervensystem}

\end{itemize}
{\mdseries
unterteilen. }

{\mdseries\color[rgb]{0.0,0.0,0.5019608}
Das Zentralnervensystem}

\begin{center}
\tablehead{}\begin{supertabular}{|m{6.649cm}|m{6.651cm}|}
\hline
Das Zentralnervensystem (ZNS) besteht aus:

\begin{itemize}
\item Gehirn

\end{itemize}
\begin{itemize}
\item {\mdseries Gro{\ss}hirn}

\item {\mdseries Kleinhirn}

\item {\mdseries Hirnstamm}

\end{itemize}
\begin{itemize}
\item R\"uckenmark

\end{itemize}
Aufgaben:

\begin{itemize}
\item Aufnahme  und Verarbeitung  von \"au{\ss}eren und internen  Reizen
(Schmerz, F\"uhlen, Sehen, H\"oren, {\dots})

\item \gNw{Bewusstsein und
{\quotedblbase}Bewusstmachung{\textquotedblleft}}

\item Abgabe  von zweckentsprechenden Impulsen  (Steuerung, Reaktion
etc.)

\end{itemize}
 &
{\centering [Warning: Draw object ignored]\par}

\\\hline
\end{supertabular}
\end{center}
{\mdseries\color{black}
%
%
ZNS - Aufgaben}

\begin{center}
\begin{minipage}{4.001cm}
{\mdseries
ZNS: ein bissl umgestellt:}

{\mdseries
Bild}

{\mdseries
ZNS Aufgaben}

{\mdseries
Anatomie}

{\mdseries
{}- Querschnitt Hirnsch\"adel}

{\mdseries
{}- Aufbau ZNS mit den versch Hirnen und Liquor}

{\mdseries
{}- Blut-Hirn-Schranke}

\end{minipage}
\end{center}
{\mdseries
\index{Gro{\ss}hirn}\gT{Gro{\ss}hirn}}

\begin{itemize}
\item {\mdseries
bewu{\ss}tes Denken}

\item {\mdseries
oberste Instanz}

\end{itemize}
{\mdseries
\index{Zwischenhirn}\gT{Zwischenhirn}}

\begin{itemize}
\item {\mdseries
Schaltinstanz und Filter R\"uckenmark - Gro{\ss}hirn}

\end{itemize}
{\mdseries
\index{Kleinhirn}\gT{Kleinhirn}}

\begin{itemize}
\item {\mdseries
Bewegungskoordination ({\quotedblbase}Makros{\textquotedblleft})}

\end{itemize}
{\mdseries
\index{Hirnstamm}\gT{Hirnstamm} und \index{verl\"angertes
R\"uckenmark}\gT{verl\"angertes R\"uckenmark}}

\begin{itemize}
\item {\mdseries
Atem-, Kreislaufzentrum}

\end{itemize}
{\mdseries
\index{R\"uckenmark}\gT{R\"uckenmark}}

\begin{itemize}
\item {\mdseries
Weiterleitung von Impulsen}

\item {\mdseries
Vom R\"uckenmark gehen die peripheren Nerven aus}

\item {\mdseries
(Muskeleigen-)Reflexe}

\end{itemize}
\subparagraph{Anatomie}
\begin{center}
\tablehead{}\begin{supertabular}{|m{4.346cm}|m{8.988cm}|}
\hline
{\mdseries\color{black} Querschnitt Gehirnsch\"adel}

\begin{itemize}
\item Kopfschwarte (Galea)

\item Knochen

\item \index{Epiduralraum}\gEs{Epiduralraum}

\item \index{Dura mater}\gT{Dura mater} (\index{Harte
Hirnhaut}harte Hirnhaut)

\item \index{Subduralraum}\gEs{Subduralraum}

\item \index{Arachnoidea}\gT{Arachnoidea}
(Spinngewebshaut)

\item
\index{Subarachnoidalraum}\gEs{Subarachnoidalraum}

\item eingebettete Gef\"a{\ss}e

\item Pia mater (weiche Hirnhaut)

\end{itemize}
 &
  [Warning: Image ignored] % Unhandled or unsupported graphics:
%\includegraphics[width=9.035cm,height=6.869cm]{AASdev20090228-img33.png}
 

\\\hline
\end{supertabular}
\end{center}
\begin{itemize}
\item \end{itemize}
{\mdseries\color{black}
Aufbau ZNS}

\begin{itemize}
\item {\mdseries
Gro{\ss}hirn (Telencephalon)}

\end{itemize}
\begin{itemize}
\item {\mdseries
\index{Graue Substanz}Graue Substanz (Rinde)}

\item {\mdseries
\index{Wei{\ss}e Substanz}Wei{\ss}e Substanz (Mark)}

\item {\mdseries
Balken (Verbindung li/re)}

\end{itemize}
\begin{itemize}
\item {\mdseries
Zwischenhirn  (Diencephalon)}

\end{itemize}
\begin{itemize}
\item {\mdseries
Thalamus, Hypothalamus (Hormone, Temperaturregulation etc.)}

\item {\mdseries
Hypophyse (Hormone)}

\end{itemize}
\begin{itemize}
\item {\mdseries
Hirnstamm (Truncus encephali)}

\item {\mdseries
Mittelhirn (Mesencephalon)}

\item {\mdseries
Br\"ucke (Pons)}

\item {\mdseries
verl\"angertes Mark (Medulla oblongata)}

\item {\mdseries
Ventrikelsystem}

\end{itemize}
\begin{itemize}
\item {\mdseries
Liquorgef\"ullte Hohlr\"aume (Ventrikel)}

\item {\mdseries
haben nichts mit den Vnetrikel im Herz zu tun!}

\end{itemize}
{\mdseries
Liquor = Gehirnwasser, in Ventrikeln gebildet}

{\mdseries
Das Hirn und das R\"uckenmark schwimmt im Liquor, dadurch ist es ein
wenig gesch\"utzt. Tritt Liquor in die Umwelt aus (Nase, Ohr, ...)  ist
das ein Zeichen f\"ur eine schwerwiegende Verletzung. }

{\mdseries
Tritt Liquor in die Umwelt aus ist das ein Zeichen f\"ur eine
schwerwiegende Verletzung. }

{\mdseries\color{black}
Blut-Hirn-Schranke }

{\mdseries
Die Blutgef\"a{\ss}wand ist im Hirn besonders dicht. Sie
\gEs{filtert} daher besonders gut und nur f\"ur
bestimmte Molek\"ule (O\gSub{2},CO\gSub{2},
Wasser,...) leicht passierbar. Sie bietet dem Hirn Schutz vor
sch\"adlichen und giftigen Substanzen. }

{\mdseries
Die Blut-Hirn-Schranke bietet dem Hirn besonderen Schutz vor
sch\"adlichen Substanzen im Blut. }

{\mdseries\color[rgb]{0.0,0.0,0.5019608}
PNS -- Peripheres Nervensystem}

\begin{center}
\tablehead{}\begin{supertabular}{|m{6.649cm}|m{6.651cm}|}
\hline
\begin{itemize}
\item Kommunikation \textit{ZNS} ${\leftarrow}$${\rightarrow}$
\textit{periphere  Organe}

\item Spinalnerven (RM, segmental)

\item Hirnnerven (12 Nervi craniales)

\item elektrische Impulse, chemische Signale (Synapsen)

\end{itemize}
Schnell: 30-70 m/sec

Arten von Nerven

\begin{itemize}
\item motorisch

\item sensorisch

\item vegetativ

\end{itemize}
\begin{itemize}
\item sympathisch

\item parasympathisch

\end{itemize}
 &
{\centering   [Warning: Image ignored]
% Unhandled or unsupported graphics:
%\includegraphics[width=5.28cm,height=4.405cm]{AASdev20090228-img34.png}
 \par}

{\mdseries\upshape R\"uckenmark mit abgehenden peripheren Nerven.
Beachte dass die Nerven geordnet nach R\"uckenmarkssegmenten abgehen. }

\\\hline
\end{supertabular}
\end{center}
\paragraph[Vegetatives (autonomes) \ Nervensystem]{Vegetatives
(autonomes)  Nervensystem}
{\mdseries\color{black}
Sympathikus  und Parasymphatikus  sind  Gegenspieler}

{\mdseries
\textbf{Sympathikus: {\quotedblbase}}Fight or flight{\textquotedblleft}
- {\quotedblbase}Kampf und Flucht{\textquotedblleft}}

{\mdseries
\textbf{Parasympathikus}: {\quotedblbase}Rest and
digest{\textquotedblleft} - {\quotedblbase}Rast und
Verdauung{\textquotedblleft}}

{\mdseries
z.B.: RR + HF}

{\mdseries
\textbf{Sympathikus}:  HF ${\uparrow}$, RR ${\uparrow}$}

{\mdseries
\textbf{Parasymphatikus}: HF ${\downarrow}$, RR ${\downarrow}$}

  [Warning: Image ignored] % Unhandled or unsupported graphics:
%\includegraphics[width=6.864cm,height=4.481cm]{AASdev20090228-img35.png}
   [Warning: Image ignored] % Unhandled or unsupported graphics:
%\includegraphics[width=6.73cm,height=4.321cm]{AASdev20090228-img36.png}
 

  [Warning: Image ignored] % Unhandled or unsupported graphics:
%\includegraphics[width=13.29cm,height=16.484cm]{AASdev20090228-img37.png}
 

{\mdseries\upshape
Nur zur Illustration: Das vegetative Nervensystem hat Einflu{\ss} auf
viele Organe.}

{\mdseries\upshape
Diese Bild NICHT LERNEN. Es dient nur der Verdeutlichung wie wichtig und
einflussreich das vegetative Nervensystem im K\"orper ist. }

\subsection{Die Haut}
{\mdseries
gr\"o{\ss}tes Sinnesorgan}

{\mdseries
ca.1/6 des K\"orpergewichtes}

{\mdseries
ca. 1,6 m{\texttwosuperior} Oberfl\"ache}

{\mdseries
Grenzfl\"ache zur Au{\ss}enwelt}

{\mdseries
Aufgaben}

\begin{itemize}
\item {\mdseries
Schutz vor \"au{\ss}eren Einfl\"ussen (mechanisch, Erreger, UV-Strahlung
etc.)}

\item {\mdseries
Schutz vor Wasserverlust/Austrocknung}

\item {\mdseries
Schutz vor W\"armeverlust /  Temperaturregulation}

\item {\mdseries
Sekretionsorgan (Schwei{\ss}, Talg)}

\item {\mdseries
Sinnesorgan (Tastsinn, Temperatursinn)}

\end{itemize}
{\mdseries\color{black}
Die Haut ist in \gEs{Schichten} aufgebaut%
%
  [Warning: Image ignored] % Unhandled or unsupported graphics:
%\includegraphics[width=12.952cm,height=9.712cm]{AASdev20090228-img38.gif}
 }

{\bfseries\color[rgb]{0.5019608,0.0,0.0}
Haut: Bild Querschnitt: muss vereinfacht werden, lat. Ausdr\"ucke weg
bzw. grau}

\begin{itemize}
\item {\mdseries
\index{Oberhaut}\textbf{Oberhaut}\textbf{
(}\gZ{Epidermis})}

\end{itemize}
\begin{itemize}
\item {\mdseries
Hornschicht, geschichtetes Plattenepithel}

\end{itemize}
\begin{itemize}
\item {\mdseries
\index{Lederhaut}\textbf{Lederhaut} (\gZ{Corium,
Dermis})}

\end{itemize}
\begin{itemize}
\item {\mdseries
Hautanhangsgebilde}

\end{itemize}
\begin{itemize}
\item {\mdseries
Schwei{\ss}dr\"usen}

\item {\mdseries
Talgdr\"usen}

\item {\mdseries
Haarfollikel, Nagelmatrix}

\end{itemize}
\begin{itemize}
\item {\mdseries
Tastk\"orperchen, freie Nervenendigungen (Schmerz)}

\item {\mdseries
Kollagenfasern (St\"utznetzwerk)}

\item {\mdseries
Arteriolen, Venolen (Versorgung der Haut, Temperaturregulation)}

\end{itemize}
\begin{itemize}
\item {\mdseries
\index{Unterhaut}\textbf{Unterhaut} (Subcutis)}

\end{itemize}
\begin{itemize}
\item {\mdseries
subkutanes Depotfett, W\"armeisolierung}

\end{itemize}
{\mdseries\color{black}
Die Hautist wichtig f\"ur die
\gEs{Thermoregulation}}

{\mdseries
Neben der Lunge dient vor allen Dingen  die Haut als wichtigstes Organ
zur \gEs{Thermoregulation}: Die Blutgef\"a{\ss}e 
der  Haut sind stark ineinander verwickelt und k\"onnen sich bei Bedarf
stark erweitern. Die stark durchblutete  Haut strahlt somit W\"arme ab.
Weiters wird durch die Verdunstungsk\"alte  des Schwei{\ss}es 
ebenfalls W\"arme abgegeben. }

\subsection{Atmung}
{\mdseries\color[rgb]{0.0,0.0,0.5019608}
Die \gEs{Atemwege} sind gegliedert}

\begin{center}
\tablehead{}\begin{supertabular}{|m{6.649cm}|m{6.651cm}|}
\hline
Obere Atemwege:

\begin{itemize}
\item {\mdseries Nasenrachenraum mit Nasennebenh\"ohlen }

\item {\mdseries Mundrachenraum}

\item {\mdseries Rachen (gemeinsamer Weg von Nahrung und Atem!)}

\item {\mdseries \index{Kehlkopf}Kehlkopf (Larynx)}

\end{itemize}
Untere Atemwege:

\begin{itemize}
\item {\mdseries \index{Luftr\"ohre}Luftr\"ohre
(\index{Trachea}Trachea)}

\item {\mdseries \index{Hauptbronchus}Hauptbronchus (li./re.)}

\item {\mdseries \index{Bronchioli}Bronchioli}

\item {\mdseries \index{Lungenbl\"aschen}Lungenbl\"aschen
(\index{Alveolen}\gT{Alveolen})}

\end{itemize}
 &
{\centering   [Warning: Image ignored]
% Unhandled or unsupported graphics:
%\includegraphics[width=5.633cm,height=7.116cm]{AASdev20090228-img39.png}
 \par}

\\\hline
\end{supertabular}
\end{center}
\begin{multicols}{2}
{\mdseries\color{black}
Obere Atemwege}

{\mdseries
Der Weg der Einatemluft beginnt in der Nase bzw. im Mund. Die Nase ist
mit Schleimhaut ausgekleidet, kn\"ocherne Vorspr\"unge sorgen f\"ur
eine Vorw\"armung der Luft. In der Schleimhaut be\-finden sich die
Sinneszellen f\"ur den Geruchssinn. }

{\mdseries
In den Sch\"adelknochen befinden sich luft\-ge\-f\"ullte Hohlr\"aume,
die so\-genannten \gT{Neben\-h\"ohlen} (Stirn-,
Kiefer\-h\"ohlen, Siebbeinzellen). Diese stehen in Verbindung mit dem
Nasenrachenraum. Entz\"undet sich die Schleimhaut spricht man von der
\gE{Neben\-h\"ohlen\-ent\-z\"undung}, welche oft
langwierig und schmerzhaft sein kann. }

{\mdseries
Die Luft gelangt nun in den \gT{Mund-Rachen\-raum} und
den \gT{Rachen}. \gE{Hier kreuzt der Weg
der Nahrung mit dem Weg der Atemluft!} Die Nahrung gelangt vom Mundraum
in die hinten gelegene Speiser\"ohre, die Luft m\"ochte vom
Nasen\-rachen\-raum in die vorne gelegene Luftr\"ohre. }

{\mdseries
Um ein Verschlucken von Nahrung in die Lunge zu verhindern
verschlie{\ss}t der am Kehlkopf gelegene \gT{Kehldeckel}
w\"ahrend des Schluckaktes die Luftr\"ohre. Dieser Reflex ist
lebenswichtig! }

{\mdseries
Im \gE{Kehlkopf}  befinden sich die
\gT{Stimmb\"ander}, welche f\"ur die Lautbildung beim
Sprechen wichtig sind.}

{\mdseries\color{black}
Untere Atemwege}

{\mdseries
In der \gT{Luftr\"ohre} (\gT{Trachea})
wird die vom Rachen und Kehlkopf kommende Luft zur Lunge
weitergeleitet. Die Luftr\"ohre teilt sich ein einen
\gE{rechten} und einen \gE{linken}
\gT{Hauptbronchus}, welche zu den jeweiligen
Lungenfl\"ugeln ziehen. Jeder Hauptbronchus verzweigt sich immer weiter
bis schlu{\ss}endlich sehr d\"unne \gT{Bronchioli} in
die \gT{Lungenbl\"aschen} (\gT{Alveolen})
m\"unden. Hier findet der Gasaustausch statt. }

{\mdseries
Die Abschnitte der Atemwege in denen kein Gasaustausch stattfindet und
der nur dem Transport dient bezeichnet man als
{\quotedblbase}\index{Totraum}\gT{Totraum}{\textquotedblleft}.}

\end{multicols}
{\bfseries\color[rgb]{0.5019608,0.0,0.0}
weiter}

\paragraph[Die Lunge ist f\"ur die Atmung zust\"andig]{Die
\gEs{Lunge} ist f\"ur die Atmung zust\"andig}
{\mdseries
Die Lunge hat 2 Fl\"ugel, welche sich jeweils in 3 (rechts) bzw. 2
(links) Lappen gliedern: }

\begin{center}
\tablehead{}\begin{supertabular}{|m{6.649cm}|m{6.651cm}|}
\hline
\multicolumn{2}{|m{13.5cm}|}{\centering
{\bfseries\upshape\color[rgb]{0.0,0.0,0.5019608} 1 Lunge}\par

}\\\hline
\centering {\bfseries\upshape\color[rgb]{0.0,0.0,0.5019608} Rechter
Fl\"ugel }\par

 &
\centering {\bfseries\upshape\color[rgb]{0.0,0.0,0.5019608} Linker
Fl\"ugel}\par

\\\hline
\begin{itemize}
\item {\mdseries 3 Lappen}

\end{itemize}
 &
\begin{itemize}
\item {\mdseries 2 Lappen}

\end{itemize}
\\\hline
\begin{itemize}
\item {\mdseries Oberlappen}

\item {\mdseries Mittellappen}

\item {\mdseries Unterlappen}

\end{itemize}
 &
\begin{itemize}
\item {\mdseries Oberlappen}

\item {\mdseries Unterlappen}

\end{itemize}
\\\hline
\end{supertabular}
\end{center}
{\mdseries
Die Oberfl\"ache ist mit dem Brustfell, der lungenseitigen
\gT{Pleura} \"uberzogen. In einem Lungefl\"ugel
verzweigt sich der Hauptbronchus in immer weitere kleinere Bronchien
bis schlie{\ss}lich die Bronchioli in die
\gT{Lungenbl\"aschen}, die \gT{Alveolen},
enden. }

\subparagraph[Die Alveolen sind die Endst\"ucke des Atemwegs und sind
f\"ur den Gasaustausch zust\"andig]{Die
\gEs{Alveolen} sind die Endst\"ucke des Atemwegs
und sind f\"ur den Gasaustausch zust\"andig}
{\mdseries
In den \gT{Lungenbl\"aschen}
(\gT{Alveolen}) hier findet der
\index{Gasaustausch}\gE{Gasaustausch} zwischen der
Einatemluft und dem Blut statt. Sauerstoff gelangt dabei aus den
Alveolen in das Blut und Kohlendiovid wird im Gegenzug vom Blut an die
Alveolen abgegeben. }

{\mdseries
Dabei besteht zwischen Luft und Blut kein direkter Kontakt, die Gase
m\"ussen Hindernisse wie die Alveolenwand und die Blutgef\"a{\ss}wand
\"uberwinden. \gE{Vergr\"o{\ss}ert sich dieser Weg
aufgrund einer Erkrankung ist die Sauerstoffaufnahme vermindert!} }

{\mdseries
Sauerstoff (O\gSub{2}):   Alveolen ${\rightarrow}$ Blut}

{\mdseries
Kohlendioxid (CO\gSub{2}):  Blut ${\rightarrow}$ Alveolen}

  [Warning: Image ignored] % Unhandled or unsupported graphics:
%\includegraphics[width=13.597cm,height=7.554cm]{AASdev20090228-img40.png}
 

\paragraph{Die Pleura}
\begin{center}
\tablehead{}\begin{supertabular}{|m{6.649cm}|m{6.651cm}|}
\hline
\begin{itemize}
\item Thoraxseitige \gT{Pleura}%
%
 (Rippenfell, fest mit  Brustwand und Zwerchfell verwachsen, kleidet
damit die Brusth\"ohle aus)

\item Lungenseitige \gT{Pleura }(Lungenfell, \"uberzieht
die Lungenoberfl\"ache)

\item Aufgrund eines Unterdrucks haften die beiden Pleuren aneinander,
sind aber gegeneinander verschieblich (wie 2 Glasplatten)

\item Geringe Menge Pleurafl\"ussigkeit ${\rightarrow}$ Adh\"asion der
beiden Bl\"atter!

\item Normalerweise keine Verwachsungen der beiden Bl\"atter

\end{itemize}
 &
{\centering   [Warning: Image ignored]
% Unhandled or unsupported graphics:
%\includegraphics[width=5.422cm,height=5.249cm]{AASdev20090228-img41.gif}
 \par}

\\\hline
\end{supertabular}
\end{center}
{\mdseries\color{black}
Welche \gE{Funktion} hat aber nun die Pleura?}

{\mdseries
Da die eine Seite mit dem Brustkorb, die andere Seite mit der Lunge
verwachsen ist und der Raum zwischen den beiden Seiten luftdicht
abgechlossen ist haften die beiden Seiten aneinander, aber sind
gegeneinander verschieblich (Es gibt ja keine Verwachsungen der beiden
eiten und die Pleurafl\"ussigkeit fungiert sogar ein bisschen als
{\quotedblbase}Schmiermittel{\textquotedblleft}). }

{\mdseries
Wenn sich nun der Brustkorb oder das Zwerchfell bei der Atmung bewegt
bewegt sich nat\"urlich das Rippenfell mit (es ist ja mit dem Brustkorb
und dem Zwerchfell verwachsen). Da der Pleurazwischenraum luftdicht ist
muss sich Zwangsl\"aufit aber auch das Lungenfell (an dem wiederum die
Lunge angewachsen ist) mitbewegen. Und wenn sich das Lungenfell
mitbewegt wird automatisch auch die Lunge wie eine Ziehharmonika
{\quotedblbase}Aufgezogen{\textquotedblleft} und Luft
{\quotedblbase}angesaugt{\textquotedblleft}. }

{\mdseries
Frage: Was w\"urde passieren wenn der Zwischenraum zwischen den
Pleurabl\"attern aufgrund einer Verletung (Z.B. ein Loch nach einer
Stichwunde in die Brust) nicht mehr luftdicht ist? Antwort: s.
\footnote{Antwort: Die Lunge w\"urde in sich zusammenfallen, da sie ja
nicht mehr aufgespannt wird. Siehe
{\quotedblbase}Pneumothorax{\textquotedblleft}.} }

{\mdseries\color[rgb]{0.0,0.0,0.5019608}
Atemmechanik%
%
}

  [Warning: Image ignored] % Unhandled or unsupported graphics:
%\includegraphics[width=13.648cm,height=7.52cm]{AASdev20090228-img42.png}
 

{\mdseries
Der wichtigste Atemmuskel ist das Zwerchfell! }

{\bfseries\color[rgb]{0.5019608,0.0,0.0}
Atemmechanik: Atemhilfsmuskulatur fehtl!}

{\mdseries\color[rgb]{0.0,0.0,0.5019608}
Luft}

\begin{center}
\tablehead{}\begin{supertabular}{|m{4.367cm}|m{4.367cm}|m{4.367cm}|}
\hline
\centering {\bfseries\upshape\color[rgb]{0.0,0.0,0.5019608}
Raumluft}\par

 &
 &
\centering {\bfseries\upshape\color[rgb]{0.0,0.0,0.5019608}
Ausatemluft}\par

\\\hline
{\mdseries \gEs{21 \%
O}\gEs{\gSub{2}}}

 &
\centering Sauerstoff\par

 &
{\mdseries \gEs{16 \%
O}\gEs{\gSub{2}}}

\\\hline
{\mdseries 79 \% N}

 &
\centering Stickstoff\par

 &
{\mdseries 79 \% N}

\\\hline
{\mdseries \gEs{0,03 \%
CO}\gEs{\gSub{2}}}

 &
\centering Kohlendioxid\par

 &
{\mdseries \gEs{4 \%
CO}\gEs{\gSub{2}}}

\\\hline
{\mdseries 1 \% Edelgase}

 &
 &
{\mdseries 1 \% Edelgase}

\\\hline
\end{supertabular}
\end{center}
{\mdseries\color[rgb]{0.0,0.0,0.5019608}
Atmungvolumina}



\begin{center}
 [Warning: Image ignored] % Unhandled or unsupported graphics:
%\includegraphics[width=10.655cm,height=6.386cm]{AASdev20090228-img43.png}

\end{center}
{\mdseries
+ 150ml in den oberen Luftwegen (Trachea, Rachenraum) ${\rightarrow}$
Dieser Teil nimmt nicht am Gasaustausch ${\rightarrow}$
\index{Totraum}\gT{Totraum} }

{\mdseries
\gNw{Atmung}}

\begin{itemize}
\item \gNw{\"au{\ss}ere  Atmung : Gasaustausch Alveolen
/Lungenkapillaren}

\item \gNw{innere  Atmung : Gasaustausch Kapillaren /Gewebe}

\end{itemize}
{\mdseries
normale Atemfrequenzen sind altersabh\"angig:}

\begin{center}
\begin{tabular}{|m{4.367cm}|m{4.367cm}|m{4.367cm}|}
\hline
 &
\centering {\bfseries\upshape\color[rgb]{0.0,0.0,0.5019608} AF}\par

 &
\centering {\bfseries\upshape\color[rgb]{0.0,0.0,0.5019608} AZV}\par

\\\hline
{\mdseries Erwachsener}

 &
{\mdseries 12 -- 16/min}

 &
{\mdseries ca.  500 ml}

\\\hline
{\mdseries Neugeborenes}

 &
{\mdseries 40/min }

 &
{\mdseries 20-40 ml}

\\\hline
{\mdseries Kleinkind}

 &
{\mdseries 25/min}

 &
{\mdseries 100-200 ml}

\\\hline
{\mdseries Schulkind}

 &
{\mdseries 20/min}

 &
{\mdseries 200-400 ml}

\\\hline
\end{tabular}
\end{center}
{\mdseries
AMV = AZV x AF, beim Erwachsenen  ca. \gEs{6-8l}}

{\mdseries
AF  ... Atemfrequenz}

{\mdseries
AMV  ... Atemminutenvolumen}

{\mdseries
AZV  ... Atemzugsvolumen ... Volumen bei einem normalen Atemzug}

\subsection{Der Verdaungstrakt}
{\mdseries
{\quotedblbase}\gT{GastroIntestinaltrakt}{\textquotedblleft}
(GIT)}

{\mdseries\color{black}
\"Ubersicht}

\begin{center}
\tablehead{}\begin{supertabular}{m{3.238cm}|m{3.674cm}|m{6.189cm}|}
\hline
\multicolumn{1}{|m{3.238cm}|}{{\mdseries
\gEs{Mundh\"ole}}

} &
 &
{\mdseries Zerkleinerung der Nahrung und Vermischung mit Speichel}

\\\hline
 &
\multicolumn{1}{m{3.674cm}}{{\mdseries Zunge}

} &
\\\hhline{~-~}
 &
\multicolumn{1}{m{3.674cm}}{{\mdseries Z\"ahne}

} &
\\\hhline{~-~}
\multicolumn{1}{|m{3.238cm}|}{{\mdseries
\gEs{Rachen}}

} &
 &
{\mdseries \textbf{Kreuzung} des Speisebreis mit dem Luftweg}

{\mdseries \textbf{Kehldeckel} verschliesst beim Schlucken die
Luftr\"ohre}

\\\hline
\multicolumn{1}{|m{3.238cm}|}{{\mdseries
\gEs{Speiser\"ohre}
(\gT{\"Osophagus})}

} &
 &
\\\hline
\multicolumn{1}{|m{3.238cm}|}{{\mdseries \gEs{Magen}}

} &
 &
{\mdseries \textbf{Magens\"aure} zerlegt Nahrung}

{\mdseries Vermischung}

{\mdseries Pf\"ortner gibt Portionsweise Speisebrei an den
Zw\"olffingerdarm ab}

\\\hline
\multicolumn{1}{|m{3.238cm}|}{{\mdseries
\gEs{D\"unndarm}}

} &
 &
\\\hline
 &
{\bfseries Zw\"olffingerdarm}

 &
{\mdseries \textbf{Gallenfl\"ussigkeit} (${\leftarrow}$Leber) und
\textbf{Bauchspeichel} (${\leftarrow}$Pankreas) wird zugesetzt}

\\\hhline{~--}
 &
{\bfseries Krummdarm}

 &
{\mdseries \textbf{Spaltung} von N\"ahrstoffen (Eiwei{\ss},
Kohlehydrate, Fett) und N\"ahrstoff\textbf{aufnahme}}

\\\hhline{~--}
 &
\multicolumn{1}{m{3.674cm}}{{\bfseries Leerdarm}

} &
\\\hhline{~-~}
\multicolumn{1}{|m{3.238cm}|}{{\mdseries
\gEs{Dickdarm} (\gT{Colon})}

} &
 &
\\\hline
 &
{\bfseries Blinddarm}

 &
{\mdseries Sackgasse}

\\\hhline{~--}
 &
{\mdseries \textbf{Wurmfortsatz} (\gT{Appendix}) }

 &
{\mdseries Anh\"angsel der Sackgasse, kann sich entz\"unden.}

\\\hhline{~--}
 &
{\mdseries \textbf{Aufsteigender} Teil}

 &
{\mdseries \textbf{Wasserentzug}, \textbf{Eindickung},
Bakterienbesiedlung}

\\\hhline{~--}
 &
\multicolumn{1}{m{3.674cm}}{{\mdseries \textbf{Querender} Teil}

} &
\\\hhline{~-~}
 &
\multicolumn{1}{m{3.674cm}}{{\mdseries \textbf{Absteigender} Teil}

} &
\\\hhline{~-~}
 &
\multicolumn{1}{m{3.674cm}}{{\mdseries S-f\"ormig gebogener Teil
(\gT{Sigma})}

} &
\\\hhline{~-~}
\multicolumn{1}{|m{3.238cm}|}{{\mdseries
\gEs{Mastdarm} (\gT{Rektum})}

} &
 &
{\mdseries {\quotedblbase}Speicherung{\textquotedblleft} des
ausscheidungsbereiten Stuhls}

\\\hline
\multicolumn{1}{|m{3.238cm}|}{{\mdseries \gEs{Anus}}

} &
 &
{\mdseries Absetzen des Stuhles}

\\\hline
\end{supertabular}
\end{center}
\begin{itemize}
\item \end{itemize}
{\mdseries\color{black}
Der Stoffwechsel}

\begin{itemize}
\item {\mdseries
Die Verdauung dient dem Stoffwechsel}

\item {\mdseries
Proze{\ss}:}

\end{itemize}
\begin{itemize}
\item {\mdseries
Nahrungsaufnahme}

\item {\mdseries
Weiterverarbeitung im GIT}

\item {\mdseries
Transport der N\"ahrstoffe und Energietr\"ager  ins Blut}

\item {\mdseries
Zellaufbau und Abbau von Abfallstoffen}

\end{itemize}
\begin{itemize}
\item {\mdseries
\"Uber die Verdauung gewinnt der K\"orper  also einerseits die n\"otigen
{\quotedblbase}Bausteine{\textquotedblleft},  andererseits auch die
n\"otige Energie um  Zellen aufbauen und in Funktion erhalten  zu
k\"onnen.}

\end{itemize}
{\mdseries\color[rgb]{0.0,0.0,0.5019608}
Mundh\"ohle}

{\mdseries
Z\"ahne, zerkleinern Nahrung}

{\mdseries
Speicheldr\"usen: produzieren Speichel um Nahrung gleitf\"ahig zu machen
und Nahrung mit Amylase (Verdauungsenzym f\"ur Zucker und St\"arke) zu
versetzen}

{\mdseries
Zunge: schiebt Nahrung in Schlund weiter, verschlie{\ss}t beim Schlucken
Kehldeckel (sonst  \gEs{Aspiration}!)}

{\mdseries\color[rgb]{0.0,0.0,0.5019608}
Die Speiser\"ohre (\"Osophagus)}

\begin{center}
\tablehead{}\begin{supertabular}{|m{6.649cm}|m{6.651cm}|}
\hline
 &
\\\hline
{\mdseries ca. 24cm langer, dreischichtiger Muskelschlauch (oberes
Drittel quergestreifte Muskulatur, untere 2/3 glatt) }

{\mdseries Verbindung Mund/Magen}

{\mdseries {\quotedblbase}schiebt{\textquotedblleft} Nahrung durch
Muskelkontraktionen (Peristaltik)
{\quotedblbase}weiter{\textquotedblleft}}

{\mdseries hier findet keine Verdauung statt}

 &
\centering   [Warning: Image ignored]
% Unhandled or unsupported graphics:
%\includegraphics[width=2.341cm,height=5.85cm]{AASdev20090228-img44.png}
    [Warning: Image ignored] % Unhandled or unsupported graphics:
%\includegraphics[width=2.103cm,height=5.707cm]{AASdev20090228-img45.png}
 \par

\\\hline
\end{supertabular}
\end{center}
{\mdseries\color[rgb]{0.0,0.0,0.5019608}
Der Magen (Gaster, Ventriculus)}

\begin{itemize}
\item 2-3 l Magensaft

\item Schleim (sch\"utzt  Magenschleimhaut)

\item Salzs\"aure (Bakterienfalle, pH 1,0)

\item Pepsin  (Enzym f\"ur Eiwei{\ss}-Verdauung)

\item Wandmuskulatur mischt den Speisebrei mit der Magens\"aure

\item portionsweise Abgabe von Speisebrei durch den Pylorus
(Magenpf\"ortner) in den Zw\"olffingerdarm

\end{itemize}
\begin{center}
\tablehead{}\begin{supertabular}{|m{6.649cm}|m{6.651cm}|}
\hline
 &
\\\hline
{\centering   [Warning: Image ignored]
% Unhandled or unsupported graphics:
%\includegraphics[width=4.116cm,height=6.196cm]{AASdev20090228-img46.png}
 \par}

 &
{\centering   [Warning: Image ignored]
% Unhandled or unsupported graphics:
%\includegraphics[width=4.746cm,height=4.44cm]{AASdev20090228-img47.png}
 \par}

\\\hline
{\mdseries\upshape Lage des Magens}

 &
{\mdseries\upshape Der Magen, der Pf\"ortner und das erste St\"uck
Zw\"olffingerdarm im Durchschnitt}

\\\hline
\end{supertabular}
\end{center}
{\mdseries\color[rgb]{0.0,0.0,0.5019608}
Zw\"olffingerdarm (Duodenum)}

  [Warning: Image ignored] % Unhandled or unsupported graphics:
%\includegraphics[width=6.774cm,height=5.079cm]{AASdev20090228-img48.png}
   [Warning: Image ignored] % Unhandled or unsupported graphics:
%\includegraphics[width=5.278cm,height=4.663cm]{AASdev20090228-img49.png}
 

\begin{itemize}
\item {\mdseries
Beginn des D\"unndarmes, ca 25 cm lang}

\item {\mdseries
Verdauungsenzyme aus Leber und Pankreas f\"ur Fett, Eiwei{\ss}, Zucker
m\"unden hier ein}

\end{itemize}
{\mdseries\color[rgb]{0.0,0.0,0.5019608}
Die \index{Leber}Leber }

{\mdseries
\gZ{{\quotedblbase}Hepar{\textquotedblleft}}}

\begin{center}
\tablehead{}\begin{supertabular}{|m{6.649cm}|m{6.651cm}|}
\hline
\begin{itemize}
\item Die Leber liegt im rechten Oberbauch und hat einen linken und
rechten Lappen. 

\item \item Abbau von Giftstoffen

\item {\quotedblbase}Verdauungsdr\"use{\textquotedblleft}

\end{itemize}
\begin{itemize}
\item produziert \index{Gallenfl\"ussigkeit}Gallenfl\"ussigkeit
({\quotedblbase}Galle{\textquotedblleft}), welche in der Gallenblase
gespeichert wird: Fettverdauung

\item Von der Leber gehen die Galleng\"ange, an den Galleng\"angen
h\"angt die 

\end{itemize}
\begin{itemize}
\item \index{Gallenblase}Gallenblase: Nur Speicher f\"ur
Gallefl\"ussigkeit

\end{itemize}
\begin{itemize}
\item speichert Fett \& Zucker

\item erh\"alt \"uber \index{Pfortader}\gT{Pfortader}
(Vena portae) n\"ahrstoffreiches Blut vom Darm zwecks Entgiftung

\end{itemize}
 &
{\centering   [Warning: Image ignored]
% Unhandled or unsupported graphics:
%\includegraphics[width=4.719cm,height=3.539cm]{AASdev20090228-img50.png}
 \par}

\\\hline
\end{supertabular}
\end{center}
  [Warning: Image ignored] % Unhandled or unsupported graphics:
%\includegraphics[width=5.469cm,height=4.037cm]{AASdev20090228-img51.png}
   [Warning: Image ignored] % Unhandled or unsupported graphics:
%\includegraphics[width=5.045cm,height=4.005cm]{AASdev20090228-img52.png}
   [Warning: Image ignored] % Unhandled or unsupported graphics:
%\includegraphics[width=2.373cm,height=4.068cm]{AASdev20090228-img53.png}
 

{\mdseries\color[rgb]{0.0,0.0,0.5019608}
\index{Bauchspeicheldr\"use}Bauchspeicheldr\"use
(\index{Pankreas}Pankreas)}

{\mdseries
15 -- 22 cm lang, Kopf und Schwanz}

{\mdseries
Ausf\"uhrungsgang m\"undet mit Gallengang in den Zw\"olffingerdarm
(Duodenum)}

{\mdseries
Aufgaben: 2 wesentlcihe Funktionen}

\begin{itemize}
\item produziert Enzyme f\"ur Verdauung ${\rightarrow}$ Zw\"offingerdarm

\item sch\"uttet Hormone ins Blut aus

\item \index{Insulin}Insulin ={\textgreater} Blutzucker ${\downarrow}$

\item Glukagon ${\rightarrow}$ Blutzucker ${\uparrow}$

\end{itemize}
Beim Diabetiker produziert die Bauchspeicheldr\"use nicht genug Insulin.




\begin{center}
 [Warning: Image ignored] % Unhandled or unsupported graphics:
%\includegraphics[width=8.804cm,height=5.081cm]{AASdev20090228-img54.jpg}

\end{center}
\begin{center}
 [Warning: Image ignored] % Unhandled or unsupported graphics:
%\includegraphics[width=5.788cm,height=4.525cm]{AASdev20090228-img55.png}

\end{center}
{\mdseries\color[rgb]{0.0,0.0,0.5019608}
\index{Krummdarm}Krumm- und \index{Leerdarm}Leerdarm }

\begin{center}
 [Warning: Image ignored] % Unhandled or unsupported graphics:
%\includegraphics[width=3.049cm,height=3.091cm]{AASdev20090228-img56.png}

\end{center}
{\mdseries
\gZ{(Jejunum, Ileum)}}

\begin{itemize}
\item {\mdseries
ca. 3 m lang}

\item {\mdseries
Oberfl\"achenvergr\"o{\ss}erung durch Darmzotten}

\item {\mdseries
\gEs{N\"ahrstoffaufnahme}}

\item {\mdseries
beweglich, am Darmgekr\"ose (Mesenterium) aufgeh\"angt}

\end{itemize}
{\mdseries\color[rgb]{0.0,0.0,0.5019608}
Dickdarm (Colon)}

{\mdseries
Der Dickdarm (Colon) ist ca. 1,5 m lang und wird in verschiedene
Abschnitte unterteilt: }

\begin{center}
 [Warning: Image ignored] % Unhandled or unsupported graphics:
%\includegraphics[width=3.751cm,height=4.328cm]{AASdev20090228-img57.png}
\captionof{figure}[Lage des Dickdarms und des Appendix]{Lage des
Dickdarms und des Appendix}

\end{center}
\begin{center}
\tablehead{}\begin{supertabular}{m{1.604cm}|m{4.8500004cm}|m{6.6460004cm}}
\hline
\multicolumn{1}{|m{1.604cm}|}{{\mdseries
\gEs{Dickdarm} (\gT{Colon})}

} &
{\mdseries teilt sich in:}

 &
\\\hline
 &
{\bfseries Blinddarm}

 &
{\mdseries Sackgasse}

\\\hhline{~--}
 &
{\mdseries \textbf{Wurmfortsatz} (\gT{Appendix}) }

 &
{\mdseries Anh\"angsel der Sackgasse, kann sich entz\"unden.}

\\\hhline{~--}
 &
{\mdseries \textbf{Aufsteigender} Teil}

 &
{\mdseries \textbf{Wasserentzug}, \textbf{Eindickung},
Bakterienbesiedlung}

\\\hhline{~--}
 &
\multicolumn{1}{m{4.8500004cm}}{{\mdseries \textbf{Querender} Teil}

} &
\\\hhline{~-~}
 &
\multicolumn{1}{m{4.8500004cm}}{{\mdseries \textbf{Absteigender} Teil}

} &
\\\hhline{~-~}
 &
\multicolumn{1}{m{4.8500004cm}}{{\mdseries S-f\"ormig gebogener Teil
(\gT{Sigma})}

} &
\\\hhline{~-~}
\end{supertabular}
\end{center}
{\mdseries
Der Leerdarm m\"undet in den Dickdarm. Diese M\"undung ist wie eine
\textit{T-Kreuzung} gebaut: \textit{Nach oben} geht es in den
\gEs{aufsteigenden Teil} des Dickdarmes,
\textit{nach unten }in den \gEs{Blinddarm}. Wie der
Name des Blinddarmes schon sagt endet er blind, er ist daher eine
Sackgasse. Am Blinddarm angeh\"angt findet man den
\gT{\textit{Wurmfortsatz}}
(\gT{Appendix}) welcher sich h\"aufig entz\"undet
(\gT{Appendizitis}). Der Wurmfortsatz erf\"ullt beim
Menschen keine nenneswerte Funktion, ausser der Arbeitsplatzsicherung
von An\"asthesisten und Chirurgen. }

{\mdseries\color[rgb]{0.0,0.0,0.5019608}
\index{Bauchfell}Bauchfell}

\begin{itemize}
\item {\mdseries
Die gesamte Bauchh\"ohle ist mit  Bauchfell
(\index{Peritoneum}\gT{Peritoneum}) ausgekleidet}

\item {\mdseries
d\"unner \"Uberzug}

\item {\mdseries
{\quotedblbase}Aufh\"angung{\textquotedblleft} der Bauchorgane }

\item {\mdseries
zus\"atzlich gro{\ss}es + kleines Netz}

\item {\mdseries
Kann sich Entz\"unden (Bauchfellentz\"undung = Peritonitis): zB bei
Darmverletzung, Bauchfelldialyse }

\end{itemize}
\begin{center}
\tablehead{}\begin{supertabular}{|m{6.649cm}|m{6.651cm}|}
\hline
{\centering   [Warning: Image ignored]
% Unhandled or unsupported graphics:
%\includegraphics[width=4.982cm,height=6.767cm]{AASdev20090228-img58.png}
 \par}

 &
{\mdseries Der Bauch {\quotedblbase}aufgeklappt: Nach oben geklappt
sieht man das gele {\quotedblbase}Gro{\ss}e Netz{\textquotedblleft},
auf die Seiten und nach unten die Bauchmuskulatur. }

{\mdseries Die hellen, fleischfarbenen Schlingen sind
D\"unndarmschlingen, die dicken braunen Schlingen sind der aufsteigende
und der quere Teil des Dickdarmes.}

\\\hline
\end{supertabular}
\end{center}
\subsection{Sonstige Bauchorgane}
{\mdseries\color[rgb]{0.0,0.0,0.5019608}
Die Milz }

{\mdseries
\gZ{(Lien, Splen)}}

\begin{center}
 [Warning: Image ignored] % Unhandled or unsupported graphics:
%\includegraphics[width=3.61cm,height=2.202cm]{AASdev20090228-img59.png}

\end{center}
\begin{itemize}
\item {\mdseries
hat nichts  mit der Verdauung  zu tun }

\begin{center}
 [Warning: Image ignored] % Unhandled or unsupported graphics:
%\includegraphics[width=4.301cm,height=6.628cm]{AASdev20090228-img60.png}

\end{center}
\item {\mdseries
12 cm, oval, \gEs{von Kapsel  umgeben}}

\item {\mdseries
im linken Oberbauch unter den Rippen gelegen}

\item {\mdseries
Teil des Immunsystems}

\item {\mdseries
entfernt alte Erythrozyten, beim Embryo auch Blut\-bildung}

\item {\mdseries
sehr gut durchblutet, Blutspeicher (Gefahr:
\index{Milzri{\ss}}Milzri{\ss})}

\end{itemize}
{\mdseries\color{black}
\gCa{Achtung:} sog.
\gEs{{\quotedblbase}}\index{Zweizeitiger
Milzriss}\gEs{Zweizeitiger
Milzriss}\gEs{{\textquotedblleft}}: Zuerst reisst
die Milz ein, die Kapsel bleib erhalten, es blutet zuerst in die Kapsel
und der Druck steigt. Nach einer gewissen Zeit reisst auch die Kapsel
und es blutet nun ungehindert aus der inneren Wunde.
\gEs{Der Patient wirkt zuerst
{\textquotedblleft}ganz normal{\textquotedblleft}, unerwartet und
pl\"otzlich verf\"allt er.} }

{\mdseries\color[rgb]{0.0,0.0,0.5019608}
Die Niere(n) (Ren)}

\begin{itemize}
\item {\mdseries
paarig angelegtes Organ}

\begin{center}
 [Warning: Image ignored] % Unhandled or unsupported graphics:
%\includegraphics[width=3.873cm,height=4.654cm]{AASdev20090228-img61.png}

\end{center}
\item {\mdseries
stark durchblutet, empfindlich auf RR- Abfall ${\rightarrow}$ Funktion
${\downarrow}$ }

\item {\mdseries
reguliert}

\end{itemize}
\begin{itemize}
\item {\mdseries
\gEs{Wasser-} und
\gEs{Elektrolythaushalt}}

\item {\mdseries
pH (S\"aure/Basen)}

\end{itemize}
\begin{itemize}
\item {\mdseries
Geht im Alter, besonders bei Diabetikern, kaputt ${\rightarrow}$
{\quotedblbase}\gEs{Chronische
Niereninsuffizienz}{\textquotedblleft} ${\rightarrow}$
{\quotedblbase}\gEs{Blutw\"asche}{\textquotedblleft}
(\index{Dialyse}\gT{Dialyse})}

\end{itemize}
{\mdseries\color{black}
Aufbau}

\begin{itemize}
\item {\mdseries
Rinde, Mark}

\item {\mdseries
Hilus}

\item {\mdseries
Gef\"a{\ss}e}

\item {\mdseries
Nierenbecken mit Anschluss an den Harntrakt}

\end{itemize}
{\bfseries\color[rgb]{0.5019608,0.0,0.0}
Niere: Anatomie und Funktion muss noch ausgearbeitet werden.}

{\mdseries\color{black}
%
%
Mikroskopischer Aufbau}

\begin{itemize}
\item {\mdseries
Nephron (ca. 1 Mio/Niere)}

\item {\mdseries
Glomerulus: Hier wird das Blut gefiltert und Plasmawasser abgepresst
(\~{}150~l~/~Tag)}

\item {\mdseries
Schleifenapparat \gZ{(Tubulus)}: Das abgepresste
Plasma\-wasser wandert den Schleifen\-apparat entlang, hier werden die
meisten wichtigen Stoffe und der Gro{\ss}\-teil des Wassers wiederum
zur\"uckgeholt, sodass am Ende nur mehr 1,5 -- 2~l Harn pro Tag in der
Harnblasen landen. }

\end{itemize}
{\mdseries\color{black}
Funktionsweise}

\begin{itemize}
\item {\mdseries
1500 l Blut pro Tag wird {\quotedblbase}bearbeitet{\textquotedblleft}}

\item {\mdseries
Glomerula: Filtration ${\rightarrow}$ 150  l
\index{Prim\"arharn}Prim\"arharn (enth\"alt Elektrolyte, etc.) wird
gewonnen und fliesst durch den Schleifenapparat}

\item {\mdseries
Schleifenapparat: {\quotedblbase}Recycling{\textquotedblleft}, Gutes
kommt zur\"uck ins Blut, der Rest fliesst weiter Richtung Harnblase
${\rightarrow}$ 1,5 -- 2  l (End-)Harn/Tag }

\end{itemize}
{\mdseries\color[rgb]{0.0,0.0,0.5019608}
Die Nebenniere(n) (Adren)}

{\mdseries
paariges, jeweils am oberen Nierenpol gelegenes Organ}

{\mdseries
Hat nichts  mit dem Harntrakt oder der eigentlichen Nierenfunktion zu
tun.}

{\mdseries
Hormonproduktion}

\begin{itemize}
\item {\mdseries
Cortison}

\item {\mdseries
Sexualhormone (Androgene)}

\item {\mdseries
\gEs{Adrenalin} , Noradrenalin}

\end{itemize}
{\mdseries
Die Nebennieren sitzen wie {\quotedblbase}M\"utzen{\textquotedblleft}
auf den Nieren (s. Illustration \ref{seq:refIllustration4} unten). }

{\mdseries\color[rgb]{0.0,0.0,0.5019608}
Harntrakt}

\begin{center}
\tablehead{}\begin{supertabular}{|m{6.649cm}|m{6.651cm}|}
\hline
{\mdseries Der Harntrakt besteht aus }

\begin{itemize}
\item {\mdseries Nieren inkl. Nierenbecken}

\item {\mdseries \index{Harnleiter}\gEs{Harnleiter} 
(Ureter ), 30 cm lang, m\"undet  in die}

\item {\mdseries \index{Blase}\gEs{Blase} 
\gZ{(Vesica )}}

\item {\mdseries \gEs{Harnr\"ohre}  (Urethra, [2642?]
25cm, [2640?] 4cm)}

\end{itemize}
{\mdseries Die Harnr\"ohre verl\"auft beim Mann von der Harnblase
\gEs{durch die Prostata} und den Penis und m\"undet
an der Eichel in die Aussenwelt. Bei einer Verg\"o{\ss}erung der
Prostata kann der Harnabflu{\ss} bis zur Unpassierbarkeit erschwert
sein.  }

{\mdseries Bei der Frau ist die Harnr\"ohre k\"urzer und m\"undet
oberhalb des Scheideneinganges. }

 &
{\centering   [Warning: Image ignored]
% Unhandled or unsupported graphics:
%\includegraphics[width=5.85cm,height=7.43cm]{AASdev20090228-img62.jpg}
 \par}

\captionof{figure}{Harntrakt mit Nieren (linke Niere im Querschnitt mit
Nierenbescken), Harnleiter, Blase und Harnr\"ohre. Auf den Nieren sieht
man wie eine Zipfelm\"utze die Nebennieren.}
\label{seq:refIllustration4}
\\\hline
\end{supertabular}
\end{center}
\subsection{Geschlechtsorgane}
{\mdseries\color[rgb]{0.0,0.0,0.5019608}
M\"annliche Geschlechtsorgane}

\begin{itemize}
\item {\mdseries
innen:}

\end{itemize}
\begin{itemize}
\item {\mdseries
Hoden (Testes)}

\end{itemize}
\begin{itemize}
\item {\mdseries
Spermienproduktion}

\item {\mdseries
Geschlechtshormon-Produktion}

\end{itemize}
\begin{itemize}
\item {\mdseries
Nebenhoden: Spermienreifung und Lagerung}

\item {\mdseries
Samenstrang, besteht aus }

\end{itemize}
\begin{itemize}
\item {\mdseries
Samenleiter}

\item {\mdseries
Nerven}

\item {\mdseries
Gef\"a{\ss}e}

\item {\mdseries
durchzieht Leistenkanal}

\end{itemize}
\begin{itemize}
\item {\mdseries
\index{Prostata}\gEs{Prostata} (Vorsteherdr\"use)}

\end{itemize}
\begin{itemize}
\item {\mdseries
produziert Samenfl\"ussigkeit (Beweglichkeit der Spermien)}

\item {\mdseries
\gEs{Vergr\"o{\ss}ert sich im Alter} und
\gEs{kann die Harnr\"ohre abdr\"ucken}, es kommt zum
Harnstau ${\rightarrow}$ h\"aufger Notfall!}

\end{itemize}
\begin{itemize}
\item {\mdseries
au{\ss}en:}

\end{itemize}
\begin{itemize}
\item {\mdseries
Glied (Penis): Schwellk\"orper, Harnsamenr\"ohre, Begattungsorgan}

\item {\mdseries
Hodensack (Skrotum): umh\"ullt Hoden und Nebenhoden}

\end{itemize}
\begin{itemize}
\item {\mdseries
Hodenstiel kann sich verdrehen ${\rightarrow}$ Notfall!}

\end{itemize}
  [Warning: Image ignored] % Unhandled or unsupported graphics:
%\includegraphics[width=13.699cm,height=8.52cm]{AASdev20090228-img63.png}
 

{\mdseries\color[rgb]{0.0,0.0,0.5019608}
Weibliche Geschlechtsorgane}

\begin{itemize}
\item {\mdseries
innen:}

\end{itemize}
\begin{itemize}
\item {\mdseries
\gEs{Eierst\"ocke} (Ovar): Reifung der Eizellen}

\item {\mdseries
\gEs{Eileiter} (Tubae): Transport des reifen Eis in
Richtung Geb\"armutter}

\item {\mdseries
\gEs{Geb\"armutter}
(\index{Uterus}\gT{Uterus}): dehnbarer Hohlmuskel, mit
Schleimhaut ausgekleidet, }

\end{itemize}
\begin{itemize}
\item {\mdseries
Schleimhaut wird regelm\"a{\ss}ig aufgebaut und wieder abgesto{\ss}en
(Menstruation, \gEs{Regelblutung})}

\item {\mdseries
hier nistet sich die befruchtete Eizelle ein. }

\item {\mdseries
Die M\"undung der Geb\"armutter in die Scheide ist der
\gT{Muttermund}. }

\end{itemize}
\begin{itemize}
\item {\mdseries
\gEs{Scheide} (Vagina)}

\end{itemize}
\begin{itemize}
\item {\mdseries
au{\ss}en: }

\end{itemize}
\begin{itemize}
\item {\mdseries
Schamlippen: go{\ss}e/\"au{\ss}ere, kleine/innere}

\item {\mdseries
Kitzler (Klitoris): entspricht der m\"annlichen Eichel}

\end{itemize}
  [Warning: Image ignored] % Unhandled or unsupported graphics:
%\includegraphics[width=13.399cm,height=8.485cm]{AASdev20090228-img64.png}
 

\paragraph{Der weibliche Zyklus}
\begin{itemize}
\item {\mdseries
Der weibliche Zyklus dauert 28 Tage}

\item {\mdseries
gerechnet von Menstruation zu Menstruation  (Monatsblutung )}

\item {\mdseries
Uterusschleimhaut wird aufgebaut und wieder abgesto{\ss}en}

\item {\mdseries
hormongesteuert}

\end{itemize}
\begin{itemize}
\item {\mdseries
\"Ostrogen, Progesteron}

\end{itemize}
\begin{itemize}
\item {\mdseries
ca. am 14. Tag: Eisprung}

\item {\mdseries
Befruchtung nein: }

\end{itemize}
\begin{itemize}
\item {\mdseries
Zyklus l\"auft normal weiter, Menstruation}

\end{itemize}
\begin{itemize}
\item {\mdseries
Befruchtung ja:}

\end{itemize}
\begin{itemize}
\item {\mdseries
Eizelle nistet sich im Uterus ein, Zyklus gestoppt}

\end{itemize}
{\mdseries
Was kann dabei schief gehen? Wenn sich die Eizelle schon im Ei\-leiter
einnistet kommt es zu einer
\index{Eileiter\-schwanger\-schaft}\gEs{Eileiter\-schwanger\-schaft}
die \gEs{lebens\-bedrohlich} ist (Eileiter rei{\ss}t
und die Frau verblutet).}

 [Warning: Image ignored] % Unhandled or unsupported graphics:
%\includegraphics[width=13.576cm,height=7.541cm]{AASdev20090228-img65.png}
\captionof{figure}[Der weibliche Zyklus. Oben: Follikelreifung. Unten:
Auf- und Abbau der Uterusschleimhaut. Am Tag 0 des Zyklus wird die
Uterusschleimheut abgestossen, es kommt zur Regelblutung, danach wird
die Schleimhaut wieder aufgebaut. Am Tag 14 erfolgt der Eispung. Kommt
es zu keiner Befruchtung wird die Schleimheut erneut am Tag 28
abgesto{\ss}en. Nach einer Befruchtung wandert die Eizelle weiter in
den Uterus und nistet sich dort ein. ]{Der weibliche Zyklus. Oben:
Follikelreifung. Unten: Auf- und Abbau der Uterusschleimhaut. Am Tag 0
des Zyklus wird die Uterusschleimheut abgestossen, es kommt zur
Regelblutung, danach wird die Schleimhaut wieder aufgebaut. Am Tag 14
erfolgt der Eispung. Kommt es zu keiner Befruchtung wird die
Schleimheut erneut am Tag 28 abgesto{\ss}en. Nach einer Befruchtung
wandert die Eizelle weiter in den Uterus und nistet sich dort ein. }


\section{[RS V] Vitalfunktionen}
{\mdseries\itshape
Sebastian Gabriel [R]; Michael Withofner}

{\bfseries\color[rgb]{0.5019608,0.0,0.0}
\gEs{Changelog:}}

{\bfseries\color[rgb]{0.5019608,0.0,0.0}
2009-02-25: Start Changelog. Kapitel mit Korr Koch korr.}

{\mdseries\color[rgb]{0.0,0.0,0.5019608}
\"Uberblick}

\begin{center}
\tablehead{}\begin{supertabular}{|m{6.649cm}|m{6.651cm}|}
\hline
{\mdseries\color{black} Vitalfunktionen I. Ordnung}

 &
{\mdseries\color{black} Vitalfunktionen II. Ordnung}

\\\hline
\begin{itemize}
\item {\mdseries Bewu{\ss}tsein}

\item {\mdseries Atmung}

\item {\mdseries Kreislauf}

\end{itemize}
{\mdseries wichtigste, grundlegende Funktionen des K\"orpers}

{\mdseries Ausfall einer dieser Funktionen }

{\mdseries bedeutet immer akute Lebensgefahr!!! }

{\mdseries mu{\ss} akutmedizinisch behandelt werden!!!}

 &
\begin{itemize}
\item {\mdseries Wasser-/Elektrolythaushalt}

\item {\mdseries S\"aure-/Basenhaushalt}

\item {\mdseries Temperaturhaushalt}

\item {\mdseries Hormonsystem}

\item {\mdseries Stoffwechsel}

\item {\mdseries Immunsystem}

\end{itemize}
\\\hline
\end{supertabular}
\end{center}
\paragraph[NACA -- Score]{NACA -- Score}
\index{NACA}{\mdseries
\gZ{National Advisory Committee for Aeronautics}}

{\mdseries
\gZ{Diente urspr\"unglich der Bewertung von 
Erkrankungen / Verletzungen und der Entscheidung welche Soldaten zuerst
auszufliegen w\"aren. Heute dient der NACA-Score der groben
statistischen Erfassung.  }}

{\mdseries
\gZ{1 bis 7:}}

\begin{center}
\begin{tabular}{m{1.201cm}m{12.1cm}}
\centering {\bfseries\upshape\color[rgb]{0.0,0.0,0.5019608}
\gZ{ NACA}}\par

 &
\\
\centering {\bfseries \gZ{I}}\par

 &
{\mdseries \gZ{ keine \"arztliche Therapie
erforderlich}}

\\
\centering {\bfseries \gZ{II}}\par

 &
{\mdseries \gZ{ ambulante Abkl\"arung}}

\\
\centering {\bfseries \gZ{III}}\par

 &
{\mdseries \gZ{ station\"are Abkl\"arung}}

\\
\centering {\bfseries \gZ{IV}}\par

 &
{\mdseries \gZ{ akute Lebensgefahr nicht
ausgeschlossen}}

\\
\centering {\bfseries \gZ{V}}\par

 &
{\mdseries \gZ{ akute Lebensgefahr}}

\\
\centering {\bfseries \gZ{VI}}\par

 &
{\mdseries \gZ{ Reanimation}}

\\
\centering {\bfseries \gZ{VII}}\par

 &
{\mdseries \gZ{ Tod}}

\\
\end{tabular}
\end{center}
\subsection{Bewusstsein}
{\mdseries
\"Uberbegriff u.a. f\"ur \gEs{Wachheit},
\gEs{Orientierung}, Aufmerksamkeit, Auffassungsgabe,
Denkverlauf und Merkf\"ahigkeit\label{ref:JRcite1Pschyrembel259}[1] . }

{\mdseries
\gZ{Die {\quotedblbase}wahre{\textquotedblleft}
Definition von Bewusstsein ist seit tausenden von Jahren ungekl\"art.
Da f\"ur die Sanit\"aterausbildung aber nur beschr\"ankt Zeit zur
Verf\"ugung steht ist die obige Definition f\"ur den Alltag erstmal
ausreichend.}}

{\mdseries
Das Bewusstsein ist besonders wichtig hinsichtlich des Schutzes des
Menschen vor Bedrohung. Ein bewusstseinsklarer Mensch kann sich gegen
Gefahren wehren, ein eingetr\"ubter oder bewusstloser Mensch kann dies
schlecht oder gar nicht. }

{\mdseries\color[rgb]{0.0,0.0,0.5019608}
Bewu{\ss}seinsst\"orung}

{\mdseries
St\"orung des Bewusstseins, besonders der
\gEs{Wachheit} und
\gEs{Orientierung}. }

{\mdseries
Orientierung}

{\mdseries
Es ist wichig zu \"Uberpr\"ufen ob der Patient orientiert ist. Man
unterscheidet dabei: }

\begin{enumerate}
\item Orientierung zur Person: Name, Alter. 

\item \"Ortliche Orientierung: Wo befindet sich der Patient gerade
(Stadt, Land, Adresse, Wohnung, Auto, etc.)?

\item Zeitliche Orientierung: Tageszeit, Wochentag, Monat, Jahreszeit,
Jahr? 

\end{enumerate}
{\mdseries\color{black}
Allgemein}

{\mdseries
Wie bereits oben erw\"ahnt ist das Bewusstsein wesentlich an der Abwehr
von Gefahren beteiligt. Je mehr das Bewusstsein gest\"ort oder
reduziert ist desto gef\"ahrlicher ist das f\"ur den Patienten. }

\begin{itemize}
\item Ab einem gewissen Punkt sind auch so wesentliche Funktionen wie
der \gT{Schluck-}, \gT{W\"urge-} bzw.
\gT{Hustenreflex} ausfallen und es kann zum verschlucken
von Mageninhalt oder Blut in die Lunge kommen
(\gEs{Aspiration}). Man spricht auch vom Ausfall der
\gT{Schutzreflexe}. 

\item Bei schweren Bewusstseinsst\"orungen kann es auch zum Erschlaffen
wichtiger Muskeln wie z.B. Der \gEs{Zunge} kommen.
F\"allt diese zur\"uck kann sie den Atemweg blockieren oder ganz
verschlie{\ss}en. 

\end{itemize}
{\mdseries
Daher sind schwere \gEs{Bewusstseinsst\"orungen}
\gEs{lebensgef\"ahrlich}!}

{\mdseries\color{black}
Ursachen}

{\mdseries
${\rightarrow}$ Sch\"adigung des zentralen Nervensystems}

\begin{center}
\tablehead{}\begin{supertabular}{|m{6.649cm}|m{6.651cm}|}
\hline
\centering {\bfseries\upshape\color[rgb]{0.0,0.0,0.5019608} prim\"are 
Ursachen}\par

 &
\centering {\bfseries\upshape\color[rgb]{0.0,0.0,0.5019608} sekund\"are 
Ursachen:}\par

\\\hline
{\mdseries direkt das ZNS betreffend}

 &
{\mdseries \"Uber Umwege das ZNS betreffend}

\\\hline
{\mdseries Verletzungen (SHT)}

{\mdseries Blutungen, Hirnschwellung}

{\mdseries Schlaganfall}

{\mdseries epileptischer Anfall}

{\mdseries Entz\"undungen}

{\mdseries \gZ{Tumor}}

 &
{\mdseries St\"orungen der O2-Versorgung :}

Atemst\"orungen  (Hypoxie)

Herz-Kreislaufst\"orungen

Elektrolytentgleisungen

{\mdseries Vergiftungen:}

Stoffwechselst\"orungen (Zuckerhaushalt, Leberversagen)

Alkohol

Medikamente, Drogen

Reizgase,  L\"osungsmittel ...

\\\hline
\end{supertabular}
\end{center}
{\mdseries\color{black}
Grade und Symptome}

{\bfseries\color[rgb]{0.5019608,0.0,0.0}
Symptome Bewusstseinsst\"orungen: Tabelle geh\"ort erg\"anzt.%
%
}

\begin{center}
\tablehead{}\begin{supertabular}{|m{3.225cm}|m{3.225cm}|m{3.225cm}|m{3.225cm}|}
\hline
bewu{\ss}tseinsklar

 &
 &
 &
\\\hline
bewu{\ss}tseinsgetr\"ubt

 &
{\mdseries desorientiert}

 &
 &
\\\hline
 &
{\mdseries delirant}

 &
 &
\\\hline
 &
agitiert

 &
 &
\\\hline
 &
somnolent

 &
{\mdseries Schl\"afrig, aber leicht erweckbar. }

 &
\\\hline
 &
sopor\"os

 &
 &
{\mdseries Gefahr!}

\\\hline
bewu{\ss}tlos

 &
Koma

 &
 &
{\mdseries Lebensgefahr! }

\\\hline
\end{supertabular}
\end{center}
{\mdseries\color{black}
Gefahren}

\begin{itemize}
\item {\mdseries
Zur\"uckfallen der Zunge ${\rightarrow}$ Erstickung!}

\item {\mdseries
\gEs{Aspiration}  (Husten- bzw. W\"urgereflex sowie 
s\"amtliche andere Schutzreflexe ausgefallen,  Muskulatur komplett
erschlafft ={\textgreater}  Erbrochenes gelangt ungehindert in die 
Lunge!)}

\end{itemize}
{\mdseries\color{black}
\index{Glasgow Coma Scale}Glasgow Coma Scale (\index{GCS}GCS)}

{\mdseries
... dient zur objektiven Einsch\"atzung der Bewu{\ss}tseinslage, Skala
mit 3 bis max. 15 Punkten.}

Bewu{\ss}tseinsklarer Patient: 15 Punkte. Toter Patient: 3 Punkte (nicht
0!)

{\mdseries
Untersucht werden: }

\begin{flushleft}
\tablehead{}\begin{supertabular}{m{3.224cm}m{0.874cm}|m{1.1459999cm}|m{7.649cm}|}
\hline
\multicolumn{1}{|m{3.224cm}|}{} &
\centering {\bfseries\upshape\color[rgb]{0.0,0.0,0.5019608} Max.}\par

 &
\centering {\bfseries\upshape\color[rgb]{0.0,0.0,0.5019608} Punkte}\par

 &
\centering {\bfseries\upshape\color[rgb]{0.0,0.0,0.5019608}
Reaktion}\par

\\\hline
\multicolumn{1}{|m{3.224cm}|}{\begin{itemize}
\item {\mdseries \gEs{Augen\"offnung}}

\end{itemize}
} &
{\mdseries 4}

 &
{\mdseries 4}

 &
{\mdseries spontan}

\\\hline
 &
 &
{\mdseries 3}

 &
{\mdseries auf Ansprache}

\\\hhline{~~--}
 &
 &
{\mdseries 2}

 &
{\mdseries Schmerzreiz}

\\\hhline{~~--}
 &
 &
{\mdseries 1}

 &
{\mdseries nicht}

\\\hhline{~~--}
\multicolumn{1}{|m{3.224cm}|}{\begin{itemize}
\item {\mdseries \gEs{Verbale Antwort}}

\end{itemize}
} &
{\mdseries 5}

 &
{\mdseries 5}

 &
{\mdseries klar}

\\\hline
 &
 &
{\mdseries 4}

 &
{\mdseries verwaschen}

\\\hhline{~~--}
 &
 &
{\mdseries 3}

 &
{\mdseries stammelnd}

\\\hhline{~~--}
 &
 &
{\mdseries 2}

 &
{\mdseries unartikuliert}

\\\hhline{~~--}
 &
 &
{\mdseries 1}

 &
{\mdseries nicht}

\\\hhline{~~--}
\multicolumn{1}{|m{3.224cm}|}{\begin{itemize}
\item {\mdseries \gEs{Motorische Reaktion}}

\end{itemize}
} &
{\mdseries 6}

 &
{\mdseries 6}

 &
{\mdseries spontan}

\\\hline
 &
 &
{\mdseries 5}

 &
{\mdseries gezielte Abwehrbewegungen}

\\\hhline{~~--}
 &
 &
{\mdseries 4}

 &
{\mdseries ungezielte Abwehrbewegungen}

\\\hhline{~~--}
 &
 &
{\mdseries 3}

 &
{\mdseries Beugung}

\\\hhline{~~--}
 &
 &
{\mdseries 2}

 &
{\mdseries Strecken}

\\\hhline{~~--}
 &
 &
{\mdseries 1}

 &
{\mdseries bewegt gar nicht}

\\\hhline{~~--}
\end{supertabular}
\end{flushleft}
{\mdseries\color{black}
Ma{\ss}nahmen bei getr\"ubten Patienten}

{\bfseries\color[rgb]{0.5019608,0.0,0.0}
Ma{\ss}nahmen bei Bewusstseinsst\"orungen: Neufassung, bitte besonders
genau durchsehen!%
%
}

\begin{itemize}
\item {\mdseries\color[rgb]{0.0,0.0,0.5019608}
BAK -- Schema  (Bewu{\ss}tsein/Atmung/Kreislauf)}

\item {\mdseries\color[rgb]{0.0,0.0,0.5019608}
Lagerung: situationsgerecht, im Zweifel \gEs{stabile
Seitenlage }}

\item {\mdseries\color[rgb]{0.0,0.0,0.5019608}
Je nach Verdachtsdiagnose: }

\end{itemize}
\begin{itemize}
\item {\mdseries\color[rgb]{0.0,0.0,0.5019608}
O2: nach Ermessen, im Zweifel ca. 6-8 l}

\item {\mdseries\color[rgb]{0.0,0.0,0.5019608}
Reanimationsbereichtschaft}

\item {\mdseries\color[rgb]{0.0,0.0,0.5019608}
Notarzt }

\end{itemize}
\begin{itemize}
\item {\mdseries\color[rgb]{0.0,0.0,0.5019608}
Engmaschiges Monitoring (Atmung, Pulsoxymetrie, RR)}

\item {\mdseries\color[rgb]{0.0,0.0,0.5019608}
W\"armeerhalt}

\item {\mdseries\color[rgb]{0.0,0.0,0.5019608}
\gEs{Ursachenforschung} -- Warum ist der Patient
eingetr\"ubt?}

\item {\mdseries\color[rgb]{0.0,0.0,0.5019608}
Neurcheck}

\end{itemize}
\begin{itemize}
\item {\mdseries\color[rgb]{0.0,0.0,0.5019608}
\gEs{Immer auch an BZ denken!}}

\end{itemize}
\begin{itemize}
\item {\mdseries\color[rgb]{0.0,0.0,0.5019608}
Traumacheck}

\end{itemize}
{\mdseries\color{black}
Ma{\ss}nahmen bei bewu{\ss}tlosen und sopor\"osen Patienten}

\begin{itemize}
\item {\mdseries\color[rgb]{0.0,0.0,0.5019608}
BAK -- Schema  (Bewu{\ss}tsein/Atmung/Kreislauf)}

\item {\mdseries\color[rgb]{0.0,0.0,0.5019608}
Lagerung: \gEs{stabile Seitenlage}}

\item {\mdseries\color[rgb]{0.0,0.0,0.5019608}
O2:  \gEs{10 l  O2 , bei
Zyanose}\footnote{\gT{Zyanose}: Bl\"auliche Verf\"arbung
der Haut infolge einer unzureichenden Sauerstoffversorgung. Besonders
auff\"allig an den Lippen, Gesicht, und
Extremit\"aten.}\gEs{  15 l}}

\item {\mdseries\color[rgb]{0.0,0.0,0.5019608}
\gEs{Reanimationsbereichtschaft!}}

\item {\mdseries\color[rgb]{0.0,0.0,0.5019608}
\gEs{Notarzt! }}

\item {\mdseries\color[rgb]{0.0,0.0,0.5019608}
Engmaschiges Monitoring (Atmung, Pulsoxymetrie, RR)}

\item {\mdseries\color[rgb]{0.0,0.0,0.5019608}
W\"armeerhalt}

\item {\mdseries\color[rgb]{0.0,0.0,0.5019608}
\gEs{Ursachenforschung} -- Warum ist der Patient
bewusstlos oder sopor\"os?}

\item {\mdseries\color[rgb]{0.0,0.0,0.5019608}
Neurcheck}

\end{itemize}
\begin{itemize}
\item {\mdseries\color[rgb]{0.0,0.0,0.5019608}
\gEs{Immer auch an BZ denken!}}

\end{itemize}
\begin{itemize}
\item {\mdseries\color[rgb]{0.0,0.0,0.5019608}
Traumacheck}

\end{itemize}
\subsection{Atmung}
{\mdseries\color{black}
Kurze Wiederholung}

{\mdseries
normale Atemfrequenz: (Erwachsener): 12-15x / min, AZV = 5-7 ml / kg KG}

{\mdseries
Atemminutenvolumen (AMV): AZV x AF}

{\mdseries
\index{Totraumventilation}\gT{Totraumventilation}:
Luftmenge in oberen  und teilweise unteren Atemwegen, die  nicht am
Gasaustausch teilnimmt. Achtung
\index{Schnappatmung}\gT{Schnappatmung}!Dabei wird nur
so wenig Luft bewegt sodass aufgrund des Totraumes keine frische Luft
in die Lungenbl\"aschen (Alveolen) kommt. H\"aufig wird die
Schnappatmung mit flacher Atmung verwechselt, aber:}

{\centering\color{black}
Schnappatmung = keine Atmung
\par}

{\mdseries\color[rgb]{0.0,0.0,0.5019608}
St\"orungen der Atmung}

{\mdseries
Ateminsuffizienz = eingeschr\"ankte Atmung ={\textgreater} O2-Mangel
(Hypoxie)   }

{\mdseries
Atemstillstand = Aussetzen d. Atmung (Apnoe)}

{\mdseries\color{black}
Ursachen:}

\begin{itemize}
\item \end{itemize}
\begin{center}
\tablehead{}\begin{supertabular}{|m{2.71cm}|m{10.59cm}|}
\hline
\centering {\bfseries\upshape\color[rgb]{0.0,0.0,0.5019608} St\"orung
des/der}\par

 &
\\\hline
O2-Angebotes

 &
\begin{itemize}
\item {\mdseries G\"arkeller (CO2 ${\uparrow}$, O2 ${\downarrow}$)}

\item {\mdseries {\quotedblbase}d\"unne Luft{\textquotedblleft} (z.B.
Hochgebirge, pO2 ${\downarrow}$)}

\item {\mdseries Defekte Gastherme (CO ${\uparrow}$m O2${\downarrow}$)}

\item {\mdseries Ertrinken}

\end{itemize}
\\\hline
O2-Diffusion

 &
\begin{itemize}
\item {\mdseries Lungen\"odem}

\item {\mdseries Lungenentz\"undung}

\item {\mdseries Lungenembolie}

\end{itemize}
\\\hline
Atemmechanik

 &
\begin{itemize}
\item {\mdseries Verlegung der }

\end{itemize}
\begin{itemize}
\item {\mdseries oberen Atemwege (Zunge, Erbrochenes, Laryngospasmus,
Glottis\"odem, Bolus)}

\item {\mdseries unteren Atemwege (Bronchitis, Asthma, Lungen\"odem)}

\end{itemize}
\begin{itemize}
\item {\mdseries Verminderung der Dehnbarkeit der Thoraxwand oder des
Lungengewebes (Rippenfraktur, Pneumothorax, Pleuraergu{\ss}, Emphysem)}

\end{itemize}
\\\hline
Neuromuskul\"aren Atemregulation

 &
\begin{itemize}
\item {\mdseries SHT (Sch\"adel Hirn Trauma)}

\item {\mdseries Schlaganfall (Insult)}

\item {\mdseries Vergiftungen}

\item {\mdseries Tumore}

\item {\mdseries hoher Querschnitt (RM durchtrennt)}

\item {\mdseries Muskelerkrankungen}

\item {\mdseries Nervenerkankung}

\item {\mdseries Stoffwechselst\"orungen}

\end{itemize}
\\\hline
\end{supertabular}
\end{center}
{\mdseries\color{black}
Symptome}

\begin{center}
\begin{tabular}{|m{3.225cm}|m{3.225cm}|m{3.225cm}|m{3.225cm}|}
\hline
\centering {\bfseries\upshape\color[rgb]{0.0,0.0,0.5019608} Gruppe}\par

 &
\centering {\bfseries\upshape\color[rgb]{0.0,0.0,0.5019608} Befund}\par

 &
\centering {\bfseries\upshape\color[rgb]{0.0,0.0,0.5019608}
Beschreibung}\par

 &
\centering {\bfseries\upshape\color[rgb]{0.0,0.0,0.5019608}
Anmerkung}\par

\\\hline
{\mdseries\color{black} Dyspnoe}

 &
{\mdseries Atemnot}

 &
{\mdseries Leitsymptom}

 &
\\\hline
{\mdseries\color{black} Atemger\"ausch}

 &
{\mdseries Brodeln}

 &
{\mdseries Blubbern, klassisch f\"ur Lunden\"odem}

 &
\\\hline
 &
{\mdseries Stridor}

 &
 &
\\\hline
 &
{\mdseries \gZ{Brummen, Giemen}}

 &
 &
\\\hline
 &
{\mdseries \gZ{Rasselger\"ausche}}

 &
 &
\\\hline
{\mdseries\color{black} Frequenz}

 &
{\mdseries Beschleunigt}

 &
{\mdseries Tachypnoe}

 &
\\\hline
 &
{\mdseries Verlangsamt}

 &
{\mdseries Bradypnoe}

 &
\\\hline
{\mdseries\color{black} Atemzugsvolumen}

 &
{\mdseries Schnappatmung}

 &
 &
\\\hline
 &
{\mdseries Flache Atmung}

 &
 &
\\\hline
 &
{\mdseries Tiefe Atmung}

 &
 &
\\\hline
{\mdseries\color{black} Hautfarbe}

 &
{\mdseries Blass}

 &
{\mdseries Normal, An\"amie, Blutverlust}

 &
\\\hline
 &
{\mdseries Rosig}

 &
{\mdseries Normal, CO-Vergiftung}

 &
\\\hline
 &
{\mdseries Bl\"aulich }

 &
{\mdseries Zyanose: Hypoxie!}

 &
\\\hline
{\mdseries\color{black} K\"orperlich}

 &
{\mdseries Einziehungen an den Rippen}

 &
{\mdseries Einsatz der Atemhilfsmuskulatur}

 &
\\\hline
 &
{\mdseries Aufst\"utzen}

 &
{\mdseries Einsatz der Atemhilfsmuskulatur}

 &
\\\hline
 &
{\mdseries Nasenfl\"ugeln}

 &
{\mdseries Atemnot, besonders bei Kleinkindern}

 &
\\\hline
 &
{\mdseries Aufrechte Position}

 &
 &
\\\hline
 &
{\mdseries {\quotedblbase}Paradoxe Atmung{\textquotedblleft}}

 &
{\mdseries Brustkorb senkt sich bei Einatmung}

 &
\\\hline
\end{tabular}
\end{center}
{\mdseries\color{black}
\textcolor[rgb]{0.7882353,0.0,0.08627451}{[F0EA?]} Gefahr der
Ateminsuffizienz}

\begin{itemize}
\item {\mdseries
Hypoxie}

\end{itemize}
\begin{itemize}
\item {\mdseries
in weiterer Folge Unterversorgung lebenswichtiger Organe wie Herz,
Hirn.}

\end{itemize}
\begin{itemize}
\item {\mdseries
Atmungsbedingte \"Ubers\"auerung (Respiratorische Azidose, s. H.IV )}

\end{itemize}
{\mdseries\color{black}
Ma{\ss}nahmen bei Atemstillstand}

\begin{itemize}
\item {\mdseries\color[rgb]{0.0,0.0,0.5019608}
Freimachen der Atemwege}

\end{itemize}
\begin{itemize}
\item {\mdseries\color[rgb]{0.0,0.0,0.5019608}
\"Offnen des Mundes (Esmarch-Handgriff)}

\item {\mdseries\color[rgb]{0.0,0.0,0.5019608}
Inspektion}

\item {\mdseries\color[rgb]{0.0,0.0,0.5019608}
fals n\"otig ausr\"aumen }

\item {\mdseries\color[rgb]{0.0,0.0,0.5019608}
(falls n\"otig) Absaugen}

\item {\mdseries\color[rgb]{0.0,0.0,0.5019608}
\"Uberstrecken der HWS}

\item {\mdseries\color[rgb]{0.0,0.0,0.5019608}
Zungengrund wird angehoben -{\textgreater} Atemwege  sind frei!}

\end{itemize}
\begin{itemize}
\item {\mdseries\color[rgb]{0.0,0.0,0.5019608}
\"Uberpr\"ufung der Atemt\"atigkeit (bei  \"uberstreckter HWS) durch}

\end{itemize}
\begin{itemize}
\item {\mdseries\color[rgb]{0.0,0.0,0.5019608}
SEHEN, H\"OREN, F\"UHLEN f\"ur mindestens 10  Sekunden!}

\end{itemize}
\begin{itemize}
\item {\mdseries\color[rgb]{0.0,0.0,0.5019608}
Sauerstoffgabe (wenn ausreichende Spontan-Atmung vorhanden)}

\item {\mdseries\color[rgb]{0.0,0.0,0.5019608}
Beatmung (keine ausreichende Spontanatmung bzw. bei Reanimation)}

\end{itemize}
\begin{itemize}
\item {\mdseries\color[rgb]{0.0,0.0,0.5019608}
Beatmungsbeutel mit Maske und Reservoir}

\end{itemize}
{\mdseries
Vgl. die entsprechenden Kapitel im Teil Erste Hilfe und
Sanit\"atstechnik bez\"uglich der genauen Durchf\"uhrung der
Ma{\ss}nahmen. }

{\mdseries
Achtung: Bei einem mit Atemstillstand vorgefundenen Patienten ist
grunds\"atzlich von einer Reanimation auszugehen. }

\subsection[Herz- / Kreislauf]{Herz- / Kreislauf}
\label{bkm:RefHeading42514202}{\mdseries
St\"orungen des Herz-/ Kreislaufsystems}

{\mdseries\color[rgb]{0.0,0.0,0.5019608}
\label{bkm:RefHeading42514602}Kreislaufstillstand}

{\mdseries
Kreislaufstillstand: Pumpversagen des Herzens aufgrund von}

{\mdseries
\gEs{Asystolie} (keine Herzt\"atigkeit)}

{\mdseries
\gEs{Kammerflimmern} (unkoordinierte
Herzt\"atigkeit)}

{\mdseries
\gEs{(Pulslose) Ventrikul\"are Tachykardie} (HF zu
schnell ${\rightarrow}$ diastolische F\"ullung nicht ausreichend
${\rightarrow}$ Auswurf ${\downarrow}$)}

{\mdseries
\gEs{elektromechanische Entkopplung} }

{\mdseries\color{black}
Ursachen}

\begin{itemize}
\item {\mdseries
Minderversorgung des Herzmuskels mit O2}

\end{itemize}
\begin{itemize}
\item {\mdseries
Herzinfarkt}

\item {\mdseries
Herzinsuffizienz}

\end{itemize}
\begin{itemize}
\item {\mdseries
St\"orungen der Atmung}

\end{itemize}
\begin{itemize}
\item {\mdseries
Asthma bronchiale}

\item {\mdseries
Lungenembolie}

\item {\mdseries
zentrale Atemregulationsst\"orungen}

\end{itemize}
\begin{itemize}
\item {\mdseries
sonstige St\"orungen}

\end{itemize}
\begin{itemize}
\item {\mdseries
Trauma}

\item {\mdseries
Vergiftungen}

\item {\mdseries
thermische Sch\"aden}

\item {\mdseries
Elektrolytverschiebungen,...}

\end{itemize}
{\mdseries\color{black}
Pathophysiologie :}

\begin{itemize}
\item {\mdseries
Aussetzen des Herzschlages ={\textgreater} Patient ist sofort pulslos!}

\item {\mdseries
Bewu{\ss}tlosigkeit (nach 15-20 sek.)}

\item {\mdseries
Atemstillstand/Schnappatmung (nach ca. 10 sek.)}

\item {\mdseries
\gEs{Irreversible Hirnsch\"aden nach 3-5 min.}, bei
K\"alte allerdings l\"anger!}

\end{itemize}
\paragraph{Tod}
{\mdseries\color{black}
\gCa{Beachte:}}

{\centering\color{black}
\gEs{[F0EA?] }\gEs{Nobody is dead
until warm and dead!}
\par}

{\mdseries\color{black}
Klinischer Tod vs. biologischer Tod}

\begin{center}
\tablehead{}\begin{supertabular}{|m{6.649cm}|m{6.651cm}|}
\hline
\centering {\bfseries\upshape\color[rgb]{0.0,0.0,0.5019608} klinischer
Tod}\par

 &
\centering {\bfseries\upshape\color[rgb]{0.0,0.0,0.5019608} biologischer
Tod}\par

\\\hline
{\mdseries keine Atmung}

{\mdseries kein Puls}

{\mdseries Reanimation mu{\ss} innerhalb der ersten  Minuten erfolgen}

{\mdseries (dann unter Umst\"anden) reversibel}

 &
{\mdseries Organtod (irreversibel) des Gehirns}

{\mdseries gleichzusetzen mit Tod des Individuums}

\\\hline
\end{supertabular}
\end{center}
{\mdseries\color{black}
Wann ist es wirklich zu sp\"at?}

{\mdseries
Todesfeststellung: Grunds\"atzlich NUR durch den Arzt!}

{\bfseries\color[rgb]{0.5019608,0.0,0.0}
Todeszeichen: Nicht wie von Roman gefordert in Tabelle mit mit sicheren
und unsicheren Todeszeichen umgebaut, halte ich f\"ur un\"ubersichtlich
und bef\"urchte Verwechslungsgefahr mit Knochenbruchzeichen.[GAB]}

{\mdseries
Unter bestimmten Voraussetzungen allerdings auch durch den Sanit\"ater
(Unter\-lassung von  Reanimationsma{\ss}nahmen):}

\begin{center}
 [Warning: Image ignored] % Unhandled or unsupported graphics:
%\includegraphics[width=4.19cm,height=3.028cm]{AASdev20090228-img66.png}

\end{center}
\begin{itemize}
\item {\mdseries
%
%
Verwesung}

\item {\mdseries
Verletzungen, die nicht mit dem Leben  vereinbar sind}

\end{itemize}
\begin{itemize}
\item {\mdseries
Kopfamputation, {\dots}}

\end{itemize}
\begin{itemize}
\item {\mdseries
(Totenflecken (nicht 100\%ig feststellbar))}

\item {\mdseries
(Leichenstarre (nicht 100\%ig feststellbar))}

\item {\mdseries
Abbruch der Reanimationsma{\ss}nahmen aufgrund von Ersch\"opfung}

\end{itemize}


\begin{center}
 [Warning: Image ignored] % Unhandled or unsupported graphics:
%\includegraphics[width=7.133cm,height=4.837cm]{AASdev20090228-img67.jpg}

\end{center}
{\mdseries\color[rgb]{0.0,0.0,0.5019608}
Herzinsuffizienz, Herzversagen}

{\mdseries
Weitere, oft auch sehr bedrohliche, St\"orungen des Herzens werden im
Kapitel 
\hyperlink{a84HerzKreislaufstrungenoutline}{Herz-Kreislaufst\"orungen} 
(I.I{}-\pageref{bkm:RefHeading42342802}:
{\quotedblbase}Herz-Kreislaufst\"orungenunten{\quotedblbase})
besprochen.}

{\mdseries\color[rgb]{0.0,0.0,0.5019608}
Temperaturregulation [F0E1?]}

{\mdseries
W\"armeproduktion durch}

\begin{itemize}
\item {\mdseries
Zellaktivit\"at}

\item {\mdseries
Stoffwechselvorg\"ange}

\item {\mdseries
Bewegung (Muskelatmung)}

\end{itemize}
{\mdseries
W\"armeabgabe \"uber:}

\begin{itemize}
\item {\mdseries
Haut}

\item {\mdseries
Schwei{\ss}}

\end{itemize}
\begin{itemize}
\item {\mdseries
Wirkung durch Verdunstungsk\"alte}

\item {\mdseries
ca. 0,5 l / Tag}

\item {\mdseries
bis max. 1,5 l / Std. bei schwerer k\"orperlicher  Arbeit}

\end{itemize}
{\mdseries
physiologische Kerntemperatur: 37{\textdegree}C}

{\mdseries
Regulation durch den Hypothalamus}

{\mdseries\color[rgb]{0.0,0.0,0.5019608}
Stoffwechsel}

{\mdseries\color{black}
 Definition }

{\mdseries
Abl\"aufe von Stoffumsetzung und  Energiegewinnung}

{\mdseries
N\"ahrstoffe werden in W\"arme und Energie  umgewandelt}

{\mdseries\color{black}
Arten}

\begin{itemize}
\item {\mdseries
Zuckerstoffwechsel (Kohlenhydrate)}

\end{itemize}
\begin{itemize}
\item {\mdseries
Energielieferant }

\end{itemize}
\begin{itemize}
\item {\mdseries
Fettstoffwechsel (Fetts\"auren, ges\"attigt und  unges\"attigt)}

\end{itemize}
\begin{itemize}
\item {\mdseries
Doppelt so hoher Energiegehalt wie  Kohlenhydrate}

\item {\mdseries
L\"osungsmittel f\"ur Vitamine}

\item {\mdseries
Bausteine der Zellmembran}

\item {\mdseries
Bestandteil von Hormonen}

\end{itemize}
\begin{itemize}
\item {\mdseries
Eiwei{\ss}stoffwechsel (Proteine)}

\end{itemize}
\begin{itemize}
\item {\mdseries
Zellmembranbausteine}

\item {\mdseries
Enzyme!}

\end{itemize}
\subsection[Wasser- und Elektro\-lyt\-haus\-halt]{%
%
Wasser- und Elektro\-lyt\-haus\-halt}
{\mdseries
Das K\"orperwasser macht ca. 60\% des K\"orpergewichtes aus und teilt
sich auf:}

\begin{itemize}
\item {\mdseries
\gT{IntraZellularRaum} (IZR), 40\% KG }

\item {\mdseries
\gT{ExtraZellularRaum} (EZR), 20\% KG }

\end{itemize}
\begin{itemize}
\item {\mdseries
\index{Zwischengewebsraum}Zwischengewebsraum, \~{}16\%:
{\quotedblbase}freier{\textquotedblleft} Raum zwischen den Zellen im
Gewebe}

\item {\mdseries
\index{Gef\"a{\ss}raum}Gef\"a{\ss}raum \~{}4\%: Raum in den
Blutgef\"a{\ss}en}

\end{itemize}
{\mdseries\color{black}
Trennung: }

{\mdseries
IZR -- EZR:    \gT{Zellmembran} (\textit{semipermeable}
Membran, siehe unten)}

{\mdseries
Zwischengewebsraum -- Gef\"a{\ss}raum:  Blutgef\"a{\ss}wand}

{\mdseries\color{black}
\index{Atome}Atome, \index{Ionen}Ionen und
\index{Elektrolyte}Elektrolyte }

{\mdseries
Atome sind die kleinsten Grundbestandteile von Marterie.}

{\mdseries
Ionen sind Atome die elektrisch geladen sind.}

{\mdseries
Geladene Ionen sind gut in Wasser l\"oslich. Wenn sie in Wasser gel\"ost
sind heissen sie {\quotedblbase}Elektrolyte{\textquotedblleft}. }

{\mdseries
Im Wasser gel\"oste, elektrisch geladene Teilchen hei{\ss}en
\gT{Elektrolyte}.}

{\mdseries
Die wichtigsten Ionen sind Natrium (Na+), Kalium (K+) (positiv geladen)
und Chlorid (Cl-)(negativ geladen). }

{\mdseries
Funktion:}

\begin{itemize}
\item {\mdseries
Osmose, Diffusion und Verschiebung von Wasser zwischen den einzelnen
R\"aumen (siehe unten)}

\end{itemize}
\begin{itemize}
\item {\mdseries
${\rightarrow}$ \gEs{Volumenregulation} des
K\"orpers, siehe unten.  }

\end{itemize}
\begin{itemize}
\item {\mdseries
Elektrische Ladung: Wie bei einer Batterie ist eine Zelle geladen (innen
Minuspol, aussen Pluspol). Das ist besonders wichtig f\"ur}

\end{itemize}
\begin{itemize}
\item {\mdseries
Herz (Erregungsleitung!)}

\item {\mdseries
Nerven}

\item {\mdseries
Muskel}

\end{itemize}
{\mdseries
\gZ{Elektrolyte sind geladene Teilchen, die  bei
Aufspaltung von S\"auren, Laugen  oder deren Salzen in w\"assriger
L\"osung  entstehen.}}

{\mdseries
\gZ{Kationen: positiv geladene Teilchen (K+ , Na+ ,Mg++
, Ca++ ,...)}}

{\mdseries
\gZ{Anionen: negativ geladene Teilchen (Cl- )}}

{\mdseries
\gZ{Ionenverteilung:}}

{\mdseries
\gZ{IZR: viel K+ , wenig Na+}}

{\mdseries
\gZ{EZR: viel Na+ , wenig K+}}

{\mdseries\color{black}
%
%
\index{Osmose}Osmose und \index{Diffusion}Diffusion}

{\bfseries\color[rgb]{0.5019608,0.0,0.0}
GAB: Osmose und Diffusion: Grunds\"atzlich wichtig und relevant,
allenfalls bessere Erkl\"arung, aber bleibt in Normalschrift, nicht
heller. [GAB]}

{\bfseries\color[rgb]{0.5019608,0.0,0.0}
KOCH: Sollte alles grau werden}

{\mdseries
Geladene Teilchen k\"onnen die Zellmembran nicht selbstst\"andig
durchdringen, Wasser dagegen schon. Man nennt die Zellmembran daher
eine {\quotedblbase}\index{semipermeable
Membran}\gT{semipermeable Membran}{\textquotedblleft}. }

{\mdseries
Die Natur m\"ochte dass in einer Fl\"ussigkeit \"uberall die gleiche
Konzentration von geladenen Teilchen herrscht. Da die Zellmembran (und
teilweise auch die Blutgef\"a{\ss}wand) nur das Wasser durchl\"asst,
nicht jedoch die geladenen Teilchen verschiebt sich nur das Wasser von
einem Raum in den anderen. Man kann sagen dass die Elekktrolyte einen
Druck auf das Wasser aus\"uben {\quotedblbase}die Seiten zu
wechseln{\textquotedblleft}. Man nennt diesen Druck den
{\quotedblbase}\index{osmotischer Druck}\gT{osmotischen
Druck}{\textquotedblleft}. Mann kann auch sagen {\quotedblbase}das
Wasser folgt dem Konzentrationsgef\"alle{\textquotedblleft}.}

\begin{center}
 [Warning: Image ignored] % Unhandled or unsupported graphics:
%\includegraphics[width=6.131cm,height=3.646cm]{AASdev20090228-img68.gif}

\end{center}
{\mdseries
\gZ{Daneben gibt es den {\quotedblbase}onkotischen
Druck{\textquotedblleft} der durch Eiwei{\ss}bestandteile ausge\"ubt
wird.}}

{\mdseries
Wie kommen die Elektrolyte in die einzelnen R\"aume wenn doch nur Wasser
die Wand passieren kann?}

{\mdseries
Daf\"ur gibt es aktive Transporter
({\quotedblbase}Pumpen{\textquotedblleft}) die die Ionen hin und her
bef\"ordern. Die Steuerung richtige Verteilung der verschiedenen Ionen
ist recht kompliziert aber lebenswichtig. Eine
{\quotedblbase}falsche{\textquotedblleft} Verteilung der Elektrolyte
kann schlimme Folgen haben (zB. Herzrhythmusst\"orungen). }

{\mdseries
Die richtige Verteilung der Elektrolyte im K\"orper ist lebenswichtig.
Elektrolyst\"orungen k\"onnen lebensgef\"ahrlich sein. }

{\mdseries\color{black}
H\"aufige Ursachen von \index{Elektrolyst\"orungen}Elektrolyst\"orungen}

\begin{itemize}
\item {\mdseries
Niereninsuffizienz (gest\"orte Ausscheidung}

\end{itemize}
\begin{itemize}
\item {\mdseries
Besonders bei Dialysepatienten}

\end{itemize}
\begin{itemize}
\item {\mdseries
Durchfall, Erbrechen }

\item {\mdseries
Medikamente}

\item {\mdseries
Vergiftung}

\end{itemize}
{\mdseries\color{black}
Wasser- und Elektrolythaushalt}

{\mdseries
Blutvolumen = Plasmawasser + zellul\"are Bestandteile
(Blutk\"orperchen)}

Beim Erwachsenen 7-8\% KG

Beim Neugeborenen / S\"augling 8-9\% KG  (70-80\%  K\"orperwasser!)

{\mdseries
Verluste bis 20\% ohne Funktionsbeeintr\"achtigung tolerierbar,
allerdings Sch\"aden in der  Mikrozirkulation}

{\mdseries
${\rightarrow}$ auch ohne klinische Schockzeichen  ist bei
Volumensverlust Volumensersatz  indiziert!}

\begin{itemize}
\item {\mdseries
osmotischer Druck}

\end{itemize}
\begin{itemize}
\item {\mdseries
Volumenregulation des K\"orpers}

\end{itemize}
\begin{itemize}
\item {\mdseries
Konzentrationsgef\"alle ${\rightarrow}$ Spannungsunterschied an den
Zellmembranen ={\textgreater} Grundlage f\"ur Erregungsleitung /
Muskelkontraktion}

\item {\mdseries
Regelung durch:}

\end{itemize}
\begin{itemize}
\item {\mdseries
Hormone \gZ{(Renin/Angiotensin/Aldosteron)}}

\item {\mdseries
Ausscheidung \"uber Niere}

\end{itemize}
{\mdseries\color[rgb]{0.0,0.0,0.5019608}
S\"aure-/Basenhaushalt }

{\mdseries
\index{S\"aure}\gT{S\"auren} = Verbindungen, die 
Wasserstoffionen (H\textsuperscript{+}) abgeben k\"onnen}

{\mdseries
\index{Lauge}\gT{Laugen} =
\index{Base}\gT{Basen} = Verbindungen, die
H\textsuperscript{+}  aufnehmen k\"onnen}

{\mdseries\color{black}
\index{pH-Wert}pH-Wert}

{\mdseries
Der pH-Wert gibt an wie sauer oder basisch eine L\"osung ist.}

{\mdseries
= Menge der freien H\textsuperscript{+}  in einer L\"osung.}

{\mdseries
pH 7 ist neutral, pH {\textgreater} 7 basisch, pH {\textless} 7 sauer}

{\mdseries
\gZ{pH normal ca. um 7,4.}}

{\mdseries
\gZ{pH {\textless} 7,35  ={\textgreater} Azidose}}

{\mdseries
\gZ{pH {\textgreater} 7,45  ={\textgreater} Alkalose}}

{\mdseries
Der richtige pH-Wert ist lebenswichtig. Der K\"orper ist da sehr
empfindlich, der pH-Wert muss sich in engen grenzen bewegen. Schon
geringe Abweichungen k\"onnen gro{\ss}e Folgen haben. }

{\mdseries
Leider gibt es im Rettungsdienst normalerweise keine M\"oglichkeit den
pH-Wert zu messen, jedoch haben wir gro{\ss}en Einflu{\ss} auf ihn
durch unsere Ma{\ss}nahmen, siehe unten. }

{\mdseries
Bei einer \"Ubers\"auerung sprechen wir von einer
\index{Azidose}\gT{Azidose}, wenn der Patient basisch
wird von einer \index{Alkalose}\gT{Alkalose}.}

\paragraph[\ Regulationsmechanismen]{ Regulationsmechanismen}
\begin{itemize}
\item {\mdseries
\textbf{Blut} mit \textbf{Puffersystem}}

\end{itemize}
\begin{itemize}
\item {\mdseries
 %
%
H\textsuperscript{+} + HCO\textsuperscript{3-}
${\leftarrow}$${\rightarrow}$ H\gSub{2 }O + CO\gSub{2
} (Formel nicht lernen!)}

\item {\mdseries
S\"aure \gZ{+ Bikarbonat} ${\leftarrow}$${\rightarrow}$
\gZ{Wasser +} CO\gSub{2}}

\end{itemize}
\begin{itemize}
\item {\mdseries
Atmet der Patient nicht (genug) staut sich das CO\gSub{2} und
bildet mit Wasser H\textsuperscript{+} ${\rightarrow}$
\textit{Atmungsbedingte \"Ubers\"auerung}}

\item {\mdseries
Atmet der Patient zu viel wird zu viel CO\gSub{2} abgeatmet und
wird von H\textsuperscript{+} und dem anderen Stoff nachgebildet
${\rightarrow}$ H\textsuperscript{+} fehlt ${\rightarrow}$
\textit{Atmungsbedingte Alkalose}}

\item {\mdseries
F\"allt zu viel H\textsuperscript{+} (S\"aure) im K\"orper an (durch den
Stoffwechsel, \textit{Stoffwechselbedingte \"Ubers\"auerung}) wird es
zu (Wasser und) CO\gSub{2} umgewandelt, das
CO\gSub{2} wird abgeatmet}

\end{itemize}
\begin{itemize}
\item {\mdseries
\textbf{Lunge} }

\end{itemize}
\begin{itemize}
\item {\mdseries
Abatmung von CO\gSub{2} ${\rightarrow}$ Einfluss auf den Puffer
(s.o.)}

\end{itemize}
\begin{itemize}
\item {\bfseries
Niere}

\end{itemize}
\begin{itemize}
\item {\mdseries
Ausscheidung (oder Wiederaufnahme)  von H\textsuperscript{+}}

\end{itemize}
\subparagraph[Atmungsbedingte \"Ubers\"auerung - Respiratorische
Azidose]{Atmungsbedingte \"Ubers\"auerung - Respiratorische Azidose}
\label{bkm:RefHeading35233321}\index{\"Ubers\"auerung}\index{Respiratorische
Azidose}{\mdseries
durch Hypoventilation / verminderten  Gasaustausch:}

{\mdseries
CO\gSub{2}  Abatmung vermindert ${\rightarrow}$ Hyperkapnie, 
Ans\"auerung des Blutes}

{\mdseries\color{black}
Ursachen}

{\mdseries
Asthma}

{\mdseries
Lungen\"odem}

{\mdseries
Atemstillstand / insuffiziente Atmung}

{\mdseries\color{black}
Therapie}

{\mdseries
Erh\"ohung des Atem-Minuten-Volumens}

\subparagraph[Stoffwechselbedingte \"Ubes\"auerung -Metabolische
Azidose]{Stoffwechselbedingte \"Ubes\"auerung -Metabolische Azidose}
\index{\"Ubes\"auerung}\index{Metabolische Azidose}{\mdseries
Zunahme von organischen S\"auren im  K\"orper}

{\mdseries\color{black}
Ursachen}

\begin{itemize}
\item {\mdseries
Nierenleistung }

\item {\mdseries
Schock}

\item {\mdseries
Durchfall}

\item {\mdseries
diabetisches Koma}

\end{itemize}
{\mdseries\color{black}
Gefahren}

\begin{itemize}
\item {\mdseries
Sch\"adigung des Zellstoffwechsels}

\item {\mdseries
Sch\"adigung des Herzmuskels}

\item {\mdseries
respiratorische und metabolische  Alkalose}

\end{itemize}
\subparagraph[Atmungsbedingte (respiratorische)
Alkalose]{Atmungsbedingte (respiratorische) Alkalose}
\index{Alkalose}{\mdseries
Hyperventilationssyndrom}

{\mdseries
Elektrolytverschiebung: Kribbeln in den Fingern, Pf\"otchenstellung,
Kr\"ampfe}

\subparagraph{Stoffwechselbedingte (metabolische) Alkalose}
{\mdseries
z.B. durch anhaltendes Erbrechen: Verlust von S\"aure}

{\mdseries\color[rgb]{0.0,0.0,0.5019608}
Zusammenfassung [F055?]}

{\mdseries
Wer zu wenig atmet wird sauer.}

{\mdseries
Wer zu viel atmet wird basisch und sp\"urt ein Kribbeln in den Fingern.
}

{\mdseries
Wer Probleme mit dem Stoffwechsel hat kann auch sauer oder basisch
werden.}

\clearpage\setcounter{page}{1}\subsection[SCHOCK ]{SCHOCK }
{\mdseries
[F0EF?][F068?][F0F0?][F0EB?][F0E9?][F0EA?][F0E1?]\textbf{\textcolor[rgb]{0.7882353,0.0,0.08627451}{[F078?]}}[F073?][F072?][F069?][F055?][F04F?][F04E?]}

{\mdseries\color{black}
Definition}

\gEs{{\quotedblbase}Lebensbedrohliche
Kreislaufst\"orung aufgrund von Unterversorgung des K\"orpers mit Blut
und Sauerstoff infolge des Versagens des
Kreislaufsystems.{\textquotedblleft} }

{\mdseries
Es k\"onnen einzelne oder mehrere Teile des Kreislaufsystems betroffen
sein \newline
(Blut, Gef\"a{\ss}e, Herz).}

{\mdseries
\gEs{Ob ein Schock oder eine Schockgefahr besteht
wird anhand von
}\gEs{Schocksymtpomen}\gEs{ und
anhand der
}\gEs{Verdachtsdiagnose}\gEs{
beurteilt.}}

{\mdseries\color{black}
Schock bedeutet:}

\begin{itemize}
\item {\mdseries
akute Kreislaufinsuffizienz = Mi{\ss}verh\"altnis
\gEs{Herzauswurf} - 
\gEs{Sauerstoffbedarf }des K\"orpers  }

\item {\mdseries
M\"ogliche Ursachen: }

\end{itemize}
\begin{itemize}
\item {\mdseries
Herz pumpt wenig}

\item {\mdseries
Zu wenig Blut vorhanden }

\item {\mdseries
Gef\"a{\ss}e sind zu weit gestellt und Blut versackt}

\end{itemize}
{\mdseries
Im Endeffekt kommt immer zu wenig Blut=Sauerstoff zu den Organen,
dadurch ist die Ge\-webedurchblutung nicht ausreichend  f\"ur die
Versorgung aller Organe mit Sauer\-stoff.}

{\mdseries
K\"orpereigene Mechanismen k\"onnen gegensteuern und versuchen den
Schock zu kompensieren (gegensteuern). Dies geht allerdings nur eine
gewisse Zeit und in einem gewissen Ausma{\ss} gut, daher ist eine
schnell  einsetzende Therapie wichtig! Bei l\"anger bestehendem Schock
kann dieser Zustand irreversibel werden. }

{\mdseries
\textbf{Kein} Schock besteht bei kurzzeitigen 
Bewu{\ss}tseinsst\"orungen wie Kollaps/Synkope (kurzfristige
Blutverteilungsst\"orung, kurzfristige zerebrale Isch\"amie, 
fl\"uchtige Symptome)}

{\mdseries\color{black}
Schock -  Kompensationsmechanismen}

{\mdseries
Der K\"orper versucht selber dagegen etwas zu tun und gegenzusteuern
({\quotedblbase}kompensieren{\textquotedblleft}). Der wichtigste
Mechanismus ist die
\index{Zentralisation}\gT{Zentralisation}. Darunter
versteht man die Blutumverteilung zugunsten lebenswichtger Organe wie
Herz, Gehirn, Lunge und Leber. Dabei werden z.B. Haut, Muskel und Darm
sowie die Extremit\"aten (=Peripherie) kaum mehr  durchblutet.
Periphere Venen kollabieren und die Rekapillarisierungszeit
verl\"angert. Wenn die Kompensationsma{\ss}nahmen  versagen spricht man
vom {\quotedblbase}\gE{dekompensierter
Schock}{\textquotedblleft}. Es besteht sp\"atestens dann akute
\gEs{Lebensgefahr}!\textbf{\textcolor[rgb]{0.7882353,0.0,0.08627451}{[F0EA?]}}
}

{\mdseries\color{black}
\gCa{Beachte:} }

\begin{center}
 [Warning: Image ignored] % Unhandled or unsupported graphics:
%\includegraphics[width=3.8cm,height=2.611cm]{AASdev20090228-img69.png}

\end{center}
{\centering\color{black}
Bei \gEs{jedem Patienten} ist zu erheben ob ein
\gEs{Schock} oder eine realistische
\gEs{Schockgefahr} besteht oder nicht bzw.
festzustellen um welche Art von Schock es sich handelt (s.u.)!
\par}

{\color{black}
\gCa{Achtung!}}

{\centering\color{black}
Jeder Schockzustand ist grunds\"atzlich ein Fall f\"ur den
\gEs{NOTARZT}!
\par}

{\mdseries\color{black}
Allgemeine Schocksymptomatik}

\begin{center}
\begin{tabular}{|m{3.9750001cm}m{9.325cm}|}
\hline
 &
\centering {\bfseries\upshape\color[rgb]{0.0,0.0,0.5019608}
Symptome}\par

\\\hline
\centering {\bfseries\color{red} \index{Schockgefahr}Schockgefahr}\par

\centering {\bfseries\color[rgb]{0.0,0.5019608,0.0} [263A?]}\par

 &
\begin{itemize}
\item {\mdseries Keine eindeutigen Anzeichen}

\item {\mdseries \gEs{Vorausschauend handeln!}}

\item {\mdseries Bedenke Verdachtsdiagnose. \newline
\gEs{Muss ich damit rechnen dass sich ein Schock
entwickelt? \newline
Besteht Schockgefahr?}}

\end{itemize}
\\\hline
\centering {\bfseries\color{red} Schock}\par

\centering {\bfseries\color{red} [E426?]}\par

 &
\begin{itemize}
\item {\mdseries HF ${\uparrow}$}

\item {\mdseries m\"a{\ss}iger RR-Abfall}

\item {\mdseries RR noch {\textgreater} HF}

\item {\mdseries neurologisch unauff\"allig, ev. Euphorisch oder
agitiert}

\item {\mdseries \gEs{Haut bla{\ss}, evt.
kaltschwei{\ss}ig}}

\item {\mdseries Venen sind kollabiert oder gestaut}

\item {\mdseries Peripherie schlecht durchblutet}

\end{itemize}
\\\hline
\centering {\bfseries\color{black} Schwerer Schock}\par

\centering {\bfseries\color{black} Vollbild}\par

\centering {\bfseries\color{black} [E427?]}\par

 &
\begin{itemize}
\item {\mdseries \gEs{Tachykardie} (HF ${\uparrow}$,
{\quotedblbase}fliehender{\textquotedblleft} Puls)}

\item {\mdseries \gEs{RR-Abfall} (schwacher Puls)}

\end{itemize}
\begin{itemize}
\item {\mdseries Herzfrequenz {\textgreater} systolischer RR}

\end{itemize}
\begin{itemize}
\item {\mdseries \gEs{Tachypnoe} (schnelle, flache
Atmung)}

\item {\mdseries neurologisch auff\"allig, getr\"ubt}

\item {\mdseries \gEs{Haut bla{\ss}}, zyanotisch
(bl\"aulich), feucht}

\item {\mdseries Venen kollabiert oder gestaut}

\item {\mdseries nur mehr schwer behebbar! ${\rightarrow}$ Schock 
m\"oglichst fr\"uhzeitig erkennen \&  behandeln}

\end{itemize}
\\\hline
\centering {\bfseries\color{black} Endstadium}\par

\centering {\bfseries\color{black} [E429?]}\par

 &
\begin{itemize}
\item {\mdseries Haut grau, klebrig}

\item {\mdseries Versagen der Herzfunktion}

\item {\mdseries RR kaum mehr me{\ss}bar}

\item {\mdseries evt. Koma}

\item {\mdseries evt. Schnappatmung}

\item {\mdseries Patient ist dabei zu sterben.}

\end{itemize}
\\\hline
\end{tabular}
\end{center}
{\mdseries\color{black}
Schockschere}

{\mdseries
Normalerweise ist die Herzfrequenz niedriger als der systolische
Blutdruck. Wenn der K\"orper versucht den Schock zu kompensieren kann
sich dieses Verh\"altnis umkehren. Man spricht dann von der
Schockschere, der Patient befindet sich dann bereits schon im Vollbild
eines Schocks! }

{\mdseries
Achtung: Kinder haben andere Normalwerte f\"ur HF und RR!}

[Warning: Draw object ignored]

{\mdseries
Wenn sich das Verh\"altnis Herzfrequenz zu Systolischem Blutdruck
umkehrt befindet sich der Patient bereits im Schock. }

{\mdseries\color[rgb]{0.0,0.0,0.5019608}
\label{bkm:RefHeading38104902}Allgemeine Ma{\ss}nahmen der
\index{Schockbek\"ampfung}Schockbek\"ampfung}

{\mdseries
Grundsatz: Der Schock muss bek\"ampft werden bevor man ihn bemerkt.}

\begin{itemize}
\item {\color[rgb]{0.0,0.0,0.5019608}
\gEs{Notarzt}!}

\item {\mdseries\color[rgb]{0.0,0.0,0.5019608}
Situationsgerechte  \gEs{Lagerung}, je nach Diagnose
und Schockart}

\item {\mdseries\color[rgb]{0.0,0.0,0.5019608}
Beengende Kleidung \"offnen}

\item {\mdseries\color[rgb]{0.0,0.0,0.5019608}
\gEs{Sauerstoffgabe} hochdosiert (ab 6l/min bei
Schockgefahr, 10l/min oder mehr bei bestehendem Schock, je nach
Schwere%
%
)}

\end{itemize}
{\bfseries\color[rgb]{0.5019608,0.0,0.0}
Freigabe: Sauerstoffgabe hochdosiert (ab 6l/min bei Schockgefahr,
10l/min oder mehr bei bestehendem Schock, je nach Schwere}

\begin{itemize}
\item {\mdseries\color[rgb]{0.0,0.0,0.5019608}
Monitoring (Pulsoxy, RR-Manschette angelegt lassen)}

\item {\mdseries\color[rgb]{0.0,0.0,0.5019608}
\textbf{Reanimationsbereitschaft} }

\item {\mdseries\color[rgb]{0.0,0.0,0.5019608}
W\"armeerhalt}

\item {\mdseries\color[rgb]{0.0,0.0,0.5019608}
Psychische Betreuung}

\item {\color[rgb]{0.0,0.0,0.5019608}
Weitere Massnahmen je nach Schockart }

\end{itemize}
{\mdseries\color[rgb]{0.0,0.0,0.5019608}
Schock - Arten}

\begin{center}
\begin{tabular}{|m{4.8430004cm}|m{8.457cm}|}
\hline
\centering {\bfseries\upshape\color[rgb]{0.0,0.0,0.5019608}
Mechanismus}\par

 &
\centering {\bfseries\upshape\color[rgb]{0.0,0.0,0.5019608}
Schockart}\par

\\\hline
{\bfseries Herz pumpt zu wenig}

 &
\begin{itemize}
\item {\mdseries \gEs{Kardiogender Schock}}

\end{itemize}
\\\hline
{\bfseries Zu wenig Blut vorhanden }

 &
\begin{itemize}
\item {\mdseries \gEs{Hypovol\"amischer Schock}}

\end{itemize}
\\\hline
{\bfseries Gef\"a{\ss}e inad\"aquat weitgestellt, Blut versackt}

 &
\begin{itemize}
\item {\mdseries \gEs{Anaphylaktischer Schock}}

\item {\mdseries \gEs{Septischer Schock, toxischer
Schock}}

\item {\mdseries \gEs{Neurogener Schock}}

\end{itemize}
\\\hline
\end{tabular}
\end{center}
\paragraph{Hypovol\"amischer Schock}
{\mdseries
Schock aufgrund eines Verlustes von Blutvolumen aus den Gef\"a{\ss}en.}

{\mdseries
Es kann dabei verloren gehen: }

\begin{itemize}
\item {\mdseries
Blut}

\item {\mdseries
Plasma}

\item {\mdseries
Wasser und Elektrolyte}

\end{itemize}
{\mdseries\color{black}
Ursachen}

\begin{center}
\tablehead{}\begin{supertabular}{|m{6.649cm}|m{6.651cm}|}
\hline
{\mdseries Fl\"ussigkeitsverlust \gEs{nach
au{\ss}en}}

 &
{\mdseries Fl\"ussigkeitsverlust \gEs{nach innen}}

\\\hline
\begin{itemize}
\item {\mdseries Blutung}

\item {\mdseries gastrointestinal (Durchfall, Diarrhoe )}

\item {\mdseries Verbrennungen (Hautbarriere versagt 
gro{\ss}\-fl\"achig)}

\item {\mdseries renal (Diabetes mellitus, Diuretika, Polyurie bei 
Nierenversagen)}

\end{itemize}
 &
\begin{itemize}
\item {\mdseries Weichteilblutung}

\end{itemize}
\begin{itemize}
\item {\mdseries Nach Trauma, rupturierte Eileiter\-schwanger\-schaft,
etc}

\end{itemize}
\begin{itemize}
\item {\mdseries Wassereinlagerung ins Gewebe}

\item {\mdseries \gZ{Aszites
({\quotedblbase}Wasserbauch{\textquotedblleft} bei Lebererkrankungen 
etc.)}}

\item {\mdseries Ileus (Darmverschlu{\ss})}

\item {\mdseries H\"amatothorax, H\"amatoperiton\"aum}

\end{itemize}
\\\hline
\end{supertabular}
\end{center}
{\mdseries
Bei akutem Fl\"ussigkeitsverlust kann ein Verlust von
\gEs{bis zu 20\% des Blutvolumens}  gut kompensiert
werden. Auf gr\"o{\ss}ere Verluste  folgt ein Absinken von Herzauswurf
und RR!}

{\mdseries\color{black}
Spezielle Ma{\ss}nahmen}

\begin{itemize}
\item {\mdseries\color[rgb]{0.0,0.0,0.5019608}
Allgemein:
\hyperlink{a761AllgemeineManahmenderSchockbekmpfungoutline}{Allgemeine
Ma{\ss}nahmen der Schockbek\"ampfung} (siehe oben) }

\item {\mdseries\color[rgb]{0.0,0.0,0.5019608}
Blutstillung  (weitere Verluste verhindern!)}

\item {\mdseries\color[rgb]{0.0,0.0,0.5019608}
ggfs. Trauma-Vorgehen beachten}

\item {\mdseries\color[rgb]{0.0,0.0,0.5019608}
%
%
Wenn m\"oglich schonende Bergung}

\item {\mdseries\color[rgb]{0.0,0.0,0.5019608}
Lagerung: grunds\"atzlich %
%
Beine hoch \newline
Wichtige Ausnahmen: {\textperiodcentered} kardiale Ursache,
{\textperiodcentered} Kopf- / Brust\-ver\-letzungen,
{\textperiodcentered} Hirndruck (SHT),  {\textperiodcentered}
Knochenbr\"uche der Beine oder des Beckens, {\textperiodcentered}
Bewusstlosigkeit\newline
${\rightarrow}$ sonst F%
%
lachlagerung oder stabile Seitenlage}

\item {\mdseries\color[rgb]{0.0,0.0,0.5019608}
evt. Aviso (Ank\"undigung) im  Krankenhaus / Schockraum}

\item {\mdseries\color[rgb]{0.0,0.0,0.5019608}
Abtransport: rasch, schonend, nicht \"uberst\"urzt}

\end{itemize}
{\bfseries\color[rgb]{0.5019608,0.0,0.0}
Begriff {\quotedblbase}Schocklagerung{\textquotedblleft} ein f\"ur alle
mal in der Ausbildung streichen, nur Hinweis dass dieser Ausdruck
veraltet ist sollten sie ihn wo aufschnappen. Stattdessen
{\quotedblbase}Beine-hoch{\textquotedblleft}-Lagerung.}

{\bfseries\color[rgb]{0.5019608,0.0,0.0}
Wenn KI gegen Beine hoch: Flachlagerung bzw. stab. Seitenlage
empfehlen?}

\paragraph{Kardiogener Schock}
{\mdseries
Schock infolge mangelnder Herzpumpleistung.}

\begin{center}
\tablehead{}\begin{supertabular}{|m{6.649cm}|m{6.651cm}|}
\hline
{\mdseries\color{black} Ursachen}

\begin{itemize}
\item {\mdseries \gEs{Herzversagen},
Herzinsuffizienz}

\end{itemize}
\begin{itemize}
\item {\mdseries Myokardinfarkt (h\"aufigste Ursache)}

\item {\mdseries Herzrhythmusst\"orungen}

\end{itemize}
\begin{itemize}
\item {\mdseries \gEs{Lungenembolie}}

\end{itemize}
 &
{\mdseries\color{black} charakterisiert durch}

\begin{itemize}
\item {\mdseries RR-Abfall}

\item {\mdseries eingeschr\"ankte Pumpleistung}

\end{itemize}
\begin{itemize}
\item {\mdseries Blut staut sich zur\"uck}

\end{itemize}
\\\hline
\end{supertabular}
\end{center}
{\mdseries\color{black}
Spezielle Ma{\ss}nahmen}

\begin{itemize}
\item {\mdseries\color[rgb]{0.0,0.0,0.5019608}
Allgemein:
\hyperlink{a761AllgemeineManahmenderSchockbekmpfungoutline}{Allgemeine
Ma{\ss}nahmen der Schockbek\"ampfung} (siehe oben) }

\item {\mdseries\color[rgb]{0.0,0.0,0.5019608}
Lagerung: \gEs{Oberk\"orper hoch!} Evt. Beine
h\"angen  lassen. }

\item {\mdseries\color[rgb]{0.0,0.0,0.5019608}
medikament\"ose Therapie durch \gEs{Notarzt} 
\gZ{(Diuretika, Katecholamine, Heparin)}}

\item {\mdseries\color[rgb]{0.0,0.0,0.5019608}
\gEs{Reanimationsbereitschaft}  }

\end{itemize}
\paragraph{Anaphylaktischer Schock}
{\mdseries
Schock aufgrund einer \"uberschie{\ss}enden allergischen Reaktion.}

{\mdseries
{\quotedblbase}Allergischer Schock{\textquotedblleft}. Ausgel\"ost durch
\"uberschie{\ss}ende Reaktion der K\"orerabwehr (Immunreaktion) auf
bestimmte Stoffe (Antigene). }

\begin{center}
 [Warning: Image ignored] % Unhandled or unsupported graphics:
%\includegraphics[width=2.641cm,height=3.837cm]{AASdev20090228-img70.jpg}

\end{center}
{\mdseries
Es werden s\"amtliche \gEs{Gef\"a{\ss}e maximal
erweitert}, zus\"atzlich steigt die
\gEs{Durchl\"assigkeit der Gef\"a{\ss}e}
\gZ{(Gef\"a{\ss}permeabilit\"at ${\uparrow}$)}
${\rightarrow}$  Plasmaverlust in den Zwischengebwebsraum
\gZ{(Interstitium)} }

{\mdseries
\gEs{RR f\"allt massiv ab!}}

{\mdseries\color{black}
M\"ogliche Antigene}

\begin{itemize}
\item {\mdseries
Pollen}

\item {\mdseries
Medikamente, Antibiotika}

\item {\mdseries
Kontrastmittel}

\item {\mdseries
tierische Gifte (Bienenstich etc.)}

\item {\mdseries
und vieles andere ...}

\end{itemize}
{\mdseries\color{black}
Symptome%
%
}

\begin{center}
\tablehead{}\begin{supertabular}{|m{3.31cm}|m{9.99cm}|}
\hline
\centering {\bfseries\upshape\color[rgb]{0.0,0.0,0.5019608} Stadium}\par

 &
\\\hline
{\bfseries Stadium I:}

 &
{\mdseries Hautreaktion: evt. Ausschlag, R\"otung;
{\quotedblbase}Quaddeln{\textquotedblleft}, Flush, \"Odeme}

{\mdseries Allgemeinsymptome: Juckreiz, Zittern, Schwindel, Kopfschmerz}

\\\hline
{\bfseries Stadium II:}

 &
{\mdseries RR ${\downarrow}$  , Tachykardie}

{\mdseries \"Ubelkeit, Erbrechen, Durchfall }

\\\hline
{\bfseries Stadium III:}

 &
{\mdseries Allgemeine Schocksymptome}

{\mdseries Bronchospasmus}

\\\hline
{\bfseries Stadium IV:}

 &
{\mdseries Atem-/Kreislaufstillstand}

{\mdseries Schwerer Schock}

\\\hline
\end{supertabular}
\end{center}
{\mdseries\color{black}
Spezielle Ma{\ss}nahmen [F0EF?][F068?][F0F0?]}

\begin{itemize}
\item {\mdseries\color[rgb]{0.0,0.0,0.5019608}
Allgemein:
\hyperlink{a761AllgemeineManahmenderSchockbekmpfungoutline}{Allgemeine
Ma{\ss}nahmen der Schockbek\"ampfung} (siehe oben) }

\item {\mdseries\color[rgb]{0.0,0.0,0.5019608}
Antigen  nach M\"oglichkeit entfernen!}

\item {\mdseries\color[rgb]{0.0,0.0,0.5019608}
Lagerung: situationsgerecht.\newline
Druckverlust im Vordergrund: Beine hoch.%
%
\newline
Atemprobleme im Vordergrund: stark erh\"ohter Oberk\"orper}

\end{itemize}
\paragraph[Septischer Schock und Toxischer Schock]{Septischer Schock und
Toxischer Schock}
\label{bkm:RefHeading41284910}{\mdseries
Siehe auch: Kap. \textit{[RS II] Hygiene }\textit{ {\guillemotright}
}\textit{Sepsis}\textit{
[}\textit{K.II}\textit{{\guillemotright}}\textit{\pageref{bkm:RefHeading41252510}}\textit{,
}\textit{unten}\textit{]}}

{\mdseries
Schock infolge einer schweren, k\"orperweiten Infektion oder aufgrund
eines Gef\"a{\ss}-beeinflussenden Giftstoffes (Toxins).}

{\mdseries
Durch bakterielle Gef\"a{\ss}wandsch\"adigung, temperaturbedingt (hohes
Fieber {\textgreater} 38{\textdegree}C) sowie aufgrund von Giftstoffen
kommt es zu Fl\"ussigkeitsaustritt in den Zwischen\-gewebs\-raum sowie
zu Gef\"a{\ss}weitstellung.}

{\mdseries\color{black}
Ursachen}

\begin{center}
\tablehead{}\begin{supertabular}{|m{6.649cm}|m{6.651cm}|}
\hline
{\mdseries Diverse Erreger und Toxine}

\begin{itemize}
\item {\mdseries Langanhaltender Entz\"undungsherd der pl\"otzlich
Erreger {\quotedblbase}ins Blut streut{\textquotedblleft} und
{\quotedblbase}explodiert{\textquotedblleft}}

\item {\mdseries Besiedelte Fremdk\"orper (Diverse medizinische Sonden,
Venflons, etc.), siehe \textit{Sepsis}\textit{
[}\textit{K.II}\textit{{\guillemotright}}\textit{\pageref{bkm:RefHeading41252510}}\textit{,
}\textit{unten}\textit{]}}

\item {\mdseries Tampons}

\item {\mdseries ...}

\end{itemize}
 &
{\mdseries\color{black} Symptome}

\begin{itemize}
\item {\mdseries Schock}

\item {\mdseries hohes Fieber oder kein Fieber}

\item {\mdseries Haut hei{\ss}, trocken, rot oder bleich, grau,
{\quotedblbase}marmoriert{\textquotedblleft}}

\end{itemize}
\\\hline
\end{supertabular}
\end{center}
{\mdseries\color{black}
Spezielle Ma{\ss}nahmen}

\begin{itemize}
\item {\mdseries\color[rgb]{0.0,0.0,0.5019608}
Allgemein:
\hyperlink{a761AllgemeineManahmenderSchockbekmpfungoutline}{Allgemeine
Ma{\ss}nahmen der Schockbek\"ampfung} (siehe oben) }

\item {\mdseries\color[rgb]{0.0,0.0,0.5019608}
Lagerung: situationsgerecht, Beine hoch}

\item {\mdseries\color[rgb]{0.0,0.0,0.5019608}
Schutz vor Unterk\"uhlung/\"Uberw\"armung}

\item {\mdseries\color[rgb]{0.0,0.0,0.5019608}
Aviso im Krankenhaus / Intensivstation
({\quotedblbase}\"Uberwachungsbett{\textquotedblleft})}

\item {\mdseries\color[rgb]{0.0,0.0,0.5019608}
\gZ{ben\"otigt Antibiotika (im KH)}}

\end{itemize}
\paragraph{Neurogener Schock}
{\mdseries
Zusammenbruch der nerv\"osen Kreislaufregulation, inad\"aquate
arterielle und ven\"ose Gef\"a{\ss}erweiterung (Vasodilatation)}

{\mdseries\color{black}
Ursachen }

\begin{itemize}
\item {\mdseries
Neurologische Erkrankungen wie Querschnitt-Verletzung des R\"uckenmarks}

\item {\mdseries
Die Gef\"a{\ss}e bekommen keine Nerven-Impulse und kontrahieren sich
nicht mehr.}

\item {\mdseries
${\rightarrow}$ Blut versackt}

\end{itemize}
{\mdseries\color{black}
Spezielle Ma{\ss}nahmen}

\begin{itemize}
\item {\mdseries\color[rgb]{0.0,0.0,0.5019608}
Allgemein:
\hyperlink{a761AllgemeineManahmenderSchockbekmpfungoutline}{Allgemeine
Ma{\ss}nahmen der Schockbek\"ampfung} (siehe oben) }

\item {\mdseries\color[rgb]{0.0,0.0,0.5019608}
Lagerung: Flachlagerung (WS-Verletzung ist Kontraindikation gg. Beine
hoch!)}

\end{itemize}


\section{[RS VII] Traumatologische Notf\"alle}
{\mdseries\itshape
Michael Withofner, Dr. med. univ.}

\begin{itemize}
\item {\mdseries
Sch\"adel-Hirn-Trauma (SHT)}

\item {\mdseries
Thorax trauma}

\item {\mdseries
Abdominal trauma / Bauchtrauma}

\item {\mdseries
Beckentrauma}

\item {\mdseries
Verletzungen der Wirbels\"aule}

\item {\mdseries
Polytrauma}

\item {\mdseries
Verletzungen der Extremit\"aten}

\item {\mdseries
Verbrennungen}

\end{itemize}
{\mdseries\color{black}
Definition Trauma}

{\mdseries
Als Trauma  bezeichnet man eine {\textperiodcentered}~Sch\"adigung,
{\textperiodcentered}~Verletzung, oder {\textperiodcentered}~Wunde, die
durch eine Gewalteinwirkung physikalisch er oder chemisch er Art
verursacht wurde.}

\subsection{Sch\"adel-Hirn-Trauma (SHT)}
{\mdseries
Folge von Gewalteinwirkung auf den  Sch\"adel}

{\mdseries\color{black}
Arten:}

\begin{itemize}
\item {\mdseries
gedecktes SHT: ohne  Er\"offnung der Dura  mater}

\item {\mdseries
offenes SHT: mit  Er\"offnung der Dura mater}

\end{itemize}
{\mdseries\color{black}
begleitende Verletzungen:}

\begin{itemize}
\item {\mdseries
Verletzung der Kopfschwarte  (Ri{\ss}quetschwunde)}

\item {\mdseries
Frakturen von}

\end{itemize}
\begin{itemize}
\item {\mdseries
Sch\"adeldach}

\item {\mdseries
Sch\"adelbasis (Blutung aus Ohren, Liquoraustritt  aus der Nase)}

\item {\mdseries
Gesichtssch\"adel}

\item {\mdseries
Unterkiefer}

\end{itemize}
{\mdseries\color[rgb]{0.0,0.0,0.5019608}
Gedecktes SHT -- Klassifikation}

\begin{itemize}
\item {\mdseries
\gT{Gehirnersch\"utterung}  {}--
\gT{Contusio cerebri}}

\item {\mdseries
\gT{Gehirnprellung}  {}-- \gT{Contusio
cerebri}}

\item {\mdseries
\gT{Gehirnquetschung}}

\end{itemize}
\paragraph[Gehirnersch\"utterung (Commotio
(cerebri))]{Gehirnersch\"utterung (Commotio (\textit{cerebri}))}
\index{Gehirnersch\"utterung}\index{Commotio (cerebri)}\begin{itemize}
\item {\mdseries
 \gEs{Bewu{\ss}tlosigkeit} /Bewu{\ss}tseinstr\"ubung
{\textless}  1h}

\item {\mdseries
\"Ubelkeit, Erbrechen, Schwindel}

\item {\mdseries
evt. fl\"uchtige neurologische Ausf\"alle}

\item {\mdseries
retrograde  \gEs{Amnesie}  ({\quotedblbase}kann sich
an den Unfallhergang nicht erinnern{\textquotedblleft} )}

\end{itemize}
\begin{itemize}
\item {\mdseries
\gZ{im
}\gZ{\textit{CT}}\gZ{/}\gZ{\textit{MR
[${\rightarrow}$}}\gZ{\textit{R}}\gZ{\textit{{}-}}\gZ{\textit{\pageref{bkm:LexComputertomographie}}}\gZ{\textit{]
}}\gZ{keine nachweisbaren Sch\"aden}}

\item {\mdseries
keine  bleibenden Sch\"aden}

\end{itemize}
{\mdseries
gedecktes SHT - Klassifikation}

\paragraph[Gehirnprellung (Contusio (cerebri))]{Gehirnprellung (Contusio
(\textit{cerebri}))}
\index{Gehirnprellung}\index{Contusio (cerebri)}\begin{itemize}
\item {\mdseries
Bewu{\ss}tlosigkeit  {\textgreater}  1h}

\item {\mdseries
Bewu{\ss}tseinstr\"ubung {\textgreater}24h}

\item {\mdseries
nachweisbare Gehirnsch\"adigung / bleibende  Ausf\"alle}

\item {\mdseries
ev. Gehirnnervenabrisse (z.B. Anosmie  bei  Riechnervabri{\ss})}

\item {\mdseries
 vegetative St\"orungen}

\end{itemize}
\paragraph[Gehirnquetschung (SHT III{\textdegree})]{Gehirnquetschung
(SHT III{\textdegree})}
\index{Gehirnquetschung}\begin{itemize}
\item {\mdseries
Bewu{\ss}tlosigkeit  {\textgreater}  24h}

\begin{center}
 [Warning: Image ignored] % Unhandled or unsupported graphics:
%\includegraphics[width=7.09cm,height=2.2cm]{AASdev20090228-img89.png}

\end{center}
\item {\mdseries
Atmungs-/ Kreislaufst\"orungen}

\item {\mdseries
Hirndruck symptomatik}

\item {\mdseries
evt. Einklemmung deAtemzentrums ${\rightarrow}$ zentraler
Atemstillsatnd}

\item {\mdseries
Behinderung der Hirndurchblutung}

\item {\mdseries
intrakranielle Blutungen}

\end{itemize}
\begin{itemize}
\item {\mdseries
\gT{epidural}}

\item {\mdseries
\gT{subdural}}

\item {\mdseries
\gT{subarachnoidal}}

\end{itemize}
\begin{itemize}
\item {\mdseries
L\"ahmungen}

\item {\mdseries
zerebrale Kr\"ampfe}

\end{itemize}
{\mdseries\color[rgb]{0.0,0.0,0.5019608}
Intrakranielle Blutungen}

\begin{itemize}
\item {\mdseries
\index{Epidurale Blutung}\gT{Epidurale Blutung} }

\end{itemize}
\begin{itemize}
\item {\mdseries
Ri{\ss} einer Meningealarterie}

\item {\mdseries
Einblutung zwischen Kalotte und Dura mater}

\item {\mdseries
{\textquotedblleft}lichtes{\textquotedblright} (luzides ) Intervall:
einige Stunden}

\end{itemize}
\begin{itemize}
\item {\mdseries
\index{Subdurale Blutung}\gT{Subdurale  Blutung} }

\end{itemize}
\begin{itemize}
\item {\mdseries
schleichender  (!) Verlauf: oft \"uber Tage}

\item {\mdseries
auch ohne  Trauma: Hypertoniker, Alkoholiker}

\end{itemize}
\begin{itemize}
\item {\mdseries
\index{Subarachnoidale Blutung}\gT{Subarachnoidale 
Blutung} (\index{SAB}\gT{SAB})}

\end{itemize}
\begin{itemize}
\item {\mdseries
z.B. geplatztes Aneurysma, Br\"uckenvenenabri{\ss} bei
Decelerationstrauma}

\end{itemize}
\begin{itemize}
\item {\mdseries
kann spontan / ohne Trauma auftreten (Aneurysma)}

\end{itemize}
\begin{itemize}
\item {\mdseries
zuerst Herd symptomatik, schneller  Verfall, 
{\quotedblbase}W\"uhlblutung{\textquotedblleft}}

\item {\mdseries
Klassisch: pl\"otzlicher, schwerer,
{\quotedblbase}\gEs{donnerschlagartiger}{\textquotedblleft}
Kopfschmerz}

\end{itemize}
{\mdseries\color[rgb]{0.0,0.0,0.5019608}
Offenes SHT}

{\mdseries
Hinweise:}

\begin{itemize}
\item {\mdseries
gro{\ss}e RQW mit tast-/sichtbaren  Knochensplittern}

\item {\mdseries
Hirnaustritt (schlechte Prognose)}

\item {\mdseries
Blutungen aus Nase / Ohren}

\item {\mdseries
Liquoraustritt, oft mit Blut vermischt 
({\quotedblbase}Tupfertest{\textquotedblleft})}

\end{itemize}
{\mdseries\color[rgb]{0.0,0.0,0.5019608}
SHT -- Allgemeine Diagnostik}

{\mdseries
Beurteilung von:}

\begin{itemize}
\item {\mdseries
A tmung / K reislauf  (BAK-Schema!!! )}

\item {\mdseries
neurologisch:}

\end{itemize}
\begin{itemize}
\item {\mdseries
B ewu{\ss}tseinslage: klar/getr\"ubt/bewu{\ss}tlos}

\item {\mdseries
Motorik/Reflexe}

\item {\mdseries
Pupillenreaktion }

\end{itemize}
\begin{itemize}
\item {\mdseries
normal: rund + seitengleich + mittelweit + bds.  direkt und indirekt
lichtreaktiv}

\item {\mdseries
abnorm: entrundet, starr, weit oder eng nicht  lichtreaktiv}

\end{itemize}
\begin{itemize}
\item {\mdseries
Traumacheck (sonstige Verletzungen ={\textgreater}  Polytrauma?)}

\item {\mdseries
Glasgow-Coma-Scale (GCS)}

\end{itemize}
{\mdseries\color{black}
Glasgow-Coma-Scale (GCS)}

{\mdseries
Bewertung von: {\textperiodcentered}~Augen\"offnung,
{\textperiodcentered}~beste verbale Reaktion,
{\textperiodcentered}~beste motorische Reaktion.}

\begin{center}
 [Warning: Image ignored] % Unhandled or unsupported graphics:
%\includegraphics[width=13.667cm,height=20.134cm]{AASdev20090228-img90.png}

\end{center}
\paragraph{Bewu{\ss}tseinsst\"orung beim SHT}
\begin{itemize}
\item {\mdseries\color[rgb]{0.0,0.0,0.5019608}
Sicherung der Vitalfunktionen}

\item {\mdseries\color[rgb]{0.0,0.0,0.5019608}
Atemwege freihalten}

\item {\mdseries\color[rgb]{0.0,0.0,0.5019608}
stabile Seitenlage bei Bewu{\ss}tlosigkeit}

\item {\mdseries\color[rgb]{0.0,0.0,0.5019608}
NOTARZT! ${\leftarrow}$${\rightarrow}$ Intubation und Beatmung ab GCS
{\textless}8 durch NA dringend  erforderlich}

\item {\mdseries\color[rgb]{0.0,0.0,0.5019608}
steriler Verband locker auf offene  Verletzungen}

\item {\mdseries\color[rgb]{0.0,0.0,0.5019608}
Knochen / Hirnteile NICHT  zur\"uckstopfen /  entfernen!}

\item {\mdseries\color[rgb]{0.0,0.0,0.5019608}
oft Kombination mit HWS-Verletzung  ={\textgreater}Stifneck    }

\end{itemize}
\begin{itemize}
\item {\mdseries
(Gefahr: hohe Querschnittl\"ahmung)}

\end{itemize}
\begin{itemize}
\item {\mdseries\color[rgb]{0.0,0.0,0.5019608}
Oberk\"orper/Kopf 30{\textdegree} erh\"oht lagern (nur wenn Patient
bewu{\ss}tseinsklar)}

\item {\mdseries\color[rgb]{0.0,0.0,0.5019608}
Sch\"adel/HWS/BWS bilden eine Linie}

\end{itemize}
\begin{itemize}
\item {\mdseries
Jedes SHT ist im Spital abzukl\"aren! (klinisch und  bildgebend)}

\end{itemize}
{\mdseries\color{black}
\gEs{Jedes SHT ist im Spital abzukl\"aren!} Es
besteht die Gefahr dass sich der Patient gegenw\"artig in der
{\quotedblbase}\gEs{luciden
Phase}{\textquotedblleft} befindet und sich sein Zustand erst in
einigen Stunden, dann aber schlagartig verschlechtert. Im Zweifelsfall
wird der Patient dort zur  Beobachtung (24h lang) aufgenommen.}

\subsection{Thoraxtrauma}
\begin{itemize}
\item {\mdseries
offenes}

\end{itemize}
\begin{itemize}
\item {\mdseries
offener Pneumothorax}

\end{itemize}
\begin{itemize}
\item {\mdseries
geschlossenes}

\end{itemize}
\begin{itemize}
\item {\mdseries
\index{Spontanpneumothorax}Spontanpneumothorax (Lungenri{\ss},
bull\"oses  Emphysem)}

\item {\mdseries
Lungenverletzung durch gebrochene Rippen }

\item {\mdseries
Bronchusverletzungen}

\end{itemize}
{\mdseries\color{black}
Thoraxtrauma -- allgemeine  Symptome}

\begin{itemize}
\item {\mdseries
Dyspnoe}

\end{itemize}
\begin{itemize}
\item {\mdseries
Tachypnoe}

\item {\mdseries
Zyanose}

\end{itemize}
\begin{itemize}
\item {\mdseries
atemabh\"angige  Schmerzen,  Inspirationsstop}

\item {\mdseries
ev. Hautemphysem (Luftansammlung in der Haut, typisches
{\quotedblbase}Knistern{\textquotedblleft})}

\item {\mdseries
Prellmarken (z.B. Lenks\"aule von PKW) }

\end{itemize}
\begin{itemize}
\item {\mdseries
={\textgreater} gr\"undlicher Traumacheck!}

\end{itemize}
{\mdseries\color{black}
Thoraxtrauma -- spezifische  Symptome}

\begin{itemize}
\item {\mdseries
instabiler Thorax}

\end{itemize}
\begin{itemize}
\item {\mdseries
Sternumfraktur}

\item {\mdseries
Serienrippenfraktur (zus. paradoxe Atmung --  Brustkorb senkt beim
Einatmen)}

\end{itemize}
\begin{itemize}
\item {\mdseries
Dyspnoe + Prellmarken +  Bluthusten  (H\"amoptysis )}

\end{itemize}
\begin{itemize}
\item {\mdseries
Lungenprellung: Gasaustausch [F0EA?] , da  Lungenabschnitte
gesch\"adigt}

\end{itemize}
\begin{itemize}
\item {\mdseries
Brustschmerz, Schock }

\end{itemize}
\begin{itemize}
\item {\mdseries
Rhythmusst\"orungen ${\rightarrow}$ Verletzung des  Herzens}

\item {\mdseries
schwerer Schock, schneller Verfall ${\rightarrow}$ Riss der  Aorta}

\item {\mdseries
Blut im Schwall aus dem Mund ${\rightarrow}$  Speiser\"ohrenverletzung}

\end{itemize}
{\mdseries\color[rgb]{0.0,0.0,0.5019608}
\index{Pneumothorax}Pneumothorax}

\begin{itemize}
\item {\mdseries
Lungen- (innerer) oder  Brustwandverletzung (\"au{\ss}erer)}

\begin{center}
 [Warning: Image ignored] % Unhandled or unsupported graphics:
%\includegraphics[width=7.086cm,height=10.686cm]{AASdev20090228-img91.png}

\end{center}
\end{itemize}
\begin{itemize}
\item {\mdseries
Luft gelangt in den Pleuraspalt}

\item {\mdseries
Unterdruck aufgehoben}

\item {\mdseries
Lunge kollabiert auf betroffener Seite}

\end{itemize}
\begin{itemize}
\item {\mdseries
{\quotedblbase}unkomplizierter{\textquotedblleft} Pneumothorax}

\end{itemize}
\begin{itemize}
\item {\mdseries
Luft kann ungehindert in Pleuraspalt und  wieder heraus}

\item {\mdseries
{\quotedblbase}nur{\textquotedblleft} Ausfall des betroffenen
Lungenfl\"ugels}

\end{itemize}
\begin{itemize}
\item {\mdseries
Ventil-, \index{Spannungspneumothorax}Spannungspneumothorax}

\end{itemize}
\begin{itemize}
\item {\mdseries
Luft kann zwar hinein, aber nicht wieder  heraus}

\item {\mdseries
Druck auf betroffener Seite steigt ${\rightarrow}$  Organe (Lunge, Herz)
werden auf die andere Seite gedr\"uckt}

\end{itemize}
\begin{itemize}
\item {\mdseries
gesunde Seite wird verdr\"angt + Abknickung  gro{\ss}er Gef\"a{\ss}e!}

\end{itemize}


\begin{center}
 [Warning: Image ignored] % Unhandled or unsupported graphics:
%\includegraphics[width=5.811cm,height=4.343cm]{AASdev20090228-img92.png}

\end{center}
{\mdseries\color{black}
Therapie}

{\mdseries
Pneumothorax, unkompliziert:}

\begin{itemize}
\item {\mdseries\color[rgb]{0.0,0.0,0.5019608}
Lagerung auf verletzte Seite (unverletzte kann  sich besser entfalten)}

\item {\mdseries\color[rgb]{0.0,0.0,0.5019608}
O2}

\item {\mdseries\color[rgb]{0.0,0.0,0.5019608}
Notarzt, legt ev. vor Ort B\"ulaudrainage}

\item {\mdseries\color[rgb]{0.0,0.0,0.5019608}
(nicht versuchen, Wunde abzudichten ={\textgreater} ev. dadurch
Spannungspneumothorax)}

\end{itemize}
{\mdseries
Ventilpneumothorax:}

\begin{itemize}
\item {\mdseries\color[rgb]{0.0,0.0,0.5019608}
Entlastung durch Drainage!}

\item {\mdseries\color[rgb]{0.0,0.0,0.5019608}
offener Pneumothorax (Mittelfellpendeln)}

\item {\mdseries\color[rgb]{0.0,0.0,0.5019608}
steril abdecken}

\end{itemize}
{\mdseries\color[rgb]{0.0,0.0,0.5019608}
Stichverletzungen}

{\mdseries
NIEMALS  HERAUSZIEHEN!!!!!!!!}

{\mdseries
steril abdecken + fixieren}

{\mdseries
verletzte Strukturen??, dient ev. als 
{\quotedblbase}St\"opsel{\textquotedblleft}??}

{\mdseries
Fremdk\"orper soll nicht tiefer rutschen}

\subsection[Bauchtrauma]{Bauchtrauma}
\index{Bauchtrauma}\begin{itemize}
\item {\mdseries
offen (perforierend),}

\end{itemize}
\begin{itemize}
\item {\mdseries
 Er\"offnung der  Bauchh\"ohle}

\item {\mdseries
Stich-, Pf\"ahlungsverletzung}

\item {\mdseries
Schnittverletzung, {\quotedblbase}Platzbauch{\textquotedblleft}}

\end{itemize}
\begin{itemize}
\item {\mdseries
geschlossen (stumpf)}

\end{itemize}
\begin{itemize}
\item {\mdseries
Schlag, Tritt}

\item {\mdseries
Prellung durch Lenkstange bei PKW-Unfall}

\end{itemize}
{\mdseries\color[rgb]{0.0,0.0,0.5019608}
Offenes Bauchtrauma}



\begin{center}
 [Warning: Image ignored] % Unhandled or unsupported graphics:
%\includegraphics[width=6.543cm,height=4.839cm]{AASdev20090228-img93.png}

\end{center}
\begin{itemize}
\item {\mdseries
ev. Austritt von Darmschlingen}

\end{itemize}
\begin{itemize}
\item {\mdseries
NICHT  zur\"uckstopfen ${\rightarrow}$ sonst 
Darmri{\ss}/Darmverschlu{\ss}}

\item {\mdseries
steril abdecken, mit Ringerl\"osung feuchthalten}

\item {\mdseries
Bauchdecke entspannen , Knierolle ${\rightarrow}$ sonst  Darmabklemmung}

\end{itemize}
\begin{itemize}
\item {\mdseries
Gefahr des Eindringens von Keimen  (Peritonitis/Sepsis!)}

\end{itemize}
\begin{itemize}
\item {\mdseries\color[rgb]{0.0,0.0,0.5019608}
steril  abdecken}

\item {\mdseries\color[rgb]{0.0,0.0,0.5019608}
Vitalparameter}

\item {\mdseries\color[rgb]{0.0,0.0,0.5019608}
O2}

\item {\mdseries\color[rgb]{0.0,0.0,0.5019608}
Notarzt}

\end{itemize}
{\mdseries\color[rgb]{0.0,0.0,0.5019608}
Geschlossenes Bauchtrauma}

{\mdseries\color{black}
Gefahren:}

\begin{itemize}
\item {\mdseries
Gef\"a{\ss}verletzungen}

\end{itemize}
\begin{itemize}
\item {\mdseries
Aorta, Vena cava,...}

\item {\mdseries
massive \gEs{innere Blutungen} ${\rightarrow}$
Volumenmangelschock}

\end{itemize}
\begin{itemize}
\item {\mdseries
Zerrei{\ss}ung  parenchymat\"oser Organe (\gEs{Milz},
\gEs{Leber})}

\end{itemize}
\begin{itemize}
\item {\mdseries
massive \gEs{innere Blutungen} ${\rightarrow}$
Volumenmangelschock}

\end{itemize}
\begin{itemize}
\item {\mdseries
Zerrei{\ss}ung  von Hohlorganen (Magen, Darm,Galleng\"ange )}

\end{itemize}
\begin{itemize}
\item {\mdseries
Austritt von Speisebrei, Magensaft, Stuhl, Gallensekret}

\end{itemize}
\begin{itemize}
\item {\mdseries
Bauchfellentz\"undung (Peritonitis) nach Std. bis  Tagen}

\item {\mdseries
Sepsis (hohe Lethalit\"at), septischer Schock}

\end{itemize}
{\mdseries\color{black}
Symptome}

\begin{itemize}
\item {\mdseries
starke  Bauchschmerzen}

\item {\mdseries
brettharte  Bauchdecke  (Abwehrspannung, D\'efense )}

\item {\mdseries
Leitsymptom: Schockzeichen}

\item {\mdseries
beim Bewu{\ss}tlosen schwierig zu  diagnostizieren}

\end{itemize}
\begin{itemize}
\item {\mdseries
gr\"undlich nach \index{Prellmarken}Prellmarken suchen!!!}

\end{itemize}
{\centering\color{black}
Immer nach \gEs{Prellmarken} suchen. Patienten
\gEs{entkleiden}!
\par}

{\mdseries
Bauchtrauma}

{\mdseries\color{black}
\index{Milzruptur}Milzruptur}

\begin{itemize}
\item {\mdseries
prim\"ar: sofort, Kapsel durchgerissen}

\item {\mdseries
sekund\"ar: Sickerblutung, Kapsel anfangs noch intakt ${\rightarrow}$
Latenz bis Ri{\ss} der Kapsel}

\end{itemize}
{\mdseries\color{black}
\gCa{Cave!} {\quotedblbase}\index{Zweizeitige
Milzruptur}\gEs{Zweizeitige
Milzruptur}{\textquotedblleft}}

{\mdseries\color{black}
\index{Aortenruptur}Aortenruptur}

\begin{itemize}
\item {\mdseries
bei VU selten, bei Sturz aus gro{\ss}er H\"ohe  h\"aufiger}

\item {\mdseries
Sehr schlechte Prognose}

\end{itemize}
{\mdseries\color{black}
\index{Nierenruptur}Nierenruptur}

\begin{itemize}
\item {\mdseries
Leitsymptom: blutiger Harn, NICHT bei  Harnleiterabri{\ss}}

\end{itemize}
{\mdseries\color{black}
Beckentrauma}

\begin{itemize}
\item {\mdseries
erhebliche Gewalteinwirkung  (Versch\"uttung, \"Uberrolltrauma)}

\item {\mdseries
massive  Blutungen (bis 5 l in wenigen  Minuten)}

\item {\mdseries
Gefahr des Ausblutens (ev. ohne sichtbare  Blutungsquelle!)}

\item {\mdseries
Mitverletzung von Organen im kleinen  Becken:}

\end{itemize}
\begin{itemize}
\item {\mdseries
Dickdarm ${\rightarrow}$ Blutung aus Anus}

\item {\mdseries
Harnsystem ${\rightarrow}$ Blut im Harn / Harn in der  Bauchh\"ohle}

\item {\mdseries
(Ausnahme: Spontanruptur der gef\"ullten  Harnblase bei stumpfem
Bauchtrauma)}

\item {\mdseries
Hoden ${\rightarrow}$ H\"amatome am Skrotum und am  Damm}

\end{itemize}
{\mdseries\color{black}
Therapie:}

\begin{itemize}
\item {\mdseries\color[rgb]{0.0,0.0,0.5019608}
Basisma{\ss}nahmen}

\item {\mdseries\color[rgb]{0.0,0.0,0.5019608}
Schockbek\"ampfung}

\item {\mdseries\color[rgb]{0.0,0.0,0.5019608}
Bei Beckenfraktur: \index{Beckengurt}\textbf{Beckengurt}, zur Not mit
G\"urtel%
%
}

\item {\mdseries\color[rgb]{0.0,0.0,0.5019608}
O2  8-10 l mindestens}

\item {\mdseries\color[rgb]{0.0,0.0,0.5019608}
Volumenzufuhr i.V. (NA)}

\item {\mdseries\color[rgb]{0.0,0.0,0.5019608}
Schmerzbek\"ampfung}

\end{itemize}
{\mdseries
Prognose h\"angt vom Ausma{\ss} des  Blutverlustes ab.}

\subsection{Wirbels\"aulenverletzungen}
\begin{itemize}
\item {\mdseries
Typische Unfallmechanismen:}

\end{itemize}
\begin{itemize}
\item {\mdseries
Sturz aus gro{\ss}er H\"ohe}

\item {\mdseries
Motorradunfall (Genickbruch)}

\item {\mdseries
eingeklemmte Personen bei VU}

\item {\mdseries
Schlag auf Kopf (Stauchung der HWS)}

\end{itemize}
\begin{itemize}
\item {\mdseries
Gefahren:}

\end{itemize}
\begin{itemize}
\item {\mdseries
Kompression/Durchtrennung des  R\"uckenmarks}

\item {\mdseries
L\"ahmungen (Parese)}

\item {\mdseries
Sensibilit\"atsst\"orungen (Par\"asthesie)}

\end{itemize}
\begin{itemize}
\item {\mdseries
inkomplett: eines von beidem}

\item {\mdseries
komplett: beides}

\end{itemize}
{\mdseries
traumatische Einblutung ins R\"uckenmark}

{\mdseries
dadurch bedingtes \"Odem ={\textgreater} RM-Kompression}

{\mdseries
Wirbels\"aulenverletzung -- Schleudertrauma, Peitschenschlag}

{\mdseries
bei Auffahrunf\"allen}

{\mdseries
Zerrung des Halsmarks (bildgebend meist nicht nachweisbar)}

{\mdseries
Schmerzen im HWS-Bereich}

{\mdseries
Kopfschmerzen}

{\mdseries
Schwindel, Tinnitus}

{\mdseries
ev. neurologische Ausf\"alle (in schweren F\"allen)}

{\mdseries
Prophylaxe: Nackenst\"utze}

{\mdseries
Therapie: Immobilisation, Stifneck, (ev. O2 )}

{\mdseries
Wirbels\"aulenverletzung -  Querschnittssydrom}

{\mdseries
inkomplett: Motorik oder Sensorik  ausgefallen}

{\mdseries
komplett: beides ausgefallen}

{\mdseries
Symptome:}

{\mdseries
schlaffe L\"ahmung, Ausfall aller Reflexe  unterhalb des betroffenen
Segmentes}

{\mdseries
Blasen-, Darml\"ahmung ={\textgreater} Inkontinenz,  \"Uberlaufblase}

{\mdseries
Ausfall der Gef\"a{\ss}-und W\"armeregulation}

{\mdseries
Wirbels\"aulenverletzungen -- spinaler  Schock}

{\mdseries
pl\"otzliche spinale Durchtrennung mit  schlagartigem Wegfall aller
zentralen  erregenden Impulse...}

{\mdseries
Gef\"a{\ss}regulation versagt}

{\mdseries
Gef\"a{\ss}tonus [F0EA?]}

{\mdseries
RR [F0EA?]  }

{\mdseries
Wirbels\"aulenverletzungen -- {\quotedblbase}hoher{\textquotedblleft} 
Querschnitt}

{\mdseries
L\"asionen auf H\"ohe C1/C2 (1. bzw. 2.  HWK)}

{\mdseries
Abtrennung des Atemzentrums in der Medulla  oblongata}

{\mdseries
Anpralltrauma (Motorrad unter Lastwagen)}

{\mdseries
{\quotedblbase}K\"opfler{\textquotedblleft} in seichtes Gew\"asser}

{\mdseries
Erh\"angen}

{\mdseries
zentrale Ateml\"ahmung (keine  Spontanatmung!)}

{\mdseries
L\"asionen oberhalb C4 (4. HWK)}

{\mdseries
Segmente, die Zwerchfell und Kehlkopf  innervieren mitbetroffen}

{\mdseries
motorische Ateml\"ahmung ={\textgreater} auf Dauer  beatmungspflichtig}

{\mdseries\color{black}
Wirbels\"aulentrauma -  Vorgehensweise}

{\mdseries
bei Verdacht  auf WS-Trauma  noch vor Ort  (bevor  der Patient bewegt 
wird) \"Uberpr\"ufung von Sensibilit\"at  und Motorik (und
Durchblutung) der Extremit\"aten !!!!!}

{\mdseries
(Merkabk\"urzung {\textquotedblleft}DMS {\textquotedblright}, Teil jedes
 Traumachecks)}

{\mdseries
Wirbels\"aulenverletzungen - Therapie}

{\mdseries
Strikte Immobilisierung (Vacuum-Matratze, Schaufeltrage, Stifneck)}

{\mdseries
Vacuummatratze immer gemeinsam mit  Stifneck}

{\mdseries
Bewegungen nur in der WS-Achse (keine  Beugung, Drehung,
Seitw\"artsneigung)}

{\mdseries
falls n\"otig, korrekte Helmabnahme (zu  zweit)}

{\mdseries
vorsichtiger L\"angszug (Stifneck anlegen)}

{\mdseries
Achtung beim Rautek-Griff: keine HWS- Stabilisierung!!!}

{\mdseries
Wirbels\"aulenverletzung -  Schocklagerung?}

{\mdseries
ja, spezielle Technik:}

{\mdseries
Patient auf Vacuummatratze + Trage}

{\mdseries
Trage schr\"ag: Kopf tief, Beine hoch}

{\mdseries
nein, wenn SHT oder kardiale  Komplikationen!!!}

{\mdseries\color{black}
Anmerkungen zur Vacuummatratze}

{\mdseries
Immer mit Schaufeltrage}

{\mdseries
Immer mit Stifneck}

{\mdseries
Umbettungen durch mehrere Personen}

\subsection{Polytrauma}
{\mdseries
Gleichzeitig entstandene Verletzungen \gEs{mehrerer
K\"orperregionen} /Organsystemen, von  denen
\gEs{wenigstens eine} allein
\gEs{oder die Kombination} mehrerer
lebensgef\"ahrlich  ist.}

\begin{itemize}
\item {\mdseries
Bauchtrauma + Oberschenkel; SHT + WS + Thoraxtrauma etc.}

\item {\mdseries
gebrochener Finger + gebrochene Zehe ${\rightarrow}$ kein Polytrauma}

\end{itemize}
{\mdseries
Problem:  Volumenmangelschock +  Ateminsuffizienz}

{\mdseries\color{black}
Polytrauma - Vorgehensweise}

\begin{itemize}
\item {\mdseries\color[rgb]{0.0,0.0,0.5019608}
Schneller Traumacheck gleichzeitig  mit  Sichern der Vitalparameter!}

\item {\mdseries\color[rgb]{0.0,0.0,0.5019608}
Bewu{\ss}tsein}

\item {\mdseries\color[rgb]{0.0,0.0,0.5019608}
Atmung}

\item {\mdseries\color[rgb]{0.0,0.0,0.5019608}
Kreislauf: Carotispuls, Pulse peripher  (orientierend!),
Kapillardurchblutung}

\item {\mdseries\color[rgb]{0.0,0.0,0.5019608}
gr\"undlicher Traumacheck erst sp\"ater}

\item {\mdseries\color[rgb]{0.0,0.0,0.5019608}
stabilisierende Ma{\ss}nahmen}

\item {\mdseries\color[rgb]{0.0,0.0,0.5019608}
Atemwege sichern (Intubation),  ven\"oser  Zugang (Notarzt!)}

\item {\mdseries\color[rgb]{0.0,0.0,0.5019608}
wenn einigerma{\ss}en stabil, gr\"undlicher  Traumacheck,
Verletzungsmuster (={\textgreater} zu  erwartende Komplikationen)}

\end{itemize}
{\mdseries\color{black}
Polytrauma beim Kind}

{\mdseries
Schockindex unbrauchbar!!!}

{\mdseries
Kinder bleiben lange stabil, verfallen  dann aber schlagartig!!!}

{\mdseries
doppelte Vorsicht (lieber eine  Notarztalarmierung zu viel, als eine zu 
wenig!)}

{\mdseries
nach Unfallhergang : Patient gilt als polytraumatisiert  bis zum Beweis 
des Gegenteils! }

{\mdseries\color{black}
Polytrauma - Unfallort}

{\mdseries
Absichern  der Unfallstelle,  m\"oglichst auff\"allig machen (Blaulicht,
 Warnblinkanlage, Scheinwerfer)}

{\mdseries
bei mehreren Verletzten: Triage}

{\mdseries
Nachalarmierung : andere  Rettungsmittel, Feuerwehr, Polizei}

{\mdseries
Eingeklemmte im Fahrzeug stabilisieren,  nicht eigenm\"achtig bergen}

{\mdseries
Crash -Rettung nur  bei  Fremdgef\"ahrdung, HWS stabilisieren!}

{\mdseries\color{black}
Polytrauma - Therapie}

\begin{itemize}
\item {\mdseries\color[rgb]{0.0,0.0,0.5019608}
Freimachen der Atemwege, 15l O2 , ggf.  Intubation}

\item {\mdseries\color[rgb]{0.0,0.0,0.5019608}
Kreislaufkontrolle, 2 gro{\ss}lumige Zug\"ange,  Volumen!}

\item {\mdseries\color[rgb]{0.0,0.0,0.5019608}
Blutstillung  (Druckverband, Abbindung wenn  indiziert)}

\item {\mdseries\color[rgb]{0.0,0.0,0.5019608}
Immobilisierung : Stifneck, Vacuummatratze,  Schaufeltrage}

\item {\mdseries\color[rgb]{0.0,0.0,0.5019608}
W\"armeschutz}

\item {\mdseries\color[rgb]{0.0,0.0,0.5019608}
BZ (Alkoholisierte)}

\item {\mdseries\color[rgb]{0.0,0.0,0.5019608}
Medikamente: }

\item {\mdseries\color[rgb]{0.0,0.0,0.5019608}
Volumssubstitute (kristalloid und kolloidal)}

\item {\mdseries\color[rgb]{0.0,0.0,0.5019608}
Schmerzmittel}

\end{itemize}
\subsection{Extremit\"atenverletzungen}
{\mdseries
Gelenke betreffend}

{\mdseries
Verstauchung (Distorsion )}

{\mdseries
Verrenkung (Luxation )}

{\mdseries
Knochenbruch (Fractur )}

\paragraph{Extremit\"atenverletzung - Verstauchung}
{\mdseries
Zerrung / \"Uberdehnung des  Kapsel-/Bandapparates}

{\mdseries\color{black}
Symptome}

{\mdseries
Schmerzen}

{\mdseries
Schwellung}

{\mdseries
Blutergu{\ss} (H\"amatom )}

{\mdseries\color{black}
Therapie}

{\mdseries
ruhigstellen}

{\mdseries
hochlagern}

{\mdseries
k\"uhlen}

{\mdseries
ad Unfallabteilung (mu{\ss} nicht sofort sein)}

\paragraph{Extremit\"atenverletzungen - Verrenkung}
{\mdseries
partielle oder vollst\"andige Verrenkung  eines Gelenks}

{\mdseries
Symptome}

{\mdseries
Schmerzen}

{\mdseries
Schwellung}

{\mdseries
Blutergu{\ss}}

{\mdseries
abnorme Gelenkstellung (federnd fixiert), bewegungsunf\"ahig}

{\mdseries\color{black}
Therapie}

{\mdseries
Stellung belassen}

{\mdseries
ruhigstellen}

{\mdseries
KEINE EINRENKVERSUCHE!!!}

\paragraph[Extremit\"atenverletzungen -
\ Knochenbr\"uche]{Extremit\"atenverletzungen -  Knochenbr\"uche}
{\mdseries\color{black}
Arten:}

\begin{itemize}
\item {\mdseries
geschlossene Fraktur}

\item {\mdseries
offene Fraktur}

\end{itemize}
{\mdseries
jeweils}

\begin{itemize}
\item {\mdseries
komplett}

\item {\mdseries
inkomplett}

\end{itemize}
{\mdseries\color{black}
Gefahren}

\begin{itemize}
\item {\mdseries
starke Blutungen}

\item {\mdseries
ev. bis Schock}

\item {\mdseries
begleitende Nervenverletzungen}

\item {\mdseries
Fettembolie (Fraktur langer R\"ohrenknochen)}

\item {\mdseries
Infektionsgefahr (Osteomyelitis bei offenen  Frakturen)}

\end{itemize}
{\mdseries
Immer gr\"undlichen Traumacheck  durchf\"uhren, um keine 
Begleitverletzungen zu \"ubersehen!!!}

{\mdseries
Extremit\"atenverletzungen -  Knochenbr\"uche}

{\mdseries\color{black}
Frakturzeichen}

\begin{center}
\begin{tabular}{|m{6.649cm}|m{6.651cm}|}
\hline
\multicolumn{2}{|m{13.5cm}|}{\centering
{\bfseries\upshape\color[rgb]{0.0,0.0,0.5019608}
\index{Frakturzeichen}Frakturzeichen}\par

}\\\hline
{\mdseries sichere}

 &
{\mdseries unsichere}

\\\hline
\begin{itemize}
\item {\mdseries abnorme Beweglichkeit
({\quotedblbase}Pseudogelenk{\textquotedblleft})}

\item {\mdseries Fehlstellung}

\item {\mdseries Krepitus (bitte NICHT absichtlich  ausprobieren!!)}

\item {\mdseries Sichtbare freie Knochenteile (= offener Bruch)}

\end{itemize}
 &
\begin{itemize}
\item {\mdseries Schmerz}

\item {\mdseries Schwellung}

\item {\mdseries H\"amatom}

\item {\mdseries gest\"orte Funktion / Bewegungsunf\"ahigkeit}

\end{itemize}
\\\hline
\end{tabular}
\end{center}
{\mdseries\color{black}
Therapie:}

{\mdseries
Traumacheck ${\rightarrow}$ je nach Bedarf}

{\mdseries
Blutstillung}

{\mdseries
Wundversorgung}

{\mdseries
Stifneck bei V.a. HWS-Verletzung}

{\mdseries
ev. NAW-Nachalarmierung (starke Schmerzen, gr\"o{\ss}erer Blutverlust,
Schock, relevanten Begleit\-verletzungen)}

{\mdseries
Ruhigstellung / Schienung (unter Zug)}

\subsection{Verbrennungen}
{\mdseries
Die Gr\"o{\ss}e des gesch\"adigten Areals wird in Prozent der
ver\-brannten K\"orperoberfl\"ache angegeben.}

{\mdseries\color{black}
Faustregel - Neunerregel}

{\mdseries
Handfl\"ache (des Patienten) = 1~\% der K\"orperoberfl\"ache (KOF)}

\begin{center}
 [Warning: Image ignored] % Unhandled or unsupported graphics:
%\includegraphics[width=6.591cm,height=7.03cm]{AASdev20090228-img94.png}

\end{center}
\begin{flushright}
\tablehead{}\begin{supertabular}{|m{2.134cm}|m{2.134cm}|m{2.134cm}|}
\hline
 &
\centering \%\par

 &
\\\hline
{\mdseries Kopf:}

 &
\centering 9\par

 &
\\\hline
{\mdseries Arm, je:}

 &
\centering 9\par

 &
\\\hline
{\mdseries Thorax vorne:}

 &
\centering 18\par

 &
\\\hline
{\mdseries Thorax hinten:}

 &
\centering 18\par

 &
\\\hline
{\mdseries Bein, je:}

 &
\centering 18\par

 &
\\\hline
{\mdseries Genital: }

 &
\centering 1\par

 &
\\\hline
 &
 &
\\\hline
\end{supertabular}
\end{flushright}
{\mdseries\color{black}
Grade}

Das Ausma{\ss} der Sch\"adigung durch Verbrennung wird in 3 Grade
eingeteilt:

\begin{flushleft}
\begin{tabular}{|m{1.836cm}|m{6.8980002cm}|m{4.3630004cm}|}
\hline
\centering {\bfseries\upshape\color[rgb]{0.0,0.0,0.5019608} Grad}\par

 &
 &
\\\hline
\centering \gEs{1. Grad}\par

 &
{\mdseries R\"otung (Juckreiz, Schmerz)}

 &
\\\hline
\centering \gEs{2. Grad}\par

 &
{\mdseries Blasenbildung, Schmerz}

{\mdseries 2a: oberfl\"achlich, Nerven intakt, }

{\mdseries 2b: tiefer, Nerven zerst\"ort}

 &
\\\hline
\centering \gEs{3. Grad}\par

 &
{\mdseries trockene Hautfetzen (wei{\ss}-gelb-schwarz), Verkohlung,
Schrumpf\-ung des Gewebes durch Hitze\-ein\-wirkung}

 &
\\\hline
\end{tabular}
\end{flushleft}
{\mdseries\color{black}
Ab \gEs{5\%} verbrannter (mind. 2. Grad) KOF bei
Kindern oder \gEs{10\%} bei Erwachsenen besteht
\gEs{Schockgefahr} (Fl\"ussigkeitsverlust)}

\begin{itemize}
\item {\mdseries
immer  auch an Begleitverletzungen  denken}

\item {\mdseries
bestimmte Unfallmechanismen}

\item {\mdseries
Traumacheck !}

\end{itemize}
{\mdseries\color{black}
Ma{\ss}nahmen}

\begin{itemize}
\item {\mdseries\color[rgb]{0.0,0.0,0.5019608}
\gEs{Eigenschutz!}}

\item {\mdseries\color[rgb]{0.0,0.0,0.5019608}
Kleiderbrand l\"oschen}

\item {\mdseries\color[rgb]{0.0,0.0,0.5019608}
Kleidung rasch entfernen (sofern nicht  eingeschmolzen!)}

\item {\mdseries\color[rgb]{0.0,0.0,0.5019608}
mind. 10 min. mit m\"a{\ss}ig kaltem Wasser  sp\"ulen --
\gEs{NICHT unterk\"uhlen!}}

\item {\mdseries\color[rgb]{0.0,0.0,0.5019608}
keimfreier, nicht verklebender Verband, locker anlegen }

\item {\mdseries\color[rgb]{0.0,0.0,0.5019608}
wenn vorhanden: Water-Jel%
%
\footnote{Streitfrage: Water-Jel: Umstritten} }

\item {\mdseries\color[rgb]{0.0,0.0,0.5019608}
Schockbek\"ampfung}

\item {\mdseries\color[rgb]{0.0,0.0,0.5019608}
NAW wenn n\"otig}

\end{itemize}
{\mdseries\color{black}
Was man besser nicht tun sollte...}

\begin{itemize}
\item {\mdseries
NEIN: Salbe, Puder, etc. auf frische  Brandwunde}

\item {\mdseries
NEIN: Blasen \"offnen}

\item {\mdseries
NEIN: festklebende Kleidung gewaltsam  entfernen}

\item {\mdseries
NEIN: aluminiumbeschichteten Verband  verwenden (Gefahr der
Hitzereflexion bei  nicht ausreichend gek\"uhlten  Verbrennungen!)}

\end{itemize}
{\mdseries\color[rgb]{0.0,0.0,0.5019608}
%
%
Inhalationstrauma}

{\mdseries\color{black}
Leitsymptome}

\begin{itemize}
\item {\mdseries
angebrannte Gesichtshaare}

\item {\mdseries
geschwollene, ru{\ss}ige Lippen}

\item {\mdseries
Ru{\ss}partikel im Sputum}

\item {\mdseries
Heiserkeit}

\item {\mdseries
Stridor}

\end{itemize}
{\mdseries
${\rightarrow}$ Rauchgasvergiftung m\"oglich!}

\subsection{Erfrierungen}
{\mdseries
...sind Gewebsnekrosen, die durch lokale  Durchblutungsst\"orungen
aufgrund von  K\"alteeinwirkung (Eis, Wind, N\"asse,  K\"alte)
entstanden sind.}

{\mdseries
Besonders geff\"ahrdet sind k\"orperferne K\"orperteile wie Zehen,
Finger, Ohren oder Nase.}

{\mdseries
Einteilung erfolgt in 4 Stadien, diese haben am Notfallort keine
entscheidende Bedeutung f\"ur die zu setzenden Ma{\ss}nahmen. Den
Schweregrad einer Erfrierung kann man erst nach Tagen voll
einsch\"atzen. Deshalb ist es wichtig, die Warnsymptome ernst zu
nehmen!}

Stadien

Symptome

zuerst 

Gef\"uhllosigkeit und Bl\"asse

starke Schmerzen und bl\"auliche Verf\"arbung

 einige Zeit sp\"ater

marmorierte Verf\"arbung der Haut mit  Blasenbildung

Taubheitsgef\"uhl, keine Empfindung bei  Ber\"uhrung

Bewegungsunf\"ahigkeit des betroffenen  K\"orperteils

starke Schmerzen

{\mdseries\color{black}
Ma{\ss}nahmen}

\begin{itemize}
\item {\mdseries\color[rgb]{0.0,0.0,0.5019608}
beengende Kleidungsst\"ucke lockern (Einschn\"urung erh\"oht Risiko
einer Erfrierung)}

\item {\mdseries\color[rgb]{0.0,0.0,0.5019608}
Keimfreier, gepolsterter Verband}

\item {\mdseries\color[rgb]{0.0,0.0,0.5019608}
Warme, gezuckerte Getr\"anke -- KEIN Alkohol}

\item {\mdseries\color[rgb]{0.0,0.0,0.5019608}
K\"orperstamm w\"armen, nie das erfrorene K\"orperteil abreiben!!}

\end{itemize}
\subsection{Unf\"alle durch Strom- und Blitzschlag}
{\mdseries
\gT{Niederspannung}: bis 1000 Volt (Steckdose:  230
Volt)}

{\mdseries
\gT{Hochspannung}: \"uber 1000 Volt, Stromst\"arke
{\textgreater}  3 Ampere (\"Uberlandleitungen) }

{\mdseries
Effekte / Verletzungsmuster je nach  Spannung (Volt) / Stromst\"arke
(Ampere)}

\begin{itemize}
\item {\mdseries
Reizeffekte (Muskel-/Nervenerregung):}

\end{itemize}
\begin{itemize}
\item {\mdseries
Verkrampfung ({\quotedblbase}kann nicht loslassen{\textquotedblleft})}

\item {\mdseries
Sensibilit\"atsst\"orungen}

\item {\mdseries
ZNS ${\rightarrow}$ Bewu{\ss}tseinsst\"orung}

\item {\mdseries
Herzrhythmusst\"orungen (Kammerflimmern, ev.  als Sp\"atfolge!)}

\end{itemize}
\begin{itemize}
\item {\mdseries
W\"armeproduktion:}

\end{itemize}
\begin{itemize}
\item {\mdseries
Strommarken (Ein-/Austrittsverbrennung)}

\item {\mdseries
Lichtbogen: oberfl\"achliche und tiefe  Verbrennungen}

\end{itemize}
\begin{itemize}
\item {\mdseries
Gleichstrom}

\item {\mdseries
Wechselstrom ${\rightarrow}$ eher  Herzrhythmusst\"orungen}

\item {\mdseries
Patient  kann Spannungsquelle ev. nicht  loslassen (Krampf der
Handmuskulatur)}

\item {\mdseries
Auch Defi stellt Strom-Gefahrenquelle  dar!!!}

\item {\mdseries
Blitzunfall}

\end{itemize}
\begin{itemize}
\item {\mdseries
mehrere Millionen (!) Volt}

\item {\mdseries
bis 100.000 Ampere}

\end{itemize}
{\mdseries\color{black}
Symptome:}

\begin{itemize}
\item {\mdseries
Bewu{\ss}tlosigkeit}

\item {\mdseries
L\"ahmungen (vor\"ubergehend)}

\item {\mdseries
Blitzfiguren auf der Haut}

\item {\mdseries
Herzrhythmusst\"orungen m\"oglich}

\end{itemize}
{\mdseries
Cave: Das eigentliche Ausma{\ss} der Verletzung ist oft
{\quotedblbase}unsichtbar{\textquotedblleft}, das der Strom im
K\"orperinneren -- also {\quotedblbase}unterirdisch{\textquotedblleft}
- viel Schaden anrichtet. }

{\mdseries\color{black}
Ma{\ss}nahmen:}

\begin{itemize}
\item {\mdseries\color[rgb]{0.0,0.0,0.5019608}
\gEs{Eigenschutz!} }

\end{itemize}
\begin{itemize}
\item {\mdseries
(Stromkreis  vorher  unterbrechen, FI bzw. bei Starkstrom durch 
Fachmann)}

\end{itemize}
\begin{itemize}
\item {\mdseries\color[rgb]{0.0,0.0,0.5019608}
Genaue  Patientenuntersuchung ${\rightarrow}$ \textbf{Eintritts}{}-  und
\textbf{Austrittsmarke}}

\item {\mdseries\color[rgb]{0.0,0.0,0.5019608}
auf ev. \textbf{Herzrhythmusst\"orungen }vorbereitet  sein
${\rightarrow}$ bis 24h Latenz}

\item {\mdseries\color[rgb]{0.0,0.0,0.5019608}
Vitalparameter ${\rightarrow}$ entsprechende  symptomatische Versorgung}

\item {\mdseries\color[rgb]{0.0,0.0,0.5019608}
NAW ${\rightarrow}$ Hospitalisierung}

\item {\mdseries\color[rgb]{0.0,0.0,0.5019608}
Stromverbrennungen wie Verbrennungen  versorgen}

\end{itemize}
\section[[RS II{]} Hygiene ]{[RS II] Hygiene }
\label{bkm:RefHeading32633609}{\mdseries\itshape
Sebastian Gabriel [*]}

\begin{center}
 [Warning: Image ignored] % Unhandled or unsupported graphics:
%\includegraphics[width=4.763cm,height=4.763cm]{AASdev20090228-img95.jpg}

\end{center}
{\bfseries\color[rgb]{0.5019608,0.0,0.0}
\gEs{Changelog:}}

{\bfseries\color[rgb]{0.5019608,0.0,0.0}
2009-02-20: Fertig f\"ur Kurs 2009/02.}

{\bfseries\color[rgb]{0.5019608,0.0,0.0}
2009-02-25: Start Changelog.}

\subsection{Grundlagen}
\begin{multicols}{2}
{\mdseries
\index{Infektion}\gT{Infektion}: Eindringen von
krankmachenden Keimen in den Organismus}

{\mdseries
\index{Pathogenit\"at}\gT{Pathogenit\"at}:
\gZ{F\"ahigkeit eines Mikroorganismus, eine Krankheit
auszul\"osen}}

{\mdseries
\index{Resistenz}\gT{Resistenz}: Widerstandsf\"ahigkeit
eines Organismus gegen die  krankmachenden  Eigenschaften eines
Erregers}

{\mdseries
\index{Inkubationszeit}\gT{Inkubationszeit}: symptomlose
Zeit zwischen Infektion (Eintritt d. Erregers in den Organismus) und
dem Auftreten erster Symptome}

{\mdseries
\index{Prodromalstadium}\gT{\textcolor[rgb]{0.7529412,0.7529412,0.7529412}{Prodromalstadium}}
 \gZ{uncharakteristische Symptome am Ende der
Inkubationszeit (Fieber, Kopfschmerz, Erbrechen, Gliederschmerzen)}}

{\mdseries
\index{Epidemie}\gT{Epidemie}: \"ortlich begrenztes,
zeitlich vermehrtes Auftreten einer  Infektionserkrankung (Grippe,
Salmonellen)}

{\mdseries
\index{Pandemie}\gT{\textcolor[rgb]{0.7529412,0.7529412,0.7529412}{Pandemie}}:
\gZ{Epidemie die sich \"uber L\"ander u. Kontinente
ausbreitet (Pest im Mittelalter)}}

{\mdseries
\index{Endemie}\gT{\textcolor[rgb]{0.7529412,0.7529412,0.7529412}{Endemie}}:
\gZ{Auf kleinere Gebiete beschr\"anktes, zeitlich
jedoch nicht begrenztes Auftreten einer Infektionserkrankung  (FSME,
Malaria, Lepra)}}

{\mdseries
\index{Morbidit\"at}\gT{\textcolor[rgb]{0.7529412,0.7529412,0.7529412}{Morbidit\"at}}:
\gZ{Anteil der an XY Erkrankten einer bestimmten
Bev\"olkerungsgruppe (z.B. 5 pro 100.000 EW)}}

{\mdseries
\index{Mortalit\"at}\gT{\textcolor[rgb]{0.7529412,0.7529412,0.7529412}{Mortalit\"at}}:
A\gZ{nteil der an XY Verstorbenen einer bestimmten
Bev\"olkerungsgruppe (z.B. 5 pro 100.000 EW)}}

{\mdseries
\index{Letalit\"at}\gT{Letalit\"at}:  Anteil der an XY
Verstorbenen beozgen auf die Gesamtzahl der Erkrankten
({\quotedblbase}1 von 5 Erkranten stirbt{\textquotedblright})}

\end{multicols}
{\mdseries\color{black}
\index{Entz\"undung}Entz\"undung}

{\mdseries
Reaktion des K\"orpers}

\begin{enumerate}
\item \gEs{R\"otung }

\item \gEs{Schwellung }

\item \gEs{\"Uberw\"armung }

\item \gEs{Schmerz }

\item \gEs{Funktionsverlust }

\end{enumerate}
{\mdseries
Ausl\"oser:}

\begin{itemize}
\item {\mdseries
Fremdk\"orper}

\item {\mdseries
Krankheitserreger}

\item {\mdseries
Mechanische Beanspruchung}

\item {\mdseries
...}

\end{itemize}
{\mdseries
Anh\"angsel
{\textquotedblleft}\index{{}-itis}\gT{{}-itis}{\textquotedblright}:
Gastritis ${\rightarrow}$
{\quotedblbase}Magenentz\"undung{\textquotedblleft}, Hepatitis
${\rightarrow}$ Leberentz\"undung, Pankreatitis ${\rightarrow}$
Entz\"undung der Bauchspeicheldr\"use, Meningitis ${\rightarrow}$
Hirnhautentz\"undung, Rhinitis ${\rightarrow}$
{\quotedblbase}Schnupfen{\textquotedblleft}, Otitis ${\rightarrow}$
Ohrenentz\"undung, ...}

\paragraph[Die wichtigsten \"Ubertragungswege]{Die wichtigsten
\"Ubertragungswege}
\begin{itemize}
\item {\mdseries
f\"akal-oral (Schmierinfektion, zB. Badesee)}

\item {\mdseries
Aerogen (Tr\"opfcheninfektion)}

\item {\mdseries
Genital (Geschlechtsverkehr)}

\item {\mdseries
\"Uber/durch die Haut}

\item {\mdseries
Diaplazent\"ar (w\"ahrend Schwangerschaft)}

\item {\mdseries
W\"ahrend der Geburt}

\item {\mdseries
Eintrittspforten: alles was irgendwie offen ist}

\end{itemize}
\begin{itemize}
\item {\mdseries
Wunden, Magen-Darm-Trakt, Atmung, Auge, Penis/Scheide, Plazenta ...}

\end{itemize}
\begin{itemize}
\item {\mdseries
Fremdk\"orper sind Wohnh\"auser und Autobahnen f\"ur Keime}

\end{itemize}
\begin{itemize}
\item {\mdseries
Harnkatheter, ven\"ose Zug\"ange, Tubus, k\"unstlicher Mageneingang,
Dialysekatheter {\dots}}

\end{itemize}
\paragraph[]{}
\paragraph[Kreuzinfektion]{Kreuzinfektion}
\begin{center}
\begin{minipage}{6.477cm}
{\mdseries
\gH{${\leftarrow}$ DAVOR f\"urchten wir uns !!!}}

\end{minipage}
\end{center}
{\mdseries
\index{Kreuzinfektion}\gT{Kreuzinfektion}:
\textbf{Patient 1} \textbf{\textcolor{magenta}{${\rightarrow}$}}
\textbf{\textit{Sanit\"ater}} \textbf{\textcolor{red}{${\rightarrow}$}}
\textbf{Patient 2}}

{\mdseries\color[rgb]{0.0,0.0,0.5019608}
Die wichtigsten Erreger}

\begin{itemize}
\item {\mdseries
Bakterien}

\end{itemize}
\begin{itemize}
\item {\mdseries
klein, im Mikroskop meist sichtbar;
{\textquotedblleft}Normale{\textquotedblright} Zellen}

\end{itemize}
\begin{itemize}
\item {\mdseries
Viren}

\end{itemize}
\begin{itemize}
\item {\mdseries
sicher klein, im Mikroskop i.d.R. nicht sichtbar, keine echten Zellen,
m\"ussen Zellen befallen}

\end{itemize}
\begin{itemize}
\item {\mdseries
Pilze}

\item {\mdseries
Tiere, Tierparasiten}

\end{itemize}
\begin{itemize}
\item {\mdseries
W\"urmer}

\end{itemize}
\begin{itemize}
\item {\mdseries
(Toxine)}

\end{itemize}
\paragraph[Bakterien]{Bakterien}
\index{Bakterien}\begin{itemize}
\item {\mdseries
{\quotedblbase}Echte{\textquotedblleft} Zellen, die leben und aktiv
etwas tun}

\end{itemize}
\begin{itemize}
\item {\mdseries
zB. Produktion von Substanzen (Toxine ...), Gasen, ...}

\end{itemize}
\begin{itemize}
\item {\mdseries
Viele Bakterien k\"onnen f\"ur den Menschen gut, egal oder auch b\"ose
sein}

\item {\mdseries
Haben die unterschiedlichsten Formen}

\end{itemize}
\begin{itemize}
\item {\mdseries
{\textquotedblleft}Kugeln{\textquotedblright}, St\"abchen, mit
Gei{\ss}eln ... }

\end{itemize}
{\mdseries\color{black}
Bakterien und Antibiotika}

\begin{itemize}
\item {\mdseries
Antibiotika t\"oten prinzipiell Bakterien, aber:}

\end{itemize}
\begin{itemize}
\item {\mdseries
Antibiotika haben ein bestimmtes Spektrum, dh. wirken nur auf bestimmte
Keime}

\item {\mdseries
Bakterien entwickeln Resistenzen gegen Antibiotika}

\item {\mdseries
Antibiotika helfen dann nicht mehr}

\end{itemize}
{\mdseries
\textbf{Resistenzen} sind ein enorm gro{\ss}es Problem in der heutigen
Medizin!}

{\mdseries
... und damit auch im Rettungs- und Krankentransport.}

{\mdseries
{\textquotedblleft}\index{Spitalskeime}\gT{Spitalskeime}{\textquotedblright}:
Keime mit {\textquotedblleft}angez\"uchteten{\textquotedblright}
Resistenzen}

{\mdseries
Unterscheide:
{\textquotedblleft}Wald-und-Wiesen{\textquotedblright}-Keime und
Spitalskeime }

{\mdseries
${\rightarrow}$ Keim ist nicht gleich Keim, auch wenn sie den gleichen
Namen haben}

{\mdseries\color{black}
Bakterien und Pathogenit\"at}

\begin{flushleft}
\tablehead{}\begin{supertabular}{|m{7.605cm}|m{5.688cm}|}
\hline
\begin{itemize}
\item {\mdseries Die B\"osartigkeit von Bakterien h\"angt u.a. ab von:}

\end{itemize}
\begin{itemize}
\item {\mdseries Welches Bakterium}

\item {\mdseries Spezies, Stamm, Resistenzen}

\item {\mdseries Wo ist die Infektion}

\item {\mdseries Wie besiedelt ist die Umgebung}

\item {\mdseries Wie ist der Immunstatus des Patienten}

\item {\mdseries Kind, Chemotherapie, AIDS, ...}

\item {\mdseries K\"orperliche Verfassung}

\item {\mdseries Alter, Ern\"ahrung, Stress, ...}

\end{itemize}
\begin{itemize}
\item {\mdseries Manche Keime machen immer krank}

\item {\mdseries Manche Keime nur unter Umst\"anden}

\item {\mdseries \gEs{Opportunisten} nur wenn der
Wirt geschw\"acht ist, und das sind viele Patienten im Rettungs- und
Krankentransport!}

\end{itemize}
 &
{\centering   [Warning: Image ignored]
% Unhandled or unsupported graphics:
%\includegraphics[width=2.858cm,height=5.001cm]{AASdev20090228-img96.jpg}
 \par}

\centering F\"ur manche Bakterien zahlt man freiwillig viel Geld, hier
am Beispiel des Lacotbacillus casei gezeigt. \par

\\\hline
\end{supertabular}
\end{flushleft}
\paragraph[Viren]{Viren}
\index{Viren}\begin{itemize}
\item {\mdseries
Keine Zellen, sehr klein, im Lichtmikroskop nicht sichtbar}

\item {\mdseries
Brauchen Wirtszelle um sich zu vermehren}

\item {\mdseries
Schleusen eigenes genetisches Programm in die Wirtszelle ein}

\item {\mdseries
\gZ{Nachweis:}}

\end{itemize}
\begin{itemize}
\item {\mdseries
\gZ{NICHT anz\"uchtbar auf N\"ahrb\"oden}}

\item {\mdseries
\gZ{Antigen- oder Antik\"orpernachweis}}

\item {\mdseries
\gZ{PCR (DNA/RNA-Nachweis)}}

\end{itemize}
\begin{itemize}
\item {\mdseries
zB: viele Verk\"uhlungen, Grippe, Feuchtblattern, viele Hepatitis-Formen
...}

\end{itemize}
\subsection{Wichtige Krankheitsbilder}
{\mdseries\color[rgb]{0.0,0.0,0.5019608}
\index{Grippe}Grippe
(\index{Influenza}\index{Influenza}\gT{Influenza})}

{\mdseries
Die Grippe ist eine \gE{virale} Erkrankung, welche je
nach Patient auch einen sehr schweren Verlauf nehmen kann und
\gEs{Komplikationen} wie Lungen- oder 
Herzmuskelentz\"undungen oder weitere Infektionen nach sich ziehen
kann. Besonders gef\"ahrdet sind \"altere und immunsupprimierte
Personen. Gegen die Grippe wird eine
\gEs{Schutzimpfung} angeboten. Diese wirkt f\"ur
ein Jahr und muss jede Saison widerholt werden. Die Impfung bietet
keinen hundertprozentigen Schutz, wird aber trotzdem f\"ur
Risikopatienten und Personen die einem erh\"oten Infektionsrisiko
ausgesetzt sind (darunter f\"allt auch Personal im Gesundheitsbereich)
empfohlen. }

{\mdseries
Da es sich bei der Grippe um eine durch Viren verursachte Krankheit
handelt sind Antibiotika wirkungslos.}

{\mdseries\color{black}
Die Grippe tritt in Wellen auf und kann t\"odlich enden}

{\mdseries
Die Grippe tritt besonders in der kalten Jahreszeit auf und tritt
regelm\"a{\ss}ig als \index{Epidemie}\gEs{Epidemie}
auf. Es kann auch zum Ausburch einer weltweiten
\index{Pandemie}\gE{Pandemie} kommen,
\gZ{z.B. Spanische Grippe (1918, 25~Mio~Tote),
Asiatische Grippe (1957, 1~Mio~Tote) und  Hongkong-Grippe (1968,
800.000~Tote)}\footnote{http://de.wikipedia.org/wiki/Influenza}\gZ{.
}}

{\mdseries
In \"Osterreich erkranken pro Jahr ca. 380.000 Personen an der
Influenza, ca. 4.500 Personen m\"ussen station\"ar behandelt werden und
ca. 120 Personen sterben an den
Folgen\footnote{http://www.bmg.gv.at/cms/site/standard.html?channel=CH0742\&doc=CMS1075899626901}.}

{\mdseries\color{black}
Symptome}

{\mdseries
Typisch f\"ur die Grippe ist ist das \gEs{rasche,
hohe anfiebern (38{\textdegree}-40{\textdegree}C) mit Sch\"uttelfrost, 
Kopf-, Hals-, Muskel-, Gelenksschmerzen, Husten und stark reduziertem
Allgemeinszustand} sowie langdauernde Abgeschlagenheit nach Ende der
Erkrankung. Bei schon vorher gschw\"achten Patienten kann der Verlauf
weniger spektakul\"ar ausshehen (kaum Fieber), das Risiko einer
Komplikationen ist allerdings nat\"urlich h\"oher. }

{\mdseries\color[rgb]{0.0,0.0,0.5019608}
\index{grippaler Infekt}grippaler Infekt,
\index{Verk\"uhlung}Verk\"uhlung, \index{Erk\"altung}Erk\"altung}

{\mdseries
Die Verk\"uhlung ist eine Infektion der oberen Luftwege mit Viren,
allerdings mit anderen als bei der {\quotedblbase}echten
Grippe{\textquotedblleft} (Influenza). Die Symptome \"ahneln denen der
Grippe, allerdings sind sie bei weitem nicht so stark ausgepr\"agt. Im
Vordergrund stehen Halsschmerzen, Husten, rinnende oder verstopfte Nase
und eine eher leicht erh\"ohte Temperatur. }

{\mdseries\color[rgb]{0.0,0.0,0.5019608}
\index{Tetanus}Tetanus (\index{Wundstarrkrampf}Wundstarrkrampf) [!!!]}

{\mdseries
\gZ{\textit{Clostridium tetani}}}

\begin{center}
 [Warning: Image ignored] % Unhandled or unsupported graphics:
%\includegraphics[width=5.984cm,height=3.539cm]{AASdev20090228-img97.gif}

\end{center}
\begin{itemize}
\item {\mdseries
Der Erreger kommt \"uberall in der Umgebung vor}

\item {\mdseries
Produziert Toxin durch welches die Skelettmuskulatur tonisch krampft}

\item {\mdseries
\gZ{Inkubation: 4-21 Tage}}

\item {\mdseries
\gZ{Vorzeichen: M\"udigkeit, Inappetenz}}

\item {\mdseries
\gEs{Klinik: Tonische Kr\"ampfe der quergestreiften
Muskulatur, oft beginnend mit Kaumuskulatur}}

\item {\mdseries
\gEs{Problem: Schmerzen, Ateml\"ahmung}}

\item {\mdseries
\gEs{Letalit\"at ca. 50\%}}

\item {\mdseries
\gEs{Impfung!}}

\end{itemize}
{\mdseries\color{black}
Tetanus \index{Impfschutz}Impfschutz}

\begin{itemize}
\item {\mdseries
Auch bei kleinen Verletzungen nach dem Impfschutz fragen und
\gEs{dokumentieren}!}

\item {\mdseries
Patient soll unverz\"uglich im Impfpass nachschauen.}

\item {\mdseries
Wenn die letzte Impfung l\"anger als \gEs{5 Jahre}
zur\"uck liegt oder der Patient noch nie geimpft wurde muss der Patient
unverz\"uglich eine Unfallabetilung aufsuchen. }

\item {\mdseries
Am Revers best\"atigen lassen!}

\end{itemize}
{\mdseries\color{black}
Eine mangelhafte Aufkl\"arung oder Dokumentation bez\"uglich Tetanus und
dessen Impfschutz ist grob fahrl\"assig und ein unverzeihbarer
Kunstfehler.}

{\mdseries\color[rgb]{0.0,0.0,0.5019608}
Tuberkulose }

{\mdseries
\gZ{Mykobacterium Tuberculosis},
{\quotedblbase}\gT{Tbc}{\textquotedblleft}}

\begin{center}
 [Warning: Image ignored] % Unhandled or unsupported graphics:
%\includegraphics[width=4.615cm,height=6.153cm]{AASdev20090228-img98.jpg}

\end{center}
\begin{itemize}
\item {\mdseries
Kann praktisch alle Organe befallen }

\item {\mdseries
Der Verlauf wird vom Organbefall bestimmt.}

\end{itemize}
\begin{itemize}
\item {\mdseries
Lunge, Skelett, Darm, Haut, Genitalien, ...}

\end{itemize}
\begin{itemize}
\item {\mdseries
Ausgangspunkt ist fast immer die \gEs{Lunge}.}

\item {\mdseries
\"Ubertragung: aerogen, oral, Hautkontakt, plazent\"ar}

\end{itemize}
\begin{itemize}
\item {\mdseries
Abh\"angig von der Herdlokalisation}

\end{itemize}
\begin{itemize}
\item {\mdseries
Inkubationszeit: 4-12 Wochen}

\item {\mdseries
Klinik:}

\end{itemize}
\begin{itemize}
\item {\mdseries
\gEs{Subfebrile Temperatur, Nachtschwei{\ss}, Husten,
evt. Bluthusten}, Lymphknotenschwellung,  Pleuritis}

\end{itemize}
\begin{itemize}
\item {\mdseries
{\textquotedblleft}\index{Offene Tb}\gT{Offene
Tb}{\textquotedblright}: Infektionsherd hat
{\textquotedblleft}Luftkontakt{\textquotedblright} ${\rightarrow}$
Tr\"opfcheninfektion m\"oglich ${\rightarrow}$ Schutzmaske}

\item {\mdseries
\gEs{Meldung an Leitstelle \& Betriebsleitung}!}

\end{itemize}
{\mdseries\color[rgb]{0.0,0.0,0.5019608}
\index{HIV}HIV - \index{Humanes Immundefizienz Virus}Humanes
Immundefizienz Virus}

\begin{itemize}
\item {\mdseries
Bef\"allt vor allem die Zellen des Abwehrsystems }

\item {\mdseries
Es vermehrt sich in ihnen, setzt sie au{\ss}er Funktion und zerst\"ort
sie schlie{\ss}lich.}

\item {\mdseries
Das k\"orpereigene Abwehrsystem kann HIV nicht aus dem K\"orper
entfernen.}

\end{itemize}
\begin{itemize}
\item {\mdseries
${\rightarrow}$Immunschw\"ache}

\end{itemize}
\begin{itemize}
\item {\mdseries
HIV verursacht AIDS (Aquired Immune Deficiency Syndrom)}

\item {\mdseries
Begleiterkrankungen: }

\end{itemize}
\begin{itemize}
\item {\mdseries
Infektionen mit normalerweise harmlosen Erregern}

\end{itemize}
\begin{itemize}
\item {\mdseries
Pilze, {\textquotedblleft}fakultative{\textquotedblright} Erreger,
{\textquotedblleft}Opportunisten{\textquotedblright}, seltene Arten der
Lungenentz\"undung}

\end{itemize}
\begin{itemize}
\item {\mdseries
Tumore}

\end{itemize}
\begin{itemize}
\item {\mdseries
\"Ubertragung}

\end{itemize}
\begin{itemize}
\item {\mdseries
K\"orperfl\"ussigkeiten}

\item {\mdseries
Blut, Blutprodukte, Sperma, Scheidensekret}

\item {\mdseries
Ungesch\"utzter Sex, Transfusionen}

\item {\mdseries
Nadelstichverletzungen}

\item {\mdseries
KEINE \"Ubertragung durch Sozialkontakte (H\"andesch\"utteln, Haushalt,
Geschirr, Toilettenanlagen)}

\end{itemize}
\begin{itemize}
\item {\mdseries
Erst nach 6-12 Wochen sind Antik\"orper nachweisbar}

\end{itemize}
\begin{itemize}
\item {\mdseries
\gEs{HIV-Test nicht
{\textquotedblleft}aktuell{\textquotedblright}, {\textquotedblleft}12
Wochen blind{\textquotedblright}}}

\end{itemize}
\begin{itemize}
\item {\mdseries
Nach ca 5-10 Jahren bricht AIDS aus}

\item {\mdseries
Therapie:}

\end{itemize}
\begin{itemize}
\item {\mdseries
Keine Heilung m\"oglich, aber deutliche Lebensverl\"angerung}

\end{itemize}
\begin{itemize}
\item {\mdseries
Prophylaxe}

\end{itemize}
\begin{itemize}
\item {\mdseries
Postexpositionsprohylaxe bei Nadelstichverletzungen}

\item {\mdseries
Kondome}

\end{itemize}
{\mdseries\color[rgb]{0.0,0.0,0.5019608}
(Virale) \index{Hepatitis}Hepatitis - versch. Erreger (A-E)}

{\mdseries
Impfung nur f\"ur Hepatitis A und B verf\"ugbar
(Twinrix{\texttrademark})}

{\mdseries
Achtung: Non-Responder m\"oglich!}

{\mdseries\color{black}
Virale Hepatitis - Klinik und Verlauf}

\begin{itemize}
\item {\mdseries
Grunds\"atzlich
{\textquotedblleft}Leberentz\"undung{\textquotedblright}}

\item {\mdseries
Vergr\"o{\ss}erte, schmerzhafte Leber }

\item {\mdseries
Ikterus (Gelbsucht) }

\item {\mdseries
Dunkler Urin (heller Sch\"uttelschaum),}

\item {\mdseries
evt. heller Stuhl (Acholie)}

\item {\mdseries
Unterscheide {\textquotedblleft}akut{\textquotedblright} und
{\textquotedblleft}chronisch{\textquotedblright}}

\end{itemize}
\begin{itemize}
\item {\mdseries
Chronisch: wenn Hepatitis nicht abheilt}

\end{itemize}
\begin{itemize}
\item {\mdseries
Ende: Leberversagen, Tod}

\end{itemize}
{\mdseries\color{black}
Hepatitis-A - Hepatitis-A-Virus (HAV)}

\begin{itemize}
\item {\mdseries
\gZ{Inkubationszeit 10 - 40 Tage, }}

\item {\mdseries
typische Reisekrankheit (Mittelmeerraum, S\"udamerika, Orient), meist
epidemisch, }

\item {\mdseries
Dauer einige Wochen, }

\item {\mdseries
kein \"Ubergang in chronischen Verlauf}

\end{itemize}
{\mdseries\color{black}
Hepatitis-B - Hepatitis-B-Virus (HBV)}

\begin{itemize}
\item {\mdseries
\gZ{Inkubationszeit 4 - 12 Wochen, }}

\item {\mdseries
meist sporadisches Auftreten, }

\item {\mdseries
\"Ubertragung durch Speichel, Blut(produkte), Sperma und andere
\gEs{K\"orperfl\"ussigkeiten}, }

\item {\mdseries
chronischer Verlauf in ca. 5 \%}

\end{itemize}
{\mdseries\color{black}
Hepatitis C - Hepatitis-C-Virus (HCV)}

\begin{itemize}
\item {\mdseries
Inkubationszeit 6 - 12 Wochen, }

\item {\mdseries
h\"aufig nach Bluttransfusion oder i.v. Drogenabusus, }

\item {\mdseries
chronischer Verlauf (in 70 -- 80~\% der F\"alle)}

\end{itemize}
{\mdseries
F\"ur den Rettungsdienst sind besonders Hep B und C wichtig: Hohe
Infektiosit\"at, chronischer Verlauf!}

{\mdseries
Gegen Hep C gibt es keine Impfung!}

{\mdseries\color[rgb]{0.0,0.0,0.5019608}
\label{bkm:RefHeading41252510}Sepsis}

{\mdseries
Siehe auch: Septischer Schock und Toxischer Schock
[H.V{\guillemotright}\pageref{bkm:RefHeading41284910}, oben]}

\begin{itemize}
\item {\mdseries
Wenn Keime und deren Produkt ausser Kontrolle geraten erfolgt Amoklauf
des K\"orpers:}

\item {\mdseries
Ungehemmte Freisetzung von Stoffen des Entz\"undungs-, Gerinnungs- und
Immunsystems}

\item {\mdseries
Auch ohne Infektion m\"oglich:}

\end{itemize}
\begin{itemize}
\item {\mdseries
\gZ{Systemic Inflammatory Response Syndrome (SIRS) }}

\end{itemize}
\begin{itemize}
\item {\mdseries
nach schwerem Trauma, Verbrennungen, Gewebshypoxie}

\item {\mdseries
Klinik abh\"angig vom Fortschritt und k\"orperlichen Abbau.}

\end{itemize}
\begin{itemize}
\item {\mdseries
Sehr stark reduzierter Allgemeinzustand}

\item {\mdseries
Fieber ({\textgreater}38{\textdegree} C) oder Hypothermie
({\textless}36{\textdegree} C)}

\item {\mdseries
Herzfrequenz {\textgreater}90/min}

\item {\mdseries
Atemfrequenz {\textgreater}20/min oder Hyperventilation}

\item {\mdseries
unkontrollierbare Gerinnung/Blutung}

\item {\mdseries
Hypotonie ${\rightarrow}$ Septischer Schock}

\item {\mdseries
Organversagen  }

\end{itemize}
{\mdseries\color[rgb]{0.0,0.0,0.5019608}
Multiresistente Keime}

{\mdseries
Vor der Entdeckung der Antibiotika waren bakterielle Infektionen sehr
gef\"urchtet, viele Menschen sind an -- nach heutiger Sicht simplen --
Infektionen gestorben. }

{\mdseries
Heute kann man zwar solche Infektionen mit Antibiotika behandeln, jedoch
k\"onnen Bakterien \gEs{Resistenzen} gegen
Antibiotika entwickeln. Die sonst wirksamen Medikamente wirken dann
nicht mehr oder nur noch sehr abgeschw\"acht. Als behandelnder Arzt
steht man dann hilflos vor dem Patienten. }

{\mdseries
\gT{Resistenzen} machen Antibiotika wirkungslos. }

{\mdseries
Es muss daher das Ziel aller sein die Ausbreitung von resistenten Keimen
zu verhindern. Dem Rettungs- und Krankentransport kommt dabei eine
\gEs{Schl\"usselrolle} zu, da es hier leicht zu
\gEs{Kreuzinfektionen} kommen kann. }

{\mdseries
\gT{Resistenzen} sind ein enorm gro{\ss}es Problem in der
heutigen Medizin!}

{\mdseries
... und damit auch im Rettungs- und Krankentransport, da es hier leicht
zu \gEs{Kreuzinfektionen} kommen kann..}

{\mdseries
{\textquotedblleft}\index{Spitalskeime}\gT{Spitalskeime}{\textquotedblright}:
Keime mit {\textquotedblleft}angez\"uchteten{\textquotedblright}
Resistenzen}

{\mdseries
Unterscheide:
{\textquotedblleft}Wald-und-Wiesen{\textquotedblright}-Keime und
Spitalskeime }

{\mdseries
Der wohl bekannteste Problemkeim ist der \gT{MRSA}
(Multiresistenter Staphylococcus aureus, siehe unten). Das liegt
schlicht daran dass er der erste war. In den letzten Jahren ist die
Resistenzentwicklung aber auch bei vielen anderen Keimen zu beobachten,
Beispiele daf\"ur sind die
\gT{\textit{ESBL}}\textit{{\textquotesingle}s} und
\gT{\textit{VRE}}\textit{{\textquotesingle}s} sowie der
\gT{\textit{VRSA}} \gZ{(Extended
Spectrum Beta-Laktamase-Bildner, Vancomycin-resistente Enterokokken,
Vancomycin-resistenter Staph. aureus)}. Das Problem der
multiresistenten Keime wird immer gr\"o{\ss}er.}

{\mdseries
Je wirkungsloser die Antibiotika werden desto wichtiger ist die Hygiene
und die \gEs{Vermeidung von (Kreuz-)Infektionen}. }

\begin{minipage}{13.67cm}
[Warning: Draw object ignored]

{\mdseries
Zeichnung \stepcounter{Drawing}{\theDrawing}: Teufelskreis der
resistenten Keime. \gEs{Oft das Bindeglied: Der
Rettungs- und Krankentransportdienst}. Quelle: }

\end{minipage}

\paragraph[MRSA - Multi-Resistenter Staphylococcus aureus]{MRSA -
Multi-Resistenter Staphylococcus aureus}
{\mdseries
\gZ{Urspr\"unglich: Methicilin-Resistenter
Staphylococcus Aureus}}

\begin{itemize}
\item {\mdseries
Staphylococcus aureus ist ein
{\textquotedblleft}Allerweltskeim{\textquotedblright}, bis zu 50\% der
Bev\"olkerung sind besiedelt}

\item {\mdseries
Vor allem vordere Nasenh\"ole besiedelt, aber auch Haut (Achselh\"ole,
Analregion), Schleimh\"aute}

\item {\mdseries
Auch im Spital h\"aufiger Keim, aber dort oft Antibiotika-Resistent:
MRSA}

\end{itemize}
\begin{itemize}
\item {\mdseries
Multi-resistenter Staph. aureus}

\item {\mdseries
{\textquotedblleft}Spitalsinfektionen{\textquotedblright}}

\item {\mdseries
schwer bis gar nicht behandelbar}

\item {\mdseries
Wunden, die nicht abheilen wollen, Sepsis, Tod ...}

\end{itemize}
\subparagraph{}
\subparagraph[MRSA-Transport - Allgemeine Ma{\ss}nahmen]{MRSA-Transport
- Allgemeine Ma{\ss}nahmen}
\label{bkm:RefHeading40262819}\begin{itemize}
\item {\mdseries
Informationen \"uber MRSA-Tr\"agerschaft}

\begin{center}
\begin{minipage}{4.941cm}
{\mdseries
\gH{ gilt im Prinzip auch f\"ur die anderen
{\quotedblbase}Problemkeime{\textquotedblleft}}}

\end{minipage}
\end{center}
\item {\mdseries
Art/Ort der MRSA-Besiedlung}

\item {\mdseries
Patientenvorbereitung und Transport}

\item {\mdseries
Wunden sind frisch verbunden und gut abgedeckt}

\item {\mdseries
Bei Besiedlung der Atemwege tr\"agt der Patient einen Mund-
/Nasenschutz}

\item {\mdseries
Vor dem Transport f\"uhrt der Patient eine hygienische
H\"andedesinfektion durch}

\item {\mdseries
Allgemeine Hygienema{\ss}nahmen}

\end{itemize}
\begin{itemize}
\item {\mdseries
Einmalhandschuhe und Schutzkittel.}

\item {\mdseries
Nach Ablegen sofortige hygienische H\"andedesinfektion  }

\end{itemize}
\begin{itemize}
\item {\mdseries
Nach Abschluss des Patiententransportes sind alle Materialien, Ger\"ate,
Instrumente und Fl\"achen, welche direkten Kontakt mit dem Patienten
hatten gem\"a{\ss} dem \"ublichen Hygieneplan zu desinfizieren.}

\end{itemize}
\begin{itemize}
\item {\mdseries
Kugelschreiber!!!}

\item {\mdseries
Alle waagerechten Oberfl\"achen des Fahrzeuginnenraumes sind mit einem
Fl\"achendesinfektionsmittel in Form einer Scheuerwischdesinfektion
desinfizierend zu reinigen.}

\item {\mdseries
Abschlie{\ss}end f\"uhrt das Einsatzpersonal eine hygienische
H\"andedesinfektion durch.}

\end{itemize}
\begin{itemize}
\item {\mdseries
Nach Abschluss der vorangegangenen Hygiene- und
Desinfektionsma{\ss}nahmen ist das Fahrzeug wieder voll einsatzbereit.}

\end{itemize}
{\mdseries
F\"ur die anderen Problemkeime (VRSA, VRE, ESBL, {\dots}) gelten
sinngem\"a{\ss} die gleichen Verfahrensweisen. }

\subsection[Desinfektion]{Desinfektion}
\index{Desinfektion}\begin{center}
\tablehead{}\begin{supertabular}{|m{6.649cm}|m{6.651cm}|}
\hline
 &
\\\hline
{\mdseries \gT{DESINFEKTION}}

{\mdseries Ziel: Keimreduktion ${\rightarrow}$ nicht infekti\"os}

{\mdseries Richtet sich gegen ausgew\"ahlte Keime}

{\mdseries Aber: nicht keimfrei! nicht steril!}

{\mdseries Chemische Substanzen}

{\mdseries Einlegebad}

 &
{\mdseries \gT{STERILISATION}}

{\mdseries Abt\"otung aller Mikroorganismen -
{\textquotedblleft}keimfrei{\textquotedblright}}

{\mdseries Dampfsterilisation (\"Uberdruck)}

{\mdseries Hei{\ss}luftsterilisation}

\\\hline
\end{supertabular}
\end{center}
{\mdseries\color{black}
\index{Hygieneplan}Hygieneplan - Hygienebeauftragter}

{\mdseries
Der Hygieneplan regelt die konkreten Hygienema{\ss}nahmen und Was Wann
Wie Womit (Wo) zu geschehen hat, siehe allgemeiner Aushang.}

{\mdseries
Der derzeitige \gEs{Hygienebeauftragte} ist
\_\_\_\_\_\_\_\_\_\_\_\_\_\_\_\_\_\_\_\_\_\_\_\_\_\_\_\_\_\_\_\_\_\_.}

{\mdseries\color[rgb]{0.0,0.0,0.5019608}
Pers\"onliche Hygiene}

\begin{itemize}
\item {\mdseries
Schmuck, Uhren, Freundschaftsarmb\"ander sind perfekte N\"ahrb\"oden
${\rightarrow}$ Vor dem Dienst ablegen!}

\begin{center}
 [Warning: Image ignored] % Unhandled or unsupported graphics:
%\includegraphics[width=3.112cm,height=3.941cm]{AASdev20090228-img99.jpg}

\end{center}
\item {\mdseries
Dienstkleidung regelm\"a{\ss}ig waschen!}

\item {\mdseries
Bei Kontamination wechseln!}

\item {\mdseries
Verwendung von Plastiksch\"urzen}

\item {\mdseries
Bei Dienstantritt}

\end{itemize}
\begin{itemize}
\item {\mdseries
H\"ande waschen}

\item {\mdseries
Nagelfalze reinigen}

\item {\mdseries
\index{Hygienische H\"andedesinfektion}Hygienische
\index{H\"andedesinfektion}H\"andedesinfektion}

\end{itemize}
{\mdseries\color[rgb]{0.0,0.0,0.5019608}
Wischdesinfektion}

\begin{itemize}
\item {\mdseries
Mit \index{Fl\"achendesinfektionsmittel}Fl\"achendesinfektionsmittel }

\item {\mdseries
Einwirkzeit lt. Packung beachten!}

\item {\mdseries
Fl\"ache muss w\"ahrend der Zeit feucht sein}

\item {\mdseries
in Autos nur Fl\"achendesinfektionsmittel ohne Alkohol!}

\end{itemize}
\begin{itemize}
\item {\mdseries
Alkohol zerst\"ort Acryl-Fl\"achen! ${\rightarrow}$
\index{Acryl-Des{\texttrademark}}\gT{Acryl-Des{\texttrademark}}
 }

\end{itemize}
\begin{itemize}
\item {\mdseries
\gEs{Von au{\ss}en in Richtung Schmutzquelle
wischen}}

\item {\mdseries
Einmalt\"ucher wechseln}

\end{itemize}
\begin{itemize}
\item {\mdseries
nicht den Schmutz \"uberall verteilen}

\end{itemize}
{\mdseries\color{black}
In den Fahrzeugen d\"urfen keine alkoholischen Desinfektionsmittel
verwendet werden da diese die Acryloberfl\"achen zerst\"oren k\"onnen!}

{\mdseries\color[rgb]{0.0,0.0,0.5019608}
Ger\"atedesinfektion}

\begin{itemize}
\item {\mdseries
Nach Kontamination:}

\end{itemize}
\begin{itemize}
\item {\mdseries
\textbf{Grobe Reinigung vor Ort}: Eintrocknen verhindern}

\begin{center}
 [Warning: Image ignored] % Unhandled or unsupported graphics:
%\includegraphics[width=4.737cm,height=2.842cm]{AASdev20090228-img100.jpg}
\captionof{figure}[Desinfektionswannen, verschiedene
Gr\"o{\ss}en]{Desinfektionswannen, verschiedene Gr\"o{\ss}en}

\end{center}
\end{itemize}
\begin{itemize}
\item {\mdseries
Wischdesinfektion}

\item {\mdseries
Ggfs. Schl\"auche durchsp\"ulen}

\item {\mdseries
Fl\"achendesinfektionsmittel in Nierenschale / Joghurtbecher, saugen
...}

\end{itemize}
\begin{itemize}
\item {\mdseries
Meldung an Leitstelle: {\textquotedblleft}Nicht EB, m\"ussen
putzen{\textquotedblright}}

\item {\mdseries
Kontaminierte Teile gegen Ersatzmaterial tauschen}

\end{itemize}
\begin{itemize}
\item {\mdseries
Zerlegen}

\item {\mdseries
Nochmals sp\"ulen}

\end{itemize}
\begin{itemize}
\item {\mdseries
Wenn Desinfektionswanne leer ${\rightarrow}$ einlegen }

\end{itemize}
\begin{itemize}
\item {\mdseries
\gEs{Vollst\"andig zerlegt} laut Anleitung}

\item {\mdseries
\gEs{Luftblasenfrei}!}

\item {\mdseries
Mit Niederhaltesieb unter die Oberfl\"ache dr\"ucken}

\item {\mdseries
Zeit auf Zettel notieren }

\end{itemize}
\begin{itemize}
\item {\mdseries
{\textquotedblleft}was, wann, von wo, von wem
eingelegt{\textquotedblright}}

\end{itemize}
\begin{itemize}
\item {\mdseries
sonst nur bereit legen, in Plastiksack}

\end{itemize}
\begin{itemize}
\item {\mdseries
{\textquotedblleft}was, von wo, von wem, warum{\textquotedblright}}

\end{itemize}
{\mdseries\color[rgb]{0.0,0.0,0.5019608}
Hygienische H\"andedesinfektion}

\begin{itemize}
\item {\mdseries
Ziel: }

\end{itemize}
\begin{itemize}
\item {\mdseries
Keimreduktion}

\end{itemize}
\begin{itemize}
\item {\mdseries
Vermeidung von Kreuzinfektionen}

\item {\mdseries
Patient1 ${\rightarrow}$ Sanit\"ater ${\rightarrow}$ Patient 2}

\item {\mdseries
{\textquotedblleft}\gEs{Unterbrechung der
Infektionskette}{\textquotedblright}}

\end{itemize}
\begin{itemize}
\item {\mdseries
Durchzuf\"uhren:}

\end{itemize}
\begin{itemize}
\item {\mdseries
Dienstantritt}

\item {\mdseries
\gEs{nach jedem Patienten!!}}

\end{itemize}
\begin{itemize}
\item {\mdseries
Womit?}

\end{itemize}
\begin{itemize}
\item {\mdseries
H\"andedesinfektionsl\"osung: Sterilium, Monopronto}

\begin{center}
 [Warning: Image ignored] % Unhandled or unsupported graphics:
%\includegraphics[width=4.001cm,height=4.001cm]{AASdev20090228-img101.jpg}

\end{center}
\end{itemize}
{\mdseries\color{black}
Vorgehen }

\begin{enumerate}
\item Ringe entfernen 

\item Eine Portion alkoholisches H\"andedesinfektionsmittel (ca. 3ml =
1~Hohlhand) aus dem Spen\-der entnehmen. 

\item mittels Ellenbogen, vollkommene Benetzung beider H\"ande und
Handgelenke 

\item Handinnenfl\"achen gegeneinander reiben

\item geschlossene Finger 

\item Handgelenke einreiben

\item Mit rechter Handinnenfl\"ache linken Handr\"ucken reiben und
umgekehrt, dabei Finger ineinander verschr\"anken.

\item Mit ineinander verschr\"ankten Fingern Handinnenfl\"achen
gegeneinander reiben. 

\begin{center}
 [Warning: Image ignored] % Unhandled or unsupported graphics:
%\includegraphics[width=3.961cm,height=3.169cm]{AASdev20090228-img102.jpg}

\end{center}
\item H\"ande ineinander verhaken und Finger gegeneinander bewegen

\item Daumen mit gegen\"uberliegender Hand vollst\"andig umschlie{\ss}en
und rotierend reiben. 

\item Daumenkuppen nicht vergessen!

\item Fingerkuppen im Handteller kreisf\"ormig einreiben bis Alkohol
verdunstet ist

\end{enumerate}
{\mdseries
 Einwirkzeit von mindestens 30 Sekunden beachten! }

[Warning: Draw object ignored]

{\mdseries\color[rgb]{0.0,0.0,0.5019608}
\"Offnen von sterilem Material}

{\mdseries
Packung {\quotedblbase}aufsch\"alen{\textquotedblleft} }

{\mdseries
  [Warning: Image ignored] % Unhandled or unsupported graphics:
%\includegraphics[width=6.656cm,height=4.993cm]{AASdev20090228-img103.jpg}
    [Warning: Image ignored] % Unhandled or unsupported graphics:
%\includegraphics[width=6.663cm,height=4.984cm]{AASdev20090228-img104.jpg}
 }

{\mdseries\color[rgb]{0.0,0.0,0.5019608}
Selbstschutz}

\begin{itemize}
\item {\mdseries
Immer Einmalhandschuhe beim Patientenkontakt tragen!}

\item {\mdseries
Handschuhe nach jedem Patienten wechseln}

\item {\mdseries
Nie mit schmutzigen Handschuhen in die Tasche greifen!}

\item {\mdseries
\gEs{Ein Helfer hat Patientenkontakt, der andere
reicht zu!}}

\item {\mdseries
Achtung {\textquotedblleft}Kugelschreiber{\textquotedblright}}

\item {\mdseries
Einmalsch\"urzen bei Bedarf}

\end{itemize}
\begin{itemize}
\item {\mdseries
evt. Infektionsschutzanz\"uge, Maske}

\end{itemize}
\begin{itemize}
\item {\mdseries
\gEs{Spitze Gegenst\"ande (Nadeln, Lanzetten, ...
}\gEs{sofort}\gEs{ in Gelbe
Box!}}

\end{itemize}
\begin{itemize}
\item {\mdseries
Ausnahme: ausdr\"uckliche Arztanweisung (Blutzuckerbestimmung aus
Nadel)}

\item {\mdseries
{\textquotedblleft}\gEs{Wer sticht, der
entsorgt!}{\textquotedblright}}

\end{itemize}
\paragraph{Kontamed, Sharp, Gelbe Box ...}
{\mdseries
Spitze Gegenst\"ande (Nadeln, Lanzetten, ...
\gEs{sofort} in Gelbe Box!}

{\mdseries
Ausnahme: ausdr\"uckliche Arztanweisung (Blutzuckerbestimmung aus
Nadel)}

{\mdseries
Wer sticht, der entsorgt!}



\begin{center}
 [Warning: Image ignored] % Unhandled or unsupported graphics:
%\includegraphics[width=3.377cm,height=4.552cm]{AASdev20090228-img105.jpg}
 [Warning: Image ignored] % Unhandled or unsupported graphics:
%\includegraphics[width=5.71cm,height=4.289cm]{AASdev20090228-img106.jpg}
\captionof{figure}[Kontamed-Boxen, verschiedene
Fabrikate.]{Kontamed-Boxen, verschiedene Fabrikate.}

\end{center}
{\mdseries
NIE Nachstopfen.}

{\mdseries
NIE ausleeren.}

{\mdseries
Wenn die Box voll ist wird sie verschlossen und abgegeben.}

{\mdseries\color[rgb]{0.0,0.0,0.5019608}
Wundreinigung}

\begin{itemize}
\item {\mdseries
Bei frischen, nicht infizierten Wunden:}

\end{itemize}
\begin{itemize}
\item {\mdseries
Reinigung von innen nach au{\ss}en}

\item {\mdseries
Ziel: keine Keime in Wunde bringen}

\end{itemize}
\begin{itemize}
\item {\mdseries
Womit?}

\end{itemize}
\begin{itemize}
\item {\mdseries
Sterilen L\"osungen oder Octenisept}

\item {\mdseries
Sterile Tupfer}

\end{itemize}
{\mdseries
KEINE Papierhandt\"ucher oder unsterile Tupfer !!!}

\begin{itemize}
\item {\mdseries
Danach steriler Verband!}

\end{itemize}
\begin{itemize}
\item {\mdseries
Bei offenen Frakturen ist keimfreiheit be\-sonders wichtig !!!}

\item {\mdseries
Gefahr einer lebenslangen Knochen\-marks\-ent\-z\"undung
(Osteomyelitis)!}

\item {\mdseries
\index{OpSite}\gT{OpSite}{}-Folie wenn vorhanden.}

\end{itemize}
 [Warning: Image ignored] % Unhandled or unsupported graphics:
%\includegraphics[width=5.742cm,height=5.658cm]{AASdev20090228-img107.png}
\captionof{figure}[Reiningung einer sauberen Wunde: Von Innen nach
Aussen, der Tupfer wird anschliessend verworfen.]{Reiningung einer
sauberen Wunde: Von Innen nach Aussen, der Tupfer wird anschliessend
verworfen.}
 [Warning: Image ignored] % Unhandled or unsupported graphics:
%\includegraphics[width=4.219cm,height=4.219cm]{AASdev20090228-img108.jpg}
\captionof{figure}[Anbringen einer OpSite(TM)-Folie: Es ist wichtig die
Folie zu spannen damit sie sich nicht selbst verklebt. ]{Anbringen
einer OpSite(TM)-Folie: Es ist wichtig die Folie zu spannen damit sie
sich nicht selbst verklebt. }


\subsection{Besondere Verhaltensweisen}
{\mdseries\color[rgb]{0.0,0.0,0.5019608}
%
%
\index{Nadelstichverletzungen}Nadelstichverletzungen}

\begin{itemize}
\item {\mdseries
Erstma{\ss}nahmen durchf\"uhren: }

\end{itemize}
\begin{itemize}
\item {\mdseries
bei Stich-bzw.Schnittverletzungen: Wunde auspressen und 10 Minuten
desinfizieren}

\item {\mdseries
bei Spritzern mit Blut oder Sekreten in Auge oder Mund: Schleimhaut mit
Wasser u./o. Schleimhautdesinfektionsmittel sp\"ulen}

\end{itemize}
\begin{itemize}
\item {\mdseries
Leitstelle informieren}

\item {\mdseries
${\rightarrow}$ \gEs{Umgehend} in ein Unfallspital
der AUVA}

\item {\mdseries
Patient: Zielspital informieren:}

\end{itemize}
\begin{itemize}
\item {\mdseries
Blut sichern f\"ur HIV- und Hepatitis-Serologie}

\end{itemize}
\begin{itemize}
\item {\mdseries
\gCa{ Unfallmeldung } (wichtig!), Blutabnahme }

\item {\mdseries
ggf. HIV-Postexpositionsprohylaxe (PEP)}

\end{itemize}
\begin{itemize}
\item {\mdseries
kurzfristige Therapie (4 Wo) }

\item {\mdseries
Einzelfall muss immer individuell beurteilt werden (Nutzen/Risiko),
starke Nebenwirkungen.}

\item {\mdseries
\gEs{Muss innerhalb von 2 Std. begonnen werden}
${\rightarrow}$ daher SOFORT in ein Unfallspital der AUVA}

\end{itemize}
\begin{itemize}
\item {\mdseries
Meldung an Betriebsleitung}

\item {\mdseries
\gEs{Unfallmeldung} (wichtig!) an die AUVA}

\end{itemize}
{\mdseries\color{black}
Unfallmeldung nicht vergessen!}

{\mdseries\color{black}
Sofort in ein Unfallspital! }

{\mdseries\color[rgb]{0.0,0.0,0.5019608}
\index{MRSA}MRSA und andere Resistente Keime}

{\mdseries
Die Ma{\ss}nahmen beim Transport von Patienten mit resistenten Keimen
werden unter {\quotedblbase}MRSA-Transport - Allgemeine
Ma{\ss}nahmen{\quotedblbase} auf Seite
K.II{\guillemotright}\pageref{bkm:RefHeading40262819} besprochen.}

\section{[RS XI] Katastrophen, Gro{\ss}schadensereignisse,
Gefahrengutunfall}
{\mdseries\itshape
Roman Koch [*]}

{\bfseries\color[rgb]{0.5019608,0.0,0.0}
\gEs{Changelog:}}

{\bfseries\color[rgb]{0.5019608,0.0,0.0}
2009-02-20: Fertig f\"ur Kurs 2009/02.}

{\bfseries\color[rgb]{0.5019608,0.0,0.0}
2009-02-25: Start Changelog.}

\subsection{Katastrophen, Gro{\ss}schadensereignisse; Unfall}
\begin{minipage}{13.7cm}
  [Warning: Image ignored] % Unhandled or unsupported graphics:
%\includegraphics[width=13.68cm,height=6.929cm]{AASdev20090228-img109.jpg}
 

{\mdseries\upshape
ICE-Ungl\"uck von Eschede am 3. Juni 1998}

 [Warning: Image ignored] % Unhandled or unsupported graphics:
%\includegraphics[width=6.212cm,height=3.933cm]{AASdev20090228-img110.jpg}
{\mdseries\upshape
Verstorbene:101}

{\mdseries\upshape
Schwerverletzte:88  Leicht- und Unverletzte:106}

{\mdseries\upshape
Foto Wikipedia, Nils Fretwurst; Lizenz: gemein frei
({\quotedblbase}public domain{\textquotedblleft})}

\end{minipage}

{\mdseries\color{black}
Definition \gT{Unfall} (Bagatelle, Notfall)}

\begin{center}
\begin{tabular}{|m{3.31cm}|m{9.99cm}|}
\hline
{\mdseries Ereignis:}

 &
{\mdseries pl\"otzlich, unerwartet}

\\\hline
{\mdseries Dauer:}

 &
{\mdseries Minuten bis 1 Std.}

\\\hline
{\mdseries Infrastruktur:}

 &
{\mdseries intakt}

\\\hline
{\mdseries Verletzte:}

 &
{\mdseries bis 14}

\\\hline
{\mdseries Erstversorgung:}

 &
{\mdseries innerhalb von Minuten durch Rettungsdienst nach Regeln der
Individual\-medizin}

\\\hline
{\mdseries Abtransport:}

 &
{\mdseries unbehindert Routinebetrieb mit \"ortlich vorhandenen Mitteln}

\\\hline
{\mdseries Definitive Versorgung:}

 &
{\mdseries in umliegende Krankenh\"auser nach Regeln der
Individualmedizin}

\\\hline
{\mdseries Material und Mannschaft:}

 &
{\mdseries Normalbetrieb Rettungs- und Krankentransport}

\\\hline
\end{tabular}
\end{center}
{\mdseries\color{black}
Definition \index{Gro{\ss}unfall}\gT{Gro{\ss}unfall}
(\index{Gro{\ss}schadensereignis}\gT{Gro{\ss}schadensereignis})}

{\mdseries
{\quotedblbase}Ein Gro{\ss}unfall liegt vor, wenn anzunehmen ist, dass
das Ereignis mit den \gNw{normalen} \"ortlichen personellen
und materiellen Kr\"aften und Mitteln nicht bew\"altigt werden kann,
aber keine erkl\"arte Katastrophensituation
vorliegt{\textquotedblleft}}

\begin{center}
\begin{tabular}{|m{3.31cm}|m{9.99cm}|}
\hline
{\mdseries Ereignis:}

 &
{\mdseries lokal begrenzt}

\\\hline
{\mdseries Dauer:}

 &
{\mdseries mehrere Stunden}

\\\hline
{\mdseries Infrastruktur:}

 &
{\mdseries intakt}

\\\hline
{\mdseries Verletzte:}

 &
{\mdseries ab 15}

\\\hline
{\mdseries Erstversorgung:}

 &
{\mdseries innerhalb von Minuten bis Stunden durch den Rettungsdienst
nach den Regeln der Katastrophenmedizin}

\\\hline
{\mdseries Abtransport:}

 &
{\mdseries relativ unbehindert nach der Erstversorgung, gen\"ugend
vorhandene Transportmittel}

\\\hline
{\mdseries Definitive Versorgung:}

 &
{\mdseries in umliegende Krankenh\"auser nach den Regeln der
Individualmedizin}

\\\hline
{\mdseries Material und Mannschaft:}

 &
{\mdseries alle verf\"ugbaren Helfer, Normalbetrieb Rettungs- und
Krankentransport inkl. Sondereinsatzgruppe und zus\"atzlichem Material}

\\\hline
\end{tabular}
\end{center}
{\mdseries\color{black}
Definition \index{Katastrophe}\gT{Katastrophe}
(au{\ss}ergew\"ohnliche Lage)}

{\mdseries
griech.:  {\quotedblbase}Umkehr der normalen
Verh\"altnisse{\textquotedblleft} }

{\mdseries
Ein \gEs{unvorhergesehenes Ereignis} mit einem
\gEs{au{\ss}ergew\"ohnlichen Schadensausma{\ss}}
das unmittelbar bevorsteht oder bereits eingetreten ist, eine
\gEs{konkrete Gefahr f\"ur Menschen},
\gZ{Tiere, Umwelt, Kulturg\"uter und Sachwerte} sowie
f\"ur die \gEs{Infrastruktur} zur Sicherstellung
der Versorgung mit lebensnotwendigen G\"utern und Dienstleistungen,
welches der \gEs{koordinierten F\"uhrung durch die
Beh\"orde} bedarf. }

\begin{center}
\begin{tabular}{|m{3.31cm}|m{9.99cm}|}
\hline
{\mdseries Ereignis:}

 &
{\mdseries (\"uber)-regionales Gebiet}

\\\hline
{\mdseries Dauer:}

 &
{\mdseries Tage, Wochen}

\\\hline
{\mdseries Infrastruktur:}

 &
{\mdseries gest\"ort, zerst\"ort}

\\\hline
{\mdseries Verletzte:}

 &
{\mdseries keine bis viele, je nach Art des Ereignisses}

\\\hline
{\mdseries Erstversorgung:}

 &
{\mdseries innerhalb von Stunden bis Tagen durch
KHD\footnotemark{}{}-Einheiten nach den Regeln der Katastrophenmedizin}

\\\hline
{\mdseries Abtransport:}

 &
{\mdseries stark behindert verz\"ogert, wenige/nicht vorhandene
Transportmittel}

\\\hline
{\mdseries Definitive Versorgung:}

 &
{\mdseries auf Verbandspl\"atzen in Behelfsspit\"alern, in Lazaretten
nach Regeln der Katastrophenmedizin}

\\\hline
{\mdseries Material und Mannschaft:}

 &
{\mdseries \"Uberlebende und verf\"ugbare Helfer, einsatzbereite
Fahrzeug, Kat-Einsatz aller Hilfsorganisationen, \"uberregionale und
internationale Hilfe}

\\\hline
\end{tabular}
\end{center}
\footnotetext{KHD: Katastrophenhilfsdienst}
{\mdseries\color{black}
Definition
\index{Individualmedizin}\gT{Individualmedizin}}

{\mdseries
Bedarfsgerechte und \gNw{bestm\"ogliche} Versorgung jedes
Patienten }

{\mdseries
Freie Arztwahl, welche auf das Vertrauensverh\"altnis zwischen Patient
und Arzt ausgerichtet ist}

{\mdseries\color{black}
Definition
\index{Katastrophenmedizin}\gT{Katastrophenmedizin}}

{\mdseries
Eine Vielzahl von Hilfsbed\"urftigen stehen wenigen \"Arzten,
Sanit\"atern und sonstigen Helfern mit oft unzureichenden Mitteln
gegen\"uber. Ziel muss es sein, unter den gegebenen Umst\"anden so
viele Leben wie m\"oglich zu retten. Mit bewussten
\gE{Verzicht auf zeit- und personalaufwendige
Maximalversorgung} des Einzelnen zugunsten der lebensrettenden
Minimalversorgung Vieler.}

{\mdseries
\gT{Katastrophenmedizin} bedeutet Massenversorgung von
Verwundeten und Kranken mit beschr\"ankten Mitteln und dem Zwang zur
Einteilung nach Dringlichkeit.}

{\mdseries\color{black}
\gNw{Bei Gro{\ss}unf\"allen und Katastrophen versucht man das
beste herauszuholen}}

\begin{center}
\tablehead{}\begin{supertabular}{|m{6.649cm}|m{6.651cm}|}
\hline
{\mdseries\color{black} \gEs{Ziele}}

\begin{itemize}
\item {\mdseries das Bestm\"ogliche}

\item {\mdseries f\"ur die gr\"o{\ss}te Anzahl}

\item {\mdseries zur rechten Zeit}

\item {\mdseries am rechten Ort}

\end{itemize}
 &
{\mdseries Die ist \gEs{nur erreichbar wenn}:}

\begin{itemize}
\item {\mdseries der Bestgeeignete auf jeder Stufe f\"uhrt}

\item {\mdseries alle Einsatzkr\"afte Ihre Aufgabe kennen }

\item {\mdseries das dazu ben\"otigte Material vorhanden, bereitgestellt
und rasch einsatzbereit ist}

\item {\mdseries klare und eindeutige F\"uhrungsstrukturen vorhanden
sind}

\item {\mdseries die F\"uhrungspyramide von allen strikt eingehalten
wird}

\item {\mdseries die Organisation eingespielt/einge\"ubt ist}

\item {\mdseries Eins\"atze geplant und vorbereitet sind}

\end{itemize}
\\\hline
{\mdseries Daf\"ur werden \gEs{Sonderausr\"ustungen}
zus\"atzlich erforderlich, da ein erh\"ohter Bedarf besteht an:}

\begin{itemize}
\item {\mdseries Material}

\item {\mdseries Medikamente}

\item {\mdseries Personal}

\end{itemize}
 &
{\mdseries Wir m\"ussen \gEs{rechnen mit}}

\begin{itemize}
\item {\mdseries einem gro{\ss}en Patientenanfall}

\item {\mdseries besonderen Verletzungsmustern}

\item {\mdseries einer l\"angeren Einsatzzeit vor Ort}

\end{itemize}
\\\hline
\end{supertabular}
\end{center}
{\mdseries\color{black}
\gNw{Es gibt verschiedene
}\gEs{Katastrophenarten}}

\begin{center}
\tablehead{}\begin{supertabular}{|m{6.649cm}|m{6.651cm}|}
\hline
{\mdseries\color{black} Naturkatastrophen}

 &
{\mdseries Erdbeben,\"Uberschwemmungen, Erdrutsche/Lawinen etc.}

\\\hline
{\mdseries\color{black} Technische Katastrophen}

 &
{\mdseries Flugzeugabsturz, Eisenbahnungl\"uck, Reaktor\-unfall,
Schiffs\-untergang etc.}

\\\hline
{\mdseries\color{black} Sekund\"arkatastrophen}

 &
{\mdseries Hungersnot, Seuchen etc.}

\\\hline
\end{supertabular}
\end{center}
{\mdseries\color{black}
\gEs{Katastrophenhilfegesetze} regeln die
Zust\"andigkeiten}

{\mdseries
Die Katastrophenhilfsgesetze \gNw{r}egeln die Aufgaben der
L\"ander, Organisationen und Zust\"andigkeiten in Ihrem Bereich.
\gEs{Ausschlie{\ss}lich die Beh\"orde erkl\"art ein
au{\ss}ergew\"ohnliches Schadensereignis zur Katastrophe} im Sinne des
Gesetzes. Beh\"ordliche Alarmpl\"ane dienen zur Festlegung der
Ma{\ss}nahmen der Katastrophenhilfe und zur Koordination der im
Anlassfall t\"atig werdenden Einrichtungen. Je nach Ausdehnung des
Schadensgebietes gibt es unterschiedliche Kompetenzen:}

\begin{center}
\begin{tabular}{|m{2.375cm}m{0.45600003cm}m{6.1590004cm}|}
\hline
{\mdseries \gE{Gemeinde }}

 &
\centering ${\rightarrow}$\par

 &
{\mdseries B\"urgermeister}

\\
{\mdseries \gE{Bezirk }}

 &
\centering ${\rightarrow}$\par

 &
{\mdseries Bezirkshauptmann}

\\
{\mdseries \gE{Bundesland }}

 &
\centering ${\rightarrow}$\par

 &
{\mdseries Landeshauptmann}

\\
{\mdseries \gE{Bundesgebiet}}

 &
\centering ${\rightarrow}$\par

 &
{\mdseries Innenminister, Bundeskanzler}

\\\hline
\end{tabular}
\end{center}
{\mdseries\color{black}
Die \gEs{Katastrophenwarnung} mittels Sirene
alarmiert die Bev\"olkerung}

{\mdseries\color[rgb]{0.0,0.0,0.5019608}
Katastrophenschutz in Wien}

\begin{center}
 [Warning: Image ignored] % Unhandled or unsupported graphics:
%\includegraphics[width=10.861cm,height=7.614cm]{AASdev20090228-img111.png}

\end{center}
\begin{center}
\tablehead{}\begin{supertabular}{|m{4.3190002cm}|m{8.981cm}|}
\hline
{\mdseries Wien ist f\"ur den Gro{\ss}\-schadens- und Katastrophenfall
gut ger\"ustet. Tausende Helfer sind dazu im Katastrophen\-schutz-Kreis
({\quotedblbase}\gT{K-Kreis}{\textquotedblleft})
org\-anisiert.}

{\mdseries Der K-Kreis ist ein Zu\-sammen\-schluss der beruflichen und
freiwilligen Hilfs- und Ein\-satz\-organisationen sowie wichtiger
Beh\"orden. Dreh\-scheibe und gemeinsame Anlaufstelle ist dabei die
Stadt Wien, Magistrats\-direktion - Katastrophen\-manage\-ment und
Sofort\-ma{\ss}\-nahmen und die Helfer Wiens.}

 &
 [Warning: Image ignored] % Unhandled or unsupported graphics:
%\includegraphics[width=9.046cm,height=13.623cm]{AASdev20090228-img112.png}
\captionof{figure}[Der K-Kreis]{Der K-Kreis}


\\\hline
\end{supertabular}
\end{center}
{\mdseries\color[rgb]{0.0,0.0,0.5019608}
Katastrophenhilfe}

{\mdseries
Hilfsma{\ss}nahmen die von uns als Sanit\"atsorganisation erwartet
werden:}

\begin{itemize}
\item {\mdseries
Bergung, Versorgung und Transport }

\item {\mdseries
Pflegerische Betreuung }

\item {\mdseries
Sorge f\"ur Unterkunft und Verpflegung }

\item {\mdseries
Karitative Betreuung von Hilfsbed\"urftigen}

\end{itemize}
\begin{center}
\tablehead{}\begin{supertabular}{|m{6.649cm}|m{6.651cm}|}
\hline
{\mdseries\color{black} Sanit\"atseinsatz}

 &
\begin{itemize}
\item {\mdseries Retten/Bergen}

\item {\mdseries Triage}

\item {\mdseries Lebens rettende Sofortma{\ss}nahmen}

\item {\mdseries Erstversorgung}

\item {\mdseries Sanit\"atshilfe}

\item Abtransport

\end{itemize}
\\\hline
{\mdseries\color{black} Pflegeeinsatz}

 &
{\mdseries Wenn die pflegerische Versorgung durch die bestehenden
Einrichtungen nicht mehr sichergestellt ist.}

\begin{itemize}
\item {\mdseries Erweiterung der vorhandenen Behandlungsm\"oglichkeiten}

\item Errichtung von Hilfskrankenh\"ausern

\end{itemize}
\\\hline
{\mdseries\color{black} Sozialeinsatz}

 &
{\mdseries Wenn Menschen infolge einer Katastrophe ihre
Selbstversorgungsm\"oglichkeit verloren haben.}

\begin{itemize}
\item {\mdseries Einrichtung/Betrieb von Notunterk\"unften/Lagern  }

\item {\mdseries Beschaffung, Zubereitung und Ausgabe von Lebensmitteln,
Kleidern, Gebrauchsgegenst\"anden  }

\item F\"ursorgliche Hilfe, Registrierung und Betreuung

\end{itemize}
\\\hline
\end{supertabular}
\end{center}
{\mdseries\color[rgb]{0.0,0.0,0.5019608}
Grunds\"atze des Einsatzes}

{\mdseries
Grundkomponenten: F\"uhrung -- Zeit -- Raum }

{\mdseries
Ziel ist es ohne \"uberst\"urztes Handeln die Ausdehnung der Krisenlage
auf benachbarte und andere R\"aume zu verhindern. Dabei wird
gem\"a{\ss} der gesetzlichen Bestimmungen von Kr\"aften der
Einsatzorganisationen, Beh\"orden und anderen Einrichtungen koordiniert
vorgegangen.}

\paragraph{F\"uhrung}
{\mdseries
Steuerndes Einwirken auf andere Menschen um ein bestimmtes Ziel zu
erreichen. }

{\mdseries
Die F\"uhrung ist ein unentbehrliches Steuerungsinstrument zur
Erf\"ullung des Einsatzauftrages
({\quotedblbase}\gE{Ziel}{\textquotedblleft}) und dem
richtigen Einsatz von \gE{Personal},
\gE{Material} und \gE{Transportmitteln} }

{\mdseries
Im Schadensraum gilt: {\quotedblbase}\gEs{Ein
einziger Verantwortlicher}{\textquotedblleft} mit Entscheidungsbefugnis
und den n\"otigen F\"uhrungsmittel
({\quotedblbase}\gE{Einheit der
F\"uhrung}{\textquotedblleft}). Dieser wirkt mittels m\"oglichst
einfacher Ma{\ss}nahmen
({\quotedblbase}\gE{Einfachheit}{\textquotedblleft}) auf
den Einsatzverlauf ein
({\quotedblbase}\gE{Handlungsfreiheit}{\textquotedblleft}).
Mittel und Mannschaften werden dabei gem\"a{\ss} ihrer Eigenart und
Leistungsf\"ahigkeit eingesetzt
({\quotedblbase}\gE{\"Okonomie der
Kr\"afte}{\textquotedblleft}).Er muss rasch auf \"Anderungen der
Situation reagieren
({\quotedblbase}\gE{Beweglichkeit}{\textquotedblleft})
und gleichzeitig vorausschauend handeln
({\quotedblbase}\gE{Reservenbildung}{\textquotedblleft}).}

{\mdseries\color{black}
{\quotedblbase}F\"uhren{\textquotedblleft} und
{\quotedblbase}Retten{\textquotedblleft} ist nicht gleichzeitig
m\"oglich!}

{\mdseries\color{black}
Sanit\"ater/NA und Not\"arzte m\"ussen
\gEs{F\"uhrungsaufgaben} \"ubernehmen}

{\mdseries
%
%
Der erste am Schadensort eintreffende Sanit\"ater muss:}

\begin{enumerate}
\item Lagemeldung

\item den Standort der \"ortlichen Einsatzleitung festlegen

\item mit bereits anwesenden Einsatzkr\"aften die ersten Ma{\ss}nahmen
zur Rettung festlegen

\item wohin Patienten im Ein

\end{enumerate}
\paragraph[Zeit]{Zeit}
{\mdseries
Jede Katastrophenereignis zeigt unabh\"angig von seiner Ursache
folgenden Zeitablauf}

\begin{center}
\begin{tabular}{|m{4.367cm}|m{4.367cm}|m{4.367cm}|}
\hline
\centering {\bfseries\upshape\color[rgb]{0.0,0.0,0.5019608} 1.
Phase}\par

 &
\centering {\bfseries\upshape\color[rgb]{0.0,0.0,0.5019608} 2.
Phase}\par

 &
\centering {\bfseries\upshape\color[rgb]{0.0,0.0,0.5019608} 3. Phase
}\par

\\\hline
{\mdseries Isolation (Minuten)}

 &
{\mdseries Retten (Stunden)}

 &
{\mdseries Wiederherstellung (Tage)}

\\\hline
{\mdseries Schaden begrenzen \newline
{}- Lage erkennen\newline
{}- Melden}

 &
{\mdseries Ordnung schaffen\newline
{}- Kommandoposten}

 &
{\mdseries Lokale Einsatzleitung\newline
{}- Verbindung\newline
{}- Lage beurteilen}

\\\hline
{\mdseries \"Uberleben}

 &
{\mdseries Weiterleben}

 &
{\mdseries Endbehandlung}

\\\hline
{\mdseries Laien:\newline
{}- Nothilfe/Selbsthilfe\newline
{}- Retten/Erste Hilfe}

 &
{\mdseries Geschulte Helfer:\newline
{}- LRSM\newline
{}- Transportf\"ahigkeit}

 &
{\mdseries Spezialisten:\newline
{}- Spital nach Dringlichkeit behandeln}

\\\hline
\multicolumn{3}{|m{13.501cm}|}{[Warning: Draw object ignored]

}\\\hline
\end{tabular}
\end{center}
\paragraph{Hilfe braucht Raum: Sanit\"atsdienstliche Organisation im
Schadensraum und in der Sanit\"atshilfsstelle (SanHist)}
{\mdseries
Der Einsatz allgemein wird unterteilt in: }

\begin{itemize}
\item Schadensraum

\item Transportraum

\item Hospitalisationsraum

\end{itemize}
{\mdseries
Der \gEs{Hospitalisationsraum} betrifft die
Transportziele und Aufnahmespit\"aler. }

{\mdseries
Der \gEs{Transportraum} umfasst in diesem
Zusammenhang alle Mittel und Wege um die Patienten vom Schadens\-raum
zun Hospitalisationsraum zu verbringen. Dieser
{\quotedblbase}allgemeine Transportraum{\textquotedblleft} ist nicht zu
ver\-wechseln mit dem Transportraum der Sanit\"atshilfestelle (s.u.). 
}

{\mdseries
Der \gEs{Schadensraum} ist der Ort des eigentlichen
(schadens-)Geschehens. Es gibt in ihm eine einheitliche
Organisations\-form f\"ur den Sanit\"atsdienst: Die
\gT{Sanit\"atshilfsstelle}
(\gT{SanHiSt}), welche aus den 3 R\"aumen
\gT{Triageraum}, \gT{Behandlungsraum} und
\gT{Transportraum} besteht. Definiert wird Sie in der
ASB\"O Rahmen\-vor\-schrift f\"ur Gro{\ss}unf\"alle und den
Durchf\"uhrungsbestimmungen f\"ur Gro{\ss}unf\"alle. }

{\centering \begin{minipage}{12.605cm}
 [Warning: Image ignored] % Unhandled or unsupported graphics:
%\includegraphics[width=12.605cm,height=14.03cm]{AASdev20090228-img113}
\captionof{figure}[Gliederung des Einsatzraumes]{Gliederung des
Einsatzraumes}
\begin{multicols}{2}
{\mdseries\upshape
\gT{\"Aussere Absperrung}: Gro{\ss}r\"aumige
Verkehrsumleitung}

{\mdseries\upshape
\gT{Innere Absperrung}: h\"alt Schaulustige fern, h\"alt
Personen in Panik zur\"uck}

{\mdseries\upshape
\gT{Sicherheitsring}: Absperrung des Schadenplatzes,
Sicherheit f\"ur Einsatzkr\"afte}

{\mdseries\upshape
\gT{Pforte}: Einbahnsystem f\"ur Einsatz-Kfz, keine
Fremdfahrzeuge, m\"oglichst nur eine Pforte }

{\mdseries\upshape
\gT{Leiter Information} (\gT{i}) f\"ur
Weiterleitung von Informationen an die Au{\ss}enwelt (Presse,
Angeh\"orige, ...)}

{\mdseries\upshape
\gT{Bergetriage} (\gT{BT}) }

{\mdseries\upshape
\gT{Sanit\"atshilfsstelle} (\gT{SanHist}):
mit Triageraum, Behandlungsraum und Transportraum;}

{\mdseries\upshape
\gT{Einsatzleiter} (\gT{EL}) }

{\mdseries\upshape
\gT{Leitender Notarzt} (\gT{LNA})}

{\mdseries\upshape
\gT{Mobile Leitstelle} (\gT{MLS}) }

{\mdseries\upshape
\gT{Material- und Meldestellet} (\gT{M})}

{\mdseries\upshape
\gT{Sammelstelle Tote} (\gT{t}) und
\gT{Unverletzte} (\gT{U}) }

\end{multicols}
\end{minipage}\par}

\subparagraph{Triageraum}
{\mdseries
Triage (frz.) = Sichtung, Auswahl}

{\mdseries
Triage ist das Begutachten aller Patienten mit geringen Mitteln und
Aufwand sowie das Einteilen in Triagegruppen mit dem Ziel m\"oglichst
viele Leben zu retten!}

{\mdseries
{\quotedblbase}Life before Limbs{\textquotedblleft} = Leben vor
Gliedma{\ss}en, Funktion vor Sch\"onheit!}

{\mdseries\color{black}
Bergetriage}

\begin{center}
\tablehead{}\begin{supertabular}{m{10.891cm}m{2.409cm}}
{\mdseries Nur wenn keine Gefahrenzone und gen\"ugend Personal/Zeit
vorhanden ist!}

\begin{itemize}
\item {\mdseries BT Team muss mind. aus einem NFS bestehen}

\item {\mdseries Gefahren beachten}

\item {\mdseries Erst Sichtung der Betroffenen und Festlegung der
Dringlichkeit}

\item {\mdseries Nicht behandeln!}

\end{itemize}
 &
{\raggedleft   [Warning: Image ignored]
% Unhandled or unsupported graphics:
%\includegraphics[width=2.396cm,height=2.149cm]{AASdev20090228-img114.png}
 \par}

\\
\end{supertabular}
\end{center}
{\mdseries\color{black}
Triage (Sichtung)}

\begin{center}
\tablehead{}\begin{supertabular}{m{10.891cm}m{2.409cm}}
\begin{itemize}
\item {\mdseries Bei Missverh\"altnis zwischen der Zahl der Betroffenen
und den Be\-handlungs\-m\"oglich\-keiten}

\item {\mdseries Schnell, mit geringem Aufwand}

\item {\mdseries Hoher Zeitdruck}

\item {\mdseries Keine detaillierte Diagnostik m\"oglich}

\item {\mdseries Triage ist dynamischer Vorgang}

\item {\mdseries St\"andige Nachtriage ist erforderlich}

\item {\mdseries Maximaler Zeitaufwand pro Patient gesamt:}

\end{itemize}
\begin{itemize}
\item {\mdseries Gehender 1 Minute}

\item {\mdseries Liegender 3 Minuten }

\end{itemize}
\begin{itemize}
\item {\mdseries Zentrale Registrierung aller Patienten mittels
Patientenleitsystem (PLS)}

\item {\mdseries Zuordnung der Patienten zu Triagegruppen}

\end{itemize}
\begin{itemize}
\item {\mdseries T1: Sofortbehandlung, lebensrettende Noteingriffe}

\item {\mdseries T2: Dringende Behandlung, Schwerverletzte}

\item {\mdseries T3: Minimalbehandlung}

\item {\mdseries T4: Abwartende Behandlung}

\end{itemize}
 &
{\raggedleft   [Warning: Image ignored]
% Unhandled or unsupported graphics:
%\includegraphics[width=2.375cm,height=2.184cm]{AASdev20090228-img115}
 \par}

\\
\end{supertabular}
\end{center}
\begin{center}
\begin{tabular}{|m{6.649cm}m{6.651cm}|}
\hline
\multicolumn{2}{|m{13.5cm}|}{\centering
{\bfseries\upshape\color[rgb]{0.0,0.0,0.5019608} Zwei Triagepl\"atze,
damit der NA abwechselnd triagieren kann}\par

}\\\hline
\multicolumn{1}{|m{6.649cm}|}{{\mdseries Vorbereitung durch RS/NFS}

{\mdseries Entkleiden}

{\mdseries Vitalparameter erheben (RR, HF, AF, SpO2)}

{\mdseries Untersuchung durch den Notarzt}

} &
{\centering   [Warning: Image ignored]
% Unhandled or unsupported graphics:
%\includegraphics[width=5.132cm,height=3.783cm]{AASdev20090228-img116}
 \par}

\\\hline
\end{tabular}
\end{center}
\subparagraph[Behandlungsraum]{Behandlungsraum}
\begin{flushleft}
\begin{tabular}{|m{3.328cm}|m{5.406cm}m{4.36cm}}
\hline
  [Warning: Image ignored] % Unhandled or unsupported graphics:
%\includegraphics[width=1.489cm,height=1.61cm]{AASdev20090228-img117.gif}
 

 &
\multicolumn{2}{m{9.966001cm}|}{{\mdseries Es werden drei
Behandlungsstellen eingerichtet}

}\\\hline
  [Warning: Image ignored] % Unhandled or unsupported graphics:
%\includegraphics[width=1.629cm,height=1.489cm]{AASdev20090228-img118}
 

 &
\multicolumn{1}{m{5.406cm}|}{{\mdseries Sofortbehandlung vor Ort}

{\mdseries T1 Lebensrettende Soforteingriffe}

} &
\multicolumn{1}{m{4.36cm}|}{{\centering   [Warning: Image ignored]
% Unhandled or unsupported graphics:
%\includegraphics[width=3.365cm,height=4.366cm]{AASdev20090228-img119.png}
 \par}

}\\\hline
  [Warning: Image ignored] % Unhandled or unsupported graphics:
%\includegraphics[width=1.669cm,height=1.578cm]{AASdev20090228-img120.png}
 

 &
{\mdseries Dringende Behandlung}

{\mdseries T2 Schwerverletzte}

 &
\\\hhline{--~}
  [Warning: Image ignored] % Unhandled or unsupported graphics:
%\includegraphics[width=1.601cm,height=1.587cm]{AASdev20090228-img121.png}
   [Warning: Image ignored] % Unhandled or unsupported graphics:
%\includegraphics[width=1.587cm,height=1.641cm]{AASdev20090228-img122.png}
 

 &
{\mdseries Warten}

{\mdseries T3 Warten Leichtverletzte}

{\mdseries T4 Warten Schwerverletzte}

 &
\\\hhline{--~}
\end{tabular}
\end{flushleft}
{\mdseries\color{black}
\gEs{Sofortbehandlung} vor Ort}

{\mdseries
Alle Patienten, welche ohne sofortige Behandlung sterben bzw. schwere
Sch\"aden erleiden}

\begin{flushleft}
\tablehead{}\begin{supertabular}{m{5.642cm}m{4.379cm}m{3.076cm}}
{\mdseries \gEs{Welche Patienten?}}

\begin{itemize}
\item {\mdseries Atemst\"orungen}

\item {\mdseries Spannungspneumothorax}

\item {\mdseries schwere \"au{\ss}ere Blutungen}

\item {\mdseries schwerer traumatischer Schock}

\end{itemize}
 &
{\mdseries \gEs{Ma{\ss}nahmen}}

\begin{itemize}
\item {\mdseries \"Uberpr\"ufung der Triage}

\item {\mdseries Volumenersatz}

\item {\mdseries Intubation}

\item {\mdseries Thoraxdrainage}

\item {\mdseries Noteingriffe}

\end{itemize}
 &
{\centering   [Warning: Image ignored]
% Unhandled or unsupported graphics:
%\includegraphics[width=2.876cm,height=2.271cm]{AASdev20090228-img123}
 \par}

\\
\end{supertabular}
\end{flushleft}
{\mdseries\color{black}
\gEs{Dringende Behandlung}}

{\mdseries
Behandlung und Herstellung der Transportf\"ahigkeit von schwer
verletzten/erkrankten Personen, eine weitere Behandlung im Krankenhaus
ist i.d.R. notwendig.}

\begin{flushleft}
\tablehead{}\begin{supertabular}{m{5.642cm}m{4.379cm}m{3.076cm}}
{\mdseries \gEs{Welche Patienten?}}

keine \gE{unmittelbare} Lebensbedrohung

\begin{itemize}
\item {\mdseries SHT}

\item {\mdseries Wirbelverletzungen mit L\"ahmungen}

\item {\mdseries Bauchtrauma, innere Blutungen}

\item {\mdseries schwere Thoraxtraumen}

\item {\mdseries gro{\ss}e Gef\"a{\ss}verletzungen}

\item {\mdseries offene Knochenbr\"uche und  Gelenks\-verletzungen}

\item {\mdseries Augenverletzungen}

\item {\mdseries gro{\ss}e Weichteilverletzungen}

\item {\mdseries Verbrennungen 15 - 30\%}

\item {\mdseries geschlossene Knochenbr\"uche und
Gelenks\-ver\-letzungen}

\end{itemize}
 &
{\mdseries \gEs{Ma{\ss}nahmen}}

\begin{itemize}
\item {\mdseries Verb\"ande}

\item {\mdseries Schockbehandlung}

\item {\mdseries Fixierung/Schienung}

\item {\mdseries Schmerzbek\"ampfung}

\item {\mdseries pflegerische Ma{\ss}nahmen}

\item {\mdseries LRSM wenn n\"otig}

\end{itemize}
 &
{\centering   [Warning: Image ignored]
% Unhandled or unsupported graphics:
%\includegraphics[width=2.775cm,height=2.44cm]{AASdev20090228-img124.png}
 \par}

\\
\end{supertabular}
\end{flushleft}
{\mdseries\color{black}
Warten \gEs{Leichtverletzte}}

{\mdseries
Alle Patienten, welche nur einer Minimalbehandlung bed\"urfen. }

\begin{flushleft}
\tablehead{}\begin{supertabular}{m{5.642cm}m{4.379cm}m{3.076cm}}
{\mdseries \gEs{Welche Patienten?}}

\begin{itemize}
\item {\mdseries Verbrennungen bis 15\%}

\item {\mdseries kleine Weichteilwunden}

\item {\mdseries einfache Knochenbr\"uche}

\item {\mdseries Prellungen, Zerrungen}

\item {\mdseries leichtes SHT}

\end{itemize}
 &
{\mdseries \gEs{Ma{\ss}nahmen}}

{\mdseries \"Uberpr\"ufung der Triage}

{\mdseries pflegerische Ma{\ss}nahmen}

{\mdseries Betreuung (KIT Team)}

{\mdseries Entlassung ambulant Behandelter zur SaStU}

 &
{\centering   [Warning: Image ignored]
% Unhandled or unsupported graphics:
%\includegraphics[width=2.677cm,height=2.299cm]{AASdev20090228-img125.png}
 \par}

\\
\end{supertabular}
\end{flushleft}
{\mdseries\color{black}
Warten \gEs{Schwerverletzte}}

{\mdseries
Alle Patienten, welche mit dem zur Verf\"ugung stehenden Mitteln
(Personal, Material) \gE{noch} nicht versorgt werden
k\"onnen . Abwartende Behandlung. }

\begin{center}
\tablehead{}\begin{supertabular}{m{5.702cm}m{4.0550003cm}m{3.3439999cm}}
{\mdseries \gEs{Welche Patienten?}}

\begin{itemize}
\item {\mdseries schweres Polytrauma}

\item {\mdseries schwerstes SHT}

\item {\mdseries offene K\"orperh\"ohlentraumen}

\item {\mdseries Mehrh\"ohlenverletzungen}

\item {\mdseries Verbrennungen \"uber 50\%}

\end{itemize}
 &
{\mdseries \gEs{Ma{\ss}nahmen}}

\begin{itemize}
\item {\mdseries \"Uberpr\"ufung der Triage}

\item {\mdseries Behandlung  (Schmerzbek\"ampfung)}

\item {\mdseries Betreuung (KIT Team)}

\item {\mdseries \"Uberwachung}

\end{itemize}
 &
{\centering   [Warning: Image ignored]
% Unhandled or unsupported graphics:
%\includegraphics[width=2.659cm,height=2.421cm]{AASdev20090228-img126.png}
 \par}

\\
\end{supertabular}
\end{center}
{\mdseries\color{black}
Transportpriorit\"at}

{\mdseries
Nach einer erfolgten Behandlung kann eine Transporpriorit\"at zugewiesen
werden. Man unterscheidet hierbei: }

\begin{itemize}
\item {\mdseries
\gEs{A}: \gEs{Hohe
Transportpriorit\"at}: Der rasche Transport zur fr\"uhzeitigen
Fachbehandlung nach \"arztlicher Notversorgung ist angesagt. }

\item {\mdseries
\gEs{B}: \gEs{Niedrige
Transportpriorit\"at}: Der Transport zur Fachbehandlung kann
verz\"ogert erfolgen. }

\end{itemize}
\subparagraph{Transportraum}
\begin{center}
\tablehead{}\begin{supertabular}{|m{2.365cm}m{5.8640003cm}m{4.8710003cm}|}
\hline
{\mdseries\color{black} Verladestelle}

 &
{\mdseries Sollte in der N\"ahe der Behandlungsstelle II errichtet
werden.}

{\mdseries Es muss protokolliert werden, welcher Patient mit welchem
Transportmittel in welches Krankenhaus eingeliefert wird.}

{\mdseries Einbahnregelung!}

 &
{\centering   [Warning: Image ignored]
% Unhandled or unsupported graphics:
%\includegraphics[width=3.373cm,height=2.484cm]{AASdev20090228-img127.png}
 \par}

\\\hline
{\mdseries\color{black} Wagenhalte\-platz}

 &
\begin{itemize}
\item {\mdseries Gro{\ss} genug}

\item {\mdseries Schr\"agparkordnung }

\item {\mdseries Sortiert nach NAW, RTW usw.}

\item {\mdseries Schl\"ussel stecken lassen}

\end{itemize}
{\centering   [Warning: Image ignored]
% Unhandled or unsupported graphics:
%\includegraphics[width=3.55cm,height=2.498cm]{AASdev20090228-img128}
 \par}

\centering Falscher Wagenhalteplatz\par

 &
{\centering   [Warning: Image ignored]
% Unhandled or unsupported graphics:
%\includegraphics[width=3.375cm,height=2.438cm]{AASdev20090228-img129}
 \par}

\\\hline
{\mdseries\color{black} Hubschrauber\-landestelle}

 &
\begin{itemize}
\item {\mdseries Die allgemeinen Vorschriften f\"ur
Hubschrauberlandepl\"atze sind zu beachten}

\item {\mdseries Freie ebene Fl\"ache }

\item {\mdseries Auf Stromleitungen achten}

\item {\mdseries Weit genug weg vom Behandlungsraum }

\end{itemize}
 &
{\centering [Warning: Draw object ignored]\par}

\\\hline
{\mdseries\color{black} Transport\-management}

 &
\begin{itemize}
\item {\mdseries Spiralf\"ormige Verteilung}

\item {\mdseries Zeitliche Staffelung}

\item {\mdseries Fachliche Zuordnung}

\item {\mdseries Katastrophe nicht ins Spital verlagern!}

\end{itemize}
 &
{\centering   [Warning: Image ignored]
% Unhandled or unsupported graphics:
%\includegraphics[width=4.954cm,height=3.88cm]{AASdev20090228-img130}
 \par}

\\\hline
\end{supertabular}
\end{center}
\subparagraph{Sammelstellen}
\begin{center}
\tablehead{}\begin{supertabular}{|m{4.367cm}m{4.367cm}m{4.367cm}|}
\hline
{\mdseries\color{black} Sammelstelle Unverletzte}

{\centering   [Warning: Image ignored]
% Unhandled or unsupported graphics:
%\includegraphics[width=2.151cm,height=2.931cm]{AASdev20090228-img131}
 \par}

 &
{\mdseries Ziel: Betreuung Unverletzter und Betroffener}

\begin{itemize}
\item {\mdseries Betreuen}

\item {\mdseries Beruhigen}

\item {\mdseries Informieren}

\item {\mdseries KIT Teams (ABW) integrieren}

\end{itemize}
 &
{\centering   [Warning: Image ignored]
% Unhandled or unsupported graphics:
%\includegraphics[width=2.01cm,height=2.154cm]{AASdev20090228-img132.png}
 \par}

\\\hline
{\mdseries\color{black} Sammelstelle Tote}

 &
{\mdseries Sollte diskret eingerichtet werden.}

{\mdseries Abgeschirmt von der Sammel\-stelle Unverletzte, dem
Be\-handlungs\-raum und der Au{\ss}en\-welt}

 &
{\centering   [Warning: Image ignored]
% Unhandled or unsupported graphics:
%\includegraphics[width=1.906cm,height=1.9cm]{AASdev20090228-img133.png}
 \par}

\\\hline
\end{supertabular}
\end{center}
\subparagraph{Material- / Meldestelle}
\begin{itemize}
\item {\mdseries
Dient zur Registrierung des Personals und Materials}

\begin{center}
 [Warning: Image ignored] % Unhandled or unsupported graphics:
%\includegraphics[width=2.176cm,height=2.126cm]{AASdev20090228-img134}

\end{center}
\item {\mdseries
Eintreffende Einheiten haben sich bei dieser zu melden}

\item {\mdseries
Steht im engen Kontakt zum Leiter SanHist und EL  }

\end{itemize}
\subparagraph{Mobile Leitstelle (MLS)}
\begin{itemize}
\item {\mdseries
Die MLS beherbergt i.d.R. die Einsatzleitung sowie wichtige
Administrations-, Telekommunkations- und Funkeinrichtungen. }

\begin{center}
 [Warning: Image ignored] % Unhandled or unsupported graphics:
%\includegraphics[width=2.902cm,height=2.915cm]{AASdev20090228-img135}

\end{center}
\item {\mdseries
Bei Bedarf werden u.U. hier Handfunkger\"ate an die Bereichsleiter
ausgegeben. }

\item {\mdseries
MLS ist die Br\"ucke zwischen verschiedenen Funkkan\"alen und
zust\"andig. }

\item {\mdseries
Wenn die MLS die Einsatzleitung beinhaltet wird sie mit einem rotem
Dreh- oder Blinklicht markiert. }

\end{itemize}
\paragraph{F\"uhrungsstruktur in der SanHist}
\begin{center}
\tablehead{}\begin{supertabular}{m{2.4459999cm}m{4.0030003cm}m{2.299cm}m{4.1520004cm}}
{\centering   [Warning: Image ignored]
% Unhandled or unsupported graphics:
%\includegraphics[width=1.654cm,height=1.624cm]{AASdev20090228-img136.png}
 \par}

 &
{\mdseries Einsatzleiter Rettungsdienst}

 &
{\centering   [Warning: Image ignored]
% Unhandled or unsupported graphics:
%\includegraphics[width=1.62cm,height=1.62cm]{AASdev20090228-img137}
 \par}

 &
{\mdseries Leiter Material-/Meldestelle}

\\
{\centering   [Warning: Image ignored]
% Unhandled or unsupported graphics:
%\includegraphics[width=1.62cm,height=1.62cm]{AASdev20090228-img138.png}
 \par}

 &
{\mdseries Leiter SanHist}

 &
{\centering   [Warning: Image ignored]
% Unhandled or unsupported graphics:
%\includegraphics[width=1.489cm,height=1.61cm]{AASdev20090228-img139.gif}
 \par}

 &
{\mdseries Leiter Behandlung}

\\
{\centering   [Warning: Image ignored]
% Unhandled or unsupported graphics:
%\includegraphics[width=1.62cm,height=1.62cm]{AASdev20090228-img140}
 \par}

 &
{\mdseries LNA Leitender Notarzt}

 &
{\centering   [Warning: Image ignored]
% Unhandled or unsupported graphics:
%\includegraphics[width=1.62cm,height=1.62cm]{AASdev20090228-img141}
 \par}

 &
{\mdseries Leiter Transport}

\\
{\centering   [Warning: Image ignored]
% Unhandled or unsupported graphics:
%\includegraphics[width=1.62cm,height=1.62cm]{AASdev20090228-img142}
 \par}

 &
{\mdseries Leiter Triageraum/Triage}

 &
{\centering   [Warning: Image ignored]
% Unhandled or unsupported graphics:
%\includegraphics[width=1.62cm,height=1.62cm]{AASdev20090228-img143.png}
 \par}

 &
{\mdseries Leiter Information}

\\
\end{supertabular}
\end{center}
{\mdseries\color[rgb]{0.0,0.0,0.5019608}
Patientenleitsystem (PLS)}

{\mdseries
Das \index{Patientenleitsystem}\gT{Patientenleitsystem}
(\index{PLS}\gT{PLS}) erm\"oglicht eine eindeutige
Kennzeichnung von Peteinten und Zuordnung von Behandlungsma{\ss}nahmen
und Befunden. Der wichtigste Bestandteil ist die
\index{Patientenleittasche}\gT{Patientenleittasche}
(\index{PLT}\gT{PLT}). }

{\mdseries
Sie besteht aus einer orangenen H\"ulle welche fortlaufend nummeriert
sind. Auf dieser Kunststoffh\"ulle k\"onnen wichtige Notizen wie
Triagegruppe, etc., vorgenommen werden. In der H\"ulle befinden sich
das Behandlungsprotokoll (blau), Idenentifikationsprotokoll (rosa),
Zusatzkarten, Kontaminationsaufkleber und ein Bogen mit gleichlautend
nummerierten Klebeetiketten. Am unteren Ende befinden sich zwei
Abschnitte welche die R\"uckverfolgbarkeit des Transports
sicherstellen. Die PLT ist mit einer Gummischnur versehen damit sie dem
Patienten umgeh\"angt werden kann. }

{\mdseries
Sp\"atestens an der Triagestelle bekommt jeder Patient oder Betroffene
{\quotedblbase}seine{\textquotedblleft} PLT. Auch Unverletzte und
Verstrorbene werden auch mittels PLS registriert.}

{\mdseries\color{black}
Zusatzkarten}

{\mdseries
Bei einer eventuellen Bergetriage bekommt jeder Patient eine
Patientenleittasche (PLT). Petienten, welche
\gE{bevorrangt} gerettet werden sollen erhalten
zus\"atzlich auf die PLT die gelbe \textbf{DRINGEND} Karte.
\gE{Tote} erhalten eine PLT mit der schwarzen
\textbf{VERSTORBEN}\textbf{ }Karte.}

{\mdseries\color{black}
Behandlungsprotokoll / Identifikationsprotokoll}

{\mdseries
Die Protokolle werden im Behandlungsraum ausgef\"ullt wenn Zeit daf\"ur
ist. Das Ausf\"ullen der Protokolle darf den Abtransport unter keinen
Umst\"anden verz\"ogern. Die Protokolle k\"onnen auch leer bleiben.}

{\mdseries\color{black}
\index{Kontaminationsaufkleber}Kontaminationsaufkleber}

{\mdseries
Wenn der Patient kontaminiert ist, wird ein Dreieck auf die PLT geklebt.
Dies soll alle nachfolgenden Personen, welche mit dem Patienten in
Kontakt treten, warnen.}

{\mdseries\color{black}
PLT-Nummer}

{\mdseries
Es erfolgt eine eindeutige Identifizierung der Patienten durch eine
einmalige Nummernkennzeichnung der Patientenleittasche. Mittels
Selbstklebenummern werden die Patienten in allen verschiedenen
Bereichen in Listen erfasst (Triage, Behandlung, Transport). Ausserdem
k\"onnen alle Effekte des Patienten damit gekennzeichnet werden.}

{\mdseries
Die Patientenleittasche \gEs{verbleibt solange am
Patienten, bis das Aufnahmeverfahren in Krankenhaus abgeschlossen ist},
d.h. bis er im Spital administrativ erfasst ist. Anschlie{\ss}end wird
die PLT der Krankengeschichte des Patienten beigef\"ugt.}

{\centering   [Warning: Image ignored]
% Unhandled or unsupported graphics:
%\includegraphics[width=13.361cm,height=7.796cm]{AASdev20090228-img144}
 \par}

\begin{center}
\begin{tabular}{m{6.649cm}m{6.651cm}}
\multicolumn{2}{m{13.5cm}}{\centering
{\bfseries\upshape\color[rgb]{0.0,0.0,0.5019608} PLT Vorderseite}\par

}\\
[Warning: Draw object ignored][Warning: Draw object ignored][Warning:
Draw object ignored]

\begin{center}
 [Warning: Image ignored] % Unhandled or unsupported graphics:
%\includegraphics[width=5.771cm,height=10.83cm]{AASdev20090228-img145}

\end{center}
 &
{\mdseries\color{black} Triagegruppen:}

{\mdseries Bei Triage nur Einteilung in I, II, III oder IV}

{\mdseries\color{black} Transportpriorit\"at:}

{\mdseries \gEs{A}: \gEs{Hohe
Transportpriorit\"at}: Der rasche Transport zur fr\"uhzeitigen
Fachbehandlung nach \"arztlicher Notversorgung ist angesagt. }

{\mdseries \gEs{B}:
\gEs{Niedrige Transportpriorit\"at}: Der Transport
zur Fachbehandlung kann verz\"ogert erfolgen. }

{\bfseries\color{red} wird sp\"ater entschieden}

{\mdseries\color{black} Zweiter Triageblock:}

\begin{itemize}
\item {\mdseries F\"ur Umtriage oder}

\item {\mdseries Spitalstriage}

\end{itemize}
\\
\end{tabular}
\end{center}
\begin{center}
\tablehead{}\begin{supertabular}{m{6.649cm}m{6.651cm}}
\multicolumn{2}{m{13.5cm}}{\centering
{\bfseries\upshape\color[rgb]{0.0,0.0,0.5019608} PLS Umtriage}\par

}\\
[Warning: Draw object ignored][Warning: Draw object ignored][Warning:
Draw object ignored][Warning: Draw object ignored][Warning: Draw object
ignored][Warning: Draw object ignored]

\begin{center}
 [Warning: Image ignored] % Unhandled or unsupported graphics:
%\includegraphics[width=4.457cm,height=9.22cm]{AASdev20090228-img146}

\end{center}
 &
{\mdseries\color{black} Umtriage innerhalb der ersten Viertelstunde}

\begin{itemize}
\item {\mdseries Mittels Pfeil von einem zum nun richtigen Triagefeld}

\item {\mdseries In der gleichen Zeile}

\end{itemize}
\\
[Warning: Draw object ignored][Warning: Draw object ignored][Warning:
Draw object ignored][Warning: Draw object ignored][Warning: Draw object
ignored]

\begin{center}
 [Warning: Image ignored] % Unhandled or unsupported graphics:
%\includegraphics[width=4.457cm,height=9.063cm]{AASdev20090228-img147}

\end{center}
 &
{\mdseries\color{black} Umtriage nach l\"angerer Zeit:}

\begin{itemize}
\item {\mdseries Neues Kreuz beim richtigen Triagefeld}

\item {\mdseries in der unteren Zeile}

\end{itemize}
\\
\end{supertabular}
\end{center}
\begin{center}
\begin{tabular}{m{6.649cm}m{6.651cm}}
\multicolumn{2}{m{13.5cm}}{\centering
{\bfseries\upshape\color[rgb]{0.0,0.0,0.5019608} PLT R\"uckseite}\par

}\\
[Warning: Draw object ignored][Warning: Draw object ignored][Warning:
Draw object ignored][Warning: Draw object ignored][Warning: Draw object
ignored][Warning: Draw object ignored][Warning: Draw object
ignored][Warning: Draw object ignored][Warning: Draw object ignored]

\begin{center}
 [Warning: Image ignored] % Unhandled or unsupported graphics:
%\includegraphics[width=6.653cm,height=12.217cm]{AASdev20090228-img148}

\end{center}
 &
{\mdseries Der Triagearzt tr\"agt auf der R\"uckseite die zu setzenden
Ma{\ss}nahmen ein:}

\begin{itemize}
\item {\mdseries Atmung (O2, Intubation, usw..)}

\item {\mdseries Kreislauf (Blutstillung usw...)}

\item {\mdseries Medikamente}

\item {\mdseries Lagerung}

\end{itemize}
Der Durchf\"uhrende (Arzt, NFS, RS) best\"atigt die Durchf\"uhrung der
Ma{\ss}nahmen im rechten Feld durch eine Paraphe

\\
 &
{\mdseries\color{black} Abriss f\"ur das Zielspital:}

{\mdseries Wird bei der Einlieferung mit der Ankunftszeit versehen und
verbleibt in der Krankengeschichte.}

\\
 &
{\mdseries\color{black} Abriss f\"ur die SanHist -- Leiter Transport:}

{\mdseries Eintragen: Zielspital, Kfz Nummer, Verladezeit}

{\mdseries Verbleibt beim Leiter Transport/Protokollf\"uhrer}

\\
\end{tabular}
\end{center}
\subsection{Gefahrengutunfall}
\begin{center}
\tablehead{}\begin{supertabular}{m{6.649cm}m{6.651cm}}
  [Warning: Image ignored] % Unhandled or unsupported graphics:
%\includegraphics[width=6.233cm,height=4.351cm]{AASdev20090228-img149.jpg}
 

 &
{\raggedleft   [Warning: Image ignored]
% Unhandled or unsupported graphics:
%\includegraphics[width=6.314cm,height=4.41cm]{AASdev20090228-img150.jpg}
 \par}

\\
\end{supertabular}
\end{center}
\begin{center}
\tablehead{}\begin{supertabular}{|m{6.649cm}|m{6.651cm}|}
\hline
{\mdseries\color{black} M\"ogliche Gefahren f\"ur den Sanit\"ater}

\begin{itemize}
\item {\mdseries Brand bzw. Explosionsgefahr}

\item {\mdseries Entweichen unter Druck stehender Gase}

\item {\mdseries Giftigkeit, Ansteckungsgefahr}

\item {\mdseries \"Atzwirkung}

\item {\mdseries Radioaktivit\"at}

\item {\mdseries Umweltgef\"ahrdung}

\end{itemize}
 &
{\mdseries\color{black} Gefahrenquellen k\"onnen entstehen bei}

\begin{itemize}
\item {\mdseries Unf\"alle bei der Herstellung}

\item {\mdseries Unf\"alle bei der Verwendung von gef\"ahrlichen
Stoffen}

\item {\mdseries Transportunf\"alle}

\end{itemize}
\\\hline
\end{supertabular}
\end{center}
{\mdseries\color{black}
Andere Gefahrenquellen:}

\begin{flushleft}
\tablehead{}\begin{supertabular}{|m{6.393cm}m{6.9030004cm}|}
\hline
{\mdseries Landwirtschaft}

 &
{\mdseries CO2 in G\"arkellern, Siloanlagen, Brunnensch\"achte}

{\mdseries Pflanzenschutzmittel}

\\\hline
{\mdseries Kl\"argruben}

 &
{\mdseries CO2 Methan, Schwefelwasserstoff}

\\\hline
{\mdseries Obstlagerhallen}

 &
{\mdseries wenig O2, viel CO2 oder N2 halten Obst frisch}

\\\hline
{\mdseries Br\"ande in Wirtschaftsgeb\"auden}

 &
{\mdseries extrem gef\"ahrliche Rauchgase}

\\\hline
{\mdseries Tiefgaragen und Tunnels}

 &
{\mdseries erh\"ohte CO- und CO2-Werte}

\\\hline
{\mdseries Fl\"ussiggas}

 &
{\mdseries Campingpl\"atze, feste und mobile K\"uchen}

\\\hline
{\mdseries Pools, Schwimmb\"ader}

 &
{\mdseries Chlor (Wasserdesinfektion)}

\\\hline
{\mdseries Industrie}

 &
{\mdseries K\"uhlmittel in gro{\ss}en K\"alteanlagen (Ammoniak)}

\\\hline
{\mdseries Speziallabors, Einrichtungen mit und zur Verarbeitung von
infekti\"osen Abf\"allen oder Tierkadaver}

 &
{\mdseries Infektionsgefahr}

\\\hline
{\mdseries Nuklearmedizinische Institute, Industrie,
Forschungseinrichtungen}

 &
{\mdseries Radioaktivit\"at}

\\\hline
\end{supertabular}
\end{flushleft}
{\mdseries\color[rgb]{0.0,0.0,0.5019608}
Kennzeichnung von gef\"ahrlichen Stoffen und Gefahrenbereichen}

\begin{center}
\begin{tabular}{|m{3.225cm}m{3.225cm}|m{3.225cm}m{3.225cm}|}
\hline
\multicolumn{2}{|m{6.65cm}|}{\centering
{\bfseries\upshape\color[rgb]{0.0,0.0,0.5019608} Hinweis auf
Gefahr}\par

} &
\multicolumn{2}{m{6.65cm}|}{\centering
{\bfseries\upshape\color[rgb]{0.0,0.0,0.5019608} Warnzeichen}\par

}\\\hline
{\centering   [Warning: Image ignored]
% Unhandled or unsupported graphics:
%\includegraphics[width=1.475cm,height=1.452cm]{AASdev20090228-img151.png}
 \par}

 &
{\mdseries Explosionsgef\"ahrlich}

 &
{\centering   [Warning: Image ignored]
% Unhandled or unsupported graphics:
%\includegraphics[width=1.541cm,height=1.426cm]{AASdev20090228-img152.png}
 \par}

 &
{\mdseries Allgemeine Gefahr}

\\\hline
{\centering   [Warning: Image ignored]
% Unhandled or unsupported graphics:
%\includegraphics[width=1.452cm,height=1.414cm]{AASdev20090228-img153.png}
 \par}

 &
{\mdseries Giftig}

 &
{\centering   [Warning: Image ignored]
% Unhandled or unsupported graphics:
%\includegraphics[width=1.673cm,height=1.414cm]{AASdev20090228-img154.png}
 \par}

 &
{\mdseries Radioaktive Stoffe, ionisierende Strahlen}

\\\hline
{\centering   [Warning: Image ignored]
% Unhandled or unsupported graphics:
%\includegraphics[width=1.497cm,height=1.611cm]{AASdev20090228-img155.png}
 \par}

 &
{\mdseries Hochentz\"undlich}

 &
{\centering   [Warning: Image ignored]
% Unhandled or unsupported graphics:
%\includegraphics[width=1.696cm,height=1.412cm]{AASdev20090228-img156.png}
 \par}

 &
{\mdseries Gasflaschen}

\\\hline
{\centering   [Warning: Image ignored]
% Unhandled or unsupported graphics:
%\includegraphics[width=1.475cm,height=1.435cm]{AASdev20090228-img157.png}
 \par}

 &
{\mdseries Reizend}

 &
{\centering   [Warning: Image ignored]
% Unhandled or unsupported graphics:
%\includegraphics[width=1.673cm,height=1.414cm]{AASdev20090228-img158.png}
 \par}

 &
{\mdseries Laserstrahlen}

\\\hline
{\centering   [Warning: Image ignored]
% Unhandled or unsupported graphics:
%\includegraphics[width=1.488cm,height=1.599cm]{AASdev20090228-img159.png}
 \par}

 &
{\mdseries \"Atzend}

 &
 &
\\\hline
\end{tabular}
\end{center}
{\mdseries\color[rgb]{0.0,0.0,0.5019608}
Transport gef\"ahrlicher G\"uter}

{\mdseries
Kennzeichnungspflicht besteht erst dann, wenn die transportierte Menge
eines Stoffes die gesetzlich erlaubte Mindestmenge \"uberschreitet.
Diese Mindestmenge gilt f\"ur jeden einzelnen Stoff und nicht f\"ur die
Gesamtmenge.}

{\mdseries
Es k\"onnen also erhebliche Mengen gef\"ahrlicher G\"uter ohne
Kennzeichnung des Fahrzeuges transportiert werden. }

{\mdseries\color{black}
Gefahrzettel}

{\mdseries
Bei St\"uckguttransporten wird jedes Versandst\"uck gekennzeichnet. An
Fahrzeugen, die kein St\"uckgut transportieren m\"ussen Gefahrenzettel
an den L\"angsseiten und an der R\"uckseite angebracht werden.}

\begin{flushleft}
\tablehead{}\begin{supertabular}{m{2.1599998cm}m{4.289cm}m{2.1599998cm}m{4.2840004cm}}
\centering {\bfseries\upshape\color[rgb]{0.0,0.0,0.5019608}
Beispiele}\par

 &
 &
 &
\\
{\centering   [Warning: Image ignored]
% Unhandled or unsupported graphics:
%\includegraphics[width=1.659cm,height=1.885cm]{AASdev20090228-img160.png}
 \par}

 &
{\mdseries Feuergef\"ahrlich (fl\"ussige Stoffe)}

 &
{\centering   [Warning: Image ignored]
% Unhandled or unsupported graphics:
%\includegraphics[width=1.856cm,height=2.096cm]{AASdev20090228-img161.png}
 \par}

 &
{\mdseries Selbstentz\"undlich}

\\
{\centering   [Warning: Image ignored]
% Unhandled or unsupported graphics:
%\includegraphics[width=1.835cm,height=1.902cm]{AASdev20090228-img162.png}
 \par}

 &
{\mdseries Feuergef\"ahrlich (feste Stoffe)}

 &
{\centering   [Warning: Image ignored]
% Unhandled or unsupported graphics:
%\includegraphics[width=1.909cm,height=2.162cm]{AASdev20090228-img163.png}
 \par}

 &
{\mdseries Entz\"undbare Gase bei Ber\"uhrung mit Wasser}

\\
\end{supertabular}
\end{flushleft}
{\mdseries\color{black}
Warntafel}

\begin{center}
\begin{tabular}{|m{3.225cm}m{3.225cm}|m{3.225cm}m{3.225cm}|}
\hline
\multicolumn{2}{|m{6.65cm}|}{\centering
{\bfseries\upshape\color[rgb]{0.0,0.0,0.5019608}
AllgemeineKennzeichnung (Sammeltransporte)}\par

} &
\multicolumn{2}{m{6.65cm}|}{\centering
{\bfseries\upshape\color[rgb]{0.0,0.0,0.5019608}
SpezielleKennzeichnung}\par

}\\\hline
{\centering [Warning: Draw object ignored]\par}

 &
 &
{\centering [Warning: Draw object ignored]\par}

 &
{\mdseries Gefahrnummer  (Kemler-Nummer)}

\\\hline
 &
 &
 &
{\mdseries Stoffnummer  (UN-Nummer)}

\\\hhline{~~~-}
\end{tabular}
\end{center}
{\mdseries\color{black}
Gefahrnummer}

\begin{flushleft}
\tablehead{}\begin{supertabular}{|m{0.97999996cm}m{12.317cm}|}
\hline
\centering 2\par

 &
{\mdseries Entweichen von Gas durch Druck oder chemische Reaktion}

\\\hline
\centering 3\par

 &
{\mdseries Entz\"undbarkeit von Fl\"ussigkeiten (D\"ampfen) und Gasen}

\\\hline
\centering 4\par

 &
{\mdseries Entz\"undbarkeit fester Stoffe}

\\\hline
\centering 5\par

 &
{\mdseries Oxidierende (brandf\"ordernde) Wirkung}

\\\hline
\centering 6\par

 &
{\mdseries Giftig}

\\\hline
\centering 7\par

 &
{\mdseries Radioaktivit\"at}

\\\hline
\centering 8\par

 &
{\mdseries \"Atzwirkung}

\\\hline
\centering 9\par

 &
{\mdseries Gefahr einer spontan heftigen Reaktion}

\\\hline
\end{supertabular}
\end{flushleft}
{\mdseries
Die Verdopplung einer Ziffer weist auf die Zunahme der entsprechenden 
Gefahr hin.}

{\mdseries
\gT{X}  vor der Nummer zur Kennzeichnung der Gefahr
${\rightarrow}$ \gEs{Stoff reagiert in 
gef\"ahrlicher Weise mit Wasser}}

{\mdseries
Wenn die Gefahr eines Stoffes ausreichend von einer einzigen Ziffer 
angegeben werden kann, wird dieser Ziffer eine Null angef\"ugt.}

{\mdseries\color{black}
Folgende Ziffernkombinationen haben eine
\gEs{besondere Bedeutung}}

\begin{flushleft}
\begin{tabular}{|m{0.97999996cm}m{12.317cm}|}
\hline
{\mdseries 22}

 &
{\mdseries tiefgek\"uhltes Gas}

\\\hline
{\mdseries X323}

 &
{\mdseries entz\"undbarer fl\"ussiger Stoff, der mit Wasser gef\"ahrlich
reagiert, wobei entz\"undbare Gase entweichen}

\\\hline
{\mdseries X333}

 &
{\mdseries selbstentz\"undliche Fl\"ussigkeit, die mit Wasser
gef\"ahrlich reagiert}

\\\hline
{\mdseries X423}

 &
{\mdseries entz\"undbarer fester Stoff, der mit Wasser gef\"ahrlich
reagiert, wobei brennbare Gase entweichen}

\\\hline
{\mdseries 44}

 &
{\mdseries entz\"undbarer fester Stoff, der sich bei erh\"ohter
Temperatur in Geschmolzenem Zustand befindet}

\\\hline
{\mdseries 539}

 &
{\mdseries entz\"undbares organisches Peroxid}

\\\hline
{\mdseries 90}

 &
{\mdseries verschiedene gef\"ahrlich Stoffe}

\\\hline
\end{tabular}
\end{flushleft}
{\mdseries\color[rgb]{0.0,0.0,0.5019608}
Einsatzrichtlinie}

{\mdseries
\textcolor{red}{G}efahr erkennen}

{\mdseries
\textcolor{red}{A}bsperrma{\ss}nahmen durchf\"uhren}

{\mdseries
\textcolor{red}{S}pezialkr\"afte anfordern}

\begin{center}
\tablehead{}\begin{supertabular}{|m{2.222cm}|m{5.1460004cm}|m{5.7320004cm}|}
\hline
{\mdseries\color{black} \textcolor{red}{G}efahr erkennen}

 &
Die Suche nach Warntafeln ist bei umgest\"urzten Fahrzeugen erschwert

 &
{\centering   [Warning: Image ignored]
% Unhandled or unsupported graphics:
%\includegraphics[width=5.137cm,height=3.747cm]{AASdev20090228-img164.jpg}
 \par}

\\\hline
{\mdseries\color{black} \textcolor{red}{A}bsperr\-ma{\ss}nahmen
durch\-f\"uhren}

 &
Die Absperrung sollte gro{\ss}r\"aumig  sein. Mit gen\"ugend Platz f\"ur
die technische Hilfeleistungen. Achtung  vor sich \"andernden
Situationen  (Windrichtung, Niederschlag)

 &
{\centering   [Warning: Image ignored]
% Unhandled or unsupported graphics:
%\includegraphics[width=3.46cm,height=3.409cm]{AASdev20090228-img165.png}
 \par}

\\\hline
{\mdseries\color{black} \textcolor{red}{S}pezial\-kr\"afte anfordern}

 &
 &
{\centering   [Warning: Image ignored]
% Unhandled or unsupported graphics:
%\includegraphics[width=5.666cm,height=3.941cm]{AASdev20090228-img166.jpg}
 \par}

\\\hline
\end{supertabular}
\end{center}
\section{Fallbeispiele und Arbeitsaufgaben}
{\mdseries\color[rgb]{0.0,0.0,0.5019608}
Fallbeispiele und Arbeitsaufgaben f\"ur BLS}

{\mdseries\color{black}
Fall }

\begin{flushleft}
\tablehead{}\begin{supertabular}{|m{2.0839999cm}|m{1.43cm}m{1.43cm}m{1.43cm}m{1.43cm}m{1.43cm}m{1.43cm}m{1.43cm}|}
\hline
\centering {\bfseries\upshape\color[rgb]{0.0,0.0,0.5019608}
Berufung}\par

 &
\multicolumn{1}{m{1.43cm}|}{\centering Typ\par

} &
\multicolumn{1}{m{1.43cm}|}{\centering {\bfseries RTW}\par

} &
\multicolumn{1}{m{1.43cm}|}{\centering Code\par

} &
\multicolumn{1}{m{1.43cm}|}{\centering {\bfseries 26-A-1}\par

} &
\multicolumn{3}{m{4.69cm}|}{}\\\hline
 &
\multicolumn{1}{m{1.43cm}|}{} &
\multicolumn{1}{m{1.43cm}|}{} &
\multicolumn{1}{m{1.43cm}|}{} &
\multicolumn{1}{m{1.43cm}|}{} &
\multicolumn{1}{m{1.43cm}|}{} &
\multicolumn{1}{m{1.43cm}|}{} &
\\\hline
\centering {\bfseries\upshape\color[rgb]{0.0,0.0,0.5019608}
Umst\"ande}\par

 &
\multicolumn{1}{m{1.43cm}|}{\centering Jahreszeit\par

} &
\multicolumn{1}{m{1.43cm}|}{\centering Tageszeit\par

} &
\multicolumn{1}{m{1.43cm}|}{\centering Wetter\par

} &
\multicolumn{1}{m{1.43cm}|}{} &
\multicolumn{1}{m{1.43cm}|}{\centering Ort\par

} &
\multicolumn{1}{m{1.43cm}|}{} &
\\\hline
 &
\multicolumn{1}{m{1.43cm}|}{\centering {\bfseries Sommer}\par

} &
\multicolumn{1}{m{1.43cm}|}{\centering {\bfseries Vormittag}\par

} &
\multicolumn{2}{m{3.06cm}|}{\centering {\bfseries schw\"ul,
hei{\ss}}\par

} &
\multicolumn{1}{m{1.43cm}|}{\centering {\bfseries Wohnung}\par

} &
\multicolumn{1}{m{1.43cm}|}{} &
\\\hhline{~-------}
\centering {\bfseries\upshape\color[rgb]{0.0,0.0,0.5019608}
Patient:}\par

 &
\multicolumn{1}{m{1.43cm}|}{\centering Alter\par

} &
\multicolumn{1}{m{1.43cm}|}{\centering Geschl.\par

} &
\multicolumn{1}{m{1.43cm}|}{\centering AZ\par

} &
\multicolumn{1}{m{1.43cm}|}{\centering EZ\par

} &
\multicolumn{1}{m{1.43cm}|}{} &
\multicolumn{1}{m{1.43cm}|}{} &
\\\hline
 &
\multicolumn{1}{m{1.43cm}|}{\centering {\bfseries 65}\par

} &
\multicolumn{1}{m{1.43cm}|}{\centering {\bfseries m}\par

} &
\multicolumn{1}{m{1.43cm}|}{\centering {\bfseries red.}\par

} &
\multicolumn{1}{m{1.43cm}|}{\centering {\bfseries gut}\par

} &
\multicolumn{1}{m{1.43cm}|}{} &
\multicolumn{1}{m{1.43cm}|}{} &
\\\hhline{~-------}
{\bfseries\color[rgb]{0.0,0.0,0.5019608} Szene}

 &
\multicolumn{7}{m{11.21cm}|}{{\bfseries Patient sitzt bla{\ss} in einem
Lehnstuhl, ansprechbar, wirkt nicht akut bedroht }

}\\\hline
{\bfseries\color[rgb]{0.0,0.0,0.5019608} Beschwerden}

 &
\multicolumn{7}{m{11.21cm}|}{{\bfseries
{\quotedblbase}Schwindel{\textquotedblleft}, erstmalig aufgetreten,
keine Schmerzen}

}\\\hline
{\bfseries\color[rgb]{0.0,0.0,0.5019608} Krankheiten}

 &
\multicolumn{7}{m{11.21cm}|}{{\bfseries {\quotedblbase}Hypertonie,
Gef\"a{\ss}verkalkung, Angina pectoris{\textquotedblleft}}

{\bfseries Blutdruckprotokoll: sie finden f\"ur die vergangenen vier
Tage folgende Eintr\"age: 170/110, 180/115, 175/113, 183/112, ...}

}\\\hline
{\bfseries\color[rgb]{0.0,0.0,0.5019608} Allergien}

 &
\multicolumn{7}{m{11.21cm}|}{{\bfseries keine bekannt}

}\\\hline
{\bfseries\color[rgb]{0.0,0.0,0.5019608} Medikamente}

 &
\multicolumn{7}{m{11.21cm}|}{{\bfseries Patient \"uberreicht
Liste:{\quotedblbase}Blutdruckmittel, aber die wirken nicht so richtig
sagt mein Arzt{\textquotedblleft}, Pat zeigt eine Sprayflasche
{\quotedblbase}Nitrolingual{\textquotedblleft} her den er gestern vom
Hausarzt {\quotedblbase}f\"ur sein Herz{\textquotedblleft} bekommen
hat.}

}\\\hline
{\bfseries\color[rgb]{0.0,0.0,0.5019608} Ereignisse}

 &
\multicolumn{7}{m{11.21cm}|}{{\bfseries {\quotedblbase}Ich habe das
erste mal diesen Nitro-Spray genommen und da ist mir damisch
word{\textquotesingle}n{\textquotedblleft}}

}\\\hline
{\bfseries\color[rgb]{0.0,0.0,0.5019608} Letzte(r)}

 &
\multicolumn{7}{m{11.21cm}|}{{\bfseries Krankenhausaufenthalt: vor 3
Jahren wg. Blinddarmentz\"undung}

{\bfseries Arztbesuch: gestern, wegen Herzschmerzen, Nitro-Spray
verschrieben bekommen}

{\bfseries Mahlzeit: Fr\"uhst\"uck, 2 Marmeladesemmeln, Kaffee}

}\\\hline
\multicolumn{8}{|m{13.494cm}|}{\centering
{\bfseries\upshape\color[rgb]{0.0,0.0,0.5019608} Erstbefund}\par

}\\\hline
Anblick

 &
\multicolumn{7}{m{11.21cm}|}{{\bfseries altersensprechend, gut
gen\"ahrt, gepflegt, jetzt schwach}

}\\\hline
Haut

 &
\multicolumn{7}{m{11.21cm}|}{{\bfseries bla{\ss}, nicht schweissig}

}\\\hline
Trauma grob

 &
\multicolumn{7}{m{11.21cm}|}{{\bfseries {}-{}-{}-}

}\\\hline
 &
\multicolumn{1}{m{1.43cm}|}{\centering RR\par

} &
\multicolumn{1}{m{1.43cm}|}{\centering HF\par

} &
\multicolumn{1}{m{1.43cm}|}{\centering Puls\par

} &
\multicolumn{1}{m{1.43cm}|}{\centering SpO2\par

} &
\multicolumn{1}{m{1.43cm}|}{\centering AF\par

} &
\multicolumn{1}{m{1.43cm}|}{\centering BZ\par

} &
\centering Temp\par

\\\hline
 &
\multicolumn{1}{m{1.43cm}|}{\centering {\bfseries 100/70}\par

} &
\multicolumn{1}{m{1.43cm}|}{\centering {\bfseries 120}\par

} &
\multicolumn{1}{m{1.43cm}|}{\centering {\bfseries rhythm.}\par

} &
\multicolumn{1}{m{1.43cm}|}{\centering {\bfseries 98\%}\par

} &
\multicolumn{1}{m{1.43cm}|}{\centering {\bfseries 12}\par

} &
\multicolumn{1}{m{1.43cm}|}{\centering {\bfseries 105}\par

} &
\centering {\bfseries 37,1}\par

\\\hline
Trauma Detail

 &
\multicolumn{7}{m{11.21cm}|}{{}-{}-{}-

}\\\hline
Neuro

 &
\multicolumn{7}{m{11.21cm}|}{{\bfseries grob neurologisch
unauff\"allig.}

}\\\hline
Ma{\ss}nahmen, Verlauf

 &
\multicolumn{1}{m{1.43cm}|}{} &
\multicolumn{5}{m{7.95cm}|}{} &
\\\hline
Transport

 &
\multicolumn{1}{m{1.43cm}|}{\centering Lagerung\par

} &
\multicolumn{5}{m{7.95cm}|}{} &
\\\hline
\"Ubergabe

 &
\multicolumn{1}{m{1.43cm}|}{\centering Abteilung:\par

} &
\multicolumn{5}{m{7.95cm}|}{} &
\\\hline
\end{supertabular}
\end{flushleft}
{\mdseries
Wie beurteilen Sie den Blutdruck des Patienten? }

{\mdseries
Welche Informationen fehlen Ihnen noch?}

{\mdseries
Wie lautet Ihre Verdachtsdiagnose? }

Welche Differentialdiagnosen ziehen Sie in Betracht?



\section{Impressum, Hinweise}
\begin{multicols}{2}
{\mdseries\color{black}
ARBEITER-SAMARITER-BUND \"OSTERREICHS GRUPPE FLORIDSDORF-DONAUSTADT}

{\mdseries
WALLENBERGGASSE 2}

{\mdseries
A-1220 WIEN}

{\mdseries
TEL. +43 1 26 33 144}

{\mdseries
FAX  +43 1 26 33 144 75}

{\mdseries
\href{http://WWW.SAMARITER.AT/}{WWW.SAMARITER.AT}}

{\mdseries
\href{mailto:OFFICE@SAMARITER.AT}{OFFICE@SAMARITER.AT}}

{\mdseries
ZVR: 317 994 238  UID: ATU 163 722 08  HG WIEN}

{\mdseries\color{black}
Betriebsgesellschaft und Rechnungsanschrift:}

ARBEITER-SAMARITER-BUND \"OSTERREICHS FLORIDSDORF-DONAUSTADT KRANKEN-,
RETTUNGSTRANSPORT UND SOZIALE DIENSTE GEM. GESMBH\newline
Kontakt und Adresse w. o.\newline
FN 287700t, HG Wien\newline
UID: ATU 63216914

Gesch\"aftsf\"uhrer: Berhnard Lehner, Ing. Alexander Prischl

\"Arztliche Leitung: Dr. Johannes Steuer, MSc,  Stv: Dr. Helmut Paul

Leitung Ausbildungszentrum: Ing. Roman Koch

Aufsichtsapotheke: Mag.pharm. Ulrike Urban, Stadtapotheke
Klosterneuburg. 

{\color{black}
Hinweise}

Alle Angaben sind rein informativ zu ver\-stehen und begr\"unden KEINE
Dienst\-an\-wei\-sung oder Lehr\-meinung. 

Wie jede Wissenschaft ist die Medizin st\"andigen Ent\-wicklungen
unterworfen. Soweit fachliche Angaben ge\-macht werden, darf der Leser
zwar darauf ver\-trau\-en, dass die Autoren gro{\ss}e Sorg\-falt
da\-rauf verwandt haben, dass diese An\-ga\-ben dem Wissens\-stand bei
Ferti\-gstel\-lung des Werkes entspricht. F\"ur Angaben, spe\-ziell
\"uber Dosierungsanweisungen und Ap\-plikations\-formen, kann jedoch
keine Ge\-w\"ahr \"ubernommen werden. Jeder Leser ist angehalten, durch
sorg\-f\"altige Pr\"u\-fung der Literatur, der Lehrmeinung so\-wie der
Beipackzettel der ver\-wendeten Pr\"aparate und gegebenenfalls nach
Konsultation eines Spezialisten fest\-zustellen, ob die gegebene
Aussage oder Empfehlung f\"ur Do\-sierungen oder die Beachtung von
Kon\-tra\-indi\-ka\-tionen gegen\"uber der Angabe in dieser
Mit\-tei\-lung abweicht. Jede Dosierung oder Ap\-pli\-kation erfolgt
auf eigene Gefahr des Be\-nutzers. Das AMS appelliert an jeden Le\-ser
ihm etwa auffallende Un\-ge\-nau\-igkei\-ten dem AMS umgehend
mitzuteilen.

Gesch\"utzte Waren\-namen und Wa\-ren\-zei\-chen werden nicht besonders
kenntlich ge\-macht. Aus dem Fehlen eines solchen Hin\-weises kann also
nicht ge\-schlossen werden, dass es sich um einen freien Waren\-namen
handelt.

\end{multicols}
{\mdseries\color{black}
Sprechstunden}

{\mdseries
Mittwoch nachmittag. }

\printindex
\listoffigures
SO LONG..........

AND THANKS FOR ALL THE FISH

\section[Literaturverzeichnis]{Literaturverzeichnis}
{\mdseries
[1] , \textbf{2002}. \textit{Pschyrembel Klinisches W\"orterbuch}. de
Gruyter, .}

{\mdseries
[2] Longmore, Wilkinson, T. C. e., \textbf{2007}. \textit{Oxford
Handbook of Clinical Medicine}. Oxford University Press, .}

\end{document}
